\documentclass[10pt,a4paper]{extarticle}

\usepackage{amsmath}
\usepackage{fontspec}
\usepackage{unicode-math}
\setmainfont{CharisSIL-Regular.ttf}[BoldFont=CharisSIL-Bold.ttf,ItalicFont=CharisSIL-Italic.ttf,BoldItalicFont=CharisSIL-BoldItalic.ttf]
\setmathfont{STIXTwoMath-Regular.otf}
\setmonofont{consola.ttf}[BoldFont=consolab.ttf,ItalicFont=consolai.ttf,Scale=0.9]
\usepackage[top=1.9cm,bottom=2.1cm,left=1.9cm,right=1.9cm]{geometry}
\usepackage{booktabs}
\usepackage{array}
\usepackage{xcolor}
\usepackage{titlesec}
\usepackage{fancyhdr}
\usepackage{microtype}
\usepackage{enumitem}
\usepackage{fvextra}
\usepackage{etoolbox}
\setlist{nosep,leftmargin=1.4em}

\definecolor{linkblue}{RGB}{0,106,150}
\definecolor{bannergray}{RGB}{229,229,229}
\usepackage[colorlinks=true,linkcolor=linkblue,urlcolor=linkblue]{hyperref}

\newcommand{\sectionname}{Chương}
\setcounter{secnumdepth}{2}
\setcounter{tocdepth}{2}
\titleformat{\section}{\Large\bfseries\color{linkblue}}{\sectionname\ \thesection.}{0.5em}{}
\titleformat{name=\section,numberless}{\Large\bfseries\color{linkblue}}{}{0em}{}
\titleformat{\subsection}{\large\bfseries}{\thesubsection.}{0.5em}{}
\titleformat{name=\subsection,numberless}{\large\bfseries}{}{0em}{}
\titleformat{\subsubsection}{\normalsize\bfseries\color{black!80}}{}{0em}{}
\titleformat{name=\subsubsection,numberless}{\normalsize\bfseries\color{black!80}}{}{0em}{}
\titleformat{\paragraph}[runin]{\small\bfseries}{}{0em}{}
\titlespacing*{\section}{0pt}{12pt}{6pt}
\titlespacing*{\subsection}{0pt}{12pt}{4pt}
\titlespacing*{\subsubsection}{0pt}{9pt}{2pt}
\titlespacing*{\paragraph}{0pt}{6pt}{6pt}
\setlength{\parskip}{3pt}
\setlength{\parindent}{0pt}
\setlength{\abovedisplayskip}{4pt}
\setlength{\belowdisplayskip}{4pt}
\emergencystretch=3em
\linespread{1.05}
\raggedbottom

\DeclareRobustCommand{\code}[1]{\texttt{\detokenize{#1}}}
\newcommand{\src}[1]{\par\noindent{\small\textit{Mã nguồn:} \texttt{\detokenize{#1}}}\par\vspace{2pt}}
\newcommand{\Sig}{\symbf{\Sigma}}
\newcommand{\vv}[1]{\symbf{#1}}

\pagestyle{fancy}
\fancyhf{}
\renewcommand{\headrulewidth}{0pt}
\fancyhead[L]{\footnotesize\textit{L.\,A.\ Kiet}}
\fancyhead[R]{\footnotesize\textit{Splat Editor: sách hướng dẫn cấu trúc chương trình}}
\fancyfoot[C]{\footnotesize\thepage}

\begin{document}
\begin{center}
\noindent\colorbox{bannergray}{\parbox{\dimexpr\textwidth-2\fboxsep}{\centering\vspace{6pt}
  {\small Đồ họa máy tính \;$\cdot$\; Trường Đại học Quốc tế, ĐHQG-HCM}\\[8pt]
  {\Huge Splat Editor}\\[6pt]
  {\Large Sách hướng dẫn cấu trúc chương trình}\\[4pt]
  {\small Giải thích từng thư mục, từng file, từng hàm, kèm mã nguồn thật và toán học}\vspace{6pt}}}
\end{center}
\vspace{4pt}

\noindent\textbf{Cách đọc cuốn sách.} Sách đi theo \emph{cấu trúc thư mục} của mã nguồn: mỗi chương là một thư mục, mỗi mục là một file, mỗi tiểu mục là một lớp hoặc một hàm.
Với mỗi hàm bạn có ba thứ: \textbf{mô tả} (chú thích của tác giả trong mã), \textbf{giải thích} (phần diễn giải thêm) và \textbf{mã nguồn nguyên văn} kèm số dòng thật
trong file. Cơ sở toán nằm ngay trong chương của file tương ứng, ở các mục \textit{Cơ sở toán}; các chỗ có ghi $\S$ trỏ tới mục đó, và Chương~\ref{ch:quickref} có bảng tra nhanh công thức--file.

\noindent\textbf{Quy ước.} Vector cột, tọa độ tay phải, trục $Y$ hướng lên, camera nhìn dọc $-z$. Ma trận $4\times4$ dùng tọa độ thuần nhất $(x,y,z,1)$.

\tableofcontents
\clearpage

\renewcommand{\sectionname}{Chương}
\section*{Tổng quan}
\addcontentsline{toc}{section}{Tổng quan: bản đồ chương trình và điểm vào}
\label{ch:overview}
\subsection*{Cấu trúc thư mục}
\begin{small}
\begin{Verbatim}[fontsize=\fontsize{6.8}{7.8}\selectfont]
SplatOpenGL/
+- main.py ................. diem vao
+- splat_editor/
|  +- __main__.py, app.py .. khoi dong, vong lap chinh (Tong quan)
|  +- core/    (Chuong 1)   toan thuan NumPy, KHONG co OpenGL
|  +- editor/  (Chuong 2)   chon, xoa, undo, co chon
|  +- io/      (Chuong 3)   doc/ghi PLY, giai ma/ma hoa
|  +- render/  (Chuong 4)   OpenGL: camera, bo ve gaussian, luong sap xep
|  |   +- shaders/ (Chuong 6)   ma GLSL .vert / .frag
|  +- ui/      (Chuong 5)   GLFW + Dear ImGui
+- report/ ................. bao cao LaTeX goc
+- guideline/report_src/ .... sach huong dan nay
\end{Verbatim}
\end{small}

\subsection*{Năm tầng và quan hệ phụ thuộc}
\begin{center}\small
\begin{tabular}{@{}llp{0.5\textwidth}@{}}\toprule
Chương & Thư mục & Phụ thuộc vào\\\midrule
1 & \code{core} & chỉ \code{io} (đọc/ghi PLY)\\
2 & \code{editor} & \code{core}\\
3 & \code{io} & \code{core.splats} (kiểu dữ liệu đích)\\
4 & \code{render} & \code{core}, \code{shaders/}, OpenGL\\
5 & \code{ui} & \code{editor}, \code{render}, \code{core}, GLFW, ImGui\\
6 & \code{render/shaders} & được nạp bởi \code{render/shader.py}\\
\bottomrule
\end{tabular}
\end{center}
\textbf{Quy tắc:} \code{core}, \code{editor} và \code{io} không biết OpenGL hay giao diện; chỉ \code{render} chạm OpenGL và chỉ \code{ui} chạm ImGui. Nhờ vậy phần logic có thể chạy và kiểm thử không cần cửa sổ.

\subsection*{Đường đi của dữ liệu}
\textbf{PLY} $\to$ \code{io} (\code{PlyFile}, \code{SplatCodec}) $\to$ \code{core} (\code{Layer}, \code{Scene}, \code{Splats}) $\to$ \code{editor} (\code{Editor}) $\to$ \code{render} (\code{Viewport}, \code{SplatRenderer}) $\to$ \textbf{shader} $\to$ màn hình.
Song song là camera, cọ, vật lý, ánh sáng, hoạt họa. Mỗi khung hình đi qua \code{Application.run} (mục dưới).
\section*{Điểm vào và vòng lặp chính}
\addcontentsline{toc}{section}{Tổng quan: điểm vào và app.py}
\subsection{\texorpdfstring{\code{main.py}}{main.py}}
\textbf{Vai trò.} Điểm vào của chương trình: chỉ import và gọi \code{splat_editor.app.main}.

\VerbatimInput[firstline=1,lastline=11,firstnumber=1,numbers=left,numbersep=3pt,frame=lines,framesep=2pt,rulecolor=\color{gray!60},breaklines=true,fontsize=\fontsize{6.6}{7.6}\selectfont]{../../main.py}
\subsection{\texorpdfstring{\code{splat_editor/__main__.py}}{splat_editor/__main__.py}}
\textbf{Vai trò.} Cho phép chạy bằng \code{python -m splat_editor}; cùng làm việc với \code{main.py}.

\VerbatimInput[firstline=1,lastline=11,firstnumber=1,numbers=left,numbersep=3pt,frame=lines,framesep=2pt,rulecolor=\color{gray!60},breaklines=true,fontsize=\fontsize{6.6}{7.6}\selectfont]{../../splat_editor/__main__.py}
\subsection{\texorpdfstring{\code{splat_editor/app.py}}{splat_editor/app.py}}
{\footnotesize 506 dòng}\par
\textbf{Vai trò.} \code{Application} ghép mọi thành phần (cửa sổ, ImGui, Editor, Viewport, Physics, Lighting, các bảng giao diện) và chạy vòng lặp chính.
Đây là nơi duy nhất biết \textbf{tất cả} các tầng; các tầng còn lại không import \code{app}. Mỗi khung hình: \code{poll_events} $\to$ \code{input.update}
$\to$ \code{update_world} $\to$ dựng ImGui $\to$ \code{viewport.sync} $\to$ \code{viewport.render} $\to$ \code{gui.draw} $\to$ \code{swap_buffers}.

\paragraph{Mô tả đầu file và cách dùng.} Phần chú thích đầu file do tác giả viết, nêu ý tưởng và ví dụ sử dụng:
\VerbatimInput[firstline=1,lastline=22,firstnumber=1,numbers=left,numbersep=3pt,frame=lines,framesep=2pt,rulecolor=\color{gray!60},breaklines=true,fontsize=\fontsize{6.2}{7.2}\selectfont]{../../splat_editor/app.py}
\subsubsection*{Khai báo mức module (dòng 23--45)}
\VerbatimInput[firstline=23,lastline=45,firstnumber=23,numbers=left,numbersep=3pt,frame=lines,framesep=2pt,rulecolor=\color{gray!60},breaklines=true,fontsize=\fontsize{6.6}{7.6}\selectfont]{../../splat_editor/app.py}
\subsubsection{\texorpdfstring{\code{class FpsCounter}}{FpsCounter}}
{\footnotesize\textit{Dòng 48--59}}\par
Các phương thức: \code{__init__}, \code{tick}.

\VerbatimInput[firstline=48,lastline=48,firstnumber=48,numbers=left,numbersep=3pt,frame=lines,framesep=2pt,rulecolor=\color{gray!60},breaklines=true,fontsize=\fontsize{6.6}{7.6}\selectfont]{../../splat_editor/app.py}
\subsubsection{\texorpdfstring{\code{FpsCounter.__init__(self, interval=0.5)}}{FpsCounter.__init__}}
{\footnotesize\textit{Dòng 49--53}}\par
\VerbatimInput[firstline=49,lastline=53,firstnumber=49,numbers=left,numbersep=3pt,frame=lines,framesep=2pt,rulecolor=\color{gray!60},breaklines=true,fontsize=\fontsize{6.6}{7.6}\selectfont]{../../splat_editor/app.py}
\subsubsection{\texorpdfstring{\code{FpsCounter.tick(self, now)}}{FpsCounter.tick}}
{\footnotesize\textit{Dòng 55--59}}\par
\VerbatimInput[firstline=55,lastline=59,firstnumber=55,numbers=left,numbersep=3pt,frame=lines,framesep=2pt,rulecolor=\color{gray!60},breaklines=true,fontsize=\fontsize{6.6}{7.6}\selectfont]{../../splat_editor/app.py}
\subsubsection{\texorpdfstring{\code{class Application}}{Application}}
{\footnotesize\textit{Dòng 62--500}}\par
Các phương thức: \code{__init__}, \code{_create_actions}, \code{open_file}, \code{add_files}, \code{open_dialog}, \code{add_dialog}, \code{save_dialog}, \code{frame_scene}, \code{_update_bounds}, \code{toggle_lighting}, \code{toggle_shadows}, \code{add_sun}, \code{add_light}, \code{pick_light}, \code{delete_key}, \code{set_mode}, \code{_spawn_player}, \code{toggle_mode}, \code{walk}, \code{cycle_view}, \code{_leave_third_person}, \code{_update_avatar}, \code{update_world}, \code{toggle_grid}, \code{toggle_sh}, \code{set_flip}, \code{exit_brush}, \code{shortcuts_window_toggle}, \code{paint}, \code{drag_speed}, \code{_draw_gui}, \code{run}.

\VerbatimInput[firstline=62,lastline=63,firstnumber=62,numbers=left,numbersep=3pt,frame=lines,framesep=2pt,rulecolor=\color{gray!60},breaklines=true,fontsize=\fontsize{6.6}{7.6}\selectfont]{../../splat_editor/app.py}
\subsubsection{\texorpdfstring{\code{Application.__init__(self, width=1500, height=900)}}{Application.__init__}}
{\footnotesize\textit{Dòng 65--88}}\par
\textbf{Giải thích.} Tạo mọi thành phần theo thứ tự phụ thuộc: cửa sổ $\to$ ImGui $\to$ Editor, Viewport, Physics, Lighting $\to$ \code{ActionRegistry} $\to$ các bảng giao diện.

\VerbatimInput[firstline=65,lastline=88,firstnumber=65,numbers=left,numbersep=3pt,frame=lines,framesep=2pt,rulecolor=\color{gray!60},breaklines=true,fontsize=\fontsize{6.6}{7.6}\selectfont]{../../splat_editor/app.py}
\subsubsection{\texorpdfstring{\code{Application._create_actions(self)}}{Application._create_actions}}
{\footnotesize\textit{Dòng 91--172}}\par
\textbf{Giải thích.} Khai báo toàn bộ hành động (mở, lưu, hoàn tác, chọn, cọ, xem, ánh sáng, chế độ\ldots) cùng phím tắt ở một nơi.

\VerbatimInput[firstline=91,lastline=172,firstnumber=91,numbers=left,numbersep=3pt,frame=lines,framesep=2pt,rulecolor=\color{gray!60},breaklines=true,fontsize=\fontsize{6.6}{7.6}\selectfont]{../../splat_editor/app.py}
\subsubsection{\texorpdfstring{\code{Application.open_file(self, path)}}{Application.open_file}}
{\footnotesize\textit{Dòng 175--177}}\par
\VerbatimInput[firstline=175,lastline=177,firstnumber=175,numbers=left,numbersep=3pt,frame=lines,framesep=2pt,rulecolor=\color{gray!60},breaklines=true,fontsize=\fontsize{6.6}{7.6}\selectfont]{../../splat_editor/app.py}
\subsubsection{\texorpdfstring{\code{Application.add_files(self, paths)}}{Application.add_files}}
{\footnotesize\textit{Dòng 179--183}}\par
\VerbatimInput[firstline=179,lastline=183,firstnumber=179,numbers=left,numbersep=3pt,frame=lines,framesep=2pt,rulecolor=\color{gray!60},breaklines=true,fontsize=\fontsize{6.6}{7.6}\selectfont]{../../splat_editor/app.py}
\subsubsection{\texorpdfstring{\code{Application.open_dialog(self)}}{Application.open_dialog}}
{\footnotesize\textit{Dòng 185--189}}\par
\VerbatimInput[firstline=185,lastline=189,firstnumber=185,numbers=left,numbersep=3pt,frame=lines,framesep=2pt,rulecolor=\color{gray!60},breaklines=true,fontsize=\fontsize{6.6}{7.6}\selectfont]{../../splat_editor/app.py}
\subsubsection{\texorpdfstring{\code{Application.add_dialog(self)}}{Application.add_dialog}}
{\footnotesize\textit{Dòng 191--194}}\par
\VerbatimInput[firstline=191,lastline=194,firstnumber=191,numbers=left,numbersep=3pt,frame=lines,framesep=2pt,rulecolor=\color{gray!60},breaklines=true,fontsize=\fontsize{6.6}{7.6}\selectfont]{../../splat_editor/app.py}
\subsubsection{\texorpdfstring{\code{Application.save_dialog(self)}}{Application.save_dialog}}
{\footnotesize\textit{Dòng 196--200}}\par
\VerbatimInput[firstline=196,lastline=200,firstnumber=196,numbers=left,numbersep=3pt,frame=lines,framesep=2pt,rulecolor=\color{gray!60},breaklines=true,fontsize=\fontsize{6.6}{7.6}\selectfont]{../../splat_editor/app.py}
\subsubsection{\texorpdfstring{\code{Application.frame_scene(self, layer=None)}}{Application.frame_scene}}
{\footnotesize\textit{Dòng 202--208}}\par
\textbf{Mô tả trong mã nguồn.} Đưa camera về nhìn toàn cảnh (hoặc chỉ một lớp).

\VerbatimInput[firstline=202,lastline=208,firstnumber=202,numbers=left,numbersep=3pt,frame=lines,framesep=2pt,rulecolor=\color{gray!60},breaklines=true,fontsize=\fontsize{6.6}{7.6}\selectfont]{../../splat_editor/app.py}
\subsubsection{\texorpdfstring{\code{Application._update_bounds(self, positions=None, reset_scale=False)}}{Application._update_bounds}}
{\footnotesize\textit{Dòng 211--224}}\par
\textbf{Mô tả trong mã nguồn.} Kích thước cảnh -\textgreater{} tầm đèn, bản đồ bóng và chiều cao người chơi mặc định.

\VerbatimInput[firstline=211,lastline=224,firstnumber=211,numbers=left,numbersep=3pt,frame=lines,framesep=2pt,rulecolor=\color{gray!60},breaklines=true,fontsize=\fontsize{6.6}{7.6}\selectfont]{../../splat_editor/app.py}
\subsubsection{\texorpdfstring{\code{Application.toggle_lighting(self)}}{Application.toggle_lighting}}
{\footnotesize\textit{Dòng 226--227}}\par
\VerbatimInput[firstline=226,lastline=227,firstnumber=226,numbers=left,numbersep=3pt,frame=lines,framesep=2pt,rulecolor=\color{gray!60},breaklines=true,fontsize=\fontsize{6.6}{7.6}\selectfont]{../../splat_editor/app.py}
\subsubsection{\texorpdfstring{\code{Application.toggle_shadows(self)}}{Application.toggle_shadows}}
{\footnotesize\textit{Dòng 229--230}}\par
\VerbatimInput[firstline=229,lastline=230,firstnumber=229,numbers=left,numbersep=3pt,frame=lines,framesep=2pt,rulecolor=\color{gray!60},breaklines=true,fontsize=\fontsize{6.6}{7.6}\selectfont]{../../splat_editor/app.py}
\subsubsection{\texorpdfstring{\code{Application.add_sun(self)}}{Application.add_sun}}
{\footnotesize\textit{Dòng 232--236}}\par
\VerbatimInput[firstline=232,lastline=236,firstnumber=232,numbers=left,numbersep=3pt,frame=lines,framesep=2pt,rulecolor=\color{gray!60},breaklines=true,fontsize=\fontsize{6.6}{7.6}\selectfont]{../../splat_editor/app.py}
\subsubsection{\texorpdfstring{\code{Application.add_light(self)}}{Application.add_light}}
{\footnotesize\textit{Dòng 238--249}}\par
\VerbatimInput[firstline=238,lastline=249,firstnumber=238,numbers=left,numbersep=3pt,frame=lines,framesep=2pt,rulecolor=\color{gray!60},breaklines=true,fontsize=\fontsize{6.6}{7.6}\selectfont]{../../splat_editor/app.py}
\subsubsection{\texorpdfstring{\code{Application.pick_light(self, x, y)}}{Application.pick_light}}
{\footnotesize\textit{Dòng 252--280}}\par
\textbf{Mô tả trong mã nguồn.} Click trong cảnh 3D: chọn khối đèn gần con trỏ nhất (và không bị cảnh che); click ra ngoài: bỏ chọn.

\VerbatimInput[firstline=252,lastline=280,firstnumber=252,numbers=left,numbersep=3pt,frame=lines,framesep=2pt,rulecolor=\color{gray!60},breaklines=true,fontsize=\fontsize{6.6}{7.6}\selectfont]{../../splat_editor/app.py}
\subsubsection{\texorpdfstring{\code{Application.delete_key(self)}}{Application.delete_key}}
{\footnotesize\textit{Dòng 282--290}}\par
\textbf{Mô tả trong mã nguồn.} Delete: xoá khối đèn đang chọn trong cảnh; không có thì xoá gaussian như trước.

\VerbatimInput[firstline=282,lastline=290,firstnumber=282,numbers=left,numbersep=3pt,frame=lines,framesep=2pt,rulecolor=\color{gray!60},breaklines=true,fontsize=\fontsize{6.6}{7.6}\selectfont]{../../splat_editor/app.py}
\subsubsection{\texorpdfstring{\code{Application.set_mode(self, mode)}}{Application.set_mode}}
{\footnotesize\textit{Dòng 292--306}}\par
\VerbatimInput[firstline=292,lastline=306,firstnumber=292,numbers=left,numbersep=3pt,frame=lines,framesep=2pt,rulecolor=\color{gray!60},breaklines=true,fontsize=\fontsize{6.6}{7.6}\selectfont]{../../splat_editor/app.py}
\subsubsection{\texorpdfstring{\code{Application._spawn_player(self)}}{Application._spawn_player}}
{\footnotesize\textit{Dòng 308--328}}\par
\textbf{Mô tả trong mã nguồn.} Đặt người chơi ở nơi có mặt đất bên dưới: chỗ camera đang đứng, rồi tới điểm camera đang nhìn, cuối cùng là giữa cảnh (camera bay ở xa trong khoảng trống thì rơi ra ngoài thế giới).

\VerbatimInput[firstline=308,lastline=328,firstnumber=308,numbers=left,numbersep=3pt,frame=lines,framesep=2pt,rulecolor=\color{gray!60},breaklines=true,fontsize=\fontsize{6.6}{7.6}\selectfont]{../../splat_editor/app.py}
\subsubsection{\texorpdfstring{\code{Application.toggle_mode(self)}}{Application.toggle_mode}}
{\footnotesize\textit{Dòng 330--331}}\par
\VerbatimInput[firstline=330,lastline=331,firstnumber=330,numbers=left,numbersep=3pt,frame=lines,framesep=2pt,rulecolor=\color{gray!60},breaklines=true,fontsize=\fontsize{6.6}{7.6}\selectfont]{../../splat_editor/app.py}
\subsubsection{\texorpdfstring{\code{Application.walk(self, forward, right, jump, sprint, sneak, dt)}}{Application.walk}}
{\footnotesize\textit{Dòng 333--357}}\par
\textbf{Mô tả trong mã nguồn.} Một bước đi bộ của chế độ sinh tồn; camera đi theo mắt người chơi.

\textbf{Giải thích.} Một bước đi bộ của chế độ sinh tồn: lấy vector vào từ phím, gọi \code{Physics.walk}, rồi đặt camera theo mắt người chơi.

\VerbatimInput[firstline=333,lastline=357,firstnumber=333,numbers=left,numbersep=3pt,frame=lines,framesep=2pt,rulecolor=\color{gray!60},breaklines=true,fontsize=\fontsize{6.6}{7.6}\selectfont]{../../splat_editor/app.py}
\subsubsection{\texorpdfstring{\code{Application.cycle_view(self)}}{Application.cycle_view}}
{\footnotesize\textit{Dòng 360--374}}\par
\textbf{Mô tả trong mã nguồn.} F5: góc nhìn thứ nhất -\textgreater{} thứ ba (sau lưng) -\textgreater{} thứ ba (trước mặt) -\textgreater{} thứ nhất.

\VerbatimInput[firstline=360,lastline=374,firstnumber=360,numbers=left,numbersep=3pt,frame=lines,framesep=2pt,rulecolor=\color{gray!60},breaklines=true,fontsize=\fontsize{6.6}{7.6}\selectfont]{../../splat_editor/app.py}
\subsubsection{\texorpdfstring{\code{Application._leave_third_person(self)}}{Application._leave_third_person}}
{\footnotesize\textit{Dòng 376--380}}\par
\VerbatimInput[firstline=376,lastline=380,firstnumber=376,numbers=left,numbersep=3pt,frame=lines,framesep=2pt,rulecolor=\color{gray!60},breaklines=true,fontsize=\fontsize{6.6}{7.6}\selectfont]{../../splat_editor/app.py}
\subsubsection{\texorpdfstring{\code{Application._update_avatar(self, feet, walked, dt)}}{Application._update_avatar}}
{\footnotesize\textit{Dòng 382--418}}\par
\textbf{Mô tả trong mã nguồn.} Hướng thân, ma trận các khớp và khoảng lùi camera (rút ngắn khi sau lưng có tường).

\VerbatimInput[firstline=382,lastline=418,firstnumber=382,numbers=left,numbersep=3pt,frame=lines,framesep=2pt,rulecolor=\color{gray!60},breaklines=true,fontsize=\fontsize{6.6}{7.6}\selectfont]{../../splat_editor/app.py}
\subsubsection{\texorpdfstring{\code{Application.update_world(self, dt)}}{Application.update_world}}
{\footnotesize\textit{Dòng 420--436}}\par
\textbf{Mô tả trong mã nguồn.} Mỗi khung hình: bản đồ va chạm (luồng nền), thời gian trong ngày.

\VerbatimInput[firstline=420,lastline=436,firstnumber=420,numbers=left,numbersep=3pt,frame=lines,framesep=2pt,rulecolor=\color{gray!60},breaklines=true,fontsize=\fontsize{6.6}{7.6}\selectfont]{../../splat_editor/app.py}
\subsubsection{\texorpdfstring{\code{Application.toggle_grid(self)}}{Application.toggle_grid}}
{\footnotesize\textit{Dòng 438--439}}\par
\VerbatimInput[firstline=438,lastline=439,firstnumber=438,numbers=left,numbersep=3pt,frame=lines,framesep=2pt,rulecolor=\color{gray!60},breaklines=true,fontsize=\fontsize{6.6}{7.6}\selectfont]{../../splat_editor/app.py}
\subsubsection{\texorpdfstring{\code{Application.toggle_sh(self)}}{Application.toggle_sh}}
{\footnotesize\textit{Dòng 441--442}}\par
\VerbatimInput[firstline=441,lastline=442,firstnumber=441,numbers=left,numbersep=3pt,frame=lines,framesep=2pt,rulecolor=\color{gray!60},breaklines=true,fontsize=\fontsize{6.6}{7.6}\selectfont]{../../splat_editor/app.py}
\subsubsection{\texorpdfstring{\code{Application.set_flip(self, flip)}}{Application.set_flip}}
{\footnotesize\textit{Dòng 444--447}}\par
\VerbatimInput[firstline=444,lastline=447,firstnumber=444,numbers=left,numbersep=3pt,frame=lines,framesep=2pt,rulecolor=\color{gray!60},breaklines=true,fontsize=\fontsize{6.6}{7.6}\selectfont]{../../splat_editor/app.py}
\subsubsection{\texorpdfstring{\code{Application.exit_brush(self)}}{Application.exit_brush}}
{\footnotesize\textit{Dòng 449--451}}\par
\VerbatimInput[firstline=449,lastline=451,firstnumber=449,numbers=left,numbersep=3pt,frame=lines,framesep=2pt,rulecolor=\color{gray!60},breaklines=true,fontsize=\fontsize{6.6}{7.6}\selectfont]{../../splat_editor/app.py}
\subsubsection{\texorpdfstring{\code{Application.shortcuts_window_toggle(self)}}{Application.shortcuts_window_toggle}}
{\footnotesize\textit{Dòng 453--454}}\par
\VerbatimInput[firstline=453,lastline=454,firstnumber=453,numbers=left,numbersep=3pt,frame=lines,framesep=2pt,rulecolor=\color{gray!60},breaklines=true,fontsize=\fontsize{6.6}{7.6}\selectfont]{../../splat_editor/app.py}
\subsubsection{\texorpdfstring{\code{Application.paint(self, x, y)}}{Application.paint}}
{\footnotesize\textit{Dòng 456--459}}\par
\VerbatimInput[firstline=456,lastline=459,firstnumber=456,numbers=left,numbersep=3pt,frame=lines,framesep=2pt,rulecolor=\color{gray!60},breaklines=true,fontsize=\fontsize{6.6}{7.6}\selectfont]{../../splat_editor/app.py}
\subsubsection{\texorpdfstring{\code{Application.drag_speed(self)}}{Application.drag_speed}}
{\footnotesize\textit{Dòng 461--463}}\par
\textbf{Mô tả trong mã nguồn.} Tốc độ kéo thanh trượt vị trí, tỉ lệ theo khoảng cách camera.

\textbf{Giải thích.} Công thức cốt lõi của hàm này:
\[
\mathrm{speed}=\max(0.005\,d,\ 10^{-4})
\]


\VerbatimInput[firstline=461,lastline=463,firstnumber=461,numbers=left,numbersep=3pt,frame=lines,framesep=2pt,rulecolor=\color{gray!60},breaklines=true,fontsize=\fontsize{6.6}{7.6}\selectfont]{../../splat_editor/app.py}
\subsubsection{\texorpdfstring{\code{Application._draw_gui(self)}}{Application._draw_gui}}
{\footnotesize\textit{Dòng 466--473}}\par
\VerbatimInput[firstline=466,lastline=473,firstnumber=466,numbers=left,numbersep=3pt,frame=lines,framesep=2pt,rulecolor=\color{gray!60},breaklines=true,fontsize=\fontsize{6.6}{7.6}\selectfont]{../../splat_editor/app.py}
\subsubsection{\texorpdfstring{\code{Application.run(self)}}{Application.run}}
{\footnotesize\textit{Dòng 475--497}}\par
\textbf{Giải thích.} \textbf{Vòng lặp chính.} Đọc ngược từ trên xuống dưới chính là thứ tự một khung hình: sự kiện $\to$ cập nhật camera/thế giới $\to$ dựng ImGui $\to$ đồng bộ GPU $\to$ vẽ cảnh $\to$ vẽ ImGui $\to$ hiển thị.

\VerbatimInput[firstline=475,lastline=497,firstnumber=475,numbers=left,numbersep=3pt,frame=lines,framesep=2pt,rulecolor=\color{gray!60},breaklines=true,fontsize=\fontsize{6.6}{7.6}\selectfont]{../../splat_editor/app.py}
\subsubsection{\texorpdfstring{\code{main(argv=None)}}{main}}
{\footnotesize\textit{Dòng 500--506}}\par
\VerbatimInput[firstline=500,lastline=506,firstnumber=500,numbers=left,numbersep=3pt,frame=lines,framesep=2pt,rulecolor=\color{gray!60},breaklines=true,fontsize=\fontsize{6.6}{7.6}\selectfont]{../../splat_editor/app.py}

\clearpage
\section{Thư mục core: logic thuần}\label{ch:core}
Thư mục \code{core} chứa \textbf{toàn bộ logic thuần}: toán học, dữ liệu gaussian, lớp (layer), cảnh (scene),
ánh sáng, vật lý, LOD, hoạt họa, lịch sử hoàn tác. Quy tắc thiết kế quan trọng nhất của thư mục này là
\textbf{không import OpenGL, GLFW hay ImGui}; mọi thứ chạy bằng NumPy nên có thể kiểm thử mà không cần cửa sổ. Các thư mục
khác phụ thuộc vào \code{core}, còn \code{core} chỉ phụ thuộc \code{io} (để đọc/ghi PLY) và không phụ thuộc ngược lại ai.

Dữ liệu chảy trong \code{core} theo chuỗi: \code{Splats} (mảng gaussian đã giải mã) $\to$ \code{Layer} (một file PLY cộng
phép biến đổi và vùng chọn) $\to$ \code{Scene} (nhiều layer). Các module còn lại (\code{lod}, \code{lighting},
\code{physics}, \code{avatar_pose}) đọc dữ liệu đó để phục vụ hiển thị và tương tác.

\subsection{\texorpdfstring{\code{splat_editor/core/math3d.py}}{splat_editor/core/math3d.py}}
{\footnotesize 135 dòng}\par
\textbf{Vai trò.} Thư viện toán 3D nhỏ, chỉ dùng \code{numpy}. Gồm ba nhóm: ma trận camera (\code{perspective}, \code{orthographic}, \code{look_at}),
phép quay (\code{euler_to_matrix}, \code{matrix_to_euler}, \code{axis_angle_matrix}) và quaternion (\code{matrix_to_quat}, \code{quat_mul}, \code{quat_to_matrix}).
Toàn bộ chương trình dùng chung các hàm này nên quy ước (vector cột, tay phải, $Y$ lên, nhìn dọc $-z$) được thống nhất ở một nơi.
Công thức chi tiết xem \S\ref{s:quat}, \S\ref{s:cam}, \S\ref{s:layer}.

\paragraph{Mô tả đầu file và cách dùng.} Phần chú thích đầu file do tác giả viết, nêu ý tưởng và ví dụ sử dụng:
\VerbatimInput[firstline=1,lastline=15,firstnumber=1,numbers=left,numbersep=3pt,frame=lines,framesep=2pt,rulecolor=\color{gray!60},breaklines=true,fontsize=\fontsize{6.2}{7.2}\selectfont]{../../splat_editor/core/math3d.py}
%% ---- merged math: math3d_quat_cam
\subsubsection{Cơ sở toán: Quaternion}\label{s:quat}
Quaternion đơn vị $q=(w,x,y,z)$ biểu diễn phép quay góc $\theta$ quanh trục $n$: $w=\cos\frac\theta2$, $(x,y,z)=\sin\frac\theta2\,n$. Dùng quaternion vì không bị gimbal, nhân gọn và dễ chuẩn hóa.

\paragraph{Quaternion sang ma trận (\code{quat_to_matrix}, giống hệt trong \code{splat_prep.vert}).}
\[
R=\begin{bmatrix}
1-2(y^2+z^2) & 2(xy-wz) & 2(xz+wy)\\
2(xy+wz) & 1-2(x^2+z^2) & 2(yz-wx)\\
2(xz-wy) & 2(yz+wx) & 1-2(x^2+y^2)
\end{bmatrix}
\]
Trong GLSL, \code{mat3(...)} điền theo \textbf{cột}, nên các số liệt kê trong shader là cột của cùng ma trận này.

\paragraph{Ma trận sang quaternion (\code{matrix_to_quat}, phương pháp Shepperd).}
Tránh chia cho số gần 0 bằng cách chọn nhánh theo phần tử lớn nhất.
\begin{itemize}
\item $\mathrm{trace}=R_{00}+R_{11}+R_{22}>0$: $s=2\sqrt{\mathrm{trace}+1}$, $w=s/4$, $x=(R_{21}-R_{12})/s$, $y=(R_{02}-R_{20})/s$, $z=(R_{10}-R_{01})/s$.
\item Ngược lại, chọn nhánh theo $R_{00}$, $R_{11}$ hoặc $R_{22}$ lớn nhất: $s=2\sqrt{1+R_{ii}-R_{jj}-R_{kk}}$, phần tử $i=s/4$, các phần tử còn lại $=(R_{ij}\pm R_{ji})/s$.
\end{itemize}

\paragraph{Nhân quaternion (\code{quat_mul}, tích Hamilton).}
\begin{align*}
w&=a_wb_w-a_xb_x-a_yb_y-a_zb_z, & x&=a_wb_x+a_xb_w+a_yb_z-a_zb_y,\\
y&=a_wb_y-a_xb_z+a_yb_w+a_zb_x, & z&=a_wb_z+a_xb_y-a_yb_x+a_zb_w.
\end{align*}
$a\otimes b$ là quay $b$ trước rồi $a$ (giống $R_aR_b$).

\paragraph{Lượng tử hóa quaternion lên GPU (\code{splat_renderer.pack}, \code{splat_prep.vert}).}
Mỗi thành phần thuộc $[-1,1]$ được lưu 8 bit: $u=\mathrm{clip}(\mathrm{round}(q\cdot127)+128,1,255)$. Shader giải bằng $q=\mathrm{normalize}((u-128)/127)$, chuẩn hóa lại để bù sai số.

\subsubsection{Cơ sở toán: Camera và phép chiếu}\label{s:cam}
Các ma trận camera dưới đây nằm trong \code{core/math3d.py}; phần camera quỹ đạo \code{OrbitCamera} (\code{render/camera.py}) được trình bày ở mục \ref{s:cam-orbit}.

\paragraph{Ma trận nhìn (\code{look_at(e, target, up)}).}
\[
f=\frac{\mathrm{target}-e}{\|\cdot\|},\quad r=\frac{f\times \mathrm{up}}{\|\cdot\|},\quad u=r\times f,\quad
V=\begin{bmatrix} r^{\top} & -r\!\cdot\! e\\ u^{\top} & -u\!\cdot\! e\\ -f^{\top} & f\!\cdot\! e\\ 0\;0\;0 & 1\end{bmatrix}.
\]
Cột 3 là $-(Re)$: đưa mắt về gốc rồi xoay hệ trục camera về $(x,y,z)=(\text{phải},\text{lên},-\text{trước})$. Khối $3\times3$ trái trên gọi là $W$ (trực giao).

\paragraph{Phép chiếu phối cảnh (\code{perspective(fov, a, n, f)}).}
Với $k=1/\tan(\mathrm{fov}/2)$:
\[
P=\begin{bmatrix}k/a&0&0&0\\0&k&0&0\\0&0&\frac{f+n}{n-f}&\frac{2fn}{n-f}\\0&0&-1&0\end{bmatrix}
\]
\begin{itemize}
\item $\mathrm{clip}=Pv$, $c_w=-v_z$ (độ sâu dương). $\mathrm{NDC}=\mathrm{clip}_{xyz}/c_w\in[-1,1]$ (\textbf{chia phối cảnh}, vật xa nhỏ đi).
\item Hàng 3: $z_{ndc}=\dfrac{(f+n)v_z+2fn}{(n-f)(-v_z)}$ ánh xạ $v_z=-n\mapsto-1$ và $v_z=-f\mapsto+1$. Quan hệ này phi tuyến, độ phân giải tập trung gần camera.
\item Camera dùng \code{near = distance*0.001}, \code{far = near*1e6} (\code{camera.projection_matrix}).
\end{itemize}

\paragraph{Phép chiếu trực giao (\code{orthographic(h, n, f)}).}
$x'=x/h$, $y'=y/h$, $z'=-\dfrac{2z}{f-n}-\dfrac{f+n}{f-n}$. Hộp $[-h,h]^2\times[n,f]$ ánh xạ vào $[-1,1]^3$. Dùng cho \textbf{bản đồ bóng} mặt trời (tia song song, không chia phối cảnh). Hàm này bị định nghĩa trùng hai lần trong file, bản sau ghi đè bản trước, cùng công thức.

\paragraph{Ví dụ số (\code{report/project3d_table.tex}).}
$p_w=(0.5,1,0.5)\to v=Vp=(0.109,0.293,-2.536)\to \mathrm{clip}=Pv=(0.197,0.709,2.341,2.536)$, $c_w=2.536\to\mathrm{NDC}=(0.078,0.279,0.923)\to$ pixel $(215.5,108.1)$ theo mục \ref{s:viewport}.

\subsubsection*{Khai báo mức module (dòng 16--18)}
\VerbatimInput[firstline=16,lastline=18,firstnumber=16,numbers=left,numbersep=3pt,frame=lines,framesep=2pt,rulecolor=\color{gray!60},breaklines=true,fontsize=\fontsize{6.6}{7.6}\selectfont]{../../splat_editor/core/math3d.py}
\subsubsection{\texorpdfstring{\code{perspective(fov_y_deg, aspect, near, far)}}{perspective}}
{\footnotesize\textit{Dòng 23--31}}\par
\textbf{Giải thích.} Dựng ma trận chiếu phối cảnh OpenGL. \code{f=1/tan(fov/2)} là tỉ lệ thu phóng; hàng cuối $(0,0,-1,0)$ làm cho $c_w=-v_z$ nên phép chia phối cảnh
sau đó cho vật xa nhỏ đi. Hai phần tử \code{m[2,2]}, \code{m[2,3]} ánh xạ $[-near,-far]$ sang $[-1,1]$ (xem \S\ref{s:cam}).

\textbf{Công thức.}
\[
k=\frac{1}{\tan(\mathrm{fov}/2)},\qquad
P=\begin{bmatrix}k/a&0&0&0\\0&k&0&0\\0&0&\dfrac{f+n}{n-f}&\dfrac{2fn}{n-f}\\0&0&-1&0\end{bmatrix}
\]
Với điểm trong không gian camera $v$: $\mathrm{clip}=Pv$, $c_w=-v_z$ và
\[
\mathrm{NDC}=\frac{\mathrm{clip}_{xyz}}{c_w},\qquad
z_{ndc}=\frac{(f+n)\,v_z+2fn}{(n-f)\,(-v_z)}\ \Rightarrow\ v_z=-n\mapsto-1,\ \ v_z=-f\mapsto+1 .
\]


\VerbatimInput[firstline=23,lastline=31,firstnumber=23,numbers=left,numbersep=3pt,frame=lines,framesep=2pt,rulecolor=\color{gray!60},breaklines=true,fontsize=\fontsize{6.6}{7.6}\selectfont]{../../splat_editor/core/math3d.py}
\subsubsection{\texorpdfstring{\code{orthographic(half_size, near, far)}}{orthographic}}
{\footnotesize\textit{Dòng 34--40}}\par
\textbf{Mô tả trong mã nguồn.} Phép chiếu trực giao đối xứng (dùng cho bản đồ bóng của mặt trời): hộp [-h,h]\textasciicircum{}2 x [near,far].

\textbf{Giải thích.} Phép chiếu trực giao dùng cho bản đồ bóng của mặt trời. \textbf{Lưu ý:} hàm này bị định nghĩa trùng hai lần trong file; bản định nghĩa sau ghi đè bản trước và hai bản giống hệt công thức,
nên không gây lỗi nhưng có thể xóa một bản cho gọn.

\textbf{Công thức.}
\[
x'=\frac{x}{h},\quad y'=\frac{y}{h},\quad z'=-\frac{2z}{f-n}-\frac{f+n}{f-n}
\qquad\big([-h,h]^2\times[n,f]\to[-1,1]^3\big)
\]


\VerbatimInput[firstline=34,lastline=40,firstnumber=34,numbers=left,numbersep=3pt,frame=lines,framesep=2pt,rulecolor=\color{gray!60},breaklines=true,fontsize=\fontsize{6.6}{7.6}\selectfont]{../../splat_editor/core/math3d.py}
\subsubsection{\texorpdfstring{\code{orthographic(half_size, near, far)}}{orthographic}}
{\footnotesize\textit{Dòng 43--49}}\par
\textbf{Mô tả trong mã nguồn.} Phép chiếu trực giao đối xứng (bản đồ bóng của mặt trời): hộp [-h,h]\textasciicircum{}2 x [near,far].

\textbf{Giải thích.} Phép chiếu trực giao dùng cho bản đồ bóng của mặt trời. \textbf{Lưu ý:} hàm này bị định nghĩa trùng hai lần trong file; bản định nghĩa sau ghi đè bản trước và hai bản giống hệt công thức,
nên không gây lỗi nhưng có thể xóa một bản cho gọn.

\textbf{Công thức.}
\[
x'=\frac{x}{h},\quad y'=\frac{y}{h},\quad z'=-\frac{2z}{f-n}-\frac{f+n}{f-n}
\qquad\big([-h,h]^2\times[n,f]\to[-1,1]^3\big)
\]


\VerbatimInput[firstline=43,lastline=49,firstnumber=43,numbers=left,numbersep=3pt,frame=lines,framesep=2pt,rulecolor=\color{gray!60},breaklines=true,fontsize=\fontsize{6.6}{7.6}\selectfont]{../../splat_editor/core/math3d.py}
\subsubsection{\texorpdfstring{\code{look_at(eye, target, up)}}{look_at}}
{\footnotesize\textit{Dòng 52--63}}\par
\textbf{Giải thích.} Dựng ma trận nhìn (view). Ba vector $r,u,-f$ là hàng của khối quay; cột tịnh tiến $-R\,e$ đưa mắt về gốc. Dòng \code{m[:3, 3] = -m[:3, :3] @ eye} chính là bước đó.

\textbf{Công thức.}
\[
f=\frac{t-e}{\|t-e\|},\quad r=\frac{f\times u_{up}}{\|f\times u_{up}\|},\quad u=r\times f,
\qquad
V=\begin{bmatrix}r^{\top}&-r\!\cdot\!e\\u^{\top}&-u\!\cdot\!e\\-f^{\top}&f\!\cdot\!e\\0\;0\;0&1\end{bmatrix}
=\begin{bmatrix}W&-We\\0&1\end{bmatrix}
\]


\VerbatimInput[firstline=52,lastline=63,firstnumber=52,numbers=left,numbersep=3pt,frame=lines,framesep=2pt,rulecolor=\color{gray!60},breaklines=true,fontsize=\fontsize{6.6}{7.6}\selectfont]{../../splat_editor/core/math3d.py}
\subsubsection{\texorpdfstring{\code{euler_to_matrix(degrees)}}{euler_to_matrix}}
{\footnotesize\textit{Dòng 68--75}}\par
\textbf{Mô tả trong mã nguồn.} Góc Euler (độ) quanh X, Y, Z -\textgreater{} ma trận xoay R = Rz * Ry * Rx.

\textbf{Giải thích.} Ghép ba ma trận quay đơn theo thứ tự $R=R_zR_yR_x$ (áp dụng $X$ trước). Nhận đầu vào là \textbf{độ}, đổi sang radian bằng \code{np.radians}.

\textbf{Công thức.}
\[
R_x=\begin{bmatrix}1&0&0\\0&c_x&-s_x\\0&s_x&c_x\end{bmatrix},\;
R_y=\begin{bmatrix}c_y&0&s_y\\0&1&0\\-s_y&0&c_y\end{bmatrix},\;
R_z=\begin{bmatrix}c_z&-s_z&0\\s_z&c_z&0\\0&0&1\end{bmatrix},\qquad R=R_zR_yR_x
\]


\VerbatimInput[firstline=68,lastline=75,firstnumber=68,numbers=left,numbersep=3pt,frame=lines,framesep=2pt,rulecolor=\color{gray!60},breaklines=true,fontsize=\fontsize{6.6}{7.6}\selectfont]{../../splat_editor/core/math3d.py}
\subsubsection{\texorpdfstring{\code{matrix_to_euler(m)}}{matrix_to_euler}}
{\footnotesize\textit{Dòng 78--87}}\par
\textbf{Mô tả trong mã nguồn.} Ngược lại của euler\_to\_matrix: ma trận R = Rz * Ry * Rx -\textgreater{} góc (độ) quanh X, Y, Z.

\textbf{Giải thích.} Phép ngược của \code{euler_to_matrix}. Nhánh \code{else} xử lý khóa gimbal ($y=\pm90^\circ$), khi đó chỉ còn xác định được tổng/hiệu của $x$ và $z$ nên gộp hết vào $z$.

\textbf{Công thức.}
\[
y=\arcsin(-R_{20}),\quad x=\operatorname{atan2}(R_{21},R_{22}),\quad z=\operatorname{atan2}(R_{10},R_{00})
\]
Khóa gimbal $|R_{20}|\ge0.999999$: $x=0$, $z=\operatorname{atan2}(-R_{01},R_{11})$.


\VerbatimInput[firstline=78,lastline=87,firstnumber=78,numbers=left,numbersep=3pt,frame=lines,framesep=2pt,rulecolor=\color{gray!60},breaklines=true,fontsize=\fontsize{6.6}{7.6}\selectfont]{../../splat_editor/core/math3d.py}
\subsubsection{\texorpdfstring{\code{axis_angle_matrix(axis, degrees)}}{axis_angle_matrix}}
{\footnotesize\textit{Dòng 90--97}}\par
\textbf{Mô tả trong mã nguồn.} Ma trận xoay `degrees` độ quanh trục `axis` (công thức Rodrigues).

\textbf{Giải thích.} Công thức Rodrigues; trục được chuẩn hóa trước. Dùng cho thao tác xoay theo trục bất kỳ (ví dụ \code{rotate_about}).

\textbf{Công thức.}
\[
R=I+\sin a\,[k]_\times+(1-\cos a)\,[k]_\times^2,\qquad
[k]_\times=\begin{bmatrix}0&-z&y\\z&0&-x\\-y&x&0\end{bmatrix}
\]


\VerbatimInput[firstline=90,lastline=97,firstnumber=90,numbers=left,numbersep=3pt,frame=lines,framesep=2pt,rulecolor=\color{gray!60},breaklines=true,fontsize=\fontsize{6.6}{7.6}\selectfont]{../../splat_editor/core/math3d.py}
\subsubsection{\texorpdfstring{\code{matrix_to_quat(m)}}{matrix_to_quat}}
{\footnotesize\textit{Dòng 100--115}}\par
\textbf{Mô tả trong mã nguồn.} Ma trận xoay 3x3 -\textgreater{} quaternion (w, x, y, z).

\textbf{Giải thích.} Phương pháp Shepperd: chọn nhánh theo phần tử lớn nhất trong (trace, $R_{00}$, $R_{11}$, $R_{22}$) để mẫu số $s$ không bao giờ gần 0.

\textbf{Công thức.}
\[
\mathrm{tr}=R_{00}+R_{11}+R_{22}>0:\quad s=2\sqrt{\mathrm{tr}+1},\;
q=\Big(\tfrac s4,\ \tfrac{R_{21}-R_{12}}{s},\ \tfrac{R_{02}-R_{20}}{s},\ \tfrac{R_{10}-R_{01}}{s}\Big)
\]
Nhánh còn lại (phần tử đường chéo lớn nhất $R_{ii}$): $s=2\sqrt{1+R_{ii}-R_{jj}-R_{kk}}$, $q_i=s/4$.


\VerbatimInput[firstline=100,lastline=115,firstnumber=100,numbers=left,numbersep=3pt,frame=lines,framesep=2pt,rulecolor=\color{gray!60},breaklines=true,fontsize=\fontsize{6.6}{7.6}\selectfont]{../../splat_editor/core/math3d.py}
\subsubsection{\texorpdfstring{\code{quat_mul(a, b)}}{quat_mul}}
{\footnotesize\textit{Dòng 118--125}}\par
\textbf{Mô tả trong mã nguồn.} Nhân quaternion a (4,) với từng quaternion trong b (N,4).

\textbf{Giải thích.} Tích Hamilton vector hóa: \code{a} là một quaternion $(4,)$ nhân với hàng loạt quaternion \code{b} $(N,4)$. Dùng khi xoay cả layer: \code{q_rot} nhân với mọi quaternion gaussian.

\textbf{Công thức.}
\[
a\otimes b=\begin{bmatrix}a_wb_w-a_xb_x-a_yb_y-a_zb_z\\a_wb_x+a_xb_w+a_yb_z-a_zb_y\\a_wb_y-a_xb_z+a_yb_w+a_zb_x\\a_wb_z+a_xb_y-a_yb_x+a_zb_w\end{bmatrix},
\qquad R(a\otimes b)=R(a)\,R(b)
\]


\VerbatimInput[firstline=118,lastline=125,firstnumber=118,numbers=left,numbersep=3pt,frame=lines,framesep=2pt,rulecolor=\color{gray!60},breaklines=true,fontsize=\fontsize{6.6}{7.6}\selectfont]{../../splat_editor/core/math3d.py}
\subsubsection{\texorpdfstring{\code{quat_to_matrix(q)}}{quat_to_matrix}}
{\footnotesize\textit{Dòng 128--135}}\par
\textbf{Mô tả trong mã nguồn.} (N,4) quaternion (w,x,y,z) -\textgreater{} (N,3,3) ma trận xoay.

\textbf{Giải thích.} Đổi hàng loạt quaternion $(N,4)$ thành $(N,3,3)$. Cùng công thức với đoạn GLSL trong \code{splat_prep.vert} (chỉ khác là shader điền theo cột).

\textbf{Công thức.}
\[
R=\begin{bmatrix}1-2(y^2+z^2)&2(xy-wz)&2(xz+wy)\\2(xy+wz)&1-2(x^2+z^2)&2(yz-wx)\\2(xz-wy)&2(yz+wx)&1-2(x^2+y^2)\end{bmatrix},
\qquad q=\Big(\cos\tfrac\theta2,\ \sin\tfrac\theta2\,\hat n\Big)
\]


\VerbatimInput[firstline=128,lastline=135,firstnumber=128,numbers=left,numbersep=3pt,frame=lines,framesep=2pt,rulecolor=\color{gray!60},breaklines=true,fontsize=\fontsize{6.6}{7.6}\selectfont]{../../splat_editor/core/math3d.py}
\subsection{\texorpdfstring{\code{splat_editor/core/transform.py}}{splat_editor/core/transform.py}}
{\footnotesize 43 dòng}\par
\textbf{Vai trò.} Phép biến đổi của một layer: vị trí, góc Euler (độ) và tỉ lệ đều. Đây là dữ liệu thuần (\code{dataclass}); \textbf{dữ liệu gaussian không bị sửa}
khi người dùng kéo/xoay/phóng layer, chỉ phép biến đổi này đổi, nhờ đó thao tác tức thời với hàng chục triệu gaussian.

\paragraph{Mô tả đầu file và cách dùng.} Phần chú thích đầu file do tác giả viết, nêu ý tưởng và ví dụ sử dụng:
\VerbatimInput[firstline=1,lastline=14,firstnumber=1,numbers=left,numbersep=3pt,frame=lines,framesep=2pt,rulecolor=\color{gray!60},breaklines=true,fontsize=\fontsize{6.2}{7.2}\selectfont]{../../splat_editor/core/transform.py}
\subsubsection*{Khai báo mức module (dòng 15--19)}
\VerbatimInput[firstline=15,lastline=19,firstnumber=15,numbers=left,numbersep=3pt,frame=lines,framesep=2pt,rulecolor=\color{gray!60},breaklines=true,fontsize=\fontsize{6.6}{7.6}\selectfont]{../../splat_editor/core/transform.py}
\subsubsection{\texorpdfstring{\code{class Transform}}{Transform}}
{\footnotesize\textit{Dòng 22--43}}\par
\textbf{Giải thích.} Ba trường: \code{position} (3 số), \code{rotation} (độ quanh X, Y, Z) và \code{scale} (một số, tỉ lệ đều). Dùng \code{default_factory} để mỗi \code{Transform} có danh sách riêng.

Các phương thức: \code{reset}, \code{rotation_matrix}, \code{translation}, \code{rotate_about}.

\VerbatimInput[firstline=22,lastline=26,firstnumber=22,numbers=left,numbersep=3pt,frame=lines,framesep=2pt,rulecolor=\color{gray!60},breaklines=true,fontsize=\fontsize{6.6}{7.6}\selectfont]{../../splat_editor/core/transform.py}
\subsubsection{\texorpdfstring{\code{Transform.reset(self)}}{Transform.reset}}
{\footnotesize\textit{Dòng 28--31}}\par
\VerbatimInput[firstline=28,lastline=31,firstnumber=28,numbers=left,numbersep=3pt,frame=lines,framesep=2pt,rulecolor=\color{gray!60},breaklines=true,fontsize=\fontsize{6.6}{7.6}\selectfont]{../../splat_editor/core/transform.py}
\subsubsection{\texorpdfstring{\code{Transform.rotation_matrix(self)}}{Transform.rotation_matrix}}
{\footnotesize\textit{Dòng 33--34}}\par
\VerbatimInput[firstline=33,lastline=34,firstnumber=33,numbers=left,numbersep=3pt,frame=lines,framesep=2pt,rulecolor=\color{gray!60},breaklines=true,fontsize=\fontsize{6.6}{7.6}\selectfont]{../../splat_editor/core/transform.py}
\subsubsection{\texorpdfstring{\code{Transform.translation(self)}}{Transform.translation}}
{\footnotesize\textit{Dòng 36--37}}\par
\VerbatimInput[firstline=36,lastline=37,firstnumber=36,numbers=left,numbersep=3pt,frame=lines,framesep=2pt,rulecolor=\color{gray!60},breaklines=true,fontsize=\fontsize{6.6}{7.6}\selectfont]{../../splat_editor/core/transform.py}
\subsubsection{\texorpdfstring{\code{Transform.rotate_about(self, delta, pivot)}}{Transform.rotate_about}}
{\footnotesize\textit{Dòng 39--43}}\par
\textbf{Mô tả trong mã nguồn.} Xoay thêm `delta` (ma trận 3x3, trong toạ độ thế giới) quanh điểm `pivot`.

\textbf{Giải thích.} Xoay thêm một ma trận \code{delta} quanh điểm \code{pivot} trong tọa độ thế giới: $R'=\Delta R$, $t'=\Delta(t-c)+c$. Góc Euler mới lấy lại bằng \code{matrix_to_euler}.

\textbf{Công thức.}
\[
R'=\Delta R,\qquad t'=\Delta\,(t-c)+c
\]


\VerbatimInput[firstline=39,lastline=43,firstnumber=39,numbers=left,numbersep=3pt,frame=lines,framesep=2pt,rulecolor=\color{gray!60},breaklines=true,fontsize=\fontsize{6.6}{7.6}\selectfont]{../../splat_editor/core/transform.py}
\subsection{\texorpdfstring{\code{splat_editor/core/splats.py}}{splat_editor/core/splats.py}}
{\footnotesize 77 dòng}\par
\textbf{Vai trò.} \code{Splats} là mảng cấu trúc-của-mảng (structure of arrays) giữ toàn bộ gaussian \textbf{đã giải mã} để vẽ: tâm, màu bậc 0, độ đục, scale,
quaternion và hệ số SH bậc cao. Cách tổ chức theo cột giúp NumPy xử lý vector hóa và giúp gửi lên GPU theo từng đoạn. Điểm đáng chú ý là trường
\code{sh_frames}: xoay SH được \textbf{hoãn} (lazy) để không phải giữ thêm một bản sao SH rất lớn.

\paragraph{Mô tả đầu file và cách dùng.} Phần chú thích đầu file do tác giả viết, nêu ý tưởng và ví dụ sử dụng:
\VerbatimInput[firstline=1,lastline=13,firstnumber=1,numbers=left,numbersep=3pt,frame=lines,framesep=2pt,rulecolor=\color{gray!60},breaklines=true,fontsize=\fontsize{6.2}{7.2}\selectfont]{../../splat_editor/core/splats.py}
\subsubsection*{Khai báo mức module (dòng 14--18)}
\VerbatimInput[firstline=14,lastline=18,firstnumber=14,numbers=left,numbersep=3pt,frame=lines,framesep=2pt,rulecolor=\color{gray!60},breaklines=true,fontsize=\fontsize{6.6}{7.6}\selectfont]{../../splat_editor/core/splats.py}
\subsubsection{\texorpdfstring{\code{class Splats}}{Splats}}
{\footnotesize\textit{Dòng 21--76}}\par
\textbf{Giải thích.} Sáu mảng song song: \code{positions} $(N,3)$, \code{colors} $(N,3)$, \code{alphas} $(N,)$, \code{scales} $(N,3)$, \code{rotations} $(N,4)$, \code{sh_rest} $(N,K,3)$.
\code{ARRAYS} liệt kê tên các mảng để \code{subset}/\code{concat} xử lý đồng loạt.

Các phương thức: \code{count}, \code{sh_count}, \code{world_sh}, \code{subset}, \code{concat}.

\VerbatimInput[firstline=21,lastline=35,firstnumber=21,numbers=left,numbersep=3pt,frame=lines,framesep=2pt,rulecolor=\color{gray!60},breaklines=true,fontsize=\fontsize{6.6}{7.6}\selectfont]{../../splat_editor/core/splats.py}
\subsubsection{\texorpdfstring{\code{Splats.count(self)}}{Splats.count}}
{\footnotesize\textit{Dòng 37--39}}\par
\VerbatimInput[firstline=37,lastline=39,firstnumber=37,numbers=left,numbersep=3pt,frame=lines,framesep=2pt,rulecolor=\color{gray!60},breaklines=true,fontsize=\fontsize{6.6}{7.6}\selectfont]{../../splat_editor/core/splats.py}
\subsubsection{\texorpdfstring{\code{Splats.sh_count(self)}}{Splats.sh_count}}
{\footnotesize\textit{Dòng 41--43}}\par
\VerbatimInput[firstline=41,lastline=43,firstnumber=41,numbers=left,numbersep=3pt,frame=lines,framesep=2pt,rulecolor=\color{gray!60},breaklines=true,fontsize=\fontsize{6.6}{7.6}\selectfont]{../../splat_editor/core/splats.py}
\subsubsection{\texorpdfstring{\code{Splats.world_sh(self, index=None, k=None)}}{Splats.world_sh}}
{\footnotesize\textit{Dòng 45--54}}\par
\textbf{Mô tả trong mã nguồn.} Hệ số SH (float32, toạ độ thế giới) của các gaussian `index` (mảng chỉ số; None: tất cả). k: chỉ lấy k hệ số đầu (mỗi bậc SH xoay độc lập nên cắt bớt bậc cao vẫn đúng).

\textbf{Giải thích.} Trả về hệ số SH đã xoay về tọa độ thế giới. Với mỗi khung \code{(start, end, rot)} trong \code{sh_frames}, chỉ những hàng nằm trong đoạn đó mới được xoay bằng \code{sh.rotate}.

\VerbatimInput[firstline=45,lastline=54,firstnumber=45,numbers=left,numbersep=3pt,frame=lines,framesep=2pt,rulecolor=\color{gray!60},breaklines=true,fontsize=\fontsize{6.6}{7.6}\selectfont]{../../splat_editor/core/splats.py}
\subsubsection{\texorpdfstring{\code{Splats.subset(self, mask)}}{Splats.subset}}
{\footnotesize\textit{Dòng 56--62}}\par
\textbf{Mô tả trong mã nguồn.} Splats mới chỉ gồm các gaussian có mask True (chỉ dùng cho dữ liệu chưa xoay).

\textbf{Giải thích.} Lọc theo mặt nạ. \code{np.compress} nhanh hơn chỉ mục boolean khoảng 2 lần với mảng lớn (SH có thể chiếm cỡ 900 MB). Có \code{assert} không cho dùng khi đã có \code{sh_frames}.

\VerbatimInput[firstline=56,lastline=62,firstnumber=56,numbers=left,numbersep=3pt,frame=lines,framesep=2pt,rulecolor=\color{gray!60},breaklines=true,fontsize=\fontsize{6.6}{7.6}\selectfont]{../../splat_editor/core/splats.py}
\subsubsection{\texorpdfstring{\code{Splats.concat(cls, parts)}}{Splats.concat}}
{\footnotesize\textit{Dòng 64--76}}\par
\textbf{Mô tả trong mã nguồn.} Gộp nhiều Splats thành một (giữ thứ tự). Lớp có ít hệ số SH hơn được thêm 0.

\textbf{Giải thích.} Gộp nhiều \code{Splats}. Layer có ít hệ số SH hơn được đệm 0 (\code{np.pad}); các khung xoay SH được dịch chỉ số bằng \code{offset}.

\VerbatimInput[firstline=64,lastline=76,firstnumber=64,numbers=left,numbersep=3pt,frame=lines,framesep=2pt,rulecolor=\color{gray!60},breaklines=true,fontsize=\fontsize{6.6}{7.6}\selectfont]{../../splat_editor/core/splats.py}
\subsection{\texorpdfstring{\code{splat_editor/core/layer.py}}{splat_editor/core/layer.py}}
{\footnotesize 187 dòng}\par
\textbf{Vai trò.} Một \code{Layer} tương ứng một file PLY đã mở, cộng phép biến đổi, vùng chọn và cờ xóa mềm. Hai ý tưởng cốt lõi:
(1) \textbf{hai hệ tọa độ} (tọa độ file và tọa độ thế giới, nối nhau bằng ma trận $M$); (2) \textbf{xóa mềm} (chỉ bật cờ \code{deleted}) nên xóa và hoàn tác gần như tức thì.
Khi lưu, dữ liệu được đổi ngược về tọa độ file để file xuất ra vẫn đúng quy ước của file gốc. Toán học: xem mục \textit{Cơ sở toán} ngay bên dưới (\S\ref{s:layer}).

\paragraph{Mô tả đầu file và cách dùng.} Phần chú thích đầu file do tác giả viết, nêu ý tưởng và ví dụ sử dụng:
\VerbatimInput[firstline=1,lastline=38,firstnumber=1,numbers=left,numbersep=3pt,frame=lines,framesep=2pt,rulecolor=\color{gray!60},breaklines=true,fontsize=\fontsize{6.2}{7.2}\selectfont]{../../splat_editor/core/layer.py}
%% ---- merged math: s:layer
\subsubsection{Cơ sở toán: Biến đổi layer}\label{s:layer}
Phần này giải thích toán của \code{core/layer.py}, cùng với \code{core/transform.py} (dữ liệu \code{Transform}) và \code{core/math3d.py} (các hàm quay, xem mục \ref{s:quat}).

\paragraph{Ma trận quay từ góc Euler (\code{euler_to_matrix}).}
\[
R_x=\begin{bmatrix}1&0&0\\0&c_x&-s_x\\0&s_x&c_x\end{bmatrix},\;
R_y=\begin{bmatrix}c_y&0&s_y\\0&1&0\\-s_y&0&c_y\end{bmatrix},\;
R_z=\begin{bmatrix}c_z&-s_z&0\\s_z&c_z&0\\0&0&1\end{bmatrix},\;
R=R_zR_yR_x
\]
Với vector cột, $X$ được áp dụng trước, rồi $Y$, rồi $Z$.

\paragraph{Ngược lại (\code{matrix_to_euler}).}
Từ $R=R_zR_yR_x$ có $R_{20}=-\sin y$, nên $y=\operatorname{asin}(-R_{20})$ (kẹp $[-1,1]$ chống lỗi làm tròn), $x=\operatorname{atan2}(R_{21},R_{22})$, $z=\operatorname{atan2}(R_{10},R_{00})$. Khi $|R_{20}|\ge0.999999$ ($y=\pm90^\circ$, \textbf{khóa gimbal}) thì $x,z$ không xác định riêng, đặt $x=0$, $z=\operatorname{atan2}(-R_{01},R_{11})$.

\paragraph{Công thức Rodrigues (\code{axis_angle_matrix}).}
Quay góc $a$ quanh trục đơn vị $(x,y,z)$, $c=\cos a$, $s=\sin a$, $t=1-c$:
\[
R=\begin{bmatrix}tx^2+c & txy-sz & txz+sy\\ txy+sz & ty^2+c & tyz-sx\\ txz-sy & tyz+sx & tz^2+c\end{bmatrix}
= I + s[k]_\times + t[k]_\times^2 .
\]

\paragraph{Ma trận thế giới của layer (\code{Layer.world_transform / world_matrix}).}
\[
p_w = s\,A\,p_f + t,\quad A=R(\mathrm{euler})\,F,\quad F=\mathrm{diag}(-1,-1,1),\quad
\Sig_w = s^2 A\Sig_f A^{\top},\quad
M=\begin{bmatrix}sA & t\\ 0 & 1\end{bmatrix}.
\]
\begin{itemize}
\item $F$ là phép \textbf{xoay $180^\circ$ quanh $Z$} ($\det=+1$, không phải phản xạ). Nó đưa dữ liệu COLMAP (Y hướng xuống) về Y hướng lên.
\item Chỉ $M$ thay đổi (dữ liệu gaussian giữ nguyên), nên kéo/xoay/zoom layer là tức thời. $M$ được gửi lên GPU thành 4 texel (\code{splat_renderer.set_transforms}, chuyển vị để ra 4 cột).
\end{itemize}

\paragraph{Xoay quanh điểm pivot (\code{Transform.rotate_about(Δ, c)}).}
$R'=\Delta R$, $\;t'=\Delta(t-c)+c$ (dời pivot về gốc, xoay, dời lại). Góc Euler mới lấy bằng \code{matrix_to_euler}$(\Delta R)$.

\paragraph{Xuất file (\code{Layer._export_block}).}
Đổi từ thế giới về file: $\mathrm{rot}'=F\cdot\mathrm{rot}$, $t'=F t$. Khi đó:
\begin{itemize}
\item vị trí $p'=(s\,\mathrm{rot}')p+t'$;
\item pháp tuyến $n'=\mathrm{rot}'\,n$ (xoay thuần, không co giãn);
\item scale (log): $\mathrm{scale}_i'=\mathrm{scale}_i+\ln s$, vì $\exp(\mathrm{scale}+\ln s)=s\exp(\mathrm{scale})$;
\item quaternion $q'=q_{\mathrm{rot}}\otimes q$;
\item SH: $\mathrm{sh}'=D(\mathrm{rot}')\,\mathrm{sh}$ (mục \ref{s:shrot}).
\end{itemize}

\paragraph{\code{Layer.world_splats} (dùng cho vật lý, bounds).}
\code{positions = sp.positions @ (s*rot).T + t}; \code{rotations = quat_mul(q_rot, q)}; \code{scales *= s}.

\subsubsection*{Khai báo mức module (dòng 39--51)}
\VerbatimInput[firstline=39,lastline=51,firstnumber=39,numbers=left,numbersep=3pt,frame=lines,framesep=2pt,rulecolor=\color{gray!60},breaklines=true,fontsize=\fontsize{6.6}{7.6}\selectfont]{../../splat_editor/core/layer.py}
\subsubsection{\texorpdfstring{\code{class Layer}}{Layer}}
{\footnotesize\textit{Dòng 54--187}}\par
\textbf{Giải thích.} Các trường quan trọng: \code{data} (dữ liệu PLY gốc, thường là memmap, để lưu lại đủ thuộc tính), \code{splats} (đã giải mã), \code{selection} và \code{deleted} (hai mặt nạ boolean),
\code{transform}, \code{visible}.

Các phương thức: \code{__init__}, \code{from_file}, \code{size}, \code{count}, \code{live_rows}, \code{selected_count}, \code{select_all}, \code{clear_selection}, \code{invert_selection}, \code{delete_selected}, \code{restore}, \code{world_transform}, \code{world_matrix}, \code{world_splats}, \code{export_blocks}, \code{export_data}, \code{_export_block}.

\VerbatimInput[firstline=54,lastline=54,firstnumber=54,numbers=left,numbersep=3pt,frame=lines,framesep=2pt,rulecolor=\color{gray!60},breaklines=true,fontsize=\fontsize{6.6}{7.6}\selectfont]{../../splat_editor/core/layer.py}
\subsubsection{\texorpdfstring{\code{Layer.__init__(self, name, data, path=None)}}{Layer.__init__}}
{\footnotesize\textit{Dòng 55--64}}\par
\VerbatimInput[firstline=55,lastline=64,firstnumber=55,numbers=left,numbersep=3pt,frame=lines,framesep=2pt,rulecolor=\color{gray!60},breaklines=true,fontsize=\fontsize{6.6}{7.6}\selectfont]{../../splat_editor/core/layer.py}
\subsubsection{\texorpdfstring{\code{Layer.from_file(cls, path)}}{Layer.from_file}}
{\footnotesize\textit{Dòng 66--68}}\par
\VerbatimInput[firstline=66,lastline=68,firstnumber=66,numbers=left,numbersep=3pt,frame=lines,framesep=2pt,rulecolor=\color{gray!60},breaklines=true,fontsize=\fontsize{6.6}{7.6}\selectfont]{../../splat_editor/core/layer.py}
\subsubsection{\texorpdfstring{\code{Layer.size(self)}}{Layer.size}}
{\footnotesize\textit{Dòng 70--73}}\par
\textbf{Mô tả trong mã nguồn.} Số hàng của mảng dữ liệu, kể cả gaussian đã xoá.

\VerbatimInput[firstline=70,lastline=73,firstnumber=70,numbers=left,numbersep=3pt,frame=lines,framesep=2pt,rulecolor=\color{gray!60},breaklines=true,fontsize=\fontsize{6.6}{7.6}\selectfont]{../../splat_editor/core/layer.py}
\subsubsection{\texorpdfstring{\code{Layer.count(self)}}{Layer.count}}
{\footnotesize\textit{Dòng 75--78}}\par
\textbf{Mô tả trong mã nguồn.} Số gaussian còn lại (chưa xoá).

\VerbatimInput[firstline=75,lastline=78,firstnumber=75,numbers=left,numbersep=3pt,frame=lines,framesep=2pt,rulecolor=\color{gray!60},breaklines=true,fontsize=\fontsize{6.6}{7.6}\selectfont]{../../splat_editor/core/layer.py}
\subsubsection{\texorpdfstring{\code{Layer.live_rows(self, rows=slice(None))}}{Layer.live_rows}}
{\footnotesize\textit{Dòng 80--83}}\par
\textbf{Mô tả trong mã nguồn.} Chỉ số (trong toàn mảng) các gaussian chưa xoá thuộc đoạn `rows`.

\VerbatimInput[firstline=80,lastline=83,firstnumber=80,numbers=left,numbersep=3pt,frame=lines,framesep=2pt,rulecolor=\color{gray!60},breaklines=true,fontsize=\fontsize{6.6}{7.6}\selectfont]{../../splat_editor/core/layer.py}
\subsubsection{\texorpdfstring{\code{Layer.selected_count(self)}}{Layer.selected_count}}
{\footnotesize\textit{Dòng 86--88}}\par
\VerbatimInput[firstline=86,lastline=88,firstnumber=86,numbers=left,numbersep=3pt,frame=lines,framesep=2pt,rulecolor=\color{gray!60},breaklines=true,fontsize=\fontsize{6.6}{7.6}\selectfont]{../../splat_editor/core/layer.py}
\subsubsection{\texorpdfstring{\code{Layer.select_all(self)}}{Layer.select_all}}
{\footnotesize\textit{Dòng 90--91}}\par
\VerbatimInput[firstline=90,lastline=91,firstnumber=90,numbers=left,numbersep=3pt,frame=lines,framesep=2pt,rulecolor=\color{gray!60},breaklines=true,fontsize=\fontsize{6.6}{7.6}\selectfont]{../../splat_editor/core/layer.py}
\subsubsection{\texorpdfstring{\code{Layer.clear_selection(self)}}{Layer.clear_selection}}
{\footnotesize\textit{Dòng 93--94}}\par
\VerbatimInput[firstline=93,lastline=94,firstnumber=93,numbers=left,numbersep=3pt,frame=lines,framesep=2pt,rulecolor=\color{gray!60},breaklines=true,fontsize=\fontsize{6.6}{7.6}\selectfont]{../../splat_editor/core/layer.py}
\subsubsection{\texorpdfstring{\code{Layer.invert_selection(self)}}{Layer.invert_selection}}
{\footnotesize\textit{Dòng 96--97}}\par
\VerbatimInput[firstline=96,lastline=97,firstnumber=96,numbers=left,numbersep=3pt,frame=lines,framesep=2pt,rulecolor=\color{gray!60},breaklines=true,fontsize=\fontsize{6.6}{7.6}\selectfont]{../../splat_editor/core/layer.py}
\subsubsection{\texorpdfstring{\code{Layer.delete_selected(self)}}{Layer.delete_selected}}
{\footnotesize\textit{Dòng 99--108}}\par
\textbf{Mô tả trong mã nguồn.} Xoá (mềm) các gaussian đang chọn. Trả về chỉ số các gaussian đã xoá (để hoàn tác), hoặc None nếu không có gì.

\textbf{Giải thích.} Xóa mềm: bật cờ \code{deleted}, bỏ chọn, tăng \code{deleted_count}, trả về chỉ số để hoàn tác. Không chép mảng nào.

\VerbatimInput[firstline=99,lastline=108,firstnumber=99,numbers=left,numbersep=3pt,frame=lines,framesep=2pt,rulecolor=\color{gray!60},breaklines=true,fontsize=\fontsize{6.6}{7.6}\selectfont]{../../splat_editor/core/layer.py}
\subsubsection{\texorpdfstring{\code{Layer.restore(self, indices)}}{Layer.restore}}
{\footnotesize\textit{Dòng 110--114}}\par
\textbf{Mô tả trong mã nguồn.} Hoàn tác delete\_selected: các gaussian `indices` hiện lại và được chọn lại như lúc xoá.

\textbf{Giải thích.} Hoàn tác xóa: tắt cờ \code{deleted}, chọn lại đúng các gaussian đó và giảm \code{deleted_count}. Độ phức tạp $O(k)$ theo số gaussian khôi phục.

\VerbatimInput[firstline=110,lastline=114,firstnumber=110,numbers=left,numbersep=3pt,frame=lines,framesep=2pt,rulecolor=\color{gray!60},breaklines=true,fontsize=\fontsize{6.6}{7.6}\selectfont]{../../splat_editor/core/layer.py}
\subsubsection{\texorpdfstring{\code{Layer.world_transform(self, flip)}}{Layer.world_transform}}
{\footnotesize\textit{Dòng 117--120}}\par
\textbf{Mô tả trong mã nguồn.} (A, t, s): p\_thế\_giới = s * A @ p\_file + t, với A = R * F là phép xoay thuần.

\textbf{Giải thích.} Trả về bộ ba $(A,t,s)$ với $A=R\cdot F$ ($F$ chỉ có khi bật \code{flip}). Mọi nơi khác (ma trận 4$\times$4, xuất file, gaussian thế giới) đều bắt đầu từ hàm này.

\textbf{Công thức.}
\[
A=R(\mathrm{euler})\,F,\qquad F=\mathrm{diag}(-1,-1,1),\qquad p_w=s\,A\,p_f+t
\]


\VerbatimInput[firstline=117,lastline=120,firstnumber=117,numbers=left,numbersep=3pt,frame=lines,framesep=2pt,rulecolor=\color{gray!60},breaklines=true,fontsize=\fontsize{6.6}{7.6}\selectfont]{../../splat_editor/core/layer.py}
\subsubsection{\texorpdfstring{\code{Layer.world_matrix(self, flip)}}{Layer.world_matrix}}
{\footnotesize\textit{Dòng 122--129}}\par
\textbf{Mô tả trong mã nguồn.} Ma trận 4x4: p\_thế\_giới = M @ (p\_file, 1). Renderer / LOD / cọ dùng ma trận này với dữ liệu toạ độ file, nên đổi phép biến đổi không phải tính lại hay gửi lại dữ liệu.

\textbf{Giải thích.} Ma trận 4$\times$4 $M=[sA\,|\,t]$. Renderer, LOD và cọ dùng $M$ với dữ liệu \textbf{ở tọa độ file}, nên đổi phép biến đổi không phải tính lại hay gửi lại dữ liệu.

\textbf{Công thức.}
\[
M=\begin{bmatrix}sA&t\\0\;0\;0&1\end{bmatrix},\qquad
\begin{bmatrix}p_w\\1\end{bmatrix}=M\begin{bmatrix}p_f\\1\end{bmatrix},\qquad
\Sig_w=s^2A\,\Sig_f\,A^{\top}
\]


\VerbatimInput[firstline=122,lastline=129,firstnumber=122,numbers=left,numbersep=3pt,frame=lines,framesep=2pt,rulecolor=\color{gray!60},breaklines=true,fontsize=\fontsize{6.6}{7.6}\selectfont]{../../splat_editor/core/layer.py}
\subsubsection{\texorpdfstring{\code{Layer.world_splats(self, flip)}}{Layer.world_splats}}
{\footnotesize\textit{Dòng 131--147}}\par
\textbf{Mô tả trong mã nguồn.} Các gaussian còn lại ở toạ độ thế giới.

\textbf{Giải thích.} Tính gaussian ở tọa độ thế giới, từng đoạn \code{CHUNK} hàng. Vị trí: $p\,(sA)^\top+t$; quaternion: \code{quat_mul(q_rot, q)}; scale nhân $s$. SH không xoay ngay mà ghi vào \code{sh_frames}.

\textbf{Công thức.}
\[
p_w=p_f\,(sA)^{\top}+t,\qquad q_w=q_{A}\otimes q_f,\qquad s_w=s\cdot s_f
\]


\VerbatimInput[firstline=131,lastline=147,firstnumber=131,numbers=left,numbersep=3pt,frame=lines,framesep=2pt,rulecolor=\color{gray!60},breaklines=true,fontsize=\fontsize{6.6}{7.6}\selectfont]{../../splat_editor/core/layer.py}
\subsubsection{\texorpdfstring{\code{Layer.export_blocks(self, flip, dtype)}}{Layer.export_blocks}}
{\footnotesize\textit{Dòng 149--154}}\par
\textbf{Mô tả trong mã nguồn.} Dữ liệu của lớp sau khi áp phép biến đổi, đổi về toạ độ file, theo kiểu `dtype`, trả về từng khối CHUNK hàng (để ghi file mà không giữ cả lớp trong RAM).

\VerbatimInput[firstline=149,lastline=154,firstnumber=149,numbers=left,numbersep=3pt,frame=lines,framesep=2pt,rulecolor=\color{gray!60},breaklines=true,fontsize=\fontsize{6.6}{7.6}\selectfont]{../../splat_editor/core/layer.py}
\subsubsection{\texorpdfstring{\code{Layer.export_data(self, flip, dtype)}}{Layer.export_data}}
{\footnotesize\textit{Dòng 156--159}}\par
\textbf{Mô tả trong mã nguồn.} Toàn bộ dữ liệu của lớp sau khi biến đổi (xem export\_blocks).

\VerbatimInput[firstline=156,lastline=159,firstnumber=156,numbers=left,numbersep=3pt,frame=lines,framesep=2pt,rulecolor=\color{gray!60},breaklines=true,fontsize=\fontsize{6.6}{7.6}\selectfont]{../../splat_editor/core/layer.py}
\subsubsection{\texorpdfstring{\code{Layer._export_block(self, rows, flip, dtype)}}{Layer._export_block}}
{\footnotesize\textit{Dòng 161--187}}\par
\textbf{Giải thích.} Xuất một khối về tọa độ file. Từng bước: đổi \code{rot, t} về tọa độ file bằng $F$; xoay vị trí, pháp tuyến, cộng $\ln s$ vào scale (vì scale lưu dạng log), nhân quaternion và xoay SH. Chi tiết toán ở \S\ref{s:layer}.

\textbf{Công thức.}
Đổi ngược về tọa độ file: $A'=FA$, $t'=Ft$. Khi đó
\[
p'=sA'p+t',\quad n'=A'n,\quad \mathrm{scale}'_i=\mathrm{scale}_i+\ln s,\quad q'=q_{A'}\otimes q,\quad \mathrm{sh}'=D(A')\,\mathrm{sh}
\]
vì $\exp(\mathrm{scale}+\ln s)=s\,e^{\mathrm{scale}}$.


\VerbatimInput[firstline=161,lastline=187,firstnumber=161,numbers=left,numbersep=3pt,frame=lines,framesep=2pt,rulecolor=\color{gray!60},breaklines=true,fontsize=\fontsize{6.6}{7.6}\selectfont]{../../splat_editor/core/layer.py}
\subsection{\texorpdfstring{\code{splat_editor/core/scene.py}}{splat_editor/core/scene.py}}
{\footnotesize 174 dòng}\par
\textbf{Vai trò.} \code{Scene} quản lý danh sách layer và \textbf{gộp các layer đang hiện thành một mảng} để vẽ, chọn và lưu. Nó cũng quyết định kiểu dữ liệu khi lưu
(\code{export_dtype}) và cung cấp các tập điểm thưa dùng cho camera (\code{world_positions}) và va chạm (\code{solid_positions}).

\paragraph{Mô tả đầu file và cách dùng.} Phần chú thích đầu file do tác giả viết, nêu ý tưởng và ví dụ sử dụng:
\VerbatimInput[firstline=1,lastline=18,firstnumber=1,numbers=left,numbersep=3pt,frame=lines,framesep=2pt,rulecolor=\color{gray!60},breaklines=true,fontsize=\fontsize{6.2}{7.2}\selectfont]{../../splat_editor/core/scene.py}
\subsubsection*{Khai báo mức module (dòng 19--22)}
\VerbatimInput[firstline=19,lastline=22,firstnumber=19,numbers=left,numbersep=3pt,frame=lines,framesep=2pt,rulecolor=\color{gray!60},breaklines=true,fontsize=\fontsize{6.6}{7.6}\selectfont]{../../splat_editor/core/scene.py}
\subsubsection{\texorpdfstring{\code{class Scene}}{Scene}}
{\footnotesize\textit{Dòng 25--174}}\par
\textbf{Giải thích.} \code{active_index} là layer đang chọn trong bảng Layers; \code{flip} là cờ lật COLMAP dùng chung cho mọi layer.

Các phương thức: \code{__init__}, \code{active_layer}, \code{is_empty}, \code{add}, \code{remove}, \code{clear}, \code{visible_layers}, \code{visible_ranges}, \code{selection_mask}, \code{deleted_mask}, \code{apply_selection}, \code{apply_selection_indices}, \code{selected_count}, \code{count}, \code{render_splats}, \code{layer_matrices}, \code{world_positions}, \code{solid_positions}, \code{world_splats}, \code{export_dtype}, \code{export_blocks}, \code{export_data}.

\VerbatimInput[firstline=25,lastline=25,firstnumber=25,numbers=left,numbersep=3pt,frame=lines,framesep=2pt,rulecolor=\color{gray!60},breaklines=true,fontsize=\fontsize{6.6}{7.6}\selectfont]{../../splat_editor/core/scene.py}
\subsubsection{\texorpdfstring{\code{Scene.__init__(self)}}{Scene.__init__}}
{\footnotesize\textit{Dòng 26--29}}\par
\VerbatimInput[firstline=26,lastline=29,firstnumber=26,numbers=left,numbersep=3pt,frame=lines,framesep=2pt,rulecolor=\color{gray!60},breaklines=true,fontsize=\fontsize{6.6}{7.6}\selectfont]{../../splat_editor/core/scene.py}
\subsubsection{\texorpdfstring{\code{Scene.active_layer(self)}}{Scene.active_layer}}
{\footnotesize\textit{Dòng 32--34}}\par
\VerbatimInput[firstline=32,lastline=34,firstnumber=32,numbers=left,numbersep=3pt,frame=lines,framesep=2pt,rulecolor=\color{gray!60},breaklines=true,fontsize=\fontsize{6.6}{7.6}\selectfont]{../../splat_editor/core/scene.py}
\subsubsection{\texorpdfstring{\code{Scene.is_empty(self)}}{Scene.is_empty}}
{\footnotesize\textit{Dòng 36--38}}\par
\VerbatimInput[firstline=36,lastline=38,firstnumber=36,numbers=left,numbersep=3pt,frame=lines,framesep=2pt,rulecolor=\color{gray!60},breaklines=true,fontsize=\fontsize{6.6}{7.6}\selectfont]{../../splat_editor/core/scene.py}
\subsubsection{\texorpdfstring{\code{Scene.add(self, layer)}}{Scene.add}}
{\footnotesize\textit{Dòng 40--43}}\par
\VerbatimInput[firstline=40,lastline=43,firstnumber=40,numbers=left,numbersep=3pt,frame=lines,framesep=2pt,rulecolor=\color{gray!60},breaklines=true,fontsize=\fontsize{6.6}{7.6}\selectfont]{../../splat_editor/core/scene.py}
\subsubsection{\texorpdfstring{\code{Scene.remove(self, index)}}{Scene.remove}}
{\footnotesize\textit{Dòng 45--48}}\par
\VerbatimInput[firstline=45,lastline=48,firstnumber=45,numbers=left,numbersep=3pt,frame=lines,framesep=2pt,rulecolor=\color{gray!60},breaklines=true,fontsize=\fontsize{6.6}{7.6}\selectfont]{../../splat_editor/core/scene.py}
\subsubsection{\texorpdfstring{\code{Scene.clear(self)}}{Scene.clear}}
{\footnotesize\textit{Dòng 50--52}}\par
\VerbatimInput[firstline=50,lastline=52,firstnumber=50,numbers=left,numbersep=3pt,frame=lines,framesep=2pt,rulecolor=\color{gray!60},breaklines=true,fontsize=\fontsize{6.6}{7.6}\selectfont]{../../splat_editor/core/scene.py}
\subsubsection{\texorpdfstring{\code{Scene.visible_layers(self)}}{Scene.visible_layers}}
{\footnotesize\textit{Dòng 54--55}}\par
\VerbatimInput[firstline=54,lastline=55,firstnumber=54,numbers=left,numbersep=3pt,frame=lines,framesep=2pt,rulecolor=\color{gray!60},breaklines=true,fontsize=\fontsize{6.6}{7.6}\selectfont]{../../splat_editor/core/scene.py}
\subsubsection{\texorpdfstring{\code{Scene.visible_ranges(self)}}{Scene.visible_ranges}}
{\footnotesize\textit{Dòng 57--64}}\par
\textbf{Mô tả trong mã nguồn.} (lớp, đầu, cuối): vị trí của mỗi lớp đang hiện trong mảng gộp của render\_splats() (kể cả gaussian đã xoá).

\textbf{Giải thích.} Sinh ra bộ ba (layer, đầu, cuối): vị trí của từng layer đang hiện trong mảng gộp. Là xương sống để chuyển giữa "chỉ số theo layer" và "chỉ số theo mảng gộp".

\VerbatimInput[firstline=57,lastline=64,firstnumber=57,numbers=left,numbersep=3pt,frame=lines,framesep=2pt,rulecolor=\color{gray!60},breaklines=true,fontsize=\fontsize{6.6}{7.6}\selectfont]{../../splat_editor/core/scene.py}
\subsubsection{\texorpdfstring{\code{Scene.selection_mask(self)}}{Scene.selection_mask}}
{\footnotesize\textit{Dòng 67--70}}\par
\textbf{Mô tả trong mã nguồn.} Vùng chọn của các lớp đang hiện, gộp theo đúng thứ tự của render\_splats().

\VerbatimInput[firstline=67,lastline=70,firstnumber=67,numbers=left,numbersep=3pt,frame=lines,framesep=2pt,rulecolor=\color{gray!60},breaklines=true,fontsize=\fontsize{6.6}{7.6}\selectfont]{../../splat_editor/core/scene.py}
\subsubsection{\texorpdfstring{\code{Scene.deleted_mask(self)}}{Scene.deleted_mask}}
{\footnotesize\textit{Dòng 72--75}}\par
\textbf{Mô tả trong mã nguồn.} Gaussian đã xoá (xoá mềm) của các lớp đang hiện, theo thứ tự của render\_splats().

\VerbatimInput[firstline=72,lastline=75,firstnumber=72,numbers=left,numbersep=3pt,frame=lines,framesep=2pt,rulecolor=\color{gray!60},breaklines=true,fontsize=\fontsize{6.6}{7.6}\selectfont]{../../splat_editor/core/scene.py}
\subsubsection{\texorpdfstring{\code{Scene.apply_selection(self, mask, value)}}{Scene.apply_selection}}
{\footnotesize\textit{Dòng 77--80}}\par
\textbf{Mô tả trong mã nguồn.} Đặt vùng chọn = value cho các gaussian có mask True (mask theo mảng gộp; bỏ qua cái đã xoá).

\VerbatimInput[firstline=77,lastline=80,firstnumber=77,numbers=left,numbersep=3pt,frame=lines,framesep=2pt,rulecolor=\color{gray!60},breaklines=true,fontsize=\fontsize{6.6}{7.6}\selectfont]{../../splat_editor/core/scene.py}
\subsubsection{\texorpdfstring{\code{Scene.apply_selection_indices(self, indices, value)}}{Scene.apply_selection_indices}}
{\footnotesize\textit{Dòng 82--92}}\par
\textbf{Mô tả trong mã nguồn.} Đặt vùng chọn = value cho các gaussian `indices` (chỉ số theo mảng gộp). Gaussian đã xoá bị bỏ qua. Trả về các chỉ số thật sự được đặt.

\VerbatimInput[firstline=82,lastline=92,firstnumber=82,numbers=left,numbersep=3pt,frame=lines,framesep=2pt,rulecolor=\color{gray!60},breaklines=true,fontsize=\fontsize{6.6}{7.6}\selectfont]{../../splat_editor/core/scene.py}
\subsubsection{\texorpdfstring{\code{Scene.selected_count(self)}}{Scene.selected_count}}
{\footnotesize\textit{Dòng 94--95}}\par
\VerbatimInput[firstline=94,lastline=95,firstnumber=94,numbers=left,numbersep=3pt,frame=lines,framesep=2pt,rulecolor=\color{gray!60},breaklines=true,fontsize=\fontsize{6.6}{7.6}\selectfont]{../../splat_editor/core/scene.py}
\subsubsection{\texorpdfstring{\code{Scene.count(self)}}{Scene.count}}
{\footnotesize\textit{Dòng 98--100}}\par
\textbf{Mô tả trong mã nguồn.} Số gaussian còn lại (chưa xoá) của các lớp đang hiện.

\VerbatimInput[firstline=98,lastline=100,firstnumber=98,numbers=left,numbersep=3pt,frame=lines,framesep=2pt,rulecolor=\color{gray!60},breaklines=true,fontsize=\fontsize{6.6}{7.6}\selectfont]{../../splat_editor/core/scene.py}
\subsubsection{\texorpdfstring{\code{Scene.render_splats(self)}}{Scene.render_splats}}
{\footnotesize\textit{Dòng 102--107}}\par
\textbf{Mô tả trong mã nguồn.} Gộp các lớp đang hiện ở toạ độ file của từng lớp (không chép nếu chỉ có một lớp), kể cả gaussian đã xoá (renderer / cọ bỏ qua chúng theo deleted\_mask()). Toạ độ thế giới = layer\_matrices()[lớp] @ toạ độ file.

\textbf{Giải thích.} Gộp layer đang hiện \textbf{ở tọa độ file} (không chép nếu chỉ có một layer). Tọa độ thế giới $=$ \code{layer_matrices()[lớp]} $\times$ tọa độ file.

\VerbatimInput[firstline=102,lastline=107,firstnumber=102,numbers=left,numbersep=3pt,frame=lines,framesep=2pt,rulecolor=\color{gray!60},breaklines=true,fontsize=\fontsize{6.6}{7.6}\selectfont]{../../splat_editor/core/scene.py}
\subsubsection{\texorpdfstring{\code{Scene.layer_matrices(self)}}{Scene.layer_matrices}}
{\footnotesize\textit{Dòng 109--111}}\par
\textbf{Mô tả trong mã nguồn.} Ma trận thế giới 4x4 của mỗi lớp đang hiện (theo thứ tự trong render\_splats).

\VerbatimInput[firstline=109,lastline=111,firstnumber=109,numbers=left,numbersep=3pt,frame=lines,framesep=2pt,rulecolor=\color{gray!60},breaklines=true,fontsize=\fontsize{6.6}{7.6}\selectfont]{../../splat_editor/core/scene.py}
\subsubsection{\texorpdfstring{\code{Scene.world_positions(self, layer=None, limit=2000000)}}{Scene.world_positions}}
{\footnotesize\textit{Dòng 113--123}}\par
\textbf{Mô tả trong mã nguồn.} Tâm gaussian (toạ độ thế giới) của một lớp / các lớp đang hiện, lấy thưa tối đa `limit` điểm (đủ để đặt camera, khỏi biến đổi cả chục triệu điểm).

\VerbatimInput[firstline=113,lastline=123,firstnumber=113,numbers=left,numbersep=3pt,frame=lines,framesep=2pt,rulecolor=\color{gray!60},breaklines=true,fontsize=\fontsize{6.6}{7.6}\selectfont]{../../splat_editor/core/scene.py}
\subsubsection{\texorpdfstring{\code{Scene.solid_positions(self, limit=3000000, min_alpha=0.35)}}{Scene.solid_positions}}
{\footnotesize\textit{Dòng 125--137}}\par
\textbf{Mô tả trong mã nguồn.} Tâm gaussian "rắn" (toạ độ thế giới) dùng làm bề mặt va chạm: bỏ gaussian đã xoá và gaussian quá mờ (thường là nhiễu / vật bay lơ lửng). Lấy thưa tối đa `limit` điểm; trả về (điểm, bước lấy thưa).

\textbf{Giải thích.} Chọn các tâm "rắn" làm bề mặt va chạm: bỏ gaussian đã xóa và gaussian quá mờ ($\alpha<0.35$), lấy thưa tối đa 3 triệu điểm bằng bước \code{step}.

\VerbatimInput[firstline=125,lastline=137,firstnumber=125,numbers=left,numbersep=3pt,frame=lines,framesep=2pt,rulecolor=\color{gray!60},breaklines=true,fontsize=\fontsize{6.6}{7.6}\selectfont]{../../splat_editor/core/scene.py}
\subsubsection{\texorpdfstring{\code{Scene.world_splats(self)}}{Scene.world_splats}}
{\footnotesize\textit{Dòng 139--142}}\par
\textbf{Mô tả trong mã nguồn.} Gộp các gaussian còn lại của các lớp đang hiện thành một Splats (toạ độ thế giới).

\VerbatimInput[firstline=139,lastline=142,firstnumber=139,numbers=left,numbersep=3pt,frame=lines,framesep=2pt,rulecolor=\color{gray!60},breaklines=true,fontsize=\fontsize{6.6}{7.6}\selectfont]{../../splat_editor/core/scene.py}
\subsubsection{\texorpdfstring{\code{Scene.export_dtype(self)}}{Scene.export_dtype}}
{\footnotesize\textit{Dòng 144--161}}\par
\textbf{Mô tả trong mã nguồn.} Kiểu dữ liệu PLY khi lưu các lớp đang hiện (None nếu không có lớp nào).

Nếu các file cùng cấu trúc thì giữ nguyên mọi thuộc tính (kể cả f\_rest\_*). Nếu khác nhau thì chuyển hết sang dạng 3DGS chuẩn (giữ f\_rest\_* khi tất cả đều có).

\textbf{Giải thích.} Chọn kiểu cột khi lưu: các file cùng cấu trúc thì giữ nguyên mọi thuộc tính; khác nhau thì chuyển sang dạng 3DGS chuẩn và chỉ giữ \code{f_rest_*} khi tất cả layer đều có.

\VerbatimInput[firstline=144,lastline=161,firstnumber=144,numbers=left,numbersep=3pt,frame=lines,framesep=2pt,rulecolor=\color{gray!60},breaklines=true,fontsize=\fontsize{6.6}{7.6}\selectfont]{../../splat_editor/core/scene.py}
\subsubsection{\texorpdfstring{\code{Scene.export_blocks(self, dtype)}}{Scene.export_blocks}}
{\footnotesize\textit{Dòng 163--166}}\par
\textbf{Mô tả trong mã nguồn.} Dữ liệu PLY của các lớp đang hiện (đã áp phép biến đổi), từng khối một.

\VerbatimInput[firstline=163,lastline=166,firstnumber=163,numbers=left,numbersep=3pt,frame=lines,framesep=2pt,rulecolor=\color{gray!60},breaklines=true,fontsize=\fontsize{6.6}{7.6}\selectfont]{../../splat_editor/core/scene.py}
\subsubsection{\texorpdfstring{\code{Scene.export_data(self)}}{Scene.export_data}}
{\footnotesize\textit{Dòng 168--174}}\par
\textbf{Mô tả trong mã nguồn.} Dữ liệu PLY của tất cả lớp đang hiện, gộp thành một mảng (None nếu không có lớp nào).

\VerbatimInput[firstline=168,lastline=174,firstnumber=168,numbers=left,numbersep=3pt,frame=lines,framesep=2pt,rulecolor=\color{gray!60},breaklines=true,fontsize=\fontsize{6.6}{7.6}\selectfont]{../../splat_editor/core/scene.py}
\subsection{\texorpdfstring{\code{splat_editor/core/sh.py}}{splat_editor/core/sh.py}}
{\footnotesize 147 dòng}\par
\textbf{Vai trò.} Spherical Harmonics (SH) bậc 1--3 cho màu thay đổi theo hướng nhìn. Module cung cấp: hàm cơ sở (\code{basis}), đánh giá màu (\code{evaluate}),
\textbf{xoay hệ số} khi xoay gaussian (\code{rotation_matrix}, \code{rotate}) và đọc/ghi cột \code{f_rest_*} của PLY. Xoay hệ số được giải bằng bình phương tối thiểu
trên 64 hướng Fibonacci thay vì công thức Wigner-D. Toán học: xem mục \textit{Cơ sở toán} ngay bên dưới (\S\ref{s:shrot}).

\paragraph{Mô tả đầu file và cách dùng.} Phần chú thích đầu file do tác giả viết, nêu ý tưởng và ví dụ sử dụng:
\VerbatimInput[firstline=1,lastline=28,firstnumber=1,numbers=left,numbersep=3pt,frame=lines,framesep=2pt,rulecolor=\color{gray!60},breaklines=true,fontsize=\fontsize{6.2}{7.2}\selectfont]{../../splat_editor/core/sh.py}
%% ---- merged math: sh
\subsubsection{Cơ sở toán: Spherical Harmonics, màu theo hướng nhìn}\label{s:sh}
Phần này liên quan đến \code{core/sh.py} và hàm \code{shColor} trong shader \code{splat_prep.vert}.

\paragraph{Công thức màu.}
\[
c(d)=\mathrm{clip}\Big(0.5+C_0 f_{dc}+\sum_{k=1}^{15}Y_k(d)\,h_k,\;0,\;1\Big),
\]
với $d$ là vector đơn vị từ camera tới gaussian (3 kênh RGB độc lập).

\paragraph{Hàm cơ sở thực (đúng quy ước dấu của 3DGS gốc).}
Với $d=(x,y,z)$, $xx=x^2$, \ldots:
\begin{itemize}
\item Bậc 1 (3 hàm): $-C_1y,\; C_1z,\; -C_1x$ với $C_1=\sqrt{3/4\pi}=0.4886025$.
\item Bậc 2 (5 hàm): $1.0925\,xy$, $-1.0925\,yz$, $0.3154\,(2zz-xx-yy)$, $-1.0925\,xz$, $0.5463\,(xx-yy)$. Các hằng số: $1.0925=\tfrac12\sqrt{15/\pi}$, $0.3154=\tfrac14\sqrt{5/\pi}$, $0.5463=\tfrac14\sqrt{15/\pi}$.
\item Bậc 3 (7 hàm): $-0.5900\,y(3xx-yy)$, $2.8906\,xyz$, $-0.4570\,y(4zz-xx-yy)$, $0.3732\,z(2zz-3xx-3yy)$, $-0.4570\,x(4zz-xx-yy)$, $1.4453\,z(xx-yy)$, $-0.5900\,x(xx-3yy)$.
\end{itemize}
Đây là các đa thức Legendre liên kết trên mặt cầu: $1+3+5+7=16=(L+1)^2$ với $L=3$. Số hệ số bậc cao $K\in\{0,3,8,15\}$ (\code{COEFFS_FOR_DEGREE}). Trong file, \code{f_rest_{c*K+k}} là hệ số $k$ của kênh $c$ (xếp theo kênh). Shader chọn bậc theo \code{uShDegree}.

\paragraph{Xoay hệ số SH khi xoay gaussian (\code{sh.rotation_matrix}, \code{sh.rotate}).}\label{s:shrot}
Khi gaussian quay $R$, hàm màu mới là $\mathrm{m\grave{a}u}'(d)=\mathrm{m\grave{a}u}(R^{\top}d)$. Mỗi bậc SH đóng kín dưới phép quay, nên tồn tại $D(R)$ ($K\times K$, khối chéo theo bậc) sao cho $\mathrm{sh}'=D\cdot\mathrm{sh}$. Thay vì dùng công thức Wigner-D, code \textbf{giải bình phương tối thiểu}: chọn 64 hướng $d_i$, cần $B(d_i)D=B(R^{\top}d_i)$, nên
\[
D=\mathrm{lstsq}\big(B(\mathrm{dirs}),\,B(\mathrm{dirs}\,R)\big)
\]
(hàng $d_i^{\top}R=(R^{\top}d_i)^{\top}$). Kết quả chính xác vì $B$ có hạng đầy đủ.
\begin{itemize}
\item \textbf{Hướng mẫu Fibonacci} (\code{_sample_dirs}): $z=1-2(i+\tfrac12)/n$, $r=\sqrt{1-z^2}$, $\varphi=(i+\tfrac12)\,\pi(3-\sqrt5)$ (góc vàng $2.39996$ rad), cho điểm phân bố gần đều trên cầu.
\item Nhân khối bằng BLAS: $(n\cdot3,K)\,@\,D^{\top}$, xử lý từng \code{CHUNK} $2^{20}$.
\end{itemize}

\paragraph{Trong shader (không xoay hệ số mỗi frame).}
$d_{\mathrm{file}}=\mathrm{normalize}(\mathrm{linear}^{\top}(p_w-e))$: đưa hướng nhìn \textbf{ngược} về tọa độ file (dùng $A^{\top}$ vì quay trực giao; scale dương bị chuẩn hóa mất). Cách này rẻ hơn xoay hệ số mỗi khung.

\subsubsection*{Khai báo mức module (dòng 29--39)}
\VerbatimInput[firstline=29,lastline=39,firstnumber=29,numbers=left,numbersep=3pt,frame=lines,framesep=2pt,rulecolor=\color{gray!60},breaklines=true,fontsize=\fontsize{6.6}{7.6}\selectfont]{../../splat_editor/core/sh.py}
\subsubsection{\texorpdfstring{\code{degree_for_coeffs(k)}}{degree_for_coeffs}}
{\footnotesize\textit{Dòng 42--47}}\par
\textbf{Mô tả trong mã nguồn.} Số hệ số bậc cao -\textgreater{} bậc SH (None nếu không hợp lệ).

\VerbatimInput[firstline=42,lastline=47,firstnumber=42,numbers=left,numbersep=3pt,frame=lines,framesep=2pt,rulecolor=\color{gray!60},breaklines=true,fontsize=\fontsize{6.6}{7.6}\selectfont]{../../splat_editor/core/sh.py}
\subsubsection{\texorpdfstring{\code{basis(dirs, k=15)}}{basis}}
{\footnotesize\textit{Dòng 50--61}}\par
\textbf{Mô tả trong mã nguồn.} (M,3) hướng đơn vị -\textgreater{} (M,k) giá trị K hàm cơ sở SH bậc cao đầu tiên.

\textbf{Giải thích.} Giá trị các hàm cơ sở SH thực tại các hướng đơn vị; đúng quy ước dấu của 3DGS gốc. Trả về $(M,k)$ với $k\le15$.

\textbf{Công thức.}
\[
Y_{1}=\big(-C_1y,\ C_1z,\ -C_1x\big),\quad C_1=\sqrt{\tfrac{3}{4\pi}}
\]
\[
Y_{2}=\big(1.0925\,xy,\ -1.0925\,yz,\ 0.3154\,(2z^2-x^2-y^2),\ -1.0925\,xz,\ 0.5463\,(x^2-y^2)\big)
\]
\[
Y_{3}=\big(-0.59\,y(3x^2-y^2),\ 2.8906\,xyz,\ -0.457\,y(4z^2-x^2-y^2),\ 0.3732\,z(2z^2-3x^2-3y^2),\ -0.457\,x(4z^2-x^2-y^2),\ 1.4453\,z(x^2-y^2),\ -0.59\,x(x^2-3y^2)\big)
\]


\VerbatimInput[firstline=50,lastline=61,firstnumber=50,numbers=left,numbersep=3pt,frame=lines,framesep=2pt,rulecolor=\color{gray!60},breaklines=true,fontsize=\fontsize{6.6}{7.6}\selectfont]{../../splat_editor/core/sh.py}
\subsubsection{\texorpdfstring{\code{evaluate(sh_rest, dirs)}}{evaluate}}
{\footnotesize\textit{Dòng 64--66}}\par
\textbf{Mô tả trong mã nguồn.} Phần màu do SH bậc cao đóng góp: (N,K,3) hệ số, (N,3) hướng -\textgreater{} (N,3).

\textbf{Giải thích.} Phần màu do SH bậc cao: \code{einsum('nk,nkc->nc')} nhân hàm cơ sở với hệ số rồi cộng theo $k$.

\textbf{Công thức.}
\[
c(d)=\mathrm{clip}\Big(0.5+C_0f_{dc}+\sum_{k=1}^{K}Y_k(d)\,h_k,\ 0,\ 1\Big),\qquad C_0=\tfrac{1}{2\sqrt\pi}
\]


\VerbatimInput[firstline=64,lastline=66,firstnumber=64,numbers=left,numbersep=3pt,frame=lines,framesep=2pt,rulecolor=\color{gray!60},breaklines=true,fontsize=\fontsize{6.6}{7.6}\selectfont]{../../splat_editor/core/sh.py}
\subsubsection{\texorpdfstring{\code{_sample_dirs(n=64)}}{_sample_dirs}}
{\footnotesize\textit{Dòng 69--75}}\par
\textbf{Mô tả trong mã nguồn.} n hướng phân bố đều trên mặt cầu (xoắn ốc Fibonacci).

\textbf{Giải thích.} 64 hướng phân bố đều trên mặt cầu bằng xoắn ốc Fibonacci (góc vàng $\pi(3-\sqrt5)$).

\textbf{Công thức.}
\[
z_i=1-\frac{2(i+\tfrac12)}{n},\quad r_i=\sqrt{1-z_i^2},\quad \varphi_i=i\,\pi(3-\sqrt5),\quad d_i=(r_i\cos\varphi_i,\ r_i\sin\varphi_i,\ z_i)
\]


\VerbatimInput[firstline=69,lastline=75,firstnumber=69,numbers=left,numbersep=3pt,frame=lines,framesep=2pt,rulecolor=\color{gray!60},breaklines=true,fontsize=\fontsize{6.6}{7.6}\selectfont]{../../splat_editor/core/sh.py}
\subsubsection*{Khai báo mức module (dòng 78--78)}
\VerbatimInput[firstline=78,lastline=78,firstnumber=78,numbers=left,numbersep=3pt,frame=lines,framesep=2pt,rulecolor=\color{gray!60},breaklines=true,fontsize=\fontsize{6.6}{7.6}\selectfont]{../../splat_editor/core/sh.py}
\subsubsection{\texorpdfstring{\code{rotation_matrix(rot, k)}}{rotation_matrix}}
{\footnotesize\textit{Dòng 81--85}}\par
\textbf{Mô tả trong mã nguồn.} Ma trận D (k,k) để xoay hệ số SH theo phép xoay 3x3 `rot`: sh' = D @ sh.

\textbf{Giải thích.} Một dòng làm cả việc: giải $B(d)\,D=B(R^\top d)$ bằng \code{np.linalg.lstsq} trên 64 hướng, ra ma trận $D$ cỡ $K\times K$.

\textbf{Công thức.}
Màu sau khi xoay gaussian: $c'(d)=c(R^\top d)$. Cần $B(d)\,D=B(R^\top d)$ với mọi $d$, nên trên 64 hướng $d_i$:
\[
D=\arg\min_{D}\ \big\|B(\mathrm{dirs})\,D-B(\mathrm{dirs}\,R)\big\|_F^2,\qquad \mathrm{sh}'=D\,\mathrm{sh}
\]


\VerbatimInput[firstline=81,lastline=85,firstnumber=81,numbers=left,numbersep=3pt,frame=lines,framesep=2pt,rulecolor=\color{gray!60},breaklines=true,fontsize=\fontsize{6.6}{7.6}\selectfont]{../../splat_editor/core/sh.py}
\subsubsection*{Khai báo mức module (dòng 88--88)}
\VerbatimInput[firstline=88,lastline=88,firstnumber=88,numbers=left,numbersep=3pt,frame=lines,framesep=2pt,rulecolor=\color{gray!60},breaklines=true,fontsize=\fontsize{6.6}{7.6}\selectfont]{../../splat_editor/core/sh.py}
\subsubsection{\texorpdfstring{\code{rotate(sh_rest, rot)}}{rotate}}
{\footnotesize\textit{Dòng 91--104}}\par
\textbf{Mô tả trong mã nguồn.} Xoay hệ số SH (N,K,3) theo phép xoay 3x3 `rot`. Trả về mảng mới cùng kiểu dữ liệu.

\textbf{Giải thích.} Xoay hàng loạt hệ số SH. Nhân ma trận 2 chiều (BLAS) từng đoạn \code{CHUNK}; nhanh hơn \code{einsum} nhiều lần khi có hàng triệu gaussian.

\VerbatimInput[firstline=91,lastline=104,firstnumber=91,numbers=left,numbersep=3pt,frame=lines,framesep=2pt,rulecolor=\color{gray!60},breaklines=true,fontsize=\fontsize{6.6}{7.6}\selectfont]{../../splat_editor/core/sh.py}
\subsubsection{\texorpdfstring{\code{column_count(names)}}{column_count}}
{\footnotesize\textit{Dòng 107--113}}\par
\textbf{Mô tả trong mã nguồn.} Số hệ số SH bậc cao K đọc được từ các cột f\_rest\_* (0 nếu không có / không hợp lệ).

\VerbatimInput[firstline=107,lastline=113,firstnumber=107,numbers=left,numbersep=3pt,frame=lines,framesep=2pt,rulecolor=\color{gray!60},breaklines=true,fontsize=\fontsize{6.6}{7.6}\selectfont]{../../splat_editor/core/sh.py}
\subsubsection{\texorpdfstring{\code{from_columns(data, names, dtype=np.float32)}}{from_columns}}
{\footnotesize\textit{Dòng 116--126}}\par
\textbf{Mô tả trong mã nguồn.} Lấy hệ số SH bậc cao (N,K,3) từ các cột f\_rest\_* của PLY (K = 0 nếu không có / không hợp lệ). Đọc từng đoạn hàng (dữ liệu có thể là memmap trên đĩa) để không tạo mảng tạm khổng lồ.

\textbf{Giải thích.} Đọc hệ số từ cột \code{f_rest_*}. Dữ liệu có thể là memmap nên đọc từng đoạn (\code{read_block}) để không tạo mảng tạm khổng lồ.

\VerbatimInput[firstline=116,lastline=126,firstnumber=116,numbers=left,numbersep=3pt,frame=lines,framesep=2pt,rulecolor=\color{gray!60},breaklines=true,fontsize=\fontsize{6.6}{7.6}\selectfont]{../../splat_editor/core/sh.py}
\subsubsection{\texorpdfstring{\code{block_coeffs(block, matrix, k)}}{block_coeffs}}
{\footnotesize\textit{Dòng 129--139}}\par
\textbf{Mô tả trong mã nguồn.} (n,K,3) hệ số SH của một khối hàng PLY (matrix: khối dạng ma trận float32, hoặc None).

\VerbatimInput[firstline=129,lastline=139,firstnumber=129,numbers=left,numbersep=3pt,frame=lines,framesep=2pt,rulecolor=\color{gray!60},breaklines=true,fontsize=\fontsize{6.6}{7.6}\selectfont]{../../splat_editor/core/sh.py}
\subsubsection{\texorpdfstring{\code{to_columns(sh_rest, out)}}{to_columns}}
{\footnotesize\textit{Dòng 142--147}}\par
\textbf{Mô tả trong mã nguồn.} Ghi hệ số SH (N,K,3) vào các cột f\_rest\_* của mảng PLY `out`.

\VerbatimInput[firstline=142,lastline=147,firstnumber=142,numbers=left,numbersep=3pt,frame=lines,framesep=2pt,rulecolor=\color{gray!60},breaklines=true,fontsize=\fontsize{6.6}{7.6}\selectfont]{../../splat_editor/core/sh.py}
\subsection{\texorpdfstring{\code{splat_editor/core/lod.py}}{splat_editor/core/lod.py}}
{\footnotesize 208 dòng}\par
\textbf{Vai trò.} Mức chi tiết (LOD) để vẽ cảnh hàng chục triệu gaussian: chia thành khối theo mã Morton, xếp gaussian trong khối theo độ quan trọng, mỗi khung hình
bỏ khối ngoài khung nhìn và chia ngân sách gaussian theo khoảng cách. Toán học: xem mục \textit{Cơ sở toán} ngay bên dưới (\S\ref{s:lod}).

\paragraph{Mô tả đầu file và cách dùng.} Phần chú thích đầu file do tác giả viết, nêu ý tưởng và ví dụ sử dụng:
\VerbatimInput[firstline=1,lastline=29,firstnumber=1,numbers=left,numbersep=3pt,frame=lines,framesep=2pt,rulecolor=\color{gray!60},breaklines=true,fontsize=\fontsize{6.2}{7.2}\selectfont]{../../splat_editor/core/lod.py}
%% ---- merged math: s:lod
\subsubsection{Cơ sở toán: LOD, Morton, Frustum}\label{s:lod}
Các công thức dưới đây được cài trong \code{core/lod.py}; phần chọn gaussian theo độ sâu được \code{render/sorter.py} gọi (xem mục \ref{s:sort}).

\paragraph{Mã Morton 3D, đường cong Z (\code{morton_codes}).}
Hộp bao theo phân vị 0.1\%--99.9\% (bỏ nhiễu): $q=\mathrm{clip}((p-lo)\cdot1023/(hi-lo),0,1023)$ (10 bit/trục). Xen kẽ bit: \code{spread(v)} đặt bit $i$ của $v$ vào vị trí $3i$ (bảng tra); $\mathrm{code}=\mathrm{spread}(q_x)\mid\mathrm{spread}(q_y)\ll1\mid\mathrm{spread}(q_z)\ll2$ (30 bit). Điểm gần nhau trong không gian có mã gần nhau, tốt cho cache.

\paragraph{Radix argsort (\code{radix_argsort}).}
Hai lượt ổn định 16 bit (thấp rồi cao) cho thứ tự đúng theo khóa 32 bit (LSD radix).

\paragraph{Khối và độ quan trọng (\code{_build_group}).}
Cắt theo Morton thành khối 1024. Trong khối xếp theo $\mathrm{importance}=\alpha\,(s_0s_1+s_1s_2+s_0s_2)$ giảm dần (tỉ lệ với diện tích bề mặt ellipsoid $\times$ độ đục). Mẹo: khóa $=65535-(16\text{ bit cao của float32 dương})$, vì float dương sắp theo thứ tự bit nên không cần log.

\paragraph{Cầu bao mỗi khối.}
$\mathrm{radius}=\tfrac12\|hi-lo\|+3\max(s)$ (nửa đường chéo hộp tâm $+3\sigma$ lớn nhất). Khi có ma trận layer: $\mathrm{center}'=sA\,\mathrm{center}+t$, $\mathrm{radius}'=\mathrm{radius}\cdot\|\text{cột 0 của }M\|=\mathrm{radius}\cdot s$.

\paragraph{Frustum culling (\code{_frustum_planes}, \code{select}).}
Với $VP=PV$, 6 mặt phẳng (Gribb--Hartmann): hàng~3 $\pm$ hàng~0/1/2 cho $(a,b,c,d)$, chuẩn hóa bằng $\|(a,b,c)\|$. Khoảng cách có dấu $\mathrm{dist}=a\,c_x+b\,c_y+c\,c_z+d$; khối nhìn thấy nếu $\mathrm{dist}>-\mathrm{radius}$ với 5 mặt (bỏ mặt xa).

\paragraph{Phân bổ ngân sách (\code{select}).}
Khối $c$ lấy $\mathrm{take}_c=\min(n_c,\lceil\lambda\,n_c/d_c\rceil)$ với $d_c=\max(\mathrm{depth}-\mathrm{radius},\epsilon)$. Tìm $\lambda$ bằng \textbf{chia đôi 40 vòng} sao cho $\sum\mathrm{take}=\mathrm{budget}$. Khối gần ($d$ nhỏ) lấy nhiều gaussian, khối xa lấy ít; ưu tiên top importance vì đã được xếp trước. Độ sâu khối: $\mathrm{depth}=-(\mathrm{center}\cdot mv_2+mv_{23})$.

\subsubsection*{Khai báo mức module (dòng 30--38)}
\VerbatimInput[firstline=30,lastline=38,firstnumber=30,numbers=left,numbersep=3pt,frame=lines,framesep=2pt,rulecolor=\color{gray!60},breaklines=true,fontsize=\fontsize{6.6}{7.6}\selectfont]{../../splat_editor/core/lod.py}
\subsubsection{\texorpdfstring{\code{_parallel(fn, n, chunk=1 << 20)}}{_parallel}}
{\footnotesize\textit{Dòng 41--44}}\par
\textbf{Mô tả trong mã nguồn.} Chạy fn(start, end) trên các đoạn [0, n) ở nhiều luồng (numpy nhả GIL).

\VerbatimInput[firstline=41,lastline=44,firstnumber=41,numbers=left,numbersep=3pt,frame=lines,framesep=2pt,rulecolor=\color{gray!60},breaklines=true,fontsize=\fontsize{6.6}{7.6}\selectfont]{../../splat_editor/core/lod.py}
\subsubsection{\texorpdfstring{\code{gather(array, index)}}{gather}}
{\footnotesize\textit{Dòng 47--55}}\par
\textbf{Mô tả trong mã nguồn.} array[index] theo trục 0, chia nhiều luồng: truy cập ngẫu nhiên bị nghẽn bởi độ trễ bộ nhớ nên nhiều luồng cùng đọc nhanh hơn \textasciitilde{}3 lần.

\VerbatimInput[firstline=47,lastline=55,firstnumber=47,numbers=left,numbersep=3pt,frame=lines,framesep=2pt,rulecolor=\color{gray!60},breaklines=true,fontsize=\fontsize{6.6}{7.6}\selectfont]{../../splat_editor/core/lod.py}
\subsubsection{\texorpdfstring{\code{_spread_table()}}{_spread_table}}
{\footnotesize\textit{Dòng 58--64}}\par
\textbf{Mô tả trong mã nguồn.} Bảng tra: số 10 bit -\textgreater{} các bit đó cách nhau 2 bit 0 (để xen kẽ thành mã Morton 3D).

\VerbatimInput[firstline=58,lastline=64,firstnumber=58,numbers=left,numbersep=3pt,frame=lines,framesep=2pt,rulecolor=\color{gray!60},breaklines=true,fontsize=\fontsize{6.6}{7.6}\selectfont]{../../splat_editor/core/lod.py}
\subsubsection*{Khai báo mức module (dòng 67--67)}
\VerbatimInput[firstline=67,lastline=67,firstnumber=67,numbers=left,numbersep=3pt,frame=lines,framesep=2pt,rulecolor=\color{gray!60},breaklines=true,fontsize=\fontsize{6.6}{7.6}\selectfont]{../../splat_editor/core/lod.py}
\subsubsection{\texorpdfstring{\code{morton_codes(positions)}}{morton_codes}}
{\footnotesize\textit{Dòng 70--84}}\par
\textbf{Mô tả trong mã nguồn.} (N,3) -\textgreater{} (N,) mã Morton 30 bit (10 bit / trục, đủ để chia khối). Dùng phân vị 0.1\%..99.9\% làm hộp bao để điểm nhiễu ở xa không làm cả cảnh dồn vào vài ô.

\textbf{Giải thích.} Mã Morton 30 bit (10 bit mỗi trục) trên hộp bao theo phân vị 0.1\%--99.9\% để bỏ nhiễu. Điểm gần nhau trong không gian có mã gần nhau.

\textbf{Công thức.}
\[
q=\mathrm{clip}\Big(\frac{p-lo}{hi-lo}\cdot1023,\ 0,\ 1023\Big),\qquad
\mathrm{code}=\mathrm{spread}(q_x)\ \mid\ \mathrm{spread}(q_y)\!\ll\!1\ \mid\ \mathrm{spread}(q_z)\!\ll\!2
\]
$\mathrm{spread}(v)$ đặt bit $i$ của $v$ vào vị trí $3i$ (mã 30 bit).


\VerbatimInput[firstline=70,lastline=84,firstnumber=70,numbers=left,numbersep=3pt,frame=lines,framesep=2pt,rulecolor=\color{gray!60},breaklines=true,fontsize=\fontsize{6.6}{7.6}\selectfont]{../../splat_editor/core/lod.py}
\subsubsection{\texorpdfstring{\code{radix_argsort(key)}}{radix_argsort}}
{\footnotesize\textit{Dòng 87--91}}\par
\textbf{Mô tả trong mã nguồn.} argsort ổn định cho khoá uint32: 2 lượt sắp theo 16 bit (numpy dùng radix sort cho số 16 bit, O(N), nhanh hơn nhiều so với argsort số 32 bit).

\textbf{Giải thích.} Sắp xếp ổn định khóa uint32 bằng hai lượt 16 bit (thấp rồi cao). NumPy dùng radix sort cho kiểu 16 bit nên nhanh hơn \code{argsort} thường.

\VerbatimInput[firstline=87,lastline=91,firstnumber=87,numbers=left,numbersep=3pt,frame=lines,framesep=2pt,rulecolor=\color{gray!60},breaklines=true,fontsize=\fontsize{6.6}{7.6}\selectfont]{../../splat_editor/core/lod.py}
\subsubsection{\texorpdfstring{\code{_build_group(positions, scales, alphas, chunk_size)}}{_build_group}}
{\footnotesize\textit{Dòng 94--132}}\par
\textbf{Mô tả trong mã nguồn.} Chỉ mục cho một nhóm gaussian (một lớp): (order, positions, sizes, box\_lo, box\_hi, radii).

\textbf{Giải thích.} Dựng chỉ mục cho một layer: sắp theo Morton, cắt thành khối 1024, trong khối sắp theo độ quan trọng $\alpha(s_0s_1+s_1s_2+s_0s_2)$ giảm dần, tính hộp bao và cầu bao ($\tfrac12\|hi-lo\|+3\max s$).

\textbf{Công thức.}
\[
\mathrm{importance}=\alpha\,(s_0s_1+s_1s_2+s_0s_2),\qquad
\mathrm{radius}=\tfrac12\|hi-lo\|+3\max(s)
\]
Với layer có ma trận $M$: $\mathrm{center}'=sA\,\mathrm{center}+t$, $\ \mathrm{radius}'=s\cdot\mathrm{radius}$.


\VerbatimInput[firstline=94,lastline=132,firstnumber=94,numbers=left,numbersep=3pt,frame=lines,framesep=2pt,rulecolor=\color{gray!60},breaklines=true,fontsize=\fontsize{6.6}{7.6}\selectfont]{../../splat_editor/core/lod.py}
\subsubsection{\texorpdfstring{\code{class LodIndex}}{LodIndex}}
{\footnotesize\textit{Dòng 135--208}}\par
\textbf{Mô tả trong mã nguồn.} Chỉ mục khối của N gaussian, chia theo nhóm `ranges` (mỗi lớp một nhóm, khối không bao giờ trộn hai lớp). Toạ độ là toạ độ riêng của từng nhóm (toạ độ file của lớp); phép biến đổi của nhóm (ma trận 4x4, đổi được bất cứ lúc nào) được áp khi chọn / chiếu, nên kéo, xoay, phóng to một lớp không phải dựng lại chỉ mục.

\textbf{Giải thích.} Giữ \code{order} (hoán vị Morton), các khối với hộp bao và bán kính, chia theo nhóm (mỗi layer một nhóm để khối không vượt qua ranh giới layer).

Các phương thức: \code{__init__}, \code{world_bounds}, \code{_frustum_planes}, \code{select}.

\VerbatimInput[firstline=135,lastline=139,firstnumber=135,numbers=left,numbersep=3pt,frame=lines,framesep=2pt,rulecolor=\color{gray!60},breaklines=true,fontsize=\fontsize{6.6}{7.6}\selectfont]{../../splat_editor/core/lod.py}
\subsubsection{\texorpdfstring{\code{LodIndex.__init__(self, positions, scales, alphas, chunk_size=CHUNK_SIZE, ranges=None)}}{LodIndex.__init__}}
{\footnotesize\textit{Dòng 141--157}}\par
\VerbatimInput[firstline=141,lastline=157,firstnumber=141,numbers=left,numbersep=3pt,frame=lines,framesep=2pt,rulecolor=\color{gray!60},breaklines=true,fontsize=\fontsize{6.6}{7.6}\selectfont]{../../splat_editor/core/lod.py}
\subsubsection{\texorpdfstring{\code{LodIndex.world_bounds(self, transforms=None)}}{LodIndex.world_bounds}}
{\footnotesize\textit{Dòng 159--166}}\par
\textbf{Mô tả trong mã nguồn.} (tâm, bán kính) hình cầu bao của các khối trong toạ độ thế giới. transforms: ma trận 4x4 (s * R | t) của từng nhóm; None: đơn vị.

\VerbatimInput[firstline=159,lastline=166,firstnumber=159,numbers=left,numbersep=3pt,frame=lines,framesep=2pt,rulecolor=\color{gray!60},breaklines=true,fontsize=\fontsize{6.6}{7.6}\selectfont]{../../splat_editor/core/lod.py}
\subsubsection{\texorpdfstring{\code{LodIndex._frustum_planes(view_proj)}}{LodIndex._frustum_planes}}
{\footnotesize\textit{Dòng 168--173}}\par
\textbf{Mô tả trong mã nguồn.} 6 mặt phẳng (a,b,c,d) của khung nhìn, pháp tuyến hướng vào trong, đã chuẩn hoá.

\textbf{Giải thích.} Tách 6 mặt phẳng khung nhìn từ $VP=PV$ theo Gribb--Hartmann: hàng 3 $\pm$ hàng 0/1/2, chuẩn hóa bằng độ dài pháp tuyến.

\textbf{Công thức.}
Với $VP=PV$ có hàng $r_0,\dots,r_3$, sáu mặt phẳng $(a,b,c,d)$ (Gribb--Hartmann):
\[
r_3\pm r_0,\quad r_3\pm r_1,\quad r_3\pm r_2,\qquad (a,b,c,d)\leftarrow\frac{(a,b,c,d)}{\|(a,b,c)\|}
\]


\VerbatimInput[firstline=168,lastline=173,firstnumber=168,numbers=left,numbersep=3pt,frame=lines,framesep=2pt,rulecolor=\color{gray!60},breaklines=true,fontsize=\fontsize{6.6}{7.6}\selectfont]{../../splat_editor/core/lod.py}
\subsubsection{\texorpdfstring{\code{LodIndex.select(self, model_view, projection, budget, transforms=None)}}{LodIndex.select}}
{\footnotesize\textit{Dòng 175--208}}\par
\textbf{Mô tả trong mã nguồn.} Vị trí trong `self.order` (tăng dần) của các gaussian cần vẽ, chưa sắp theo độ sâu. Chỉ số gaussian gốc = self.order[kết quả].

\textbf{Giải thích.} Hàm quan trọng nhất của LOD. Bỏ khối ngoài khung nhìn (khoảng cách có dấu $>-$bán kính), tìm $\lambda$ bằng chia đôi 40 vòng sao cho tổng gaussian lấy bằng \code{budget}; khối gần lấy nhiều, khối xa lấy ít.

\textbf{Công thức.}
Khối nhìn thấy nếu $a\,c_x+b\,c_y+c\,c_z+d>-\mathrm{radius}$ với 5 mặt. Số gaussian lấy từ khối $c$:
\[
\mathrm{take}_c=\min\Big(n_c,\ \Big\lceil\frac{\lambda\,n_c}{d_c}\Big\rceil\Big),\quad d_c=\max(\mathrm{depth}_c-\mathrm{radius}_c,\ \epsilon),\quad \sum_c\mathrm{take}_c=\mathrm{budget}
\]
$\lambda$ tìm bằng chia đôi 40 vòng.


\VerbatimInput[firstline=175,lastline=208,firstnumber=175,numbers=left,numbersep=3pt,frame=lines,framesep=2pt,rulecolor=\color{gray!60},breaklines=true,fontsize=\fontsize{6.6}{7.6}\selectfont]{../../splat_editor/core/lod.py}
\subsection{\texorpdfstring{\code{splat_editor/core/lighting.py}}{splat_editor/core/lighting.py}}
{\footnotesize 247 dòng}\par
\textbf{Vai trò.} Trạng thái ánh sáng: ambient, một mặt trời, tối đa 8 khối đèn điểm, chu kỳ ngày--đêm dạng máy trạng thái và bản đồ bóng. Không có OpenGL:
\code{Lighting.params()} đổi trạng thái thành một \code{dict} để renderer gửi vào shader. Toán học: xem mục \textit{Cơ sở toán} ngay bên dưới (\S\ref{s:phong}).

\paragraph{Mô tả đầu file và cách dùng.} Phần chú thích đầu file do tác giả viết, nêu ý tưởng và ví dụ sử dụng:
\VerbatimInput[firstline=1,lastline=20,firstnumber=1,numbers=left,numbersep=3pt,frame=lines,framesep=2pt,rulecolor=\color{gray!60},breaklines=true,fontsize=\fontsize{6.2}{7.2}\selectfont]{../../splat_editor/core/lighting.py}
%% ---- merged math: s:sun
\subsubsection{Cơ sở toán: Mặt trời, ngày--đêm và bản đồ bóng}\label{s:sun}
Phần này gồm quỹ đạo mặt trời (\code{Lighting.sun_direction}), máy trạng thái \code{DayCycle} và thiết lập bản đồ bóng (\code{Lighting.shadow_setup}). Phần dùng bản đồ bóng trong shader (\code{shadowFactor}, \code{shadow.vert}) xem mục \ref{s:phong} (\code{splat_prep.vert}) và \code{shadow.vert}. Mô hình Phong và pháp tuyến gaussian cũng nằm ở mục \ref{s:phong}.

\paragraph{Mặt trời (\code{Lighting.sun_direction}).}
\begin{align*}
\mathrm{angle}&=\tfrac{\mathrm{hour}-6}{12}\pi\quad(\text{giờ }6\to0,\;12\to\tfrac\pi2,\;18\to\pi)\\
\mathrm{elevation}&=\sin(\mathrm{angle})\cdot65^\circ\quad(\text{âm ban đêm})\\
\mathrm{azimuth}&=\mathrm{azimuth}_0+(\mathrm{angle}-\tfrac\pi2)\\
\mathrm{dir}&=(\cos el\cos az,\;\sin el,\;\cos el\sin az)\quad(\text{cầu}\to\text{Descartes, đơn vị})
\end{align*}
Chu kỳ: \code{hour += dt*24/day_length (mod 24)}, \code{day_length = 36 s}, tức 1.5 s/giờ.
\begin{itemize}
\item $\mathrm{daylight}=\mathrm{smoothstep}(-0.08,0.3,\mathrm{dir}_y)$.
\item Màu mặt trời: $\mathrm{mix}((1,.5,.25),(1,.96,.88),\mathrm{smoothstep}(0,.35,s))\cdot I\cdot\mathrm{smoothstep}(-.02,.2,s)$ (cam lúc thấp, trắng lúc cao).
\item Ambient: $\mathrm{base}\cdot\mathrm{mix}((.12,.16,.32),1,\mathrm{daylight})$ (xanh tối ban đêm).
\item Bầu trời: $\mathrm{mix}(\text{đêm},\text{ngày},\mathrm{daylight})$ rồi trộn thêm cam với trọng số $0.55\exp(-(s/0.18)^2)$, một bump Gauss quanh đường chân trời (hoàng hôn).
\item $\mathrm{mix}(a,b,t)=a+(b-a)t$.
\end{itemize}

\paragraph{Máy trạng thái ngày--đêm (\code{DayCycle}).}
Mốc bắt đầu: DAWN 5h, DAY 7h, DUSK 17h, NIGHT 19h. \code{phase_for(h)}: $h\in[\mathrm{start},\mathrm{end})$, hoặc nếu $\mathrm{start}>\mathrm{end}$ (qua nửa đêm) thì $h\ge\mathrm{start}\lor h<\mathrm{end}$. Sự kiện \code{on_enter} bật/tắt đèn \code{auto_night}.

\paragraph{Bản đồ bóng (\code{Lighting.shadow_setup}, \code{shadow.vert}, \code{shadowFactor}).}
\begin{itemize}
\item $r=1.1\cdot\mathrm{radius}$ cảnh; mắt đèn $\mathrm{eye}=\mathrm{center}+d\cdot2r$; $\mathrm{up}=(0,1,0)$ (hoặc $(0,0,1)$ nếu $|d_y|\ge0.99$ để tránh $f\parallel\mathrm{up}$).
\item $VP=\mathrm{Ortho}(r,0.01r,4r)\cdot\mathrm{LookAt}(\mathrm{eye},\mathrm{center})$; độ phân giải $2048^2$.
\item $\mathrm{texel}=2r/2048$; $\mathrm{bias}=3\,\mathrm{texel}/(3.99r)$ (độ lệch độ sâu quy về NDC $[0,1]$; miền $z$ dài $3.99r$).
\item Điểm tra: $lp=VP\cdot(p+n\cdot2\,\mathrm{texel})$ (đẩy dọc pháp tuyến chống tự bóng); $s=lp_{xyz}/lp_w\cdot0.5+0.5$ (NDC sang $[0,1]$ để lấy mẫu texture); nằm ngoài $[0,1]$ thì coi là sáng.
\item \textbf{PCF $3\times3$:} $S=\frac19\sum_{i,j=-1}^{1}\big[s_z-\mathrm{bias}\le D(s_{xy}+(i,j)/2048)\big]$.
\item Chỉ gaussian có $\alpha\ge0.3$ và chưa xóa mới được vẽ vào depth map (dạng điểm).
\end{itemize}

\subsubsection*{Khai báo mức module (dòng 22--30)}
\VerbatimInput[firstline=21,lastline=29,firstnumber=21,numbers=left,numbersep=3pt,frame=lines,framesep=2pt,rulecolor=\color{gray!60},breaklines=true,fontsize=\fontsize{6.6}{7.6}\selectfont]{../../splat_editor/core/lighting.py}
\subsubsection{\texorpdfstring{\code{class DayPhase(Enum)}}{DayPhase}}
{\footnotesize\textit{Dòng 32--36}}\par
\VerbatimInput[firstline=32,lastline=36,firstnumber=32,numbers=left,numbersep=3pt,frame=lines,framesep=2pt,rulecolor=\color{gray!60},breaklines=true,fontsize=\fontsize{6.6}{7.6}\selectfont]{../../splat_editor/core/lighting.py}
\subsubsection{\texorpdfstring{\code{class DayCycle}}{DayCycle}}
{\footnotesize\textit{Dòng 40--65}}\par
\textbf{Mô tả trong mã nguồn.} Máy trạng thái ngày đêm. Mỗi trạng thái bắt đầu ở một giờ cố định; chuyển trạng thái khi giờ chạm mốc của trạng thái kế tiếp (hoặc nhảy thẳng tới trạng thái đúng nếu kéo thanh trượt giờ).

\textbf{Giải thích.} Máy trạng thái 4 pha: NIGHT $\to$ DAWN (5h) $\to$ DAY (7h) $\to$ DUSK (17h) $\to$ NIGHT (19h). Khi chuyển pha gọi sự kiện \code{on_enter} (bật/tắt đèn \code{auto_night}).

\textbf{Công thức.}
\[
\mathrm{phase}(h)=\begin{cases}\text{DAWN}&5\le h<7\\\text{DAY}&7\le h<17\\\text{DUSK}&17\le h<19\\\text{NIGHT}&h\ge19\ \text{hoặc}\ h<5\end{cases}
\]


Các phương thức: \code{__init__}, \code{phase_for}, \code{update}.

\VerbatimInput[firstline=39,lastline=45,firstnumber=39,numbers=left,numbersep=3pt,frame=lines,framesep=2pt,rulecolor=\color{gray!60},breaklines=true,fontsize=\fontsize{6.6}{7.6}\selectfont]{../../splat_editor/core/lighting.py}
\subsubsection{\texorpdfstring{\code{DayCycle.__init__(self, hour=10.0)}}{DayCycle.__init__}}
{\footnotesize\textit{Dòng 47--48}}\par
\VerbatimInput[firstline=47,lastline=48,firstnumber=47,numbers=left,numbersep=3pt,frame=lines,framesep=2pt,rulecolor=\color{gray!60},breaklines=true,fontsize=\fontsize{6.6}{7.6}\selectfont]{../../splat_editor/core/lighting.py}
\subsubsection{\texorpdfstring{\code{DayCycle.phase_for(cls, hour)}}{DayCycle.phase_for}}
{\footnotesize\textit{Dòng 50--56}}\par
\VerbatimInput[firstline=50,lastline=56,firstnumber=50,numbers=left,numbersep=3pt,frame=lines,framesep=2pt,rulecolor=\color{gray!60},breaklines=true,fontsize=\fontsize{6.6}{7.6}\selectfont]{../../splat_editor/core/lighting.py}
\subsubsection{\texorpdfstring{\code{DayCycle.update(self, hour)}}{DayCycle.update}}
{\footnotesize\textit{Dòng 58--64}}\par
\textbf{Mô tả trong mã nguồn.} Trả về trạng thái mới nếu vừa chuyển, ngược lại None (sự kiện on\_enter).

\VerbatimInput[firstline=58,lastline=64,firstnumber=58,numbers=left,numbersep=3pt,frame=lines,framesep=2pt,rulecolor=\color{gray!60},breaklines=true,fontsize=\fontsize{6.6}{7.6}\selectfont]{../../splat_editor/core/lighting.py}
\subsubsection{\texorpdfstring{\code{class PointLight}}{PointLight}}
{\footnotesize\textit{Dòng 67--75}}\par
\VerbatimInput[firstline=67,lastline=75,firstnumber=67,numbers=left,numbersep=3pt,frame=lines,framesep=2pt,rulecolor=\color{gray!60},breaklines=true,fontsize=\fontsize{6.6}{7.6}\selectfont]{../../splat_editor/core/lighting.py}
\subsubsection{\texorpdfstring{\code{_smooth(edge0, edge1, x)}}{_smooth}}
{\footnotesize\textit{Dòng 78--80}}\par
\VerbatimInput[firstline=78,lastline=80,firstnumber=78,numbers=left,numbersep=3pt,frame=lines,framesep=2pt,rulecolor=\color{gray!60},breaklines=true,fontsize=\fontsize{6.6}{7.6}\selectfont]{../../splat_editor/core/lighting.py}
\subsubsection{\texorpdfstring{\code{_mix(a, b, t)}}{_mix}}
{\footnotesize\textit{Dòng 83--84}}\par
\VerbatimInput[firstline=83,lastline=84,firstnumber=83,numbers=left,numbersep=3pt,frame=lines,framesep=2pt,rulecolor=\color{gray!60},breaklines=true,fontsize=\fontsize{6.6}{7.6}\selectfont]{../../splat_editor/core/lighting.py}
\subsubsection{\texorpdfstring{\code{class Lighting}}{Lighting}}
{\footnotesize\textit{Dòng 91--276}}\par
Các phương thức: \code{__init__}, \code{set_bounds}, \code{add_sun}, \code{remove_sun}, \code{add_light}, \code{remove_light}, \code{active_light}, \code{sun_direction}, \code{daylight}, \code{sun_color}, \code{sun_color_for_gizmo}, \code{ambient_color}, \code{sky_color}, \code{update}, \code{_on_enter}, \code{shadow_setup}, \code{point_lights}, \code{params}.

\VerbatimInput[firstline=87,lastline=88,firstnumber=87,numbers=left,numbersep=3pt,frame=lines,framesep=2pt,rulecolor=\color{gray!60},breaklines=true,fontsize=\fontsize{6.6}{7.6}\selectfont]{../../splat_editor/core/lighting.py}
\subsubsection{\texorpdfstring{\code{Lighting.__init__(self)}}{Lighting.__init__}}
{\footnotesize\textit{Dòng 90--112}}\par
\VerbatimInput[firstline=90,lastline=112,firstnumber=90,numbers=left,numbersep=3pt,frame=lines,framesep=2pt,rulecolor=\color{gray!60},breaklines=true,fontsize=\fontsize{6.6}{7.6}\selectfont]{../../splat_editor/core/lighting.py}
\subsubsection{\texorpdfstring{\code{Lighting.set_bounds(self, center, radius)}}{Lighting.set_bounds}}
{\footnotesize\textit{Dòng 115--117}}\par
\VerbatimInput[firstline=115,lastline=117,firstnumber=115,numbers=left,numbersep=3pt,frame=lines,framesep=2pt,rulecolor=\color{gray!60},breaklines=true,fontsize=\fontsize{6.6}{7.6}\selectfont]{../../splat_editor/core/lighting.py}
\subsubsection{\texorpdfstring{\code{Lighting.add_sun(self)}}{Lighting.add_sun}}
{\footnotesize\textit{Dòng 119--120}}\par
\VerbatimInput[firstline=119,lastline=120,firstnumber=119,numbers=left,numbersep=3pt,frame=lines,framesep=2pt,rulecolor=\color{gray!60},breaklines=true,fontsize=\fontsize{6.6}{7.6}\selectfont]{../../splat_editor/core/lighting.py}
\subsubsection{\texorpdfstring{\code{Lighting.remove_sun(self)}}{Lighting.remove_sun}}
{\footnotesize\textit{Dòng 122--123}}\par
\VerbatimInput[firstline=122,lastline=123,firstnumber=122,numbers=left,numbersep=3pt,frame=lines,framesep=2pt,rulecolor=\color{gray!60},breaklines=true,fontsize=\fontsize{6.6}{7.6}\selectfont]{../../splat_editor/core/lighting.py}
\subsubsection{\texorpdfstring{\code{Lighting.add_light(self, position, light_range=None)}}{Lighting.add_light}}
{\footnotesize\textit{Dòng 125--133}}\par
\VerbatimInput[firstline=125,lastline=133,firstnumber=125,numbers=left,numbersep=3pt,frame=lines,framesep=2pt,rulecolor=\color{gray!60},breaklines=true,fontsize=\fontsize{6.6}{7.6}\selectfont]{../../splat_editor/core/lighting.py}
\subsubsection{\texorpdfstring{\code{Lighting.remove_light(self, index)}}{Lighting.remove_light}}
{\footnotesize\textit{Dòng 135--139}}\par
\VerbatimInput[firstline=135,lastline=139,firstnumber=135,numbers=left,numbersep=3pt,frame=lines,framesep=2pt,rulecolor=\color{gray!60},breaklines=true,fontsize=\fontsize{6.6}{7.6}\selectfont]{../../splat_editor/core/lighting.py}
\subsubsection{\texorpdfstring{\code{Lighting.active_light(self)}}{Lighting.active_light}}
{\footnotesize\textit{Dòng 141--143}}\par
\VerbatimInput[firstline=141,lastline=143,firstnumber=141,numbers=left,numbersep=3pt,frame=lines,framesep=2pt,rulecolor=\color{gray!60},breaklines=true,fontsize=\fontsize{6.6}{7.6}\selectfont]{../../splat_editor/core/lighting.py}
\subsubsection{\texorpdfstring{\code{Lighting.sun_direction(self)}}{Lighting.sun_direction}}
{\footnotesize\textit{Dòng 146--152}}\par
\textbf{Mô tả trong mã nguồn.} Vector đơn vị từ cảnh hướng về mặt trời. Giờ 6: mọc, 12: trên cao nhất, 18: lặn.

\textbf{Giải thích.} Tọa độ cầu $\to$ Descartes: góc giờ $\tfrac{hour-6}{12}\pi$, độ cao $\sin(angle)\cdot65^\circ$. Giờ 6 mọc, 12 ở đỉnh, 18 lặn.

\textbf{Công thức.}
\[
\theta=\frac{h-6}{12}\pi,\quad el=65^\circ\sin\theta,\quad az=az_0+\Big(\theta-\frac\pi2\Big),\quad
\hat d=\begin{bmatrix}\cos el\cos az\\\sin el\\\cos el\sin az\end{bmatrix}
\]


\VerbatimInput[firstline=146,lastline=152,firstnumber=146,numbers=left,numbersep=3pt,frame=lines,framesep=2pt,rulecolor=\color{gray!60},breaklines=true,fontsize=\fontsize{6.6}{7.6}\selectfont]{../../splat_editor/core/lighting.py}
\subsubsection{\texorpdfstring{\code{Lighting.daylight(self)}}{Lighting.daylight}}
{\footnotesize\textit{Dòng 154--156}}\par
\textbf{Mô tả trong mã nguồn.} 0 (đêm) .. 1 (ban ngày), mượt quanh lúc mặt trời qua đường chân trời.

\VerbatimInput[firstline=154,lastline=156,firstnumber=154,numbers=left,numbersep=3pt,frame=lines,framesep=2pt,rulecolor=\color{gray!60},breaklines=true,fontsize=\fontsize{6.6}{7.6}\selectfont]{../../splat_editor/core/lighting.py}
\subsubsection{\texorpdfstring{\code{Lighting.sun_color(self)}}{Lighting.sun_color}}
{\footnotesize\textit{Dòng 158--162}}\par
\VerbatimInput[firstline=158,lastline=162,firstnumber=158,numbers=left,numbersep=3pt,frame=lines,framesep=2pt,rulecolor=\color{gray!60},breaklines=true,fontsize=\fontsize{6.6}{7.6}\selectfont]{../../splat_editor/core/lighting.py}
\subsubsection{\texorpdfstring{\code{Lighting.sun_color_for_gizmo(self)}}{Lighting.sun_color_for_gizmo}}
{\footnotesize\textit{Dòng 164--166}}\par
\textbf{Mô tả trong mã nguồn.} Màu quả cầu mặt trời (luôn sáng rõ, đỏ cam khi gần đường chân trời).

\VerbatimInput[firstline=164,lastline=166,firstnumber=164,numbers=left,numbersep=3pt,frame=lines,framesep=2pt,rulecolor=\color{gray!60},breaklines=true,fontsize=\fontsize{6.6}{7.6}\selectfont]{../../splat_editor/core/lighting.py}
\subsubsection{\texorpdfstring{\code{Lighting.ambient_color(self)}}{Lighting.ambient_color}}
{\footnotesize\textit{Dòng 168--173}}\par
\VerbatimInput[firstline=168,lastline=173,firstnumber=168,numbers=left,numbersep=3pt,frame=lines,framesep=2pt,rulecolor=\color{gray!60},breaklines=true,fontsize=\fontsize{6.6}{7.6}\selectfont]{../../splat_editor/core/lighting.py}
\subsubsection{\texorpdfstring{\code{Lighting.sky_color(self)}}{Lighting.sky_color}}
{\footnotesize\textit{Dòng 175--179}}\par
\textbf{Giải thích.} Màu trời: trộn đêm/ngày theo \code{daylight}, rồi thêm màu cam hoàng hôn với trọng số Gauss quanh đường chân trời.

\textbf{Công thức.}
\[
\mathrm{daylight}=\mathrm{smoothstep}(-0.08,0.3,\hat d_y),\qquad
c_{sky}=\mathrm{mix}(c_{night},c_{day},\mathrm{daylight})+0.55\,e^{-(s/0.18)^2}\,c_{dusk}
\]
$\mathrm{smoothstep}(a,b,x)=t^2(3-2t)$, $t=\mathrm{clip}\big(\tfrac{x-a}{b-a},0,1\big)$.


\VerbatimInput[firstline=175,lastline=179,firstnumber=175,numbers=left,numbersep=3pt,frame=lines,framesep=2pt,rulecolor=\color{gray!60},breaklines=true,fontsize=\fontsize{6.6}{7.6}\selectfont]{../../splat_editor/core/lighting.py}
\subsubsection{\texorpdfstring{\code{Lighting.update(self, dt)}}{Lighting.update}}
{\footnotesize\textit{Dòng 184--188}}\par
\textbf{Mô tả trong mã nguồn.} Tiến thời gian, chạy máy trạng thái ngày đêm.

\textbf{Giải thích.} Mỗi khung hình: tăng giờ ($24/day\_length$ giờ mỗi giây), chạy máy trạng thái ngày đêm.

\textbf{Công thức.}
\[
h\leftarrow(h+\Delta t\cdot24/T_{day})\bmod24,\qquad T_{day}=36\ \mathrm{s}\ \Rightarrow\ 1.5\ \mathrm{s/giờ}
\]


\VerbatimInput[firstline=184,lastline=188,firstnumber=184,numbers=left,numbersep=3pt,frame=lines,framesep=2pt,rulecolor=\color{gray!60},breaklines=true,fontsize=\fontsize{6.6}{7.6}\selectfont]{../../splat_editor/core/lighting.py}
\subsubsection{\texorpdfstring{\code{Lighting._on_enter(self, phase)}}{Lighting._on_enter}}
{\footnotesize\textit{Dòng 190--201}}\par
\VerbatimInput[firstline=190,lastline=201,firstnumber=190,numbers=left,numbersep=3pt,frame=lines,framesep=2pt,rulecolor=\color{gray!60},breaklines=true,fontsize=\fontsize{6.6}{7.6}\selectfont]{../../splat_editor/core/lighting.py}
\subsubsection{\texorpdfstring{\code{Lighting.shadow_setup(self)}}{Lighting.shadow_setup}}
{\footnotesize\textit{Dòng 204--213}}\par
\textbf{Mô tả trong mã nguồn.} (ma trận ánh sáng, kích thước texel theo thế giới, độ lệch độ sâu) cho bản đồ bóng trực giao.

\textbf{Giải thích.} Dựng ma trận ánh sáng cho bản đồ bóng trực giao: bán kính $r=1.1\times$bán kính cảnh, mắt đèn $=$ tâm $+d\cdot2r$; trả về $(VP,\ texel,\ bias)$.

\textbf{Công thức.}
\[
r=1.1\,R_{scene},\quad e=c+\hat d\cdot2r,\quad VP=\mathrm{Ortho}(r,\,0.01r,\,4r)\cdot\mathrm{LookAt}(e,c),
\]
\[
\mathrm{texel}=\frac{2r}{2048},\qquad \mathrm{bias}=\frac{3\,\mathrm{texel}}{3.99\,r}
\]


\VerbatimInput[firstline=204,lastline=213,firstnumber=204,numbers=left,numbersep=3pt,frame=lines,framesep=2pt,rulecolor=\color{gray!60},breaklines=true,fontsize=\fontsize{6.6}{7.6}\selectfont]{../../splat_editor/core/lighting.py}
\subsubsection{\texorpdfstring{\code{Lighting.point_lights(self)}}{Lighting.point_lights}}
{\footnotesize\textit{Dòng 216--217}}\par
\VerbatimInput[firstline=216,lastline=217,firstnumber=216,numbers=left,numbersep=3pt,frame=lines,framesep=2pt,rulecolor=\color{gray!60},breaklines=true,fontsize=\fontsize{6.6}{7.6}\selectfont]{../../splat_editor/core/lighting.py}
\subsubsection{\texorpdfstring{\code{Lighting.params(self)}}{Lighting.params}}
{\footnotesize\textit{Dòng 219--247}}\par
\textbf{Mô tả trong mã nguồn.} Trạng thái ánh sáng dạng dict cho renderer; `signature` đổi khi cần tính lại màu gaussian.

\textbf{Giải thích.} Đóng gói trạng thái thành \code{dict} cho renderer. Trường \code{signature} đổi khi phải tính lại màu của mọi gaussian (pass 1).

\VerbatimInput[firstline=219,lastline=247,firstnumber=219,numbers=left,numbersep=3pt,frame=lines,framesep=2pt,rulecolor=\color{gray!60},breaklines=true,fontsize=\fontsize{6.6}{7.6}\selectfont]{../../splat_editor/core/lighting.py}
\subsection{\texorpdfstring{\code{splat_editor/core/physics.py}}{splat_editor/core/physics.py}}
{\footnotesize 320 dòng}\par
\textbf{Vai trò.} Vật lý cho chế độ sinh tồn: voxel hóa tâm gaussian thành bản đồ va chạm, người chơi (trọng lực, nhảy, leo bậc, trượt tường).
Truy vấn cột dùng khóa 63 bit và \code{np.searchsorted} nên rất rẻ. Toán học: xem mục \textit{Cơ sở toán} ngay bên dưới (\S\ref{s:phys}).

\paragraph{Mô tả đầu file và cách dùng.} Phần chú thích đầu file do tác giả viết, nêu ý tưởng và ví dụ sử dụng:
\VerbatimInput[firstline=1,lastline=22,firstnumber=1,numbers=left,numbersep=3pt,frame=lines,framesep=2pt,rulecolor=\color{gray!60},breaklines=true,fontsize=\fontsize{6.2}{7.2}\selectfont]{../../splat_editor/core/physics.py}
%% ---- merged math: s:phys
\subsubsection{Cơ sở toán: Vật lý và va chạm}\label{s:phys}

\paragraph{Voxel hóa (\code{CollisionWorld.build}).}
$\mathrm{cell}=\lfloor(p-\mathrm{origin})/v\rfloor$, $\mathrm{origin}=\min-v$; $v=\mathrm{height}/8$; kẹp $[0,2^{21}-1]$. \textbf{Khóa 63 bit:} $\mathrm{key}=(c_x\ll42)\mid(c_z\ll21)\mid c_y$ (mỗi trục 21 bit). $c_y$ ở bit thấp nên các ô cùng cột $(c_x,c_z)$ nằm liền nhau trong mảng đã sắp. Ô đặc nếu có ít nhất \code{min_count} gaussian (3/2/1 tùy mức lấy mẫu). \code{floor_y} là phân vị 1\% của $y$ (sàn an toàn).

\paragraph{Truy vấn cột (\code{_highest}).}
$\mathrm{query}=\mathrm{key}(i_x,i_z,i_y)$; $\mathrm{idx}=\mathrm{searchsorted}(\mathrm{keys},\mathrm{query},\text{`right'})-1$, tức khóa lớn nhất $\le\mathrm{query}$; hợp lệ nếu $\mathrm{found}\gg21=\mathrm{query}\gg21$ (cùng cột). Trả về $i_y$ cao nhất $\le y$. Độ cao mặt trên ô: $\mathrm{origin}_y+(i_y+1)v$. Độ phức tạp $O(\log N)$ mỗi truy vấn, đã vector hóa.

\paragraph{Tia kiểm tra tầm nhìn (\code{segment_blocked}).}
Lấy mẫu dọc đoạn với bước $v/2$: $n=\lfloor L/(v/2)\rfloor+1$, $t=\mathrm{linspace}(0,1,n)$; điểm $=\mathrm{start}+(\mathrm{end}-\mathrm{start})t$; tra khóa; trúng ô đặc thì bị chặn.

\paragraph{Người chơi.}
Kích thước tỉ lệ với chiều cao $H$ (2.74): mắt $0.92H$, bán kính $0.14H$, bước leo $0.35H$, nhảy $0.75H$, $g=12H$, tốc độ đi $2.5H$ ($\times1.7$ chạy, $\times0.4$ rình).
\begin{itemize}
\item \textbf{Va chạm cột trụ:} 9 điểm vòng \code{_RING} $=$ tâm + 8 điểm $(\cos\frac{k\pi}{4},\sin\frac{k\pi}{4})\cdot\mathrm{bán\ kính}$.
\item \textbf{Trượt tường:} thử $(x+dx,z+dz)$; nếu bị chặn thử chỉ $x$; rồi chỉ $z$.
\item \textbf{Nhảy:} $v_y=\sqrt{2g\,h_{\mathrm{nhảy}}}$ (từ $v^2=2gh$).
\item \textbf{Euler:} $v_y\leftarrow v_y-g\Delta t$, $y\leftarrow y+v_y\Delta t$ ($\Delta t$ kẹp $\le0.05$); chạm đất khi $v_y\le0\ \land\ y_{\mathrm{mới}}\le\mathrm{ground}$.
\item \textbf{Làm mượt mắt:} $\mathrm{eye}_y\mathrel{+}=(\mathrm{target}-\mathrm{eye}_y)\min(18\Delta t,1)$ (lọc mũ, che bậc thang).
\end{itemize}

\subsubsection*{Khai báo mức module (dòng 23--31)}
\VerbatimInput[firstline=23,lastline=31,firstnumber=23,numbers=left,numbersep=3pt,frame=lines,framesep=2pt,rulecolor=\color{gray!60},breaklines=true,fontsize=\fontsize{6.6}{7.6}\selectfont]{../../splat_editor/core/physics.py}
\subsubsection{\texorpdfstring{\code{class CollisionWorld}}{CollisionWorld}}
{\footnotesize\textit{Dòng 34--112}}\par
\textbf{Mô tả trong mã nguồn.} Tập các voxel đặc, dựng một lần từ vị trí gaussian (toạ độ thế giới).

Các phương thức: \code{__init__}, \code{ready}, \code{build}, \code{_highest}, \code{column_top}, \code{segment_blocked}, \code{blocked}.

\VerbatimInput[firstline=34,lastline=35,firstnumber=34,numbers=left,numbersep=3pt,frame=lines,framesep=2pt,rulecolor=\color{gray!60},breaklines=true,fontsize=\fontsize{6.6}{7.6}\selectfont]{../../splat_editor/core/physics.py}
\subsubsection{\texorpdfstring{\code{CollisionWorld.__init__(self)}}{CollisionWorld.__init__}}
{\footnotesize\textit{Dòng 37--44}}\par
\VerbatimInput[firstline=37,lastline=44,firstnumber=37,numbers=left,numbersep=3pt,frame=lines,framesep=2pt,rulecolor=\color{gray!60},breaklines=true,fontsize=\fontsize{6.6}{7.6}\selectfont]{../../splat_editor/core/physics.py}
\subsubsection{\texorpdfstring{\code{CollisionWorld.ready(self)}}{CollisionWorld.ready}}
{\footnotesize\textit{Dòng 46--48}}\par
\VerbatimInput[firstline=46,lastline=48,firstnumber=46,numbers=left,numbersep=3pt,frame=lines,framesep=2pt,rulecolor=\color{gray!60},breaklines=true,fontsize=\fontsize{6.6}{7.6}\selectfont]{../../splat_editor/core/physics.py}
\subsubsection{\texorpdfstring{\code{CollisionWorld.build(self, positions, voxel, min_count=1)}}{CollisionWorld.build}}
{\footnotesize\textit{Dòng 50--69}}\par
\textbf{Mô tả trong mã nguồn.} positions (N,3) float; voxel: cạnh ô; min\_count: số gaussian tối thiểu để ô là đặc.

\textbf{Giải thích.} Voxel hóa: \code{cell = floor((p-origin)/v)}, khóa 63 bit $(c_x\ll42)|(c_z\ll21)|c_y$, đếm gaussian mỗi ô và giữ những ô có ít nhất \code{min_count}. Mảng khóa được sắp để dùng \code{searchsorted}.

\textbf{Công thức.}
\[
v=\frac{H}{8},\quad \mathrm{cell}=\Big\lfloor\frac{p-\mathrm{origin}}{v}\Big\rfloor,\quad
\mathrm{key}=(c_x\ll42)\ \mid\ (c_z\ll21)\ \mid\ c_y
\]
Ô đặc nếu có $\ge\mathrm{min\_count}$ gaussian rắn.


\VerbatimInput[firstline=50,lastline=69,firstnumber=50,numbers=left,numbersep=3pt,frame=lines,framesep=2pt,rulecolor=\color{gray!60},breaklines=true,fontsize=\fontsize{6.6}{7.6}\selectfont]{../../splat_editor/core/physics.py}
\subsubsection{\texorpdfstring{\code{CollisionWorld._highest(self, x, z, y)}}{CollisionWorld._highest}}
{\footnotesize\textit{Dòng 71--82}}\par
\textbf{Mô tả trong mã nguồn.} Chỉ số iy của ô đặc cao nhất tại cột (x, z) mà iy \textless{}= ô chứa độ cao y; -1 nếu không có.

\textbf{Giải thích.} Tìm ô đặc cao nhất trong cột không vượt quá độ cao cho trước, bằng một lần \code{searchsorted} và kiểm tra cùng cột (\code{found>>21 == query>>21}). $O(\log N)$, vector hóa.

\textbf{Công thức.}
\[
\mathrm{idx}=\mathrm{searchsorted}(\mathrm{keys},\,q,\,\text{`right'})-1,\qquad
\text{hợp lệ nếu }\ \mathrm{keys}[\mathrm{idx}]\gg21=q\gg21,\qquad y_{top}=\mathrm{origin}_y+(i_y+1)\,v
\]


\VerbatimInput[firstline=71,lastline=82,firstnumber=71,numbers=left,numbersep=3pt,frame=lines,framesep=2pt,rulecolor=\color{gray!60},breaklines=true,fontsize=\fontsize{6.6}{7.6}\selectfont]{../../splat_editor/core/physics.py}
\subsubsection{\texorpdfstring{\code{CollisionWorld.column_top(self, x, z, y_max)}}{CollisionWorld.column_top}}
{\footnotesize\textit{Dòng 84--87}}\par
\textbf{Mô tả trong mã nguồn.} Độ cao mặt trên của ô đặc cao nhất không vượt quá y\_max tại mỗi cột (-inf: không có).

\VerbatimInput[firstline=84,lastline=87,firstnumber=84,numbers=left,numbersep=3pt,frame=lines,framesep=2pt,rulecolor=\color{gray!60},breaklines=true,fontsize=\fontsize{6.6}{7.6}\selectfont]{../../splat_editor/core/physics.py}
\subsubsection{\texorpdfstring{\code{CollisionWorld.segment_blocked(self, start, end, skip=1.0, skip_end=0.0)}}{CollisionWorld.segment_blocked}}
{\footnotesize\textit{Dòng 89--106}}\par
\textbf{Mô tả trong mã nguồn.} Đoạn thẳng start→end có đi xuyên qua ô đặc nào không (bỏ qua `skip` ô đầu và `skip\_end` ô cuối: đầu = ngay tại mắt, cuối = vật nằm sát / trong mặt đất như khối đèn đang nằm trên sàn).

\textbf{Giải thích.} Tia kiểm tra tầm nhìn: lấy mẫu dọc đoạn với bước $v/2$ rồi tra voxel. Dùng cho camera ngôi thứ ba (lùi khi sau lưng có tường) và chọn đèn.

\textbf{Công thức.}
\[
n=\Big\lfloor\frac{L}{v/2}\Big\rfloor+1,\qquad p_k=\mathrm{start}+\frac{k}{n-1}(\mathrm{end}-\mathrm{start}),\quad k=0,\dots,n-1
\]


\VerbatimInput[firstline=89,lastline=106,firstnumber=89,numbers=left,numbersep=3pt,frame=lines,framesep=2pt,rulecolor=\color{gray!60},breaklines=true,fontsize=\fontsize{6.6}{7.6}\selectfont]{../../splat_editor/core/physics.py}
\subsubsection{\texorpdfstring{\code{CollisionWorld.blocked(self, x, z, y0, y1)}}{CollisionWorld.blocked}}
{\footnotesize\textit{Dòng 108--112}}\par
\textbf{Mô tả trong mã nguồn.} Có ô đặc nào nằm trong khoảng độ cao [y0, y1] tại mỗi cột không.

\VerbatimInput[firstline=108,lastline=112,firstnumber=108,numbers=left,numbersep=3pt,frame=lines,framesep=2pt,rulecolor=\color{gray!60},breaklines=true,fontsize=\fontsize{6.6}{7.6}\selectfont]{../../splat_editor/core/physics.py}
\subsubsection{\texorpdfstring{\code{class Player}}{Player}}
{\footnotesize\textit{Dòng 115--141}}\par
\textbf{Mô tả trong mã nguồn.} Người chơi: một cột đứng có bán kính, chân ở `feet`. Mắt ở trên chân.

Các phương thức: \code{__init__}, \code{place}, \code{eye}.

\VerbatimInput[firstline=115,lastline=116,firstnumber=115,numbers=left,numbersep=3pt,frame=lines,framesep=2pt,rulecolor=\color{gray!60},breaklines=true,fontsize=\fontsize{6.6}{7.6}\selectfont]{../../splat_editor/core/physics.py}
\subsubsection{\texorpdfstring{\code{Player.__init__(self)}}{Player.__init__}}
{\footnotesize\textit{Dòng 118--124}}\par
\VerbatimInput[firstline=118,lastline=124,firstnumber=118,numbers=left,numbersep=3pt,frame=lines,framesep=2pt,rulecolor=\color{gray!60},breaklines=true,fontsize=\fontsize{6.6}{7.6}\selectfont]{../../splat_editor/core/physics.py}
\subsubsection{\texorpdfstring{\code{Player.place(self, eye, eye_height)}}{Player.place}}
{\footnotesize\textit{Dòng 126--132}}\par
\VerbatimInput[firstline=126,lastline=132,firstnumber=126,numbers=left,numbersep=3pt,frame=lines,framesep=2pt,rulecolor=\color{gray!60},breaklines=true,fontsize=\fontsize{6.6}{7.6}\selectfont]{../../splat_editor/core/physics.py}
\subsubsection{\texorpdfstring{\code{Player.eye(self, eye_height, dt=None)}}{Player.eye}}
{\footnotesize\textit{Dòng 134--141}}\par
\textbf{Mô tả trong mã nguồn.} Vị trí mắt; dt != None: làm mượt độ cao khi đang đứng trên mặt đất.

\VerbatimInput[firstline=134,lastline=141,firstnumber=134,numbers=left,numbersep=3pt,frame=lines,framesep=2pt,rulecolor=\color{gray!60},breaklines=true,fontsize=\fontsize{6.6}{7.6}\selectfont]{../../splat_editor/core/physics.py}
\subsubsection{\texorpdfstring{\code{class Physics}}{Physics}}
{\footnotesize\textit{Dòng 144--320}}\par
Các phương thức: \code{__init__}, \code{eye_height}, \code{gravity}, \code{walk_speed}, \code{set_scale}, \code{set_height}, \code{invalidate}, \code{building}, \code{ready}, \code{has_world}, \code{ensure_world}, \code{drop_height}, \code{visibility}, \code{walk}.

\VerbatimInput[firstline=144,lastline=147,firstnumber=144,numbers=left,numbersep=3pt,frame=lines,framesep=2pt,rulecolor=\color{gray!60},breaklines=true,fontsize=\fontsize{6.6}{7.6}\selectfont]{../../splat_editor/core/physics.py}
\subsubsection{\texorpdfstring{\code{Physics.__init__(self)}}{Physics.__init__}}
{\footnotesize\textit{Dòng 149--162}}\par
\VerbatimInput[firstline=149,lastline=162,firstnumber=149,numbers=left,numbersep=3pt,frame=lines,framesep=2pt,rulecolor=\color{gray!60},breaklines=true,fontsize=\fontsize{6.6}{7.6}\selectfont]{../../splat_editor/core/physics.py}
\subsubsection{\texorpdfstring{\code{Physics.eye_height(self)}}{Physics.eye_height}}
{\footnotesize\textit{Dòng 165--167}}\par
\VerbatimInput[firstline=165,lastline=167,firstnumber=165,numbers=left,numbersep=3pt,frame=lines,framesep=2pt,rulecolor=\color{gray!60},breaklines=true,fontsize=\fontsize{6.6}{7.6}\selectfont]{../../splat_editor/core/physics.py}
\subsubsection{\texorpdfstring{\code{Physics.gravity(self)}}{Physics.gravity}}
{\footnotesize\textit{Dòng 169--171}}\par
\VerbatimInput[firstline=169,lastline=171,firstnumber=169,numbers=left,numbersep=3pt,frame=lines,framesep=2pt,rulecolor=\color{gray!60},breaklines=true,fontsize=\fontsize{6.6}{7.6}\selectfont]{../../splat_editor/core/physics.py}
\subsubsection{\texorpdfstring{\code{Physics.walk_speed(self)}}{Physics.walk_speed}}
{\footnotesize\textit{Dòng 173--175}}\par
\VerbatimInput[firstline=173,lastline=175,firstnumber=173,numbers=left,numbersep=3pt,frame=lines,framesep=2pt,rulecolor=\color{gray!60},breaklines=true,fontsize=\fontsize{6.6}{7.6}\selectfont]{../../splat_editor/core/physics.py}
\subsubsection{\texorpdfstring{\code{Physics.set_scale(self, radius)}}{Physics.set_scale}}
{\footnotesize\textit{Dòng 177--179}}\par
\textbf{Mô tả trong mã nguồn.} Giữ chiều cao mặc định cố định (2.74); chỉnh lại được trong giao diện.

\VerbatimInput[firstline=177,lastline=179,firstnumber=177,numbers=left,numbersep=3pt,frame=lines,framesep=2pt,rulecolor=\color{gray!60},breaklines=true,fontsize=\fontsize{6.6}{7.6}\selectfont]{../../splat_editor/core/physics.py}
\subsubsection{\texorpdfstring{\code{Physics.set_height(self, height)}}{Physics.set_height}}
{\footnotesize\textit{Dòng 181--184}}\par
\VerbatimInput[firstline=181,lastline=184,firstnumber=181,numbers=left,numbersep=3pt,frame=lines,framesep=2pt,rulecolor=\color{gray!60},breaklines=true,fontsize=\fontsize{6.6}{7.6}\selectfont]{../../splat_editor/core/physics.py}
\subsubsection{\texorpdfstring{\code{Physics.invalidate(self)}}{Physics.invalidate}}
{\footnotesize\textit{Dòng 187--190}}\par
\textbf{Mô tả trong mã nguồn.} Cảnh đổi (mở / xoá / kéo lớp): dựng lại bản đồ va chạm khi cần (chờ yên 0.4 giây).

\VerbatimInput[firstline=187,lastline=190,firstnumber=187,numbers=left,numbersep=3pt,frame=lines,framesep=2pt,rulecolor=\color{gray!60},breaklines=true,fontsize=\fontsize{6.6}{7.6}\selectfont]{../../splat_editor/core/physics.py}
\subsubsection{\texorpdfstring{\code{Physics.building(self)}}{Physics.building}}
{\footnotesize\textit{Dòng 192--194}}\par
\VerbatimInput[firstline=192,lastline=194,firstnumber=192,numbers=left,numbersep=3pt,frame=lines,framesep=2pt,rulecolor=\color{gray!60},breaklines=true,fontsize=\fontsize{6.6}{7.6}\selectfont]{../../splat_editor/core/physics.py}
\subsubsection{\texorpdfstring{\code{Physics.ready(self)}}{Physics.ready}}
{\footnotesize\textit{Dòng 196--198}}\par
\VerbatimInput[firstline=196,lastline=198,firstnumber=196,numbers=left,numbersep=3pt,frame=lines,framesep=2pt,rulecolor=\color{gray!60},breaklines=true,fontsize=\fontsize{6.6}{7.6}\selectfont]{../../splat_editor/core/physics.py}
\subsubsection{\texorpdfstring{\code{Physics.has_world(self)}}{Physics.has_world}}
{\footnotesize\textit{Dòng 200--202}}\par
\VerbatimInput[firstline=200,lastline=202,firstnumber=200,numbers=left,numbersep=3pt,frame=lines,framesep=2pt,rulecolor=\color{gray!60},breaklines=true,fontsize=\fontsize{6.6}{7.6}\selectfont]{../../splat_editor/core/physics.py}
\subsubsection{\texorpdfstring{\code{Physics.ensure_world(self, scene, needed)}}{Physics.ensure_world}}
{\footnotesize\textit{Dòng 204--228}}\par
\textbf{Mô tả trong mã nguồn.} Gọi mỗi khung hình. Dựng bản đồ ở luồng nền nếu cần mà đang cũ.

\textbf{Giải thích.} Mỗi khung hình: nếu bản đồ va chạm cũ và cảnh đã yên 0.4 giây thì dựng lại \textbf{ở luồng nền}, tránh giật hình.

\VerbatimInput[firstline=204,lastline=228,firstnumber=204,numbers=left,numbersep=3pt,frame=lines,framesep=2pt,rulecolor=\color{gray!60},breaklines=true,fontsize=\fontsize{6.6}{7.6}\selectfont]{../../splat_editor/core/physics.py}
\subsubsection{\texorpdfstring{\code{Physics.drop_height(self, x, z, y_from, radius=0.0)}}{Physics.drop_height}}
{\footnotesize\textit{Dòng 231--242}}\par
\textbf{Mô tả trong mã nguồn.} Độ cao mặt đất dưới (x, z) tính từ y\_from trở xuống; None nếu không có.

\VerbatimInput[firstline=231,lastline=242,firstnumber=231,numbers=left,numbersep=3pt,frame=lines,framesep=2pt,rulecolor=\color{gray!60},breaklines=true,fontsize=\fontsize{6.6}{7.6}\selectfont]{../../splat_editor/core/physics.py}
\subsubsection{\texorpdfstring{\code{Physics.visibility(self, eye, targets, skip_end=0.0)}}{Physics.visibility}}
{\footnotesize\textit{Dòng 244--249}}\par
\textbf{Mô tả trong mã nguồn.} Tỉ lệ các điểm `targets` nhìn thấy được từ `eye` (không bị cảnh che); 1.0 nếu chưa có bản đồ.

\VerbatimInput[firstline=244,lastline=249,firstnumber=244,numbers=left,numbersep=3pt,frame=lines,framesep=2pt,rulecolor=\color{gray!60},breaklines=true,fontsize=\fontsize{6.6}{7.6}\selectfont]{../../splat_editor/core/physics.py}
\subsubsection{\texorpdfstring{\code{Physics.walk(self, forward, right, yaw, jump, sprint, sneak, dt)}}{Physics.walk}}
{\footnotesize\textit{Dòng 252--320}}\par
\textbf{Mô tả trong mã nguồn.} Đi bộ: W/S/A/D theo hướng nhìn `yaw`, nhảy, chạy nhanh, đi chậm. Trả về False nếu chưa có bản đồ va chạm (người chơi đứng yên)

\textbf{Giải thích.} Đi bộ kiểu Minecraft: hướng theo \code{yaw}, trượt tường (thử $x{+}z$, rồi chỉ $x$, rồi chỉ $z$), va chạm cột trụ bằng 9 điểm vòng, nhảy $v_y=\sqrt{2gh}$, trọng lực Euler, leo bậc.

\textbf{Công thức.}
\[
g=12H,\quad v_{jump}=\sqrt{2g\,h_{jump}},\quad h_{jump}=0.75H
\]
\[
v_y\leftarrow v_y-g\Delta t,\quad y\leftarrow y+v_y\Delta t,\quad \mathrm{eye}_y\mathrel{+}=(\mathrm{target}-\mathrm{eye}_y)\min(18\Delta t,1)
\]
Va chạm cột trụ: 9 điểm $(0,0)$ và $r\,(\cos\tfrac{k\pi}4,\sin\tfrac{k\pi}4)$, $k=0..7$.


\VerbatimInput[firstline=252,lastline=320,firstnumber=252,numbers=left,numbersep=3pt,frame=lines,framesep=2pt,rulecolor=\color{gray!60},breaklines=true,fontsize=\fontsize{6.6}{7.6}\selectfont]{../../splat_editor/core/physics.py}
\subsection{\texorpdfstring{\code{splat_editor/core/avatar_pose.py}}{splat_editor/core/avatar_pose.py}}
{\footnotesize 172 dòng}\par
\textbf{Vai trò.} Tính \textbf{góc các khớp} của nhân vật (nhịp đi, cúi, ngồi, cử chỉ) và dựng ma trận model cho 6 khối cơ thể. Không đụng OpenGL; việc vẽ nằm ở
\code{render/avatar_renderer.py}. Toán học: xem mục \textit{Cơ sở toán} ngay bên dưới (\S\ref{s:anim}).

\paragraph{Mô tả đầu file và cách dùng.} Phần chú thích đầu file do tác giả viết, nêu ý tưởng và ví dụ sử dụng:
\VerbatimInput[firstline=1,lastline=18,firstnumber=1,numbers=left,numbersep=3pt,frame=lines,framesep=2pt,rulecolor=\color{gray!60},breaklines=true,fontsize=\fontsize{6.2}{7.2}\selectfont]{../../splat_editor/core/avatar_pose.py}
%% ---- merged math: s:anim
\subsubsection{Cơ sở toán: Hoạt họa nhân vật}\label{s:anim}
Các công thức dưới đây được cài đặt trong \code{core/avatar_pose.py}; việc vẽ 6 khối bằng các ma trận này nằm ở \code{render/avatar_renderer.py}.
\begin{itemize}
\item Bộ lọc tiến về mục tiêu: $\mathrm{approach}(v,t,\mathrm{rate},dt)=v+(t-v)\min(1,dt\cdot\mathrm{rate})$ (cùng dạng quán tính camera bay).
\item \textbf{Nhịp đi:} $\mathrm{walked}=\mathrm{distance}/\mathrm{height}$; $\mathrm{phase}\mathrel{+}=\mathrm{walked}\cdot\mathrm{STRIDE}\cdot1.75$ (pha tỉ lệ \textbf{quãng đường}, không phải thời gian, nên chân không trượt); $\mathrm{swing}=\sin(\mathrm{phase})\cdot60^\circ\cdot\mathrm{amount}$; chân trái $+\mathrm{swing}$, chân phải $-\mathrm{swing}$, tay ngược pha (đối xứng chéo).
\item Biên độ mục tiêu $=\min(1,\mathrm{speed}/(4.3/1.75))$.
\item \textbf{Thở:} $\mathrm{idle}=2.2\sin(1.13t)+0.8\sin(1.91t)$ (hai tần số không công ước nên không lặp cứng).
\item \textbf{Cử chỉ:} zombie $-85^\circ+3\sin(2.4t)$; cheer $\pm(158+6\sin7t)$; wave $142+20\sin9t$; pha trộn tuyến tính \code{_mix_pair}.
\item Ngồi: chân $-90^\circ$, hạ thân 10 đơn vị skin; $\mathrm{stance}=1-0.48\,\mathrm{sit}-0.12\,\mathrm{crouch}$.
\end{itemize}

\paragraph{Cây phân cấp (\code{limb_matrices}).}
\begin{align*}
\mathrm{root}&=T(\mathrm{feet})\,R_y(\mathrm{facing}+180^\circ)\,S(\mathrm{scale})\,T(0,-\mathrm{drop},0)\\
\mathrm{upper}&=\mathrm{root}\cdot T(0,12,0)\,R_x(\mathrm{lean})\,T(0,-12,0)\quad(\text{xoay quanh hông})\\
\mathrm{leg}&=\mathrm{root}\cdot T(0,12,0)\,R_x(a)\,T(0,-12,0)\\
\mathrm{arm}&=\mathrm{upper}\cdot T(p_x,24,0)\,R_z(a_1)\,R_x(a_0)\,T(-p_x,-24,0)\quad(\text{quanh vai, }p_x=\pm6)\\
\mathrm{head}&=\mathrm{upper}\cdot T(0,24,0)\,R_y(\mathrm{yaw})\,R_x(\mathrm{pitch})\,T(0,-24,0)\quad(\text{quanh cổ})
\end{align*}
Mẫu chung: $M_{\mathrm{con}}=M_{\mathrm{cha}}\cdot T(p)\,R\,T(-p)$ (xoay quanh điểm $p$: dời $p$ về gốc, xoay, dời lại). $\mathrm{scale}=\mathrm{chiều\ cao}/32$. Ma trận quay trục \code{_rot(axis, deg)} điền $c,-s,s,c$ vào cặp chỉ số $(1,2)$ cho $X$, $(2,0)$ cho $Y$, $(0,1)$ cho $Z$.

\subsubsection*{Khai báo mức module (dòng 19--31)}
\VerbatimInput[firstline=19,lastline=31,firstnumber=19,numbers=left,numbersep=3pt,frame=lines,framesep=2pt,rulecolor=\color{gray!60},breaklines=true,fontsize=\fontsize{6.6}{7.6}\selectfont]{../../splat_editor/core/avatar_pose.py}
\subsubsection{\texorpdfstring{\code{_approach(value, target, rate, dt)}}{_approach}}
{\footnotesize\textit{Dòng 34--35}}\par
\VerbatimInput[firstline=34,lastline=35,firstnumber=34,numbers=left,numbersep=3pt,frame=lines,framesep=2pt,rulecolor=\color{gray!60},breaklines=true,fontsize=\fontsize{6.6}{7.6}\selectfont]{../../splat_editor/core/avatar_pose.py}
\subsubsection{\texorpdfstring{\code{class AvatarPose}}{AvatarPose}}
{\footnotesize\textit{Dòng 38--124}}\par
Các phương thức: \code{__init__}, \code{toggle_sit}, \code{stand}, \code{toggle_gesture}, \code{update}, \code{stance}, \code{joints}.

\VerbatimInput[firstline=38,lastline=38,firstnumber=38,numbers=left,numbersep=3pt,frame=lines,framesep=2pt,rulecolor=\color{gray!60},breaklines=true,fontsize=\fontsize{6.6}{7.6}\selectfont]{../../splat_editor/core/avatar_pose.py}
\subsubsection{\texorpdfstring{\code{AvatarPose.__init__(self)}}{AvatarPose.__init__}}
{\footnotesize\textit{Dòng 39--47}}\par
\VerbatimInput[firstline=39,lastline=47,firstnumber=39,numbers=left,numbersep=3pt,frame=lines,framesep=2pt,rulecolor=\color{gray!60},breaklines=true,fontsize=\fontsize{6.6}{7.6}\selectfont]{../../splat_editor/core/avatar_pose.py}
\subsubsection{\texorpdfstring{\code{AvatarPose.toggle_sit(self)}}{AvatarPose.toggle_sit}}
{\footnotesize\textit{Dòng 50--53}}\par
\VerbatimInput[firstline=50,lastline=53,firstnumber=50,numbers=left,numbersep=3pt,frame=lines,framesep=2pt,rulecolor=\color{gray!60},breaklines=true,fontsize=\fontsize{6.6}{7.6}\selectfont]{../../splat_editor/core/avatar_pose.py}
\subsubsection{\texorpdfstring{\code{AvatarPose.stand(self)}}{AvatarPose.stand}}
{\footnotesize\textit{Dòng 55--56}}\par
\VerbatimInput[firstline=55,lastline=56,firstnumber=55,numbers=left,numbersep=3pt,frame=lines,framesep=2pt,rulecolor=\color{gray!60},breaklines=true,fontsize=\fontsize{6.6}{7.6}\selectfont]{../../splat_editor/core/avatar_pose.py}
\subsubsection{\texorpdfstring{\code{AvatarPose.toggle_gesture(self, name)}}{AvatarPose.toggle_gesture}}
{\footnotesize\textit{Dòng 58--61}}\par
\VerbatimInput[firstline=58,lastline=61,firstnumber=58,numbers=left,numbersep=3pt,frame=lines,framesep=2pt,rulecolor=\color{gray!60},breaklines=true,fontsize=\fontsize{6.6}{7.6}\selectfont]{../../splat_editor/core/avatar_pose.py}
\subsubsection{\texorpdfstring{\code{AvatarPose.update(self, dt, distance, crouching, height=1.0)}}{AvatarPose.update}}
{\footnotesize\textit{Dòng 64--79}}\par
\textbf{Mô tả trong mã nguồn.} distance: quãng đường ngang đã đi (đơn vị cảnh); height: chiều cao người (cùng đơn vị) để nhịp bước và biên độ không phụ thuộc kích thước cảnh.

\textbf{Giải thích.} Pha đi bộ tỉ lệ với \textbf{quãng đường} đi được (không phải thời gian) nên chân không trượt; biên độ tiến dần về mục tiêu bằng bộ lọc mũ \code{_approach}.

\textbf{Công thức.}
\[
\mathrm{approach}(v,t,\mathrm{rate},dt)=v+(t-v)\min(1,\,dt\cdot\mathrm{rate}),\qquad
\mathrm{phase}\mathrel{+}=\frac{\mathrm{dist}}{H}\cdot\mathrm{STRIDE}\cdot1.75
\]
\[
\mathrm{swing}=60^\circ\sin(\mathrm{phase})\cdot\mathrm{amount},\qquad
\mathrm{idle}=2.2\sin(1.13t)+0.8\sin(1.91t)
\]


\VerbatimInput[firstline=64,lastline=79,firstnumber=64,numbers=left,numbersep=3pt,frame=lines,framesep=2pt,rulecolor=\color{gray!60},breaklines=true,fontsize=\fontsize{6.6}{7.6}\selectfont]{../../splat_editor/core/avatar_pose.py}
\subsubsection{\texorpdfstring{\code{AvatarPose.stance(self)}}{AvatarPose.stance}}
{\footnotesize\textit{Dòng 81--84}}\par
\textbf{Mô tả trong mã nguồn.} Hệ số chiều cao mắt so với lúc đứng (ngồi / cúi thì thấp hơn).

\VerbatimInput[firstline=81,lastline=84,firstnumber=81,numbers=left,numbersep=3pt,frame=lines,framesep=2pt,rulecolor=\color{gray!60},breaklines=true,fontsize=\fontsize{6.6}{7.6}\selectfont]{../../splat_editor/core/avatar_pose.py}
\subsubsection{\texorpdfstring{\code{AvatarPose.joints(self)}}{AvatarPose.joints}}
{\footnotesize\textit{Dòng 87--124}}\par
\textbf{Mô tả trong mã nguồn.} Góc (độ) của từng khớp. Tay: (quay quanh X, quay quanh Z).

\textbf{Giải thích.} Trả về góc (độ) từng khớp: swing chân/tay, thở (hai tần số không công ước), cử chỉ (zombie, cheer, wave), ngồi và cúi.

\VerbatimInput[firstline=87,lastline=124,firstnumber=87,numbers=left,numbersep=3pt,frame=lines,framesep=2pt,rulecolor=\color{gray!60},breaklines=true,fontsize=\fontsize{6.6}{7.6}\selectfont]{../../splat_editor/core/avatar_pose.py}
\subsubsection{\texorpdfstring{\code{_mix_pair(a, b, k)}}{_mix_pair}}
{\footnotesize\textit{Dòng 127--128}}\par
\VerbatimInput[firstline=127,lastline=128,firstnumber=127,numbers=left,numbersep=3pt,frame=lines,framesep=2pt,rulecolor=\color{gray!60},breaklines=true,fontsize=\fontsize{6.6}{7.6}\selectfont]{../../splat_editor/core/avatar_pose.py}
\subsubsection*{Khai báo mức module (dòng 132--135)}
\VerbatimInput[firstline=132,lastline=135,firstnumber=132,numbers=left,numbersep=3pt,frame=lines,framesep=2pt,rulecolor=\color{gray!60},breaklines=true,fontsize=\fontsize{6.6}{7.6}\selectfont]{../../splat_editor/core/avatar_pose.py}
\subsubsection{\texorpdfstring{\code{_translate(x, y, z)}}{_translate}}
{\footnotesize\textit{Dòng 138--141}}\par
\VerbatimInput[firstline=138,lastline=141,firstnumber=138,numbers=left,numbersep=3pt,frame=lines,framesep=2pt,rulecolor=\color{gray!60},breaklines=true,fontsize=\fontsize{6.6}{7.6}\selectfont]{../../splat_editor/core/avatar_pose.py}
\subsubsection{\texorpdfstring{\code{_rot(axis, degrees)}}{_rot}}
{\footnotesize\textit{Dòng 144--149}}\par
\VerbatimInput[firstline=144,lastline=149,firstnumber=144,numbers=left,numbersep=3pt,frame=lines,framesep=2pt,rulecolor=\color{gray!60},breaklines=true,fontsize=\fontsize{6.6}{7.6}\selectfont]{../../splat_editor/core/avatar_pose.py}
\subsubsection{\texorpdfstring{\code{limb_matrices(feet, facing, head_yaw, head_pitch, joints, scale)}}{limb_matrices}}
{\footnotesize\textit{Dòng 152--172}}\par
\textbf{Mô tả trong mã nguồn.} Ma trận model (4x4) của 6 khối theo thứ tự LIMBS.

feet: vị trí bàn chân; facing: yaw camera (rad) của thân; head\_yaw: lệch của đầu so với thân (rad); head\_pitch: pitch camera (rad, dương = nhìn xuống); scale: kích thước một điểm ảnh skin (= chiều cao / 32).

\textbf{Giải thích.} Dựng ma trận model 4$\times$4 cho 6 khối theo cây phân cấp $M_{con}=M_{cha}\,T(p)\,R\,T(-p)$ (xoay quanh khớp $p$). Xem \S\ref{s:anim}.

\textbf{Công thức.}
\[
M_{con}=M_{cha}\,T(p)\,R\,T(-p)
\]
\[
\mathrm{root}=T(\mathrm{feet})R_y(\mathrm{facing}+180^\circ)S\,T(0,-\mathrm{drop},0),\quad
\mathrm{arm}=\mathrm{upper}\,T(\pm6,24,0)\,R_z(a_1)R_x(a_0)\,T(\mp6,-24,0)
\]


\VerbatimInput[firstline=152,lastline=172,firstnumber=152,numbers=left,numbersep=3pt,frame=lines,framesep=2pt,rulecolor=\color{gray!60},breaklines=true,fontsize=\fontsize{6.6}{7.6}\selectfont]{../../splat_editor/core/avatar_pose.py}
\subsection{\texorpdfstring{\code{splat_editor/core/history.py}}{splat_editor/core/history.py}}
{\footnotesize 83 dòng}\par
\textbf{Vai trò.} Lịch sử hoàn tác theo mẫu \textbf{Command}: mỗi thao tác đã thực hiện là một đối tượng biết tự hoàn tác. \code{History} giữ tối đa 20 bước. Vì xóa là xóa mềm
nên mỗi lệnh chỉ cần giữ danh sách chỉ số đã xóa.

\paragraph{Mô tả đầu file và cách dùng.} Phần chú thích đầu file do tác giả viết, nêu ý tưởng và ví dụ sử dụng:
\VerbatimInput[firstline=1,lastline=17,firstnumber=1,numbers=left,numbersep=3pt,frame=lines,framesep=2pt,rulecolor=\color{gray!60},breaklines=true,fontsize=\fontsize{6.2}{7.2}\selectfont]{../../splat_editor/core/history.py}
%% ---- merged math: undo
\subsubsection{Cơ sở toán: Undo}\label{s:math-undo}
Mẫu Command, tối đa 20 bước. Xóa mềm bằng bit \code{deleted}; hoàn tác là lật lại các chỉ số ($O(k)$). Trạng thái GPU 1 byte/gaussian: bit 0 = chọn, bit 1 = xóa (\code{STATE_*}). Màu chọn: $\mathrm{mix}(\mathrm{color},\mathrm{selectColor},0.7\cdot\mathrm{selected})$.

\subsubsection*{Khai báo mức module (dòng 18--18)}
\VerbatimInput[firstline=18,lastline=18,firstnumber=18,numbers=left,numbersep=3pt,frame=lines,framesep=2pt,rulecolor=\color{gray!60},breaklines=true,fontsize=\fontsize{6.6}{7.6}\selectfont]{../../splat_editor/core/history.py}
\subsubsection{\texorpdfstring{\code{class Command(ABC)}}{Command}}
{\footnotesize\textit{Dòng 21--34}}\par
\textbf{Mô tả trong mã nguồn.} Một thao tác đã thực hiện, có thể hoàn tác.

\textbf{Giải thích.} Lớp trừu tượng: mỗi thao tác phải có \code{description} và \code{undo()}. Muốn thêm thao tác hoàn tác được chỉ cần kế thừa lớp này.

Các phương thức: \code{description}, \code{undo}, \code{release}.

\VerbatimInput[firstline=21,lastline=22,firstnumber=21,numbers=left,numbersep=3pt,frame=lines,framesep=2pt,rulecolor=\color{gray!60},breaklines=true,fontsize=\fontsize{6.6}{7.6}\selectfont]{../../splat_editor/core/history.py}
\subsubsection{\texorpdfstring{\code{Command.description(self)}}{Command.description}}
{\footnotesize\textit{Dòng 24--27}}\par
\VerbatimInput[firstline=24,lastline=27,firstnumber=24,numbers=left,numbersep=3pt,frame=lines,framesep=2pt,rulecolor=\color{gray!60},breaklines=true,fontsize=\fontsize{6.6}{7.6}\selectfont]{../../splat_editor/core/history.py}
\subsubsection{\texorpdfstring{\code{Command.undo(self)}}{Command.undo}}
{\footnotesize\textit{Dòng 29--31}}\par
\VerbatimInput[firstline=29,lastline=31,firstnumber=29,numbers=left,numbersep=3pt,frame=lines,framesep=2pt,rulecolor=\color{gray!60},breaklines=true,fontsize=\fontsize{6.6}{7.6}\selectfont]{../../splat_editor/core/history.py}
\subsubsection{\texorpdfstring{\code{Command.release(self, load)}}{Command.release}}
{\footnotesize\textit{Dòng 33--34}}\par
\textbf{Mô tả trong mã nguồn.} Thay mọi mảng dữ liệu `data` đang giữ bằng load(data) (xem Editor.\_release\_file).

\VerbatimInput[firstline=33,lastline=34,firstnumber=33,numbers=left,numbersep=3pt,frame=lines,framesep=2pt,rulecolor=\color{gray!60},breaklines=true,fontsize=\fontsize{6.6}{7.6}\selectfont]{../../splat_editor/core/history.py}
\subsubsection{\texorpdfstring{\code{class DeleteGaussiansCommand(Command)}}{DeleteGaussiansCommand}}
{\footnotesize\textit{Dòng 37--52}}\par
Các phương thức: \code{__init__}, \code{count}, \code{description}, \code{undo}.

\VerbatimInput[firstline=37,lastline=37,firstnumber=37,numbers=left,numbersep=3pt,frame=lines,framesep=2pt,rulecolor=\color{gray!60},breaklines=true,fontsize=\fontsize{6.6}{7.6}\selectfont]{../../splat_editor/core/history.py}
\subsubsection{\texorpdfstring{\code{DeleteGaussiansCommand.__init__(self, snapshots)}}{DeleteGaussiansCommand.__init__}}
{\footnotesize\textit{Dòng 38--39}}\par
\VerbatimInput[firstline=38,lastline=39,firstnumber=38,numbers=left,numbersep=3pt,frame=lines,framesep=2pt,rulecolor=\color{gray!60},breaklines=true,fontsize=\fontsize{6.6}{7.6}\selectfont]{../../splat_editor/core/history.py}
\subsubsection{\texorpdfstring{\code{DeleteGaussiansCommand.count(self)}}{DeleteGaussiansCommand.count}}
{\footnotesize\textit{Dòng 41--44}}\par
\textbf{Mô tả trong mã nguồn.} Số gaussian đã xoá.

\VerbatimInput[firstline=41,lastline=44,firstnumber=41,numbers=left,numbersep=3pt,frame=lines,framesep=2pt,rulecolor=\color{gray!60},breaklines=true,fontsize=\fontsize{6.6}{7.6}\selectfont]{../../splat_editor/core/history.py}
\subsubsection{\texorpdfstring{\code{DeleteGaussiansCommand.description(self)}}{DeleteGaussiansCommand.description}}
{\footnotesize\textit{Dòng 46--48}}\par
\VerbatimInput[firstline=46,lastline=48,firstnumber=46,numbers=left,numbersep=3pt,frame=lines,framesep=2pt,rulecolor=\color{gray!60},breaklines=true,fontsize=\fontsize{6.6}{7.6}\selectfont]{../../splat_editor/core/history.py}
\subsubsection{\texorpdfstring{\code{DeleteGaussiansCommand.undo(self)}}{DeleteGaussiansCommand.undo}}
{\footnotesize\textit{Dòng 50--52}}\par
\VerbatimInput[firstline=50,lastline=52,firstnumber=50,numbers=left,numbersep=3pt,frame=lines,framesep=2pt,rulecolor=\color{gray!60},breaklines=true,fontsize=\fontsize{6.6}{7.6}\selectfont]{../../splat_editor/core/history.py}
\subsubsection{\texorpdfstring{\code{class History}}{History}}
{\footnotesize\textit{Dòng 55--83}}\par
\textbf{Giải thích.} Ngăn xếp tối đa \code{max_steps} (mặc định 20). \code{push} cắt bỏ phần cũ nhất bằng \code{del self._stack[:-max_steps]}.

Các phương thức: \code{__init__}, \code{__len__}, \code{__iter__}, \code{can_undo}, \code{push}, \code{undo}, \code{clear}.

\VerbatimInput[firstline=55,lastline=55,firstnumber=55,numbers=left,numbersep=3pt,frame=lines,framesep=2pt,rulecolor=\color{gray!60},breaklines=true,fontsize=\fontsize{6.6}{7.6}\selectfont]{../../splat_editor/core/history.py}
\subsubsection{\texorpdfstring{\code{History.__init__(self, max_steps=20)}}{History.__init__}}
{\footnotesize\textit{Dòng 56--58}}\par
\VerbatimInput[firstline=56,lastline=58,firstnumber=56,numbers=left,numbersep=3pt,frame=lines,framesep=2pt,rulecolor=\color{gray!60},breaklines=true,fontsize=\fontsize{6.6}{7.6}\selectfont]{../../splat_editor/core/history.py}
\subsubsection{\texorpdfstring{\code{History.__len__(self)}}{History.__len__}}
{\footnotesize\textit{Dòng 60--61}}\par
\VerbatimInput[firstline=60,lastline=61,firstnumber=60,numbers=left,numbersep=3pt,frame=lines,framesep=2pt,rulecolor=\color{gray!60},breaklines=true,fontsize=\fontsize{6.6}{7.6}\selectfont]{../../splat_editor/core/history.py}
\subsubsection{\texorpdfstring{\code{History.__iter__(self)}}{History.__iter__}}
{\footnotesize\textit{Dòng 63--64}}\par
\VerbatimInput[firstline=63,lastline=64,firstnumber=63,numbers=left,numbersep=3pt,frame=lines,framesep=2pt,rulecolor=\color{gray!60},breaklines=true,fontsize=\fontsize{6.6}{7.6}\selectfont]{../../splat_editor/core/history.py}
\subsubsection{\texorpdfstring{\code{History.can_undo(self)}}{History.can_undo}}
{\footnotesize\textit{Dòng 66--68}}\par
\VerbatimInput[firstline=66,lastline=68,firstnumber=66,numbers=left,numbersep=3pt,frame=lines,framesep=2pt,rulecolor=\color{gray!60},breaklines=true,fontsize=\fontsize{6.6}{7.6}\selectfont]{../../splat_editor/core/history.py}
\subsubsection{\texorpdfstring{\code{History.push(self, command)}}{History.push}}
{\footnotesize\textit{Dòng 70--72}}\par
\VerbatimInput[firstline=70,lastline=72,firstnumber=70,numbers=left,numbersep=3pt,frame=lines,framesep=2pt,rulecolor=\color{gray!60},breaklines=true,fontsize=\fontsize{6.6}{7.6}\selectfont]{../../splat_editor/core/history.py}
\subsubsection{\texorpdfstring{\code{History.undo(self)}}{History.undo}}
{\footnotesize\textit{Dòng 74--80}}\par
\textbf{Mô tả trong mã nguồn.} Hoàn tác thao tác gần nhất. Trả về Command đã hoàn tác hoặc None.

\VerbatimInput[firstline=74,lastline=80,firstnumber=74,numbers=left,numbersep=3pt,frame=lines,framesep=2pt,rulecolor=\color{gray!60},breaklines=true,fontsize=\fontsize{6.6}{7.6}\selectfont]{../../splat_editor/core/history.py}
\subsubsection{\texorpdfstring{\code{History.clear(self)}}{History.clear}}
{\footnotesize\textit{Dòng 82--83}}\par
\VerbatimInput[firstline=82,lastline=83,firstnumber=82,numbers=left,numbersep=3pt,frame=lines,framesep=2pt,rulecolor=\color{gray!60},breaklines=true,fontsize=\fontsize{6.6}{7.6}\selectfont]{../../splat_editor/core/history.py}
\subsection{\texorpdfstring{\code{splat_editor/core/rect.py}}{splat_editor/core/rect.py}}
{\footnotesize 33 dòng}\par
\textbf{Vai trò.} Hình chữ nhật trên màn hình (tọa độ cửa sổ, gốc trên-trái), dùng cho bố cục giao diện, cọ chọn và ma trận chiếu.

\paragraph{Mô tả đầu file và cách dùng.} Phần chú thích đầu file do tác giả viết, nêu ý tưởng và ví dụ sử dụng:
\VerbatimInput[firstline=1,lastline=10,firstnumber=1,numbers=left,numbersep=3pt,frame=lines,framesep=2pt,rulecolor=\color{gray!60},breaklines=true,fontsize=\fontsize{6.2}{7.2}\selectfont]{../../splat_editor/core/rect.py}
\subsubsection*{Khai báo mức module (dòng 11--11)}
\VerbatimInput[firstline=11,lastline=11,firstnumber=11,numbers=left,numbersep=3pt,frame=lines,framesep=2pt,rulecolor=\color{gray!60},breaklines=true,fontsize=\fontsize{6.6}{7.6}\selectfont]{../../splat_editor/core/rect.py}
\subsubsection{\texorpdfstring{\code{class Rect(NamedTuple)}}{Rect}}
{\footnotesize\textit{Dòng 14--33}}\par
Các phương thức: \code{is_empty}, \code{aspect}, \code{center}, \code{contains}.

\VerbatimInput[firstline=14,lastline=18,firstnumber=14,numbers=left,numbersep=3pt,frame=lines,framesep=2pt,rulecolor=\color{gray!60},breaklines=true,fontsize=\fontsize{6.6}{7.6}\selectfont]{../../splat_editor/core/rect.py}
\subsubsection{\texorpdfstring{\code{Rect.is_empty(self)}}{Rect.is_empty}}
{\footnotesize\textit{Dòng 20--22}}\par
\VerbatimInput[firstline=20,lastline=22,firstnumber=20,numbers=left,numbersep=3pt,frame=lines,framesep=2pt,rulecolor=\color{gray!60},breaklines=true,fontsize=\fontsize{6.6}{7.6}\selectfont]{../../splat_editor/core/rect.py}
\subsubsection{\texorpdfstring{\code{Rect.aspect(self)}}{Rect.aspect}}
{\footnotesize\textit{Dòng 24--26}}\par
\VerbatimInput[firstline=24,lastline=26,firstnumber=24,numbers=left,numbersep=3pt,frame=lines,framesep=2pt,rulecolor=\color{gray!60},breaklines=true,fontsize=\fontsize{6.6}{7.6}\selectfont]{../../splat_editor/core/rect.py}
\subsubsection{\texorpdfstring{\code{Rect.center(self)}}{Rect.center}}
{\footnotesize\textit{Dòng 28--30}}\par
\VerbatimInput[firstline=28,lastline=30,firstnumber=28,numbers=left,numbersep=3pt,frame=lines,framesep=2pt,rulecolor=\color{gray!60},breaklines=true,fontsize=\fontsize{6.6}{7.6}\selectfont]{../../splat_editor/core/rect.py}
\subsubsection{\texorpdfstring{\code{Rect.contains(self, x, y)}}{Rect.contains}}
{\footnotesize\textit{Dòng 32--33}}\par
\VerbatimInput[firstline=32,lastline=33,firstnumber=32,numbers=left,numbersep=3pt,frame=lines,framesep=2pt,rulecolor=\color{gray!60},breaklines=true,fontsize=\fontsize{6.6}{7.6}\selectfont]{../../splat_editor/core/rect.py}

\clearpage
\section{Thư mục editor: thao tác chỉnh sửa}\label{ch:editor}
Thư mục \code{editor} là tầng \textbf{thao tác chỉnh sửa}. Nó dùng \code{Scene}, \code{History} và \code{LodIndex}
của \code{core} nhưng cũng \textbf{không biết gì về OpenGL hay giao diện}. Hai file: \code{editor.py} (lớp \code{Editor}: mở/lưu file, chọn,
xóa, hoàn tác, phát tín hiệu "dữ liệu đã đổi" bằng số phiên bản) và \code{brush.py} (cọ chọn: tìm gaussian nằm trong nét tô trên màn hình).

\subsection{\texorpdfstring{\code{splat_editor/editor/editor.py}}{splat_editor/editor/editor.py}}
{\footnotesize 251 dòng}\par
\textbf{Vai trò.} \code{Editor} là cửa ngõ duy nhất của mọi thao tác chỉnh sửa. Cơ chế quan trọng là \textbf{số phiên bản}: mỗi khi dữ liệu đổi, Editor tăng
\code{geometry_version}, \code{deleted_version}, \code{selection_version} hoặc \code{transform_version}; \code{Viewport.sync} so sánh các số này để biết phần nào cần gửi lại GPU.
Nhờ vậy xóa/hoàn tác chỉ gửi mảng trạng thái 1 byte/gaussian, kéo layer chỉ gửi ma trận 4$\times$4.

\paragraph{Mô tả đầu file và cách dùng.} Phần chú thích đầu file do tác giả viết, nêu ý tưởng và ví dụ sử dụng:
\VerbatimInput[firstline=1,lastline=24,firstnumber=1,numbers=left,numbersep=3pt,frame=lines,framesep=2pt,rulecolor=\color{gray!60},breaklines=true,fontsize=\fontsize{6.2}{7.2}\selectfont]{../../splat_editor/editor/editor.py}
\subsubsection*{Khai báo mức module (dòng 25--34)}
\VerbatimInput[firstline=25,lastline=34,firstnumber=25,numbers=left,numbersep=3pt,frame=lines,framesep=2pt,rulecolor=\color{gray!60},breaklines=true,fontsize=\fontsize{6.6}{7.6}\selectfont]{../../splat_editor/editor/editor.py}
\subsubsection{\texorpdfstring{\code{class Editor}}{Editor}}
{\footnotesize\textit{Dòng 37--251}}\par
Các phương thức: \code{__init__}, \code{splats}, \code{lod}, \code{layer_matrices}, \code{total_count}, \code{notify_geometry_changed}, \code{notify_deleted_changed}, \code{notify_transform_changed}, \code{notify_selection_changed}, \code{selection_changes_since}, \code{set_status}, \code{_load_layer}, \code{open_file}, \code{add_file}, \code{can_save}, \code{suggested_save_name}, \code{save_file}, \code{_release_file}, \code{remove_layer}, \code{select_all}, \code{select_none}, \code{invert_selection}, \code{paint}, \code{delete_selected}, \code{delete}, \code{undo}.

\VerbatimInput[firstline=37,lastline=37,firstnumber=37,numbers=left,numbersep=3pt,frame=lines,framesep=2pt,rulecolor=\color{gray!60},breaklines=true,fontsize=\fontsize{6.6}{7.6}\selectfont]{../../splat_editor/editor/editor.py}
\subsubsection{\texorpdfstring{\code{Editor.__init__(self, max_undo=20)}}{Editor.__init__}}
{\footnotesize\textit{Dòng 38--51}}\par
\VerbatimInput[firstline=38,lastline=51,firstnumber=38,numbers=left,numbersep=3pt,frame=lines,framesep=2pt,rulecolor=\color{gray!60},breaklines=true,fontsize=\fontsize{6.6}{7.6}\selectfont]{../../splat_editor/editor/editor.py}
\subsubsection{\texorpdfstring{\code{Editor.splats(self)}}{Editor.splats}}
{\footnotesize\textit{Dòng 54--61}}\par
\textbf{Mô tả trong mã nguồn.} Tất cả lớp đang hiện, gộp lại, ở toạ độ file của từng lớp (tính lại khi dữ liệu đổi). Toạ độ thế giới = layer\_matrices()[lớp] @ toạ độ file.

\textbf{Giải thích.} Thuộc tính lười: gộp layer đang hiện thành một \code{Splats} và chỉ tính lại khi \code{geometry_version} đổi.

\VerbatimInput[firstline=54,lastline=61,firstnumber=54,numbers=left,numbersep=3pt,frame=lines,framesep=2pt,rulecolor=\color{gray!60},breaklines=true,fontsize=\fontsize{6.6}{7.6}\selectfont]{../../splat_editor/editor/editor.py}
\subsubsection{\texorpdfstring{\code{Editor.lod(self)}}{Editor.lod}}
{\footnotesize\textit{Dòng 63--72}}\par
\textbf{Mô tả trong mã nguồn.} Chỉ mục không gian (khối + LOD) của self.splats, mỗi lớp một nhóm; dùng chung cho renderer và cọ chọn. Chỉ dựng lại khi dữ liệu đổi, không phải khi kéo / xoay lớp.

\textbf{Giải thích.} Chỉ mục LOD dùng chung cho renderer, bộ sắp xếp và cọ; dựng lại cùng lúc với \code{splats}, mỗi layer một nhóm.

\VerbatimInput[firstline=63,lastline=72,firstnumber=63,numbers=left,numbersep=3pt,frame=lines,framesep=2pt,rulecolor=\color{gray!60},breaklines=true,fontsize=\fontsize{6.6}{7.6}\selectfont]{../../splat_editor/editor/editor.py}
\subsubsection{\texorpdfstring{\code{Editor.layer_matrices(self)}}{Editor.layer_matrices}}
{\footnotesize\textit{Dòng 74--75}}\par
\VerbatimInput[firstline=74,lastline=75,firstnumber=74,numbers=left,numbersep=3pt,frame=lines,framesep=2pt,rulecolor=\color{gray!60},breaklines=true,fontsize=\fontsize{6.6}{7.6}\selectfont]{../../splat_editor/editor/editor.py}
\subsubsection{\texorpdfstring{\code{Editor.total_count(self)}}{Editor.total_count}}
{\footnotesize\textit{Dòng 77--80}}\par
\textbf{Mô tả trong mã nguồn.} Số gaussian còn lại (chưa xoá) của các lớp đang hiện.

\VerbatimInput[firstline=77,lastline=80,firstnumber=77,numbers=left,numbersep=3pt,frame=lines,framesep=2pt,rulecolor=\color{gray!60},breaklines=true,fontsize=\fontsize{6.6}{7.6}\selectfont]{../../splat_editor/editor/editor.py}
\subsubsection{\texorpdfstring{\code{Editor.notify_geometry_changed(self)}}{Editor.notify_geometry_changed}}
{\footnotesize\textit{Dòng 82--86}}\par
\textbf{Mô tả trong mã nguồn.} Dữ liệu gaussian đổi (mở / xoá / ẩn lớp...): phải dựng lại và gửi lại lên GPU.

\VerbatimInput[firstline=82,lastline=86,firstnumber=82,numbers=left,numbersep=3pt,frame=lines,framesep=2pt,rulecolor=\color{gray!60},breaklines=true,fontsize=\fontsize{6.6}{7.6}\selectfont]{../../splat_editor/editor/editor.py}
\subsubsection{\texorpdfstring{\code{Editor.notify_deleted_changed(self)}}{Editor.notify_deleted_changed}}
{\footnotesize\textit{Dòng 88--91}}\par
\textbf{Mô tả trong mã nguồn.} Tập gaussian đã xoá đổi (xoá / hoàn tác): renderer chỉ gửi lại cờ, không gửi lại dữ liệu.

\VerbatimInput[firstline=88,lastline=91,firstnumber=88,numbers=left,numbersep=3pt,frame=lines,framesep=2pt,rulecolor=\color{gray!60},breaklines=true,fontsize=\fontsize{6.6}{7.6}\selectfont]{../../splat_editor/editor/editor.py}
\subsubsection{\texorpdfstring{\code{Editor.notify_transform_changed(self)}}{Editor.notify_transform_changed}}
{\footnotesize\textit{Dòng 93--95}}\par
\textbf{Mô tả trong mã nguồn.} Chỉ phép biến đổi của lớp (hoặc lật trục) đổi: renderer chỉ cần cập nhật ma trận.

\VerbatimInput[firstline=93,lastline=95,firstnumber=93,numbers=left,numbersep=3pt,frame=lines,framesep=2pt,rulecolor=\color{gray!60},breaklines=true,fontsize=\fontsize{6.6}{7.6}\selectfont]{../../splat_editor/editor/editor.py}
\subsubsection{\texorpdfstring{\code{Editor.notify_selection_changed(self, indices=None, value=None)}}{Editor.notify_selection_changed}}
{\footnotesize\textit{Dòng 97--104}}\par
\textbf{Mô tả trong mã nguồn.} indices / value: thay đổi chỉ gồm `vùng chọn[indices] = value` (để gửi GPU từng phần).

\textbf{Giải thích.} Ghi lại thay đổi vùng chọn dạng \code{(chỉ số, giá trị)} để GPU chỉ nhận phần đã tô thay vì cả mảng.

\VerbatimInput[firstline=97,lastline=104,firstnumber=97,numbers=left,numbersep=3pt,frame=lines,framesep=2pt,rulecolor=\color{gray!60},breaklines=true,fontsize=\fontsize{6.6}{7.6}\selectfont]{../../splat_editor/editor/editor.py}
\subsubsection{\texorpdfstring{\code{Editor.selection_changes_since(self, version)}}{Editor.selection_changes_since}}
{\footnotesize\textit{Dòng 106--111}}\par
\textbf{Mô tả trong mã nguồn.} [(chỉ số, giá trị)] đã tô kể từ `version`, hoặc None nếu có thay đổi khác (phải gửi lại cả).

\VerbatimInput[firstline=106,lastline=111,firstnumber=106,numbers=left,numbersep=3pt,frame=lines,framesep=2pt,rulecolor=\color{gray!60},breaklines=true,fontsize=\fontsize{6.6}{7.6}\selectfont]{../../splat_editor/editor/editor.py}
\subsubsection{\texorpdfstring{\code{Editor.set_status(self, text)}}{Editor.set_status}}
{\footnotesize\textit{Dòng 113--115}}\par
\VerbatimInput[firstline=113,lastline=115,firstnumber=113,numbers=left,numbersep=3pt,frame=lines,framesep=2pt,rulecolor=\color{gray!60},breaklines=true,fontsize=\fontsize{6.6}{7.6}\selectfont]{../../splat_editor/editor/editor.py}
\subsubsection{\texorpdfstring{\code{Editor._load_layer(self, path)}}{Editor._load_layer}}
{\footnotesize\textit{Dòng 118--123}}\par
\VerbatimInput[firstline=118,lastline=123,firstnumber=118,numbers=left,numbersep=3pt,frame=lines,framesep=2pt,rulecolor=\color{gray!60},breaklines=true,fontsize=\fontsize{6.6}{7.6}\selectfont]{../../splat_editor/editor/editor.py}
\subsubsection{\texorpdfstring{\code{Editor.open_file(self, path)}}{Editor.open_file}}
{\footnotesize\textit{Dòng 125--135}}\par
\textbf{Mô tả trong mã nguồn.} Mở file mới, thay thế toàn bộ cảnh. Mở lỗi thì giữ nguyên cảnh cũ.

\VerbatimInput[firstline=125,lastline=135,firstnumber=125,numbers=left,numbersep=3pt,frame=lines,framesep=2pt,rulecolor=\color{gray!60},breaklines=true,fontsize=\fontsize{6.6}{7.6}\selectfont]{../../splat_editor/editor/editor.py}
\subsubsection{\texorpdfstring{\code{Editor.add_file(self, path)}}{Editor.add_file}}
{\footnotesize\textit{Dòng 137--145}}\par
\textbf{Mô tả trong mã nguồn.} Thêm file vào cảnh đang có. Trả về lớp mới hoặc None.

\VerbatimInput[firstline=137,lastline=145,firstnumber=137,numbers=left,numbersep=3pt,frame=lines,framesep=2pt,rulecolor=\color{gray!60},breaklines=true,fontsize=\fontsize{6.6}{7.6}\selectfont]{../../splat_editor/editor/editor.py}
\subsubsection{\texorpdfstring{\code{Editor.can_save(self)}}{Editor.can_save}}
{\footnotesize\textit{Dòng 147--148}}\par
\VerbatimInput[firstline=147,lastline=148,firstnumber=147,numbers=left,numbersep=3pt,frame=lines,framesep=2pt,rulecolor=\color{gray!60},breaklines=true,fontsize=\fontsize{6.6}{7.6}\selectfont]{../../splat_editor/editor/editor.py}
\subsubsection{\texorpdfstring{\code{Editor.suggested_save_name(self)}}{Editor.suggested_save_name}}
{\footnotesize\textit{Dòng 150--152}}\par
\VerbatimInput[firstline=150,lastline=152,firstnumber=150,numbers=left,numbersep=3pt,frame=lines,framesep=2pt,rulecolor=\color{gray!60},breaklines=true,fontsize=\fontsize{6.6}{7.6}\selectfont]{../../splat_editor/editor/editor.py}
\subsubsection{\texorpdfstring{\code{Editor.save_file(self, path)}}{Editor.save_file}}
{\footnotesize\textit{Dòng 154--168}}\par
\VerbatimInput[firstline=154,lastline=168,firstnumber=154,numbers=left,numbersep=3pt,frame=lines,framesep=2pt,rulecolor=\color{gray!60},breaklines=true,fontsize=\fontsize{6.6}{7.6}\selectfont]{../../splat_editor/editor/editor.py}
\subsubsection{\texorpdfstring{\code{Editor._release_file(self, path)}}{Editor._release_file}}
{\footnotesize\textit{Dòng 170--184}}\par
\textbf{Mô tả trong mã nguồn.} File PLY đang mở được ánh xạ thẳng vào bộ nhớ (memmap); Windows không cho ghi đè file đang ánh xạ. Nếu lưu đè lên chính nó thì chép dữ liệu đó vào RAM trước.

\textbf{Giải thích.} Windows không cho ghi đè file đang được memmap, nên trước khi lưu đè phải nạp dữ liệu vào RAM (qua \code{load}) và thả memmap.

\VerbatimInput[firstline=170,lastline=184,firstnumber=170,numbers=left,numbersep=3pt,frame=lines,framesep=2pt,rulecolor=\color{gray!60},breaklines=true,fontsize=\fontsize{6.6}{7.6}\selectfont]{../../splat_editor/editor/editor.py}
\subsubsection{\texorpdfstring{\code{Editor.remove_layer(self, index=None)}}{Editor.remove_layer}}
{\footnotesize\textit{Dòng 187--192}}\par
\VerbatimInput[firstline=187,lastline=192,firstnumber=187,numbers=left,numbersep=3pt,frame=lines,framesep=2pt,rulecolor=\color{gray!60},breaklines=true,fontsize=\fontsize{6.6}{7.6}\selectfont]{../../splat_editor/editor/editor.py}
\subsubsection{\texorpdfstring{\code{Editor.select_all(self)}}{Editor.select_all}}
{\footnotesize\textit{Dòng 195--198}}\par
\VerbatimInput[firstline=195,lastline=198,firstnumber=195,numbers=left,numbersep=3pt,frame=lines,framesep=2pt,rulecolor=\color{gray!60},breaklines=true,fontsize=\fontsize{6.6}{7.6}\selectfont]{../../splat_editor/editor/editor.py}
\subsubsection{\texorpdfstring{\code{Editor.select_none(self)}}{Editor.select_none}}
{\footnotesize\textit{Dòng 200--203}}\par
\VerbatimInput[firstline=200,lastline=203,firstnumber=200,numbers=left,numbersep=3pt,frame=lines,framesep=2pt,rulecolor=\color{gray!60},breaklines=true,fontsize=\fontsize{6.6}{7.6}\selectfont]{../../splat_editor/editor/editor.py}
\subsubsection{\texorpdfstring{\code{Editor.invert_selection(self)}}{Editor.invert_selection}}
{\footnotesize\textit{Dòng 205--208}}\par
\VerbatimInput[firstline=205,lastline=208,firstnumber=205,numbers=left,numbersep=3pt,frame=lines,framesep=2pt,rulecolor=\color{gray!60},breaklines=true,fontsize=\fontsize{6.6}{7.6}\selectfont]{../../splat_editor/editor/editor.py}
\subsubsection{\texorpdfstring{\code{Editor.paint(self, x, y, view_proj, rect)}}{Editor.paint}}
{\footnotesize\textit{Dòng 210--221}}\par
\textbf{Mô tả trong mã nguồn.} Tô cọ từ vị trí lần trước tới (x, y): chọn hoặc bỏ chọn gaussian.

\textbf{Giải thích.} Tô cọ từ vị trí trước đến vị trí hiện tại: gọi \code{BrushTool.hit} để lấy chỉ số rồi \code{Scene.apply_selection_indices}.

\VerbatimInput[firstline=210,lastline=221,firstnumber=210,numbers=left,numbersep=3pt,frame=lines,framesep=2pt,rulecolor=\color{gray!60},breaklines=true,fontsize=\fontsize{6.6}{7.6}\selectfont]{../../splat_editor/editor/editor.py}
\subsubsection{\texorpdfstring{\code{Editor.delete_selected(self)}}{Editor.delete_selected}}
{\footnotesize\textit{Dòng 224--236}}\par
\textbf{Giải thích.} Xóa mềm mọi layer đang có vùng chọn, đẩy một \code{DeleteGaussiansCommand} vào lịch sử và tăng \code{deleted_version}.

\VerbatimInput[firstline=224,lastline=236,firstnumber=224,numbers=left,numbersep=3pt,frame=lines,framesep=2pt,rulecolor=\color{gray!60},breaklines=true,fontsize=\fontsize{6.6}{7.6}\selectfont]{../../splat_editor/editor/editor.py}
\subsubsection{\texorpdfstring{\code{Editor.delete(self)}}{Editor.delete}}
{\footnotesize\textit{Dòng 238--243}}\par
\textbf{Mô tả trong mã nguồn.} Phím Delete: xoá gaussian đang chọn nếu có, nếu không thì xoá lớp đang chọn.

\VerbatimInput[firstline=238,lastline=243,firstnumber=238,numbers=left,numbersep=3pt,frame=lines,framesep=2pt,rulecolor=\color{gray!60},breaklines=true,fontsize=\fontsize{6.6}{7.6}\selectfont]{../../splat_editor/editor/editor.py}
\subsubsection{\texorpdfstring{\code{Editor.undo(self)}}{Editor.undo}}
{\footnotesize\textit{Dòng 245--251}}\par
\VerbatimInput[firstline=245,lastline=251,firstnumber=245,numbers=left,numbersep=3pt,frame=lines,framesep=2pt,rulecolor=\color{gray!60},breaklines=true,fontsize=\fontsize{6.6}{7.6}\selectfont]{../../splat_editor/editor/editor.py}
\subsection{\texorpdfstring{\code{splat_editor/editor/brush.py}}{splat_editor/editor/brush.py}}
{\footnotesize 140 dòng}\par
\textbf{Vai trò.} Cọ chọn xuyên độ sâu: gaussian nào có tâm chiếu vào nét tô đều được chọn. Để nhanh, cọ tận dụng hộp bao của các khối LOD:
chiếu 8 đỉnh hộp ra hình chữ nhật, khối nằm trọn trong nét thì lấy cả khối, khối chỉ chạm thì mới kiểm từng gaussian. Toán học: xem mục \textit{Cơ sở toán} ngay bên dưới (\S\ref{s:brush}).

\paragraph{Mô tả đầu file và cách dùng.} Phần chú thích đầu file do tác giả viết, nêu ý tưởng và ví dụ sử dụng:
\VerbatimInput[firstline=1,lastline=20,firstnumber=1,numbers=left,numbersep=3pt,frame=lines,framesep=2pt,rulecolor=\color{gray!60},breaklines=true,fontsize=\fontsize{6.2}{7.2}\selectfont]{../../splat_editor/editor/brush.py}
%% ---- merged math: s:viewport
\subsubsection{Cơ sở toán: biến đổi viewport (\texorpdfstring{\code{brush._to_screen}}{brush._to_screen})}\label{s:viewport}
\[
x_{px}=x_0+\tfrac{n_x+1}{2}W,\qquad y_{px}=y_0+\tfrac{1-n_y}{2}H\ \text{(đảo $y$: màn hình có $y$ hướng xuống)}.
\]

%% ---- merged math: s:brush
\subsubsection{Cơ sở toán: Cọ chọn}\label{s:brush}

\paragraph{Khoảng cách điểm--đoạn, capsule (\code{_segment_dist2}).}
$\tau=\mathrm{clip}\!\big((p-a)\cdot(b-a)/\|b-a\|^2,0,1\big)$, $d^2=\|p-(a+\tau(b-a))\|^2$; chọn nếu $d^2\le r^2$. Nếu $a=b$ thì $\tau=0$. Capsule là tập điểm cách đoạn không quá $r$ (hình viên thuốc, lồi).

\paragraph{Chiếu tâm ra pixel.}
$\mathrm{clip}=M_{[:3,:3]}p+M_{[:3,3]}$ với $M=VP\cdot M_{layer}$; chỉ xét $w>10^{-6}$ (trước camera); $x=\mathrm{left}+(c_x/w+1)/2\cdot W$, $y=\mathrm{top}+(1-c_y/w)/2\cdot H$.

\paragraph{Tăng tốc theo khối.}
8 đỉnh hộp bao khối (\code{_CORNERS}: bit $i$ của chỉ số chọn lo/hi) được chiếu thành hình chữ nhật bao (hộp lồi và phép chiếu bảo toàn tính lồi khi $w>0$). Khối chạm nét nếu hình chữ nhật giao hộp $[\min(a,b)-r,\max(a,b)+r]$. Khối \textbf{nằm trọn trong nét} nếu cả 4 góc hình chữ nhật có $d^2\le r^2$ (capsule lồi nên cả hình chữ nhật nằm trong), khi đó lấy cả khối mà không kiểm từng gaussian. Khối cắt mặt phẳng camera (dấu $w$ lẫn lộn) luôn được kiểm. Cọ \textbf{xuyên sâu} (không test độ sâu) như SuperSplat.

\subsubsection*{Khai báo mức module (dòng 21--28)}
\VerbatimInput[firstline=21,lastline=28,firstnumber=21,numbers=left,numbersep=3pt,frame=lines,framesep=2pt,rulecolor=\color{gray!60},breaklines=true,fontsize=\fontsize{6.6}{7.6}\selectfont]{../../splat_editor/editor/brush.py}
\subsubsection{\texorpdfstring{\code{class BrushTool}}{BrushTool}}
{\footnotesize\textit{Dòng 31--140}}\par
Các phương thức: \code{__init__}, \code{is_painting}, \code{toggle}, \code{resize}, \code{begin_stroke}, \code{end_stroke}, \code{_to_screen}, \code{_segment_dist2}, \code{_candidate_chunks}, \code{hit}.

\VerbatimInput[firstline=31,lastline=33,firstnumber=31,numbers=left,numbersep=3pt,frame=lines,framesep=2pt,rulecolor=\color{gray!60},breaklines=true,fontsize=\fontsize{6.6}{7.6}\selectfont]{../../splat_editor/editor/brush.py}
\subsubsection{\texorpdfstring{\code{BrushTool.__init__(self)}}{BrushTool.__init__}}
{\footnotesize\textit{Dòng 35--39}}\par
\VerbatimInput[firstline=35,lastline=39,firstnumber=35,numbers=left,numbersep=3pt,frame=lines,framesep=2pt,rulecolor=\color{gray!60},breaklines=true,fontsize=\fontsize{6.6}{7.6}\selectfont]{../../splat_editor/editor/brush.py}
\subsubsection{\texorpdfstring{\code{BrushTool.is_painting(self)}}{BrushTool.is_painting}}
{\footnotesize\textit{Dòng 42--44}}\par
\VerbatimInput[firstline=42,lastline=44,firstnumber=42,numbers=left,numbersep=3pt,frame=lines,framesep=2pt,rulecolor=\color{gray!60},breaklines=true,fontsize=\fontsize{6.6}{7.6}\selectfont]{../../splat_editor/editor/brush.py}
\subsubsection{\texorpdfstring{\code{BrushTool.toggle(self)}}{BrushTool.toggle}}
{\footnotesize\textit{Dòng 46--47}}\par
\VerbatimInput[firstline=46,lastline=47,firstnumber=46,numbers=left,numbersep=3pt,frame=lines,framesep=2pt,rulecolor=\color{gray!60},breaklines=true,fontsize=\fontsize{6.6}{7.6}\selectfont]{../../splat_editor/editor/brush.py}
\subsubsection{\texorpdfstring{\code{BrushTool.resize(self, factor)}}{BrushTool.resize}}
{\footnotesize\textit{Dòng 49--50}}\par
\VerbatimInput[firstline=49,lastline=50,firstnumber=49,numbers=left,numbersep=3pt,frame=lines,framesep=2pt,rulecolor=\color{gray!60},breaklines=true,fontsize=\fontsize{6.6}{7.6}\selectfont]{../../splat_editor/editor/brush.py}
\subsubsection{\texorpdfstring{\code{BrushTool.begin_stroke(self, select)}}{BrushTool.begin_stroke}}
{\footnotesize\textit{Dòng 52--55}}\par
\textbf{Mô tả trong mã nguồn.} select=True: tô để chọn; False: tô để bỏ chọn.

\VerbatimInput[firstline=52,lastline=55,firstnumber=52,numbers=left,numbersep=3pt,frame=lines,framesep=2pt,rulecolor=\color{gray!60},breaklines=true,fontsize=\fontsize{6.6}{7.6}\selectfont]{../../splat_editor/editor/brush.py}
\subsubsection{\texorpdfstring{\code{BrushTool.end_stroke(self)}}{BrushTool.end_stroke}}
{\footnotesize\textit{Dòng 57--58}}\par
\VerbatimInput[firstline=57,lastline=58,firstnumber=57,numbers=left,numbersep=3pt,frame=lines,framesep=2pt,rulecolor=\color{gray!60},breaklines=true,fontsize=\fontsize{6.6}{7.6}\selectfont]{../../splat_editor/editor/brush.py}
\subsubsection{\texorpdfstring{\code{BrushTool._to_screen(clip, viewport)}}{BrushTool._to_screen}}
{\footnotesize\textit{Dòng 61--66}}\par
\textbf{Giải thích.} Chiếu tâm gaussian ra pixel bằng $VP\cdot M_{layer}$; chỉ giữ điểm có $w>10^{-6}$ (trước camera).

\textbf{Công thức.}
\[
\mathrm{clip}=M_{[:3,:3]}\,p+M_{[:3,3]},\quad M=VP\cdot M_{layer},\qquad
x_{px}=x_0+\frac{c_x/w+1}{2}W,\quad y_{px}=y_0+\frac{1-c_y/w}{2}H
\]


\VerbatimInput[firstline=61,lastline=66,firstnumber=61,numbers=left,numbersep=3pt,frame=lines,framesep=2pt,rulecolor=\color{gray!60},breaklines=true,fontsize=\fontsize{6.6}{7.6}\selectfont]{../../splat_editor/editor/brush.py}
\subsubsection{\texorpdfstring{\code{BrushTool._segment_dist2(p, a, b)}}{BrushTool._segment_dist2}}
{\footnotesize\textit{Dòng 68--75}}\par
\textbf{Mô tả trong mã nguồn.} Bình phương khoảng cách từ các điểm p (...,2) tới đoạn thẳng a -\textgreater{} b.

\textbf{Giải thích.} Bình phương khoảng cách điểm--đoạn (capsule): chiếu điểm lên đường thẳng, kẹp tham số $\tau\in[0,1]$.

\textbf{Công thức.}
\[
\tau=\mathrm{clip}\Big(\frac{(p-a)\cdot(b-a)}{\|b-a\|^2},0,1\Big),\qquad d^2=\big\|p-\big(a+\tau(b-a)\big)\big\|^2\le r^2
\]


\VerbatimInput[firstline=68,lastline=75,firstnumber=68,numbers=left,numbersep=3pt,frame=lines,framesep=2pt,rulecolor=\color{gray!60},breaklines=true,fontsize=\fontsize{6.6}{7.6}\selectfont]{../../splat_editor/editor/brush.py}
\subsubsection{\texorpdfstring{\code{BrushTool._candidate_chunks(self, lod, matrices, viewport, a, b)}}{BrushTool._candidate_chunks}}
{\footnotesize\textit{Dòng 77--97}}\par
\textbf{Mô tả trong mã nguồn.} (khối phải kiểm tra từng gaussian, khối nằm trọn trong nét cọ) theo hình chữ nhật bao hình chiếu hộp bao của khối. matrices: (số nhóm, 4, 4) chiếu * thế giới của từng nhóm.

\textbf{Giải thích.} Với mỗi khối, chiếu 8 đỉnh hộp bao thành hình chữ nhật; trả về hai nhóm: khối cần kiểm từng gaussian và khối nằm \textbf{trọn} trong nét cọ.

\textbf{Công thức.}
Hộp bao khối có 8 đỉnh $v_i\in\{lo,hi\}^3$; hình chữ nhật bao $=[\min_i x_i,\max_i x_i]\times[\min_i y_i,\max_i y_i]$.
Khối nằm \textbf{trọn} trong nét nếu $d^2(\text{4 góc})\le r^2$ (capsule lồi nên cả hình chữ nhật nằm trong).


\VerbatimInput[firstline=77,lastline=97,firstnumber=77,numbers=left,numbersep=3pt,frame=lines,framesep=2pt,rulecolor=\color{gray!60},breaklines=true,fontsize=\fontsize{6.6}{7.6}\selectfont]{../../splat_editor/editor/brush.py}
\subsubsection{\texorpdfstring{\code{BrushTool.hit(self, lod, view_proj, viewport, start, end, transforms=None)}}{BrushTool.hit}}
{\footnotesize\textit{Dòng 99--140}}\par
\textbf{Mô tả trong mã nguồn.} Chỉ số các gaussian có tâm cách đoạn thẳng start -\textgreater{} end (pixel) không quá `radius`. transforms: ma trận thế giới 4x4 của từng nhóm trong lod (None: đơn vị).

\textbf{Giải thích.} Hàm chính của cọ: lấy chỉ số gaussian có tâm cách đoạn nét tô không quá bán kính. Khối nằm trọn thì lấy cả khối; khối cắt mặt phẳng camera luôn bị kiểm; phần kiểm chia nhiều luồng.

\VerbatimInput[firstline=99,lastline=140,firstnumber=99,numbers=left,numbersep=3pt,frame=lines,framesep=2pt,rulecolor=\color{gray!60},breaklines=true,fontsize=\fontsize{6.6}{7.6}\selectfont]{../../splat_editor/editor/brush.py}

\clearpage
\section{Thư mục io: đọc ghi và giải mã}\label{ch:io}
Thư mục \code{io} là tầng \textbf{đọc/ghi và biến đổi định dạng} giữa file PLY và cấu trúc \code{Splats}. \code{ply.py}
lo phần byte (header, kiểu dữ liệu, ánh xạ bộ nhớ); \code{codec.py} lo phần ý nghĩa (sigmoid, exp, hằng số $C_0$, chuẩn hóa quaternion).
Tách đôi như vậy để thay đổi định dạng lưu mà không đụng tới logic hiển thị.

\subsection{\texorpdfstring{\code{splat_editor/io/ply.py}}{splat_editor/io/ply.py}}
{\footnotesize 134 dòng}\par
\textbf{Vai trò.} Đọc/ghi PLY. File \code{binary_little_endian} được ánh xạ bộ nhớ (\code{np.memmap}) nên mở file 4 GB gần như tức thì và không chiếm RAM.
Mỗi property là một cột của mảng cấu trúc NumPy (\code{data['x']}, \code{data['opacity']}, \ldots). Ghi file theo từng khối để không giữ cả layer trong RAM.

\paragraph{Mô tả đầu file và cách dùng.} Phần chú thích đầu file do tác giả viết, nêu ý tưởng và ví dụ sử dụng:
\VerbatimInput[firstline=1,lastline=22,firstnumber=1,numbers=left,numbersep=3pt,frame=lines,framesep=2pt,rulecolor=\color{gray!60},breaklines=true,fontsize=\fontsize{6.2}{7.2}\selectfont]{../../splat_editor/io/ply.py}
\subsubsection*{Khai báo mức module (dòng 23--25)}
\VerbatimInput[firstline=23,lastline=25,firstnumber=23,numbers=left,numbersep=3pt,frame=lines,framesep=2pt,rulecolor=\color{gray!60},breaklines=true,fontsize=\fontsize{6.6}{7.6}\selectfont]{../../splat_editor/io/ply.py}
\subsubsection{\texorpdfstring{\code{read_block(data, start, end)}}{read_block}}
{\footnotesize\textit{Dòng 28--36}}\par
\textbf{Mô tả trong mã nguồn.} Hàng [start, end) của `data` (có thể là memmap): chép liền một mạch vào RAM rồi trả về (khối, ma trận float32 (n, số cột) nếu mọi cột đều là float32 little endian, ngược lại None). Đọc cột qua ma trận nhanh hơn nhiều so với truy cập từng trường của mảng có cấu trúc.

\textbf{Giải thích.} Chép hàng $[start,end)$ của memmap liền một mạch vào RAM, trả về khối cùng dạng ma trận float32 nếu mọi cột cùng kiểu (nhanh hơn truy cập từng cột).

\VerbatimInput[firstline=28,lastline=36,firstnumber=28,numbers=left,numbersep=3pt,frame=lines,framesep=2pt,rulecolor=\color{gray!60},breaklines=true,fontsize=\fontsize{6.6}{7.6}\selectfont]{../../splat_editor/io/ply.py}
\subsubsection{\texorpdfstring{\code{class PlyError(ValueError)}}{PlyError}}
{\footnotesize\textit{Dòng 39--40}}\par
\textbf{Mô tả trong mã nguồn.} File PLY sai định dạng hoặc không được hỗ trợ.

\VerbatimInput[firstline=39,lastline=40,firstnumber=39,numbers=left,numbersep=3pt,frame=lines,framesep=2pt,rulecolor=\color{gray!60},breaklines=true,fontsize=\fontsize{6.6}{7.6}\selectfont]{../../splat_editor/io/ply.py}
\subsubsection{\texorpdfstring{\code{class PlyFile}}{PlyFile}}
{\footnotesize\textit{Dòng 43--134}}\par
Các phương thức: \code{read}, \code{_read_header}, \code{_read_ascii}, \code{_read_binary}, \code{write}, \code{write_blocks}.

\VerbatimInput[firstline=43,lastline=54,firstnumber=43,numbers=left,numbersep=3pt,frame=lines,framesep=2pt,rulecolor=\color{gray!60},breaklines=true,fontsize=\fontsize{6.6}{7.6}\selectfont]{../../splat_editor/io/ply.py}
\subsubsection{\texorpdfstring{\code{PlyFile.read(cls, path)}}{PlyFile.read}}
{\footnotesize\textit{Dòng 55--67}}\par
\textbf{Mô tả trong mã nguồn.} Đọc phần 'element vertex' của file PLY. Hỗ trợ ascii và binary (little/big endian).

\VerbatimInput[firstline=55,lastline=67,firstnumber=55,numbers=left,numbersep=3pt,frame=lines,framesep=2pt,rulecolor=\color{gray!60},breaklines=true,fontsize=\fontsize{6.6}{7.6}\selectfont]{../../splat_editor/io/ply.py}
\subsubsection{\texorpdfstring{\code{PlyFile._read_header(cls, f)}}{PlyFile._read_header}}
{\footnotesize\textit{Dòng 69--92}}\par
\textbf{Mô tả trong mã nguồn.} -\textgreater{} (định dạng, [(tên element, số lượng, [(tên property, kiểu numpy)])])

\textbf{Giải thích.} Phân tích header: định dạng (ascii/binary little/big endian), danh sách element, property và kiểu NumPy tương ứng.

\VerbatimInput[firstline=69,lastline=92,firstnumber=69,numbers=left,numbersep=3pt,frame=lines,framesep=2pt,rulecolor=\color{gray!60},breaklines=true,fontsize=\fontsize{6.6}{7.6}\selectfont]{../../splat_editor/io/ply.py}
\subsubsection{\texorpdfstring{\code{PlyFile._read_ascii(f, count, props)}}{PlyFile._read_ascii}}
{\footnotesize\textit{Dòng 94--101}}\par
\VerbatimInput[firstline=94,lastline=101,firstnumber=94,numbers=left,numbersep=3pt,frame=lines,framesep=2pt,rulecolor=\color{gray!60},breaklines=true,fontsize=\fontsize{6.6}{7.6}\selectfont]{../../splat_editor/io/ply.py}
\subsubsection{\texorpdfstring{\code{PlyFile._read_binary(f, count, props, endian)}}{PlyFile._read_binary}}
{\footnotesize\textit{Dòng 103--114}}\par
\textbf{Giải thích.} Với little endian dùng \code{np.memmap} (không đọc cả file). Offset dữ liệu lấy từ độ dài header.

\VerbatimInput[firstline=103,lastline=114,firstnumber=103,numbers=left,numbersep=3pt,frame=lines,framesep=2pt,rulecolor=\color{gray!60},breaklines=true,fontsize=\fontsize{6.6}{7.6}\selectfont]{../../splat_editor/io/ply.py}
\subsubsection{\texorpdfstring{\code{PlyFile.write(cls, path, data)}}{PlyFile.write}}
{\footnotesize\textit{Dòng 117--120}}\par
\textbf{Mô tả trong mã nguồn.} Ghi mảng có cấu trúc (giống kết quả read) ra file PLY binary\_little\_endian.

\VerbatimInput[firstline=117,lastline=120,firstnumber=117,numbers=left,numbersep=3pt,frame=lines,framesep=2pt,rulecolor=\color{gray!60},breaklines=true,fontsize=\fontsize{6.6}{7.6}\selectfont]{../../splat_editor/io/ply.py}
\subsubsection{\texorpdfstring{\code{PlyFile.write_blocks(cls, path, dtype, count, blocks)}}{PlyFile.write_blocks}}
{\footnotesize\textit{Dòng 122--134}}\par
\textbf{Mô tả trong mã nguồn.} Ghi `count` gaussian kiểu `dtype` từ các khối (mảng có cấu trúc) lần lượt, để không phải gộp cả cảnh thành một mảng lớn trong RAM.

\textbf{Giải thích.} Ghi từng khối lần lượt để không phải giữ toàn bộ layer trong RAM khi lưu file hàng chục triệu gaussian.

\VerbatimInput[firstline=122,lastline=134,firstnumber=122,numbers=left,numbersep=3pt,frame=lines,framesep=2pt,rulecolor=\color{gray!60},breaklines=true,fontsize=\fontsize{6.6}{7.6}\selectfont]{../../splat_editor/io/ply.py}
\subsection{\texorpdfstring{\code{splat_editor/io/codec.py}}{splat_editor/io/codec.py}}
{\footnotesize 122 dòng}\par
\textbf{Vai trò.} \code{SplatCodec} chuyển giữa các cột "thô" của file 3DGS và \code{Splats} đã giải mã (sigmoid cho $\alpha$, exp cho scale, $C_0$ cho màu, chuẩn hóa quaternion),
và chiều ngược lại khi lưu. Toán học: xem mục \textit{Cơ sở toán} ngay bên dưới (\S\ref{s:codec}).

\paragraph{Mô tả đầu file và cách dùng.} Phần chú thích đầu file do tác giả viết, nêu ý tưởng và ví dụ sử dụng:
\VerbatimInput[firstline=1,lastline=11,firstnumber=1,numbers=left,numbersep=3pt,frame=lines,framesep=2pt,rulecolor=\color{gray!60},breaklines=true,fontsize=\fontsize{6.2}{7.2}\selectfont]{../../splat_editor/io/codec.py}
%% ---- merged math: s:codec
\subsubsection{Cơ sở toán: Giải mã và mã hóa file PLY}\label{s:codec}
Một gaussian gồm tâm $\mu$, màu $c$, độ đục $\alpha$, kích thước $s$, quay $q$. File lưu giá trị ``thô'':

\begin{center}\small
\begin{tabular}{@{}lll@{}}\toprule
Thuộc tính & Giải mã & Lý do\\\midrule
\code{x,y,z} & $\mu=(x,y,z)$ & \\
\code{f_dc_0..2} & $c = 0.5 + C_0\,f_{dc}$,\; $C_0=\tfrac{1}{2\sqrt\pi}=0.28209479$ & $C_0=Y_{00}$ (SH bậc 0)\\
\code{opacity} & $\alpha = 1/(1+e^{-\mathrm{opacity}})$ & tối ưu logit để $\alpha\in(0,1)$\\
\code{scale_i} & $s_i = \exp(\mathrm{scale}_i)$ & tối ưu log để $s>0$\\
\code{rot_0..3} & $q\leftarrow q/\|q\|$; nếu $\|q\|=0$ thì $(1,0,0,0)$ & quaternion phải đơn vị\\\bottomrule
\end{tabular}
\end{center}

\textbf{Mã hóa ngược:} $f_{dc}=(c-0.5)/C_0$, $\mathrm{scale}=\ln\max(s,10^{-12})$, $\mathrm{opacity}=\ln\frac{a}{1-a}$ với $a=\mathrm{clip}(\alpha,10^{-6},1-10^{-6})$ để tránh $\ln 0$.

\textbf{Point cloud} (không có \code{scale_0}): bán kính mỗi điểm $=\tfrac12 (V/N)^{1/3}$, với $V$ là tích ba cạnh hộp bao. Đó là ước lượng khoảng cách trung bình giữa các điểm. Màu \code{uchar} chia 255.

\textbf{Dung lượng:} $3+3+3+45+1+3+4=62$ float32 $=248$ byte/gaussian (SH bậc 3: 15 hệ số $\times$ 3 kênh $=45$).

\subsubsection*{Khai báo mức module (dòng 12--25)}
\VerbatimInput[firstline=12,lastline=25,firstnumber=12,numbers=left,numbersep=3pt,frame=lines,framesep=2pt,rulecolor=\color{gray!60},breaklines=true,fontsize=\fontsize{6.6}{7.6}\selectfont]{../../splat_editor/io/codec.py}
\subsubsection{\texorpdfstring{\code{class SplatCodec}}{SplatCodec}}
{\footnotesize\textit{Dòng 28--122}}\par
Các phương thức: \code{decode}, \code{_decode_block}, \code{encode}.

\VerbatimInput[firstline=28,lastline=28,firstnumber=28,numbers=left,numbersep=3pt,frame=lines,framesep=2pt,rulecolor=\color{gray!60},breaklines=true,fontsize=\fontsize{6.6}{7.6}\selectfont]{../../splat_editor/io/codec.py}
\subsubsection{\texorpdfstring{\code{SplatCodec.decode(cls, data, sh_dtype=np.float16)}}{SplatCodec.decode}}
{\footnotesize\textit{Dòng 29--63}}\par
\textbf{Mô tả trong mã nguồn.} Chuyển các cột của file 3DGS PLY thành Splats.

File 3DGS chuẩn có: x y z, f\_dc\_0..2 (màu), f\_rest\_* (màu theo hướng nhìn, tuỳ chọn), opacity (logit), scale\_0..2 (log), rot\_0..3. Nếu chỉ là point cloud (x y z red green blue) thì tự tạo gaussian nhỏ hình cầu.

Giải mã từng đoạn CHUNK hàng (dữ liệu có thể là memmap trên đĩa) để không tạo mảng tạm bằng cả file. Hệ số SH chỉ dùng để hiển thị nên lưu float16 (đỡ một nửa RAM); khi lưu file vẫn đọc lại từ dữ liệu gốc nên không mất độ chính xác.

\textbf{Giải thích.} Giải mã từng đoạn (\code{_decode_block}) thành các mảng \code{Splats}; điểm đám mây (không có \code{scale_0}) được gán bán kính ước lượng.

\VerbatimInput[firstline=29,lastline=63,firstnumber=29,numbers=left,numbersep=3pt,frame=lines,framesep=2pt,rulecolor=\color{gray!60},breaklines=true,fontsize=\fontsize{6.6}{7.6}\selectfont]{../../splat_editor/io/codec.py}
\subsubsection{\texorpdfstring{\code{SplatCodec._decode_block(data, names, dtype)}}{SplatCodec._decode_block}}
{\footnotesize\textit{Dòng 65--98}}\par
\textbf{Mô tả trong mã nguồn.} (tâm, màu, alpha, scale, quaternion) của một đoạn hàng (scale = 1 nếu là point cloud).

\textbf{Giải thích.} Công thức giải mã: $c=0.5+C_0f_{dc}$, $\alpha=\mathrm{sigmoid}$, $s=\exp$, chuẩn hóa quaternion (rỗng $\to(1,0,0,0)$). Bán kính point cloud $=\tfrac12(V/N)^{1/3}$.

\textbf{Công thức.}
\[
c=0.5+C_0\,f_{dc},\quad \alpha=\frac{1}{1+e^{-\mathrm{opacity}}},\quad s_i=e^{\mathrm{scale}_i},\quad q\leftarrow\frac{q}{\|q\|}
\]
Point cloud: bán kính $r=\tfrac12\sqrt[3]{V/N}$, $V$ là thể tích hộp bao, $N$ số điểm.


\VerbatimInput[firstline=65,lastline=98,firstnumber=65,numbers=left,numbersep=3pt,frame=lines,framesep=2pt,rulecolor=\color{gray!60},breaklines=true,fontsize=\fontsize{6.6}{7.6}\selectfont]{../../splat_editor/io/codec.py}
\subsubsection{\texorpdfstring{\code{SplatCodec.encode(source, splats, dtype)}}{SplatCodec.encode}}
{\footnotesize\textit{Dòng 100--122}}\par
\textbf{Mô tả trong mã nguồn.} Đổi dữ liệu PLY `source` sang kiểu 3DGS chuẩn `dtype`.

Cột nào có sẵn trong file thì chép nguyên, cột thiếu thì tính từ `splats` (ví dụ point cloud chỉ có red/green/blue -\textgreater{} tính f\_dc, opacity, scale, rot).

\textbf{Giải thích.} Chiều ngược khi lưu: $f_{dc}=(c-0.5)/C_0$, $\ln$ cho scale, logit cho độ đục (kẹp $[10^{-6},1-10^{-6}]$ tránh $\ln0$).

\textbf{Công thức.}
\[
f_{dc}=\frac{c-0.5}{C_0},\quad \mathrm{scale}=\ln\max(s,10^{-12}),\quad
\mathrm{opacity}=\ln\frac{a}{1-a},\ \ a=\mathrm{clip}(\alpha,10^{-6},1-10^{-6})
\]


\VerbatimInput[firstline=100,lastline=122,firstnumber=100,numbers=left,numbersep=3pt,frame=lines,framesep=2pt,rulecolor=\color{gray!60},breaklines=true,fontsize=\fontsize{6.6}{7.6}\selectfont]{../../splat_editor/io/codec.py}

\clearpage
\section{Thư mục render: OpenGL}\label{ch:render}
Thư mục \code{render} là nơi duy nhất \textbf{nói chuyện trực tiếp với OpenGL 3.3}. Nó nhận dữ liệu từ \code{Editor}/\code{Scene}
rồi biến thành hình ảnh. Gồm: bọc shader (\code{shader.py}), camera (\code{camera.py}), bộ vẽ gaussian hai pass (\code{splat_renderer.py}),
luồng nền chọn và sắp xếp (\code{sorter.py}), khung nhìn ghép mọi thứ (\code{viewport.py}), các bộ vẽ phụ (lưới sàn, đèn, nhân vật) và thư mục
\code{shaders/} chứa mã GLSL. Pipeline hiển thị gaussian gồm 5 bước: (1) đóng gói dữ liệu lên GPU, (2) luồng nền chọn LOD và sắp xếp xa $\to$ gần,
(3) pass 1 tính ellipse và màu bằng transform feedback, (4) pass 2 vẽ hình vuông và trộn alpha, (5) chỉnh ngân sách theo thời gian GPU.

\subsection{\texorpdfstring{\code{splat_editor/render/shader.py}}{splat_editor/render/shader.py}}
{\footnotesize 110 dòng}\par
\textbf{Vai trò.} Bọc chương trình shader OpenGL: nạp \code{.vert/.frag} từ thư mục \code{shaders/}, biên dịch, liên kết, cache vị trí uniform, và các hàm \code{set_*}.
Hỗ trợ chương trình chỉ có vertex shader dùng \textbf{transform feedback} (pass 1 của bộ vẽ gaussian).

\paragraph{Mô tả đầu file và cách dùng.} Phần chú thích đầu file do tác giả viết, nêu ý tưởng và ví dụ sử dụng:
\VerbatimInput[firstline=1,lastline=24,firstnumber=1,numbers=left,numbersep=3pt,frame=lines,framesep=2pt,rulecolor=\color{gray!60},breaklines=true,fontsize=\fontsize{6.2}{7.2}\selectfont]{../../splat_editor/render/shader.py}
\subsubsection*{Khai báo mức module (dòng 25--31)}
\VerbatimInput[firstline=25,lastline=31,firstnumber=25,numbers=left,numbersep=3pt,frame=lines,framesep=2pt,rulecolor=\color{gray!60},breaklines=true,fontsize=\fontsize{6.6}{7.6}\selectfont]{../../splat_editor/render/shader.py}
\subsubsection{\texorpdfstring{\code{class ShaderError(RuntimeError)}}{ShaderError}}
{\footnotesize\textit{Dòng 34--35}}\par
\VerbatimInput[firstline=34,lastline=35,firstnumber=34,numbers=left,numbersep=3pt,frame=lines,framesep=2pt,rulecolor=\color{gray!60},breaklines=true,fontsize=\fontsize{6.6}{7.6}\selectfont]{../../splat_editor/render/shader.py}
\subsubsection{\texorpdfstring{\code{class ShaderProgram}}{ShaderProgram}}
{\footnotesize\textit{Dòng 38--110}}\par
Các phương thức: \code{__init__}, \code{load}, \code{_compile}, \code{use}, \code{location}, \code{set_int}, \code{set_float}, \code{set_vec2}, \code{set_vec3}, \code{set_vec3_array}, \code{set_float_array}, \code{set_mat3}, \code{set_mat4}.

\VerbatimInput[firstline=38,lastline=38,firstnumber=38,numbers=left,numbersep=3pt,frame=lines,framesep=2pt,rulecolor=\color{gray!60},breaklines=true,fontsize=\fontsize{6.6}{7.6}\selectfont]{../../splat_editor/render/shader.py}
\subsubsection{\texorpdfstring{\code{ShaderProgram.__init__(self, vertex_src, fragment_src=None, feedback_varyings=None)}}{ShaderProgram.__init__}}
{\footnotesize\textit{Dòng 39--58}}\par
\textbf{Mô tả trong mã nguồn.} fragment\_src=None + feedback\_varyings: chương trình chỉ chạy vertex shader và ghi kết quả vào buffer (transform feedback, xen kẽ theo thứ tự `feedback\_varyings`).

\textbf{Giải thích.} Biên dịch vertex (và fragment nếu có) shader, liên kết chương trình. Tham số \code{feedback_varyings} bật transform feedback: chương trình chỉ chạy vertex shader và ghi đầu ra vào buffer.

\VerbatimInput[firstline=39,lastline=58,firstnumber=39,numbers=left,numbersep=3pt,frame=lines,framesep=2pt,rulecolor=\color{gray!60},breaklines=true,fontsize=\fontsize{6.6}{7.6}\selectfont]{../../splat_editor/render/shader.py}
\subsubsection{\texorpdfstring{\code{ShaderProgram.load(cls, name, fragment=True, feedback_varyings=None)}}{ShaderProgram.load}}
{\footnotesize\textit{Dòng 60--65}}\par
\textbf{Mô tả trong mã nguồn.} Nạp shaders/\textless{}name\textgreater{}.vert và shaders/\textless{}name\textgreater{}.frag (fragment=False: chỉ .vert).

\VerbatimInput[firstline=60,lastline=65,firstnumber=60,numbers=left,numbersep=3pt,frame=lines,framesep=2pt,rulecolor=\color{gray!60},breaklines=true,fontsize=\fontsize{6.6}{7.6}\selectfont]{../../splat_editor/render/shader.py}
\subsubsection{\texorpdfstring{\code{ShaderProgram._compile(src, kind)}}{ShaderProgram._compile}}
{\footnotesize\textit{Dòng 67--74}}\par
\VerbatimInput[firstline=67,lastline=74,firstnumber=67,numbers=left,numbersep=3pt,frame=lines,framesep=2pt,rulecolor=\color{gray!60},breaklines=true,fontsize=\fontsize{6.6}{7.6}\selectfont]{../../splat_editor/render/shader.py}
\subsubsection{\texorpdfstring{\code{ShaderProgram.use(self)}}{ShaderProgram.use}}
{\footnotesize\textit{Dòng 76--77}}\par
\VerbatimInput[firstline=76,lastline=77,firstnumber=76,numbers=left,numbersep=3pt,frame=lines,framesep=2pt,rulecolor=\color{gray!60},breaklines=true,fontsize=\fontsize{6.6}{7.6}\selectfont]{../../splat_editor/render/shader.py}
\subsubsection{\texorpdfstring{\code{ShaderProgram.location(self, name)}}{ShaderProgram.location}}
{\footnotesize\textit{Dòng 79--82}}\par
\VerbatimInput[firstline=79,lastline=82,firstnumber=79,numbers=left,numbersep=3pt,frame=lines,framesep=2pt,rulecolor=\color{gray!60},breaklines=true,fontsize=\fontsize{6.6}{7.6}\selectfont]{../../splat_editor/render/shader.py}
\subsubsection{\texorpdfstring{\code{ShaderProgram.set_int(self, name, value)}}{ShaderProgram.set_int}}
{\footnotesize\textit{Dòng 85--86}}\par
\VerbatimInput[firstline=85,lastline=86,firstnumber=85,numbers=left,numbersep=3pt,frame=lines,framesep=2pt,rulecolor=\color{gray!60},breaklines=true,fontsize=\fontsize{6.6}{7.6}\selectfont]{../../splat_editor/render/shader.py}
\subsubsection{\texorpdfstring{\code{ShaderProgram.set_float(self, name, value)}}{ShaderProgram.set_float}}
{\footnotesize\textit{Dòng 88--89}}\par
\VerbatimInput[firstline=88,lastline=89,firstnumber=88,numbers=left,numbersep=3pt,frame=lines,framesep=2pt,rulecolor=\color{gray!60},breaklines=true,fontsize=\fontsize{6.6}{7.6}\selectfont]{../../splat_editor/render/shader.py}
\subsubsection{\texorpdfstring{\code{ShaderProgram.set_vec2(self, name, x, y)}}{ShaderProgram.set_vec2}}
{\footnotesize\textit{Dòng 91--92}}\par
\VerbatimInput[firstline=91,lastline=92,firstnumber=91,numbers=left,numbersep=3pt,frame=lines,framesep=2pt,rulecolor=\color{gray!60},breaklines=true,fontsize=\fontsize{6.6}{7.6}\selectfont]{../../splat_editor/render/shader.py}
\subsubsection{\texorpdfstring{\code{ShaderProgram.set_vec3(self, name, v)}}{ShaderProgram.set_vec3}}
{\footnotesize\textit{Dòng 94--95}}\par
\VerbatimInput[firstline=94,lastline=95,firstnumber=94,numbers=left,numbersep=3pt,frame=lines,framesep=2pt,rulecolor=\color{gray!60},breaklines=true,fontsize=\fontsize{6.6}{7.6}\selectfont]{../../splat_editor/render/shader.py}
\subsubsection{\texorpdfstring{\code{ShaderProgram.set_vec3_array(self, name, values)}}{ShaderProgram.set_vec3_array}}
{\footnotesize\textit{Dòng 97--99}}\par
\VerbatimInput[firstline=97,lastline=99,firstnumber=97,numbers=left,numbersep=3pt,frame=lines,framesep=2pt,rulecolor=\color{gray!60},breaklines=true,fontsize=\fontsize{6.6}{7.6}\selectfont]{../../splat_editor/render/shader.py}
\subsubsection{\texorpdfstring{\code{ShaderProgram.set_float_array(self, name, values)}}{ShaderProgram.set_float_array}}
{\footnotesize\textit{Dòng 101--103}}\par
\VerbatimInput[firstline=101,lastline=103,firstnumber=101,numbers=left,numbersep=3pt,frame=lines,framesep=2pt,rulecolor=\color{gray!60},breaklines=true,fontsize=\fontsize{6.6}{7.6}\selectfont]{../../splat_editor/render/shader.py}
\subsubsection{\texorpdfstring{\code{ShaderProgram.set_mat3(self, name, m)}}{ShaderProgram.set_mat3}}
{\footnotesize\textit{Dòng 105--107}}\par
\textbf{Mô tả trong mã nguồn.} m: ma trận numpy theo hàng (row-major), tự chuyển vị khi gửi lên.

\VerbatimInput[firstline=105,lastline=107,firstnumber=105,numbers=left,numbersep=3pt,frame=lines,framesep=2pt,rulecolor=\color{gray!60},breaklines=true,fontsize=\fontsize{6.6}{7.6}\selectfont]{../../splat_editor/render/shader.py}
\subsubsection{\texorpdfstring{\code{ShaderProgram.set_mat4(self, name, m)}}{ShaderProgram.set_mat4}}
{\footnotesize\textit{Dòng 109--110}}\par
\VerbatimInput[firstline=109,lastline=110,firstnumber=109,numbers=left,numbersep=3pt,frame=lines,framesep=2pt,rulecolor=\color{gray!60},breaklines=true,fontsize=\fontsize{6.6}{7.6}\selectfont]{../../splat_editor/render/shader.py}
\subsection{\texorpdfstring{\code{splat_editor/render/camera.py}}{splat_editor/render/camera.py}}
{\footnotesize 123 dòng}\par
\textbf{Vai trò.} Camera quỹ đạo (\code{OrbitCamera}): luôn nhìn vào \code{target}, quay quanh bằng yaw/pitch. Các thao tác: \code{orbit}, \code{pan}, \code{look}, \code{fly}, \code{frame}.
Toán học: xem mục \textit{Cơ sở toán} ngay bên dưới (\S\ref{s:cam-orbit}; ma trận \code{look_at}/\code{perspective} ở \S\ref{s:cam}).

\paragraph{Mô tả đầu file và cách dùng.} Phần chú thích đầu file do tác giả viết, nêu ý tưởng và ví dụ sử dụng:
\VerbatimInput[firstline=1,lastline=15,firstnumber=1,numbers=left,numbersep=3pt,frame=lines,framesep=2pt,rulecolor=\color{gray!60},breaklines=true,fontsize=\fontsize{6.2}{7.2}\selectfont]{../../splat_editor/render/camera.py}
%% ---- merged math: camera_orbit
\subsubsection{Cơ sở toán: Camera quỹ đạo}\label{s:cam-orbit}
Các ma trận \code{look_at}, \code{perspective}, \code{orthographic} được giải thích ở mục \ref{s:cam} (file \code{core/math3d.py}). Dưới đây là toán của \code{OrbitCamera}.
\begin{itemize}
\item \textbf{Vị trí mắt} (tọa độ cầu): $e=t+d\,(\cos\varphi\sin\theta,\;\sin\varphi,\;\cos\varphi\cos\theta)$ với $\theta$ là yaw, $\varphi$ là pitch, $d$ là distance (\code{position}).
\item Pitch kẹp $\pm1.55$ rad ($\approx\pm88.8^\circ$) để $f\times\mathrm{up}$ không suy biến (\code{PITCH_LIMIT}).
\item \code{orbit}: yaw $-=dx\cdot0.005$, pitch $+=dy\cdot0.005$ (rad/pixel).
\item \textbf{Pan} (\code{pan}): chiều cao nhìn thấy tại khoảng cách $d$ là $2d\tan(\mathrm{fov}/2)$, trải trên $H$ pixel, nên $\mathrm{units\_per\_pixel}=2d\tan(\mathrm{fov}/2)/H$. Khi đó $\mathrm{target}\mathrel{+}=(-dx\cdot\mathrm{right}+dy\cdot\mathrm{up})\cdot\mathrm{units\_per\_pixel}$, với right, up là hàng 0 và 1 của $V$.
\item \textbf{Look} (\code{look}): giữ mắt cố định, đổi yaw/pitch, rồi \code{target = eye - (position_moi - target)} để mắt không dịch.
\item \textbf{Bay} (\code{fly}): hướng ngang $=\mathrm{forward}\cdot(-\sin\theta,0,-\cos\theta)+\mathrm{right}\cdot(\cos\theta,0,-\sin\theta)$, chuẩn hóa (đi chéo không nhanh hơn), $y=\mathrm{up}$. Vận tốc mong muốn $v^*=\text{hướng}\cdot\text{tốc độ}\cdot(2.5\text{ nếu chạy})$. \textbf{Lọc mũ:} $v\leftarrow v+(v^*-v)\min(10\Delta t,1)$; $\mathrm{target}\mathrel{+}=v\Delta t$ (tích phân Euler). Dừng hẳn khi $|v|<10^{-3}\cdot\text{tốc độ}$.
\item \textbf{Frame} (\code{frame}): tâm là median các điểm; bán kính $R$ là phân vị 90\% khoảng cách tới tâm (bỏ nhiễu); $d=1.3R/\sin(\mathrm{fov}/2)$ (hình cầu bán kính $R$ vừa khung dọc: $R=d\sin(\mathrm{fov}/2)$); tốc độ bay $=0.5R$.
\item \textbf{Ngôi thứ ba:} mode 1: $\mathrm{cam}=\mathrm{eye}-\mathrm{forward}\cdot D$; mode 2: $\mathrm{cam}=\mathrm{eye}+\mathrm{forward}\cdot D$ nhìn ngược lại. $D$ co lại nếu tia mắt$\to$camera chạm voxel (\code{segment_blocked}, mục \ref{s:phys}).
\item $\mathrm{forward}()=(\mathrm{target}-\mathrm{position})/\mathrm{distance}$.
\end{itemize}

\subsubsection*{Khai báo mức module (dòng 16--20)}
\VerbatimInput[firstline=16,lastline=20,firstnumber=16,numbers=left,numbersep=3pt,frame=lines,framesep=2pt,rulecolor=\color{gray!60},breaklines=true,fontsize=\fontsize{6.6}{7.6}\selectfont]{../../splat_editor/render/camera.py}
\subsubsection{\texorpdfstring{\code{class OrbitCamera}}{OrbitCamera}}
{\footnotesize\textit{Dòng 23--123}}\par
Các phương thức: \code{__init__}, \code{reset}, \code{position}, \code{view_matrix}, \code{projection_matrix}, \code{orbit}, \code{pan}, \code{look}, \code{forward}, \code{set_eye}, \code{fly}, \code{frame}.

\VerbatimInput[firstline=23,lastline=27,firstnumber=23,numbers=left,numbersep=3pt,frame=lines,framesep=2pt,rulecolor=\color{gray!60},breaklines=true,fontsize=\fontsize{6.6}{7.6}\selectfont]{../../splat_editor/render/camera.py}
\subsubsection{\texorpdfstring{\code{OrbitCamera.__init__(self)}}{OrbitCamera.__init__}}
{\footnotesize\textit{Dòng 29--31}}\par
\VerbatimInput[firstline=29,lastline=31,firstnumber=29,numbers=left,numbersep=3pt,frame=lines,framesep=2pt,rulecolor=\color{gray!60},breaklines=true,fontsize=\fontsize{6.6}{7.6}\selectfont]{../../splat_editor/render/camera.py}
\subsubsection{\texorpdfstring{\code{OrbitCamera.reset(self)}}{OrbitCamera.reset}}
{\footnotesize\textit{Dòng 33--41}}\par
\VerbatimInput[firstline=33,lastline=41,firstnumber=33,numbers=left,numbersep=3pt,frame=lines,framesep=2pt,rulecolor=\color{gray!60},breaklines=true,fontsize=\fontsize{6.6}{7.6}\selectfont]{../../splat_editor/render/camera.py}
\subsubsection{\texorpdfstring{\code{OrbitCamera.position(self)}}{OrbitCamera.position}}
{\footnotesize\textit{Dòng 44--47}}\par
\VerbatimInput[firstline=44,lastline=47,firstnumber=44,numbers=left,numbersep=3pt,frame=lines,framesep=2pt,rulecolor=\color{gray!60},breaklines=true,fontsize=\fontsize{6.6}{7.6}\selectfont]{../../splat_editor/render/camera.py}
\subsubsection{\texorpdfstring{\code{OrbitCamera.view_matrix(self)}}{OrbitCamera.view_matrix}}
{\footnotesize\textit{Dòng 49--60}}\par
\textbf{Giải thích.} Gọi \code{look_at(position, target, up)}; vị trí mắt tính từ tọa độ cầu $e=t+d(\cos\varphi\sin\theta,\sin\varphi,\cos\varphi\cos\theta)$.

\textbf{Công thức.}
\[
e=t+d\begin{bmatrix}\cos\varphi\sin\theta\\\sin\varphi\\\cos\varphi\cos\theta\end{bmatrix},\qquad
V=\mathrm{look\_at}(e,t,\hat y),\qquad |\varphi|\le1.55\ \mathrm{rad}
\]


\VerbatimInput[firstline=49,lastline=60,firstnumber=49,numbers=left,numbersep=3pt,frame=lines,framesep=2pt,rulecolor=\color{gray!60},breaklines=true,fontsize=\fontsize{6.6}{7.6}\selectfont]{../../splat_editor/render/camera.py}
\subsubsection{\texorpdfstring{\code{OrbitCamera.projection_matrix(self, aspect)}}{OrbitCamera.projection_matrix}}
{\footnotesize\textit{Dòng 62--64}}\par
\textbf{Giải thích.} Chiếu phối cảnh với \code{near = distance*0.001}, \code{far = near*10^6} để zoom từ rất gần đến rất xa mà không cắt.

\textbf{Công thức.}
\[
n=0.001\,d,\qquad f=10^{6}\,n
\]


\VerbatimInput[firstline=62,lastline=64,firstnumber=62,numbers=left,numbersep=3pt,frame=lines,framesep=2pt,rulecolor=\color{gray!60},breaklines=true,fontsize=\fontsize{6.6}{7.6}\selectfont]{../../splat_editor/render/camera.py}
\subsubsection{\texorpdfstring{\code{OrbitCamera.orbit(self, dx, dy)}}{OrbitCamera.orbit}}
{\footnotesize\textit{Dòng 67--70}}\par
\textbf{Mô tả trong mã nguồn.} dx, dy: số pixel chuột di chuyển.

\textbf{Giải thích.} Công thức cốt lõi của hàm này:
\[
\theta\leftarrow\theta-0.005\,dx,\qquad \varphi\leftarrow\mathrm{clip}(\varphi+0.005\,dy,\ \pm1.55)
\]


\VerbatimInput[firstline=67,lastline=70,firstnumber=67,numbers=left,numbersep=3pt,frame=lines,framesep=2pt,rulecolor=\color{gray!60},breaklines=true,fontsize=\fontsize{6.6}{7.6}\selectfont]{../../splat_editor/render/camera.py}
\subsubsection{\texorpdfstring{\code{OrbitCamera.pan(self, dx, dy, viewport_height)}}{OrbitCamera.pan}}
{\footnotesize\textit{Dòng 72--77}}\par
\textbf{Giải thích.} Kéo camera: 1 pixel ứng với $2d\tan(fov/2)/H$ đơn vị thế giới tại khoảng cách $d$, di chuyển theo hai trục right/up của view.

\textbf{Công thức.}
\[
u_{px}=\frac{2d\tan(\mathrm{fov}/2)}{H},\qquad t\leftarrow t+\big(-dx\,\hat r+dy\,\hat u\big)\,u_{px}
\]


\VerbatimInput[firstline=72,lastline=77,firstnumber=72,numbers=left,numbersep=3pt,frame=lines,framesep=2pt,rulecolor=\color{gray!60},breaklines=true,fontsize=\fontsize{6.6}{7.6}\selectfont]{../../splat_editor/render/camera.py}
\subsubsection{\texorpdfstring{\code{OrbitCamera.look(self, dx, dy)}}{OrbitCamera.look}}
{\footnotesize\textit{Dòng 79--84}}\par
\textbf{Mô tả trong mã nguồn.} Nhìn quanh tại chỗ (kiểu góc nhìn thứ nhất): giữ nguyên vị trí, xoay hướng nhìn.

\VerbatimInput[firstline=79,lastline=84,firstnumber=79,numbers=left,numbersep=3pt,frame=lines,framesep=2pt,rulecolor=\color{gray!60},breaklines=true,fontsize=\fontsize{6.6}{7.6}\selectfont]{../../splat_editor/render/camera.py}
\subsubsection{\texorpdfstring{\code{OrbitCamera.forward(self)}}{OrbitCamera.forward}}
{\footnotesize\textit{Dòng 86--88}}\par
\textbf{Mô tả trong mã nguồn.} Vector đơn vị hướng nhìn.

\VerbatimInput[firstline=86,lastline=88,firstnumber=86,numbers=left,numbersep=3pt,frame=lines,framesep=2pt,rulecolor=\color{gray!60},breaklines=true,fontsize=\fontsize{6.6}{7.6}\selectfont]{../../splat_editor/render/camera.py}
\subsubsection{\texorpdfstring{\code{OrbitCamera.set_eye(self, eye)}}{OrbitCamera.set_eye}}
{\footnotesize\textit{Dòng 90--92}}\par
\textbf{Mô tả trong mã nguồn.} Đặt mắt camera tại `eye`, giữ nguyên hướng nhìn (góc nhìn thứ nhất của chế độ sinh tồn).

\VerbatimInput[firstline=90,lastline=92,firstnumber=90,numbers=left,numbersep=3pt,frame=lines,framesep=2pt,rulecolor=\color{gray!60},breaklines=true,fontsize=\fontsize{6.6}{7.6}\selectfont]{../../splat_editor/render/camera.py}
\subsubsection{\texorpdfstring{\code{OrbitCamera.fly(self, forward, right, up, dt, sprint=False)}}{OrbitCamera.fly}}
{\footnotesize\textit{Dòng 94--109}}\par
\textbf{Mô tả trong mã nguồn.} Bay kiểu Minecraft sáng tạo: W/S A/D đi ngang theo hướng nhìn, lên / xuống theo trục Y.

Vận tốc tiến dần về vận tốc mong muốn nên có quán tính nhẹ khi bắt đầu / dừng.

\textbf{Giải thích.} Bay kiểu Minecraft: hướng ngang chuẩn hóa, lọc mũ vận tốc $v\leftarrow v+(v^*-v)\min(10\Delta t,1)$ rồi tích phân Euler.

\textbf{Công thức.}
\[
v^*=\hat w\cdot\mathrm{speed}\cdot\begin{cases}2.5&\text{chạy}\\1&\text{đi}\end{cases},\quad
v\leftarrow v+(v^*-v)\min(10\Delta t,1),\quad t\leftarrow t+v\Delta t
\]
\[
\hat w=\mathrm{normalize}\big(f_{in}\,(-\sin\theta,0,-\cos\theta)+r_{in}\,(\cos\theta,0,-\sin\theta)\big)
\]


\VerbatimInput[firstline=94,lastline=109,firstnumber=94,numbers=left,numbersep=3pt,frame=lines,framesep=2pt,rulecolor=\color{gray!60},breaklines=true,fontsize=\fontsize{6.6}{7.6}\selectfont]{../../splat_editor/render/camera.py}
\subsubsection{\texorpdfstring{\code{OrbitCamera.frame(self, points)}}{OrbitCamera.frame}}
{\footnotesize\textit{Dòng 111--123}}\par
\textbf{Mô tả trong mã nguồn.} Đặt camera để nhìn thấy toàn bộ đám điểm.

\textbf{Giải thích.} Đặt camera nhìn toàn bộ đám điểm: tâm là median, bán kính là phân vị 90\%, khoảng cách $1.3R/\sin(fov/2)$.

\textbf{Công thức.}
\[
t=\mathrm{median}(p_i),\quad R=P_{90}\big(\|p_i-t\|\big),\quad d=\frac{1.3\,R}{\sin(\mathrm{fov}/2)},\quad \mathrm{speed}=0.5R
\]


\VerbatimInput[firstline=111,lastline=123,firstnumber=111,numbers=left,numbersep=3pt,frame=lines,framesep=2pt,rulecolor=\color{gray!60},breaklines=true,fontsize=\fontsize{6.6}{7.6}\selectfont]{../../splat_editor/render/camera.py}
\subsection{\texorpdfstring{\code{splat_editor/render/sorter.py}}{splat_editor/render/sorter.py}}
{\footnotesize 140 dòng}\par
\textbf{Vai trò.} Luồng nền \code{DepthSorter}: nhận yêu cầu (ma trận camera và ngân sách), chọn gaussian theo LOD rồi sắp xếp \textbf{xa $\to$ gần} bằng khóa uint16.
Luồng chính không bao giờ phải chờ; trong lúc chờ renderer dùng thứ tự cũ. Toán học: xem mục \textit{Cơ sở toán} ngay bên dưới (\S\ref{s:sort}).

\paragraph{Mô tả đầu file và cách dùng.} Phần chú thích đầu file do tác giả viết, nêu ý tưởng và ví dụ sử dụng:
\VerbatimInput[firstline=1,lastline=28,firstnumber=1,numbers=left,numbersep=3pt,frame=lines,framesep=2pt,rulecolor=\color{gray!60},breaklines=true,fontsize=\fontsize{6.2}{7.2}\selectfont]{../../splat_editor/render/sorter.py}
%% ---- merged math: s:sort
\subsubsection{Cơ sở toán: Sắp xếp theo độ sâu}\label{s:sort}
\begin{itemize}
\item Độ sâu camera của tâm: $z_{cam}=(MV\cdot M_{layer})[2,:]\cdot(p,1)$. Camera nhìn $-z$ nên $z_{cam}<0$ là phía trước, càng âm càng xa. Giữ $\mathrm{depth}<0$.
\item Khóa 16 bit: $\mathrm{key}=\lfloor(\mathrm{depth}-lo)\cdot65535/(hi-lo)\rfloor$ (lượng tử đều trên $[lo,hi]$); \code{argsort(uint16, stable)} là radix sort $O(N)$. Depth (âm) tăng dần tương ứng xa $\to$ gần. Lượng tử 16 bit đủ vì chỉ cần thứ tự gần đúng.
\item Chạy ở luồng nền; kết quả cũ bị bỏ nếu \code{version} đã đổi.
\end{itemize}

\subsubsection*{Khai báo mức module (dòng 29--33)}
\VerbatimInput[firstline=29,lastline=33,firstnumber=29,numbers=left,numbersep=3pt,frame=lines,framesep=2pt,rulecolor=\color{gray!60},breaklines=true,fontsize=\fontsize{6.6}{7.6}\selectfont]{../../splat_editor/render/sorter.py}
\subsubsection{\texorpdfstring{\code{class DepthSorter}}{DepthSorter}}
{\footnotesize\textit{Dòng 36--140}}\par
Các phương thức: \code{__init__}, \code{set_splats}, \code{request}, \code{take_result}, \code{close}, \code{_run}, \code{_select}, \code{_depth_rows}, \code{sort}.

\VerbatimInput[firstline=36,lastline=36,firstnumber=36,numbers=left,numbersep=3pt,frame=lines,framesep=2pt,rulecolor=\color{gray!60},breaklines=true,fontsize=\fontsize{6.6}{7.6}\selectfont]{../../splat_editor/render/sorter.py}
\subsubsection{\texorpdfstring{\code{DepthSorter.__init__(self)}}{DepthSorter.__init__}}
{\footnotesize\textit{Dòng 37--47}}\par
\VerbatimInput[firstline=37,lastline=47,firstnumber=37,numbers=left,numbersep=3pt,frame=lines,framesep=2pt,rulecolor=\color{gray!60},breaklines=true,fontsize=\fontsize{6.6}{7.6}\selectfont]{../../splat_editor/render/sorter.py}
\subsubsection{\texorpdfstring{\code{DepthSorter.set_splats(self, lod, state=None)}}{DepthSorter.set_splats}}
{\footnotesize\textit{Dòng 49--58}}\par
\textbf{Mô tả trong mã nguồn.} lod: LodIndex (tâm gaussian theo thứ tự vị trí: lod.positions). state: mảng trạng thái (uint8) theo vị trí, renderer sửa tại chỗ; None: không bỏ gì.

\VerbatimInput[firstline=49,lastline=58,firstnumber=49,numbers=left,numbersep=3pt,frame=lines,framesep=2pt,rulecolor=\color{gray!60},breaklines=true,fontsize=\fontsize{6.6}{7.6}\selectfont]{../../splat_editor/render/sorter.py}
\subsubsection{\texorpdfstring{\code{DepthSorter.request(self, model_view, projection, budget, transforms=None)}}{DepthSorter.request}}
{\footnotesize\textit{Dòng 60--66}}\par
\textbf{Mô tả trong mã nguồn.} Yêu cầu chọn + sắp xếp. Yêu cầu chưa xử lý bị thay bằng yêu cầu mới nhất. budget: số gaussian tối đa (None: không giới hạn); transforms: ma trận thế giới mỗi lớp.

\textbf{Giải thích.} Gửi yêu cầu cho luồng nền; yêu cầu chưa xử lý bị thay bằng yêu cầu mới nhất (không xếp hàng), nên luồng nền luôn làm việc với camera mới nhất.

\VerbatimInput[firstline=60,lastline=66,firstnumber=60,numbers=left,numbersep=3pt,frame=lines,framesep=2pt,rulecolor=\color{gray!60},breaklines=true,fontsize=\fontsize{6.6}{7.6}\selectfont]{../../splat_editor/render/sorter.py}
\subsubsection{\texorpdfstring{\code{DepthSorter.take_result(self)}}{DepthSorter.take_result}}
{\footnotesize\textit{Dòng 68--74}}\par
\textbf{Mô tả trong mã nguồn.} (slots, ranks) nếu có kết quả mới, ngược lại None.

\VerbatimInput[firstline=68,lastline=74,firstnumber=68,numbers=left,numbersep=3pt,frame=lines,framesep=2pt,rulecolor=\color{gray!60},breaklines=true,fontsize=\fontsize{6.6}{7.6}\selectfont]{../../splat_editor/render/sorter.py}
\subsubsection{\texorpdfstring{\code{DepthSorter.close(self)}}{DepthSorter.close}}
{\footnotesize\textit{Dòng 76--79}}\par
\VerbatimInput[firstline=76,lastline=79,firstnumber=76,numbers=left,numbersep=3pt,frame=lines,framesep=2pt,rulecolor=\color{gray!60},breaklines=true,fontsize=\fontsize{6.6}{7.6}\selectfont]{../../splat_editor/render/sorter.py}
\subsubsection{\texorpdfstring{\code{DepthSorter._run(self)}}{DepthSorter._run}}
{\footnotesize\textit{Dòng 81--94}}\par
\VerbatimInput[firstline=81,lastline=94,firstnumber=81,numbers=left,numbersep=3pt,frame=lines,framesep=2pt,rulecolor=\color{gray!60},breaklines=true,fontsize=\fontsize{6.6}{7.6}\selectfont]{../../splat_editor/render/sorter.py}
\subsubsection{\texorpdfstring{\code{DepthSorter._select(lod, state, model_view, projection, budget, transforms)}}{DepthSorter._select}}
{\footnotesize\textit{Dòng 96--108}}\par
\textbf{Mô tả trong mã nguồn.} lod.select rồi bỏ gaussian đã xoá. Nếu bị giới hạn bởi budget mà bỏ đi nhiều, chọn lại với budget nới theo tỉ lệ còn lại để vẫn vẽ gần đủ budget gaussian.

\VerbatimInput[firstline=96,lastline=108,firstnumber=96,numbers=left,numbersep=3pt,frame=lines,framesep=2pt,rulecolor=\color{gray!60},breaklines=true,fontsize=\fontsize{6.6}{7.6}\selectfont]{../../splat_editor/render/sorter.py}
\subsubsection{\texorpdfstring{\code{DepthSorter._depth_rows(model_view, transforms, lod, slots)}}{DepthSorter._depth_rows}}
{\footnotesize\textit{Dòng 110--118}}\par
\textbf{Mô tả trong mã nguồn.} Hàng z của ma trận model-view * thế giới cho từng đoạn `slots` (tăng dần) theo lớp: [(số phần tử, hàng (4,))].

\VerbatimInput[firstline=110,lastline=118,firstnumber=110,numbers=left,numbersep=3pt,frame=lines,framesep=2pt,rulecolor=\color{gray!60},breaklines=true,fontsize=\fontsize{6.6}{7.6}\selectfont]{../../splat_editor/render/sorter.py}
\subsubsection{\texorpdfstring{\code{DepthSorter.sort(positions, depth_rows)}}{DepthSorter.sort}}
{\footnotesize\textit{Dòng 120--140}}\par
\textbf{Mô tả trong mã nguồn.} Chỉ số các tâm nằm trước camera, xếp xa -\textgreater{} gần. depth\_rows: ma trận model-view 4x4, hoặc [(số phần tử, hàng z)] cho từng đoạn liên tiếp.

\textbf{Giải thích.} Tính độ sâu camera $z=(MV\cdot M)[2,:]\cdot(p,1)$, giữ điểm trước camera ($z<0$), lượng tử thành uint16 và \code{argsort} ổn định (radix) để ra thứ tự xa $\to$ gần.

\textbf{Công thức.}
\[
z_{cam}=(MV\,M_{layer})_{[2,:]}\cdot(p,1)\ \ (<0\text{ là trước camera}),\qquad
\mathrm{key}=\Big\lfloor\frac{z-lo}{hi-lo}\cdot65535\Big\rfloor
\]
Sắp \code{key} tăng dần $=$ xa $\to$ gần (radix sort $O(N)$).


\VerbatimInput[firstline=120,lastline=140,firstnumber=120,numbers=left,numbersep=3pt,frame=lines,framesep=2pt,rulecolor=\color{gray!60},breaklines=true,fontsize=\fontsize{6.6}{7.6}\selectfont]{../../splat_editor/render/sorter.py}
\subsection{\texorpdfstring{\code{splat_editor/render/splat_renderer.py}}{splat_editor/render/splat_renderer.py}}
{\footnotesize 504 dòng}\par
\textbf{Vai trò.} Trái tim của hiển thị: \code{SplatRenderer} đóng gói gaussian thành 32 byte, đưa lên GPU, chạy pass 1 (transform feedback) và pass 2 (vẽ hình vuông, trộn alpha),
quản lý vùng chọn/xóa, bản đồ bóng và tự chỉnh ngân sách. Toán học: xem mục \textit{Cơ sở toán} ngay bên dưới (\S\ref{s:ewa}, \S\ref{s:blend}).

\paragraph{Mô tả đầu file và cách dùng.} Phần chú thích đầu file do tác giả viết, nêu ý tưởng và ví dụ sử dụng:
\VerbatimInput[firstline=1,lastline=34,firstnumber=1,numbers=left,numbersep=3pt,frame=lines,framesep=2pt,rulecolor=\color{gray!60},breaklines=true,fontsize=\fontsize{6.2}{7.2}\selectfont]{../../splat_editor/render/splat_renderer.py}
%% ---- merged math: s:budget
\subsubsection{Cơ sở toán: tự chỉnh ngân sách (\texorpdfstring{\code{SplatRenderer.adapt_budget}}{SplatRenderer.adapt_budget})}\label{s:budget}
Đo thời gian GPU bằng query. $\mathrm{cost}=\mathrm{median}(\mathrm{ms}/\mathrm{count})$; $\mathrm{wanted}=\mathrm{target\_ms}/\mathrm{cost}$. Nếu lệch dưới 10\% thì giữ nguyên; ngược lại $\mathrm{step}=\mathrm{clip}(\mathrm{wanted}/\mathrm{current},0.75,1.33)^{1/2}$ (căn bậc hai nên đổi tối đa khoảng 15\%/lần, làm trơn để tránh dao động); $\mathrm{budget}=\mathrm{clip}(\mathrm{current}\cdot\mathrm{step},\,\mathrm{MIN\_BUDGET},\,\mathrm{count})$.

%% ---- merged math: quantization
\subsubsection{Cơ sở toán: Lượng tử hóa và đóng gói GPU}\label{s:quant}
Phần này giải thích cách \code{SplatRenderer.pack} nén dữ liệu gaussian, và cách shader \code{splat_prep.vert} giải nén (xem thêm mục shader tương ứng).
\begin{itemize}
\item \textbf{Half-float (IEEE-754 16 bit)}, \code{halfToFloat}: $e$ là bit 10--14, $m$ là bit 0--9, bit 15 là dấu. Nếu $e=0$ (subnormal) thì $v=m\cdot2^{-24}$; nếu $e\ne0$ thì $v=(m+1024)\cdot2^{e-25}=(1+m/1024)\cdot2^{e-15}$. Dùng cho màu, $\alpha$, scale để tiết kiệm bộ nhớ.
\item \textbf{SH 8 bit:} $\mathrm{scale\_sh}=\mathrm{percentile}_{99.9}(|\mathrm{coeff}|)/127$; lưu $u=\mathrm{clip}(\mathrm{round}(\mathrm{coeff}/\mathrm{scale\_sh})+128,1,255)$; giải $\mathrm{coeff}=(u-128)\cdot\mathrm{scale\_sh}$ (shader: \code{(texel*255 - 128) * uShScale}).
\item \textbf{Màu RGBA8} (\code{packColor}): $u=\mathrm{clamp}(c,0,1)\cdot255+0.5$ (làm tròn).
\item 1 gaussian $=8$ uint32 (2 texel RGBA32UI): xyz (float), màu rg/b+$\alpha$, scale xyz (half)+layer id, quaternion $4\times8$ bit.
\end{itemize}

\subsubsection*{Khai báo mức module (dòng 35--54)}
\VerbatimInput[firstline=35,lastline=54,firstnumber=35,numbers=left,numbersep=3pt,frame=lines,framesep=2pt,rulecolor=\color{gray!60},breaklines=true,fontsize=\fontsize{6.6}{7.6}\selectfont]{../../splat_editor/render/splat_renderer.py}
\subsubsection{\texorpdfstring{\code{class SplatRenderer}}{SplatRenderer}}
{\footnotesize\textit{Dòng 57--504}}\par
Các phương thức: \code{__init__}, \code{pack}, \code{_upload_chunked}, \code{set_splats}, \code{_upload_sh}, \code{_upload_state}, \code{_set_state_bit}, \code{set_selection}, \code{set_deleted}, \code{update_selection}, \code{set_transforms}, \code{set_lighting}, \code{_upload_order}, \code{_update_order}, \code{_prepare}, \code{_set_light_uniforms}, \code{_create_shadow_map}, \code{_render_shadow_map}, \code{_collect_timings}, \code{adapt_budget}, \code{draw}, \code{_draw_splats}.

\VerbatimInput[firstline=57,lastline=58,firstnumber=57,numbers=left,numbersep=3pt,frame=lines,framesep=2pt,rulecolor=\color{gray!60},breaklines=true,fontsize=\fontsize{6.6}{7.6}\selectfont]{../../splat_editor/render/splat_renderer.py}
\subsubsection{\texorpdfstring{\code{SplatRenderer.__init__(self)}}{SplatRenderer.__init__}}
{\footnotesize\textit{Dòng 60--127}}\par
\VerbatimInput[firstline=60,lastline=127,firstnumber=60,numbers=left,numbersep=3pt,frame=lines,framesep=2pt,rulecolor=\color{gray!60},breaklines=true,fontsize=\fontsize{6.6}{7.6}\selectfont]{../../splat_editor/render/splat_renderer.py}
\subsubsection{\texorpdfstring{\code{SplatRenderer.pack(splats, index=slice(None), layers=0)}}{SplatRenderer.pack}}
{\footnotesize\textit{Dòng 130--147}}\par
\textbf{Mô tả trong mã nguồn.} Gaussian `index` -\textgreater{} mảng (n,8) uint32 = 2 texel RGBA32UI / gaussian: texel 0: x y z (float32, toạ độ file của lớp), r g (half) texel 1: b alpha (half), sx sy (half), sz (half) | số thứ tự lớp (16 bit), quaternion w x y z (4 x snorm8).

\textbf{Giải thích.} Đóng gói gaussian thành 8 uint32 (2 texel RGBA32UI): tâm float32; màu, alpha và scale half-float; quaternion 4$\times$8 bit; id layer. Tổng 32 byte/gaussian.

\textbf{Công thức.}
\[
32\ \mathrm{byte}=\underbrace{12}_{\text{tâm float32}}+\underbrace{6}_{\text{rgb half}}+\underbrace{2}_{\alpha\ \text{half}}+\underbrace{6}_{\text{scale half}}+\underbrace{2}_{\text{layer}}+\underbrace{4}_{\text{quaternion }4\times8\mathrm{bit}}
\]
\[
u_q=\mathrm{clip}\big(\mathrm{round}(127\,q)+128,\ 1,\ 255\big),\qquad q=\mathrm{normalize}\Big(\frac{u_q-128}{127}\Big)
\]


\VerbatimInput[firstline=130,lastline=147,firstnumber=130,numbers=left,numbersep=3pt,frame=lines,framesep=2pt,rulecolor=\color{gray!60},breaklines=true,fontsize=\fontsize{6.6}{7.6}\selectfont]{../../splat_editor/render/splat_renderer.py}
\subsubsection{\texorpdfstring{\code{SplatRenderer._upload_chunked(buffer, count, bytes_per_item, make_chunk)}}{SplatRenderer._upload_chunked}}
{\footnotesize\textit{Dòng 149--169}}\par
\textbf{Mô tả trong mã nguồn.} Cấp buffer rồi gửi dữ liệu từng đoạn (make\_chunk(start, end) -\textgreater{} mảng numpy). Các đoạn được chuẩn bị song song ở nhiều luồng (numpy nhả GIL), luồng chính chỉ gửi lên GPU.

\VerbatimInput[firstline=149,lastline=169,firstnumber=149,numbers=left,numbersep=3pt,frame=lines,framesep=2pt,rulecolor=\color{gray!60},breaklines=true,fontsize=\fontsize{6.6}{7.6}\selectfont]{../../splat_editor/render/splat_renderer.py}
\subsubsection{\texorpdfstring{\code{SplatRenderer.set_splats(self, splats, lod=None)}}{SplatRenderer.set_splats}}
{\footnotesize\textit{Dòng 172--197}}\par
\textbf{Mô tả trong mã nguồn.} lod: LodIndex của splats (None: tự dựng).

\textbf{Giải thích.} Gửi dữ liệu lên GPU theo thứ tự Morton của \code{LodIndex}, từng đoạn \code{UPLOAD_CHUNK}; nếu SH vượt \code{SH_VRAM_BUDGET} thì tự hạ bậc.

\VerbatimInput[firstline=172,lastline=197,firstnumber=172,numbers=left,numbersep=3pt,frame=lines,framesep=2pt,rulecolor=\color{gray!60},breaklines=true,fontsize=\fontsize{6.6}{7.6}\selectfont]{../../splat_editor/render/splat_renderer.py}
\subsubsection{\texorpdfstring{\code{SplatRenderer._upload_sh(self, splats, order)}}{SplatRenderer._upload_sh}}
{\footnotesize\textit{Dòng 199--223}}\par
\textbf{Mô tả trong mã nguồn.} (N,K,3) hệ số SH -\textgreater{} K' texel RGBA8 / gaussian (lượng tử 8 bit, thành phần A bỏ trống). K' \textless{}= K: hạ bậc SH nếu không đủ SH\_VRAM\_BUDGET. Hệ số giữ ở toạ độ file: shader xoay hướng nhìn ngược về toạ độ file thay vì xoay hệ số.

\textbf{Giải thích.} Lượng tử hóa SH về 8 bit: $u=\mathrm{clip}(\mathrm{round}(c/scale\_sh)+128,1,255)$ với \code{scale_sh} là phân vị 99.9\% chia 127.

\textbf{Công thức.}
\[
\mathrm{scale\_sh}=\frac{P_{99.9}(|h|)}{127},\qquad u=\mathrm{clip}\Big(\mathrm{round}\Big(\frac{h}{\mathrm{scale\_sh}}\Big)+128,\ 1,\ 255\Big),\qquad h\approx(u-128)\,\mathrm{scale\_sh}
\]


\VerbatimInput[firstline=199,lastline=223,firstnumber=199,numbers=left,numbersep=3pt,frame=lines,framesep=2pt,rulecolor=\color{gray!60},breaklines=true,fontsize=\fontsize{6.6}{7.6}\selectfont]{../../splat_editor/render/splat_renderer.py}
\subsubsection{\texorpdfstring{\code{SplatRenderer._upload_state(self, resize=False, lo=0, hi=None)}}{SplatRenderer._upload_state}}
{\footnotesize\textit{Dòng 225--235}}\par
\textbf{Mô tả trong mã nguồn.} Gửi self.\_state[lo:hi] lên GPU (resize: cấp lại cả buffer).

\VerbatimInput[firstline=225,lastline=235,firstnumber=225,numbers=left,numbersep=3pt,frame=lines,framesep=2pt,rulecolor=\color{gray!60},breaklines=true,fontsize=\fontsize{6.6}{7.6}\selectfont]{../../splat_editor/render/splat_renderer.py}
\subsubsection{\texorpdfstring{\code{SplatRenderer._set_state_bit(self, bit, mask)}}{SplatRenderer._set_state_bit}}
{\footnotesize\textit{Dòng 237--244}}\par
\textbf{Mô tả trong mã nguồn.} Đặt bit `bit` của mọi gaussian theo mask (N,) bool theo chỉ số gaussian gốc. Sửa tại chỗ: luồng sắp xếp dùng chung mảng này.

\textbf{Giải thích.} Đặt bit trạng thái (chọn/xóa) theo mặt nạ, sửa tại chỗ trên mảng 1 byte/gaussian rồi chỉ gửi đoạn bị đổi.

\VerbatimInput[firstline=237,lastline=244,firstnumber=237,numbers=left,numbersep=3pt,frame=lines,framesep=2pt,rulecolor=\color{gray!60},breaklines=true,fontsize=\fontsize{6.6}{7.6}\selectfont]{../../splat_editor/render/splat_renderer.py}
\subsubsection{\texorpdfstring{\code{SplatRenderer.set_selection(self, mask)}}{SplatRenderer.set_selection}}
{\footnotesize\textit{Dòng 246--248}}\par
\textbf{Mô tả trong mã nguồn.} mask: (N,) bool theo chỉ số gaussian gốc, gaussian nào đang được chọn.

\VerbatimInput[firstline=246,lastline=248,firstnumber=246,numbers=left,numbersep=3pt,frame=lines,framesep=2pt,rulecolor=\color{gray!60},breaklines=true,fontsize=\fontsize{6.6}{7.6}\selectfont]{../../splat_editor/render/splat_renderer.py}
\subsubsection{\texorpdfstring{\code{SplatRenderer.set_deleted(self, mask)}}{SplatRenderer.set_deleted}}
{\footnotesize\textit{Dòng 250--254}}\par
\textbf{Mô tả trong mã nguồn.} mask: (N,) bool theo chỉ số gaussian gốc, gaussian nào đã bị xoá (xoá mềm). Rẻ: chỉ gửi 1 byte / gaussian rồi chọn + sắp xếp lại, không gửi lại dữ liệu gaussian.

\VerbatimInput[firstline=250,lastline=254,firstnumber=250,numbers=left,numbersep=3pt,frame=lines,framesep=2pt,rulecolor=\color{gray!60},breaklines=true,fontsize=\fontsize{6.6}{7.6}\selectfont]{../../splat_editor/render/splat_renderer.py}
\subsubsection{\texorpdfstring{\code{SplatRenderer.update_selection(self, indices, value)}}{SplatRenderer.update_selection}}
{\footnotesize\textit{Dòng 256--267}}\par
\textbf{Mô tả trong mã nguồn.} Đặt vùng chọn = value cho các gaussian `indices` (chỉ số gốc). Chỉ gửi lên GPU đoạn bị đổi: nét cọ chạm các gaussian gần nhau, mà gần nhau trong không gian thì cũng gần nhau trên GPU (thứ tự Morton), nên đoạn này thường rất nhỏ.

\VerbatimInput[firstline=256,lastline=267,firstnumber=256,numbers=left,numbersep=3pt,frame=lines,framesep=2pt,rulecolor=\color{gray!60},breaklines=true,fontsize=\fontsize{6.6}{7.6}\selectfont]{../../splat_editor/render/splat_renderer.py}
\subsubsection{\texorpdfstring{\code{SplatRenderer.set_transforms(self, matrices)}}{SplatRenderer.set_transforms}}
{\footnotesize\textit{Dòng 269--278}}\par
\textbf{Mô tả trong mã nguồn.} Ma trận thế giới 4x4 của từng lớp (thứ tự như trong splats). Rất rẻ: kéo / xoay / phóng to lớp chỉ gọi hàm này, không gửi lại dữ liệu gaussian.

\textbf{Giải thích.} Gửi ma trận thế giới 4$\times$4 của mỗi layer làm 4 texel. Rất rẻ: kéo/xoay/phóng layer không phải gửi lại dữ liệu gaussian.

\textbf{Công thức.}
Mỗi layer gửi 4 texel $=$ 4 cột của $M=\begin{bmatrix}sA&t\\0&1\end{bmatrix}$ (chuyển vị trước khi gửi vì numpy lưu theo hàng).


\VerbatimInput[firstline=269,lastline=278,firstnumber=269,numbers=left,numbersep=3pt,frame=lines,framesep=2pt,rulecolor=\color{gray!60},breaklines=true,fontsize=\fontsize{6.6}{7.6}\selectfont]{../../splat_editor/render/splat_renderer.py}
\subsubsection{\texorpdfstring{\code{SplatRenderer.set_lighting(self, params)}}{SplatRenderer.set_lighting}}
{\footnotesize\textit{Dòng 280--282}}\par
\textbf{Mô tả trong mã nguồn.} params: Lighting.params(). Đổi chữ ký -\textgreater{} pass 1 phải tính lại màu của mọi gaussian.

\VerbatimInput[firstline=280,lastline=282,firstnumber=280,numbers=left,numbersep=3pt,frame=lines,framesep=2pt,rulecolor=\color{gray!60},breaklines=true,fontsize=\fontsize{6.6}{7.6}\selectfont]{../../splat_editor/render/splat_renderer.py}
\subsubsection{\texorpdfstring{\code{SplatRenderer._upload_order(self, slots, ranks)}}{SplatRenderer._upload_order}}
{\footnotesize\textit{Dòng 285--299}}\par
\VerbatimInput[firstline=285,lastline=299,firstnumber=285,numbers=left,numbersep=3pt,frame=lines,framesep=2pt,rulecolor=\color{gray!60},breaklines=true,fontsize=\fontsize{6.6}{7.6}\selectfont]{../../splat_editor/render/splat_renderer.py}
\subsubsection{\texorpdfstring{\code{SplatRenderer._update_order(self, model_view, projection)}}{SplatRenderer._update_order}}
{\footnotesize\textit{Dòng 301--308}}\par
\VerbatimInput[firstline=301,lastline=308,firstnumber=301,numbers=left,numbersep=3pt,frame=lines,framesep=2pt,rulecolor=\color{gray!60},breaklines=true,fontsize=\fontsize{6.6}{7.6}\selectfont]{../../splat_editor/render/splat_renderer.py}
\subsubsection{\texorpdfstring{\code{SplatRenderer._prepare(self, model_view, projection, width, height)}}{SplatRenderer._prepare}}
{\footnotesize\textit{Dòng 311--350}}\par
\textbf{Mô tả trong mã nguồn.} Pass 1: tính ellipse + màu cho mọi gaussian được chọn, ghi vào feedback\_buffer.

\textbf{Giải thích.} \textbf{Pass 1}: bật transform feedback, vẽ \code{GL_POINTS} cho các gaussian được chọn; \code{splat_prep.vert} tính ellipse + màu và ghi 32 byte vào \code{feedback_buffer}. Camera đứng yên thì bỏ qua pass này.

\textbf{Công thức.}
\[
\Sigma'=J\,W\,\Sigma\,W^{\top}J^{\top},\qquad \Sigma=R\,S\,S^{\top}R^{\top}
\]
(chi tiết trong shader \code{splat_prep.vert}, Chương~\ref{ch:shader}).


\VerbatimInput[firstline=311,lastline=350,firstnumber=311,numbers=left,numbersep=3pt,frame=lines,framesep=2pt,rulecolor=\color{gray!60},breaklines=true,fontsize=\fontsize{6.6}{7.6}\selectfont]{../../splat_editor/render/splat_renderer.py}
\subsubsection{\texorpdfstring{\code{SplatRenderer._set_light_uniforms(self, s, light)}}{SplatRenderer._set_light_uniforms}}
{\footnotesize\textit{Dòng 356--387}}\par
\VerbatimInput[firstline=356,lastline=387,firstnumber=356,numbers=left,numbersep=3pt,frame=lines,framesep=2pt,rulecolor=\color{gray!60},breaklines=true,fontsize=\fontsize{6.6}{7.6}\selectfont]{../../splat_editor/render/splat_renderer.py}
\subsubsection{\texorpdfstring{\code{SplatRenderer._create_shadow_map(self)}}{SplatRenderer._create_shadow_map}}
{\footnotesize\textit{Dòng 389--403}}\par
\VerbatimInput[firstline=389,lastline=403,firstnumber=389,numbers=left,numbersep=3pt,frame=lines,framesep=2pt,rulecolor=\color{gray!60},breaklines=true,fontsize=\fontsize{6.6}{7.6}\selectfont]{../../splat_editor/render/splat_renderer.py}
\subsubsection{\texorpdfstring{\code{SplatRenderer._render_shadow_map(self, light)}}{SplatRenderer._render_shadow_map}}
{\footnotesize\textit{Dòng 405--442}}\par
\textbf{Mô tả trong mã nguồn.} Shadow mapping: vẽ tâm các gaussian đã chọn vào bản đồ độ sâu nhìn từ mặt trời. Chỉ vẽ lại khi mặt trời hoặc tập gaussian được chọn đổi (không vẽ lại mỗi lần camera nhúc nhích nhẹ).

\textbf{Giải thích.} Bản đồ bóng: vẽ tâm các gaussian ($\alpha\ge0.3$, chưa xóa) vào bản đồ độ sâu $2048^2$ nhìn từ mặt trời; chỉ vẽ lại khi ánh sáng hoặc cảnh đổi.

\textbf{Công thức.}
\[
\text{texel}=\frac{2r}{\mathrm{SHADOW\_SIZE}},\qquad \mathrm{SHADOW\_SIZE}=2048
\]


\VerbatimInput[firstline=405,lastline=442,firstnumber=405,numbers=left,numbersep=3pt,frame=lines,framesep=2pt,rulecolor=\color{gray!60},breaklines=true,fontsize=\fontsize{6.6}{7.6}\selectfont]{../../splat_editor/render/splat_renderer.py}
\subsubsection{\texorpdfstring{\code{SplatRenderer._collect_timings(self)}}{SplatRenderer._collect_timings}}
{\footnotesize\textit{Dòng 445--455}}\par
\VerbatimInput[firstline=445,lastline=455,firstnumber=445,numbers=left,numbersep=3pt,frame=lines,framesep=2pt,rulecolor=\color{gray!60},breaklines=true,fontsize=\fontsize{6.6}{7.6}\selectfont]{../../splat_editor/render/splat_renderer.py}
\subsubsection{\texorpdfstring{\code{SplatRenderer.adapt_budget(self, target_ms)}}{SplatRenderer.adapt_budget}}
{\footnotesize\textit{Dòng 457--470}}\par
\textbf{Mô tả trong mã nguồn.} Chỉnh budget để thời gian GPU lúc camera di chuyển (khung hình có chạy pass 1, nặng nhất) gần target\_ms. Ước lượng chi phí mỗi gaussian = ms / số gaussian thật sự xử lý ở khung hình đó (không dùng budget: kết quả sắp xếp đến trễ vài khung hình), rồi đổi budget từ từ.

\textbf{Giải thích.} Tự chỉnh ngân sách: $\mathrm{wanted}=\mathrm{target\_ms}/\mathrm{cost}$, bước $=\mathrm{clip}(\mathrm{wanted}/\mathrm{current},0.75,1.33)^{1/2}$ (căn bậc hai làm trơn, đổi tối đa khoảng 15\%/lần).

\textbf{Công thức.}
\[
\mathrm{cost}=\mathrm{median}\Big(\frac{ms}{count}\Big),\quad \mathrm{wanted}=\frac{\mathrm{target\_ms}}{\mathrm{cost}},\quad
\mathrm{step}=\mathrm{clip}\Big(\frac{\mathrm{wanted}}{\mathrm{current}},0.75,1.33\Big)^{1/2}
\]
\[
\mathrm{budget}\leftarrow\mathrm{clip}(\mathrm{current}\cdot\mathrm{step},\ \mathrm{MIN\_BUDGET},\ \mathrm{count})
\]


\VerbatimInput[firstline=457,lastline=470,firstnumber=457,numbers=left,numbersep=3pt,frame=lines,framesep=2pt,rulecolor=\color{gray!60},breaklines=true,fontsize=\fontsize{6.6}{7.6}\selectfont]{../../splat_editor/render/splat_renderer.py}
\subsubsection{\texorpdfstring{\code{SplatRenderer.draw(self, model_view, projection, width, height)}}{SplatRenderer.draw}}
{\footnotesize\textit{Dòng 472--490}}\par
\VerbatimInput[firstline=472,lastline=490,firstnumber=472,numbers=left,numbersep=3pt,frame=lines,framesep=2pt,rulecolor=\color{gray!60},breaklines=true,fontsize=\fontsize{6.6}{7.6}\selectfont]{../../splat_editor/render/splat_renderer.py}
\subsubsection{\texorpdfstring{\code{SplatRenderer._draw_splats(self)}}{SplatRenderer._draw_splats}}
{\footnotesize\textit{Dòng 492--504}}\par
\textbf{Mô tả trong mã nguồn.} Pass 2: vẽ hình vuông xa -\textgreater{} gần.

\textbf{Giải thích.} \textbf{Pass 2}: \code{glDrawArraysInstanced} vẽ mỗi gaussian một hình vuông (4 đỉnh) theo thứ tự xa $\to$ gần; trộn \code{GL_ONE, GL_ONE_MINUS_SRC_ALPHA} (premultiplied), tắt depth test.

\textbf{Công thức.}
\[
C_i=c_i\,a_i+(1-a_i)\,C_{i-1}\qquad\Longleftrightarrow\qquad C=\mathrm{src.rgb}+(1-\mathrm{src.a})\,C_{dst}
\]
Không giao hoán, nên phải vẽ xa $\to$ gần; tắt depth test.


\VerbatimInput[firstline=492,lastline=504,firstnumber=492,numbers=left,numbersep=3pt,frame=lines,framesep=2pt,rulecolor=\color{gray!60},breaklines=true,fontsize=\fontsize{6.6}{7.6}\selectfont]{../../splat_editor/render/splat_renderer.py}
\subsection{\texorpdfstring{\code{splat_editor/render/viewport.py}}{splat_editor/render/viewport.py}}
{\footnotesize 212 dòng}\par
\textbf{Vai trò.} \code{Viewport} ghép camera, lưới sàn, bộ vẽ gaussian, đèn và nhân vật vào một hình chữ nhật của cửa sổ. Hỗ trợ vẽ ở độ phân giải thấp hơn (\code{render_scale})
qua framebuffer riêng và các mức chất lượng (\code{QUALITY_PRESETS}).

\paragraph{Mô tả đầu file và cách dùng.} Phần chú thích đầu file do tác giả viết, nêu ý tưởng và ví dụ sử dụng:
\VerbatimInput[firstline=1,lastline=19,firstnumber=1,numbers=left,numbersep=3pt,frame=lines,framesep=2pt,rulecolor=\color{gray!60},breaklines=true,fontsize=\fontsize{6.2}{7.2}\selectfont]{../../splat_editor/render/viewport.py}
\subsubsection*{Khai báo mức module (dòng 20--29)}
\VerbatimInput[firstline=20,lastline=29,firstnumber=20,numbers=left,numbersep=3pt,frame=lines,framesep=2pt,rulecolor=\color{gray!60},breaklines=true,fontsize=\fontsize{6.6}{7.6}\selectfont]{../../splat_editor/render/viewport.py}
\subsubsection{\texorpdfstring{\code{class QualityPreset(NamedTuple)}}{QualityPreset}}
{\footnotesize\textit{Dòng 32--36}}\par
\VerbatimInput[firstline=32,lastline=36,firstnumber=32,numbers=left,numbersep=3pt,frame=lines,framesep=2pt,rulecolor=\color{gray!60},breaklines=true,fontsize=\fontsize{6.6}{7.6}\selectfont]{../../splat_editor/render/viewport.py}
\subsubsection*{Khai báo mức module (dòng 39--47)}
\VerbatimInput[firstline=39,lastline=47,firstnumber=39,numbers=left,numbersep=3pt,frame=lines,framesep=2pt,rulecolor=\color{gray!60},breaklines=true,fontsize=\fontsize{6.6}{7.6}\selectfont]{../../splat_editor/render/viewport.py}
\subsubsection{\texorpdfstring{\code{class Viewport}}{Viewport}}
{\footnotesize\textit{Dòng 50--212}}\par
Các phương thức: \code{__init__}, \code{sync}, \code{set_quality}, \code{view_projection}, \code{frame}, \code{render}, \code{sky}, \code{_draw_scene}, \code{_bind_offscreen}.

\VerbatimInput[firstline=50,lastline=51,firstnumber=50,numbers=left,numbersep=3pt,frame=lines,framesep=2pt,rulecolor=\color{gray!60},breaklines=true,fontsize=\fontsize{6.6}{7.6}\selectfont]{../../splat_editor/render/viewport.py}
\subsubsection{\texorpdfstring{\code{Viewport.__init__(self)}}{Viewport.__init__}}
{\footnotesize\textit{Dòng 53--83}}\par
\VerbatimInput[firstline=53,lastline=83,firstnumber=53,numbers=left,numbersep=3pt,frame=lines,framesep=2pt,rulecolor=\color{gray!60},breaklines=true,fontsize=\fontsize{6.6}{7.6}\selectfont]{../../splat_editor/render/viewport.py}
\subsubsection{\texorpdfstring{\code{Viewport.sync(self, editor)}}{Viewport.sync}}
{\footnotesize\textit{Dòng 85--108}}\par
\textbf{Mô tả trong mã nguồn.} Gửi lại gaussian / vùng chọn lên GPU nếu Editor đã thay đổi.

\textbf{Giải thích.} So sánh số phiên bản của Editor với số đã gửi: \code{geometry} $\to$ gửi lại dữ liệu; \code{deleted} $\to$ gửi mảng cờ; \code{selection} $\to$ gửi phần thay đổi; \code{transform} $\to$ gửi ma trận.

\VerbatimInput[firstline=85,lastline=108,firstnumber=85,numbers=left,numbersep=3pt,frame=lines,framesep=2pt,rulecolor=\color{gray!60},breaklines=true,fontsize=\fontsize{6.6}{7.6}\selectfont]{../../splat_editor/render/viewport.py}
\subsubsection{\texorpdfstring{\code{Viewport.set_quality(self, name)}}{Viewport.set_quality}}
{\footnotesize\textit{Dòng 110--116}}\par
\VerbatimInput[firstline=110,lastline=116,firstnumber=110,numbers=left,numbersep=3pt,frame=lines,framesep=2pt,rulecolor=\color{gray!60},breaklines=true,fontsize=\fontsize{6.6}{7.6}\selectfont]{../../splat_editor/render/viewport.py}
\subsubsection{\texorpdfstring{\code{Viewport.view_projection(self, rect)}}{Viewport.view_projection}}
{\footnotesize\textit{Dòng 118--119}}\par
\VerbatimInput[firstline=118,lastline=119,firstnumber=118,numbers=left,numbersep=3pt,frame=lines,framesep=2pt,rulecolor=\color{gray!60},breaklines=true,fontsize=\fontsize{6.6}{7.6}\selectfont]{../../splat_editor/render/viewport.py}
\subsubsection{\texorpdfstring{\code{Viewport.frame(self, positions)}}{Viewport.frame}}
{\footnotesize\textit{Dòng 121--126}}\par
\textbf{Mô tả trong mã nguồn.} Đưa camera về nhìn toàn bộ các điểm `positions` (toạ độ thế giới; None: đặt lại camera).

\VerbatimInput[firstline=121,lastline=126,firstnumber=121,numbers=left,numbersep=3pt,frame=lines,framesep=2pt,rulecolor=\color{gray!60},breaklines=true,fontsize=\fontsize{6.6}{7.6}\selectfont]{../../splat_editor/render/viewport.py}
\subsubsection{\texorpdfstring{\code{Viewport.render(self, rect, window_size, framebuffer_size)}}{Viewport.render}}
{\footnotesize\textit{Dòng 128--170}}\par
\textbf{Mô tả trong mã nguồn.} rect: vùng 3D theo toạ độ cửa sổ (gốc trên-trái), được đổi sang pixel framebuffer.

\textbf{Giải thích.} Vẽ một khung hình vào \code{rect}: chọn độ phân giải theo trạng thái chuyển động, bind framebuffer riêng khi \code{render_scale<1}, vẽ nền, mặt trời, gaussian, lưới, đèn, nhân vật rồi phóng lên khung nhìn.

\textbf{Công thức.}
\[
x_{px}=x_0+\frac{n_x+1}{2}W,\qquad y_{px}=y_0+\frac{1-n_y}{2}H
\]


\VerbatimInput[firstline=128,lastline=170,firstnumber=128,numbers=left,numbersep=3pt,frame=lines,framesep=2pt,rulecolor=\color{gray!60},breaklines=true,fontsize=\fontsize{6.6}{7.6}\selectfont]{../../splat_editor/render/viewport.py}
\subsubsection{\texorpdfstring{\code{Viewport.sky(self)}}{Viewport.sky}}
{\footnotesize\textit{Dòng 172--177}}\par
\textbf{Mô tả trong mã nguồn.} Màu nền: đổi theo giờ trong ngày khi có mặt trời, ngược lại là màu người dùng chọn.

\VerbatimInput[firstline=172,lastline=177,firstnumber=172,numbers=left,numbersep=3pt,frame=lines,framesep=2pt,rulecolor=\color{gray!60},breaklines=true,fontsize=\fontsize{6.6}{7.6}\selectfont]{../../splat_editor/render/viewport.py}
\subsubsection{\texorpdfstring{\code{Viewport._draw_scene(self, view, projection, w, h)}}{Viewport._draw_scene}}
{\footnotesize\textit{Dòng 179--193}}\par
\VerbatimInput[firstline=179,lastline=193,firstnumber=179,numbers=left,numbersep=3pt,frame=lines,framesep=2pt,rulecolor=\color{gray!60},breaklines=true,fontsize=\fontsize{6.6}{7.6}\selectfont]{../../splat_editor/render/viewport.py}
\subsubsection{\texorpdfstring{\code{Viewport._bind_offscreen(self, w, h)}}{Viewport._bind_offscreen}}
{\footnotesize\textit{Dòng 195--212}}\par
\VerbatimInput[firstline=195,lastline=212,firstnumber=195,numbers=left,numbersep=3pt,frame=lines,framesep=2pt,rulecolor=\color{gray!60},breaklines=true,fontsize=\fontsize{6.6}{7.6}\selectfont]{../../splat_editor/render/viewport.py}
\subsection{\texorpdfstring{\code{splat_editor/render/grid_renderer.py}}{splat_editor/render/grid_renderer.py}}
{\footnotesize 51 dòng}\par
\textbf{Vai trò.} Lưới sàn vô hạn: một tam giác phủ toàn màn hình, fragment shader bắn tia từ camera tới mặt phẳng $y=h$ và tô đường lưới.

\paragraph{Mô tả đầu file và cách dùng.} Phần chú thích đầu file do tác giả viết, nêu ý tưởng và ví dụ sử dụng:
\VerbatimInput[firstline=1,lastline=14,firstnumber=1,numbers=left,numbersep=3pt,frame=lines,framesep=2pt,rulecolor=\color{gray!60},breaklines=true,fontsize=\fontsize{6.2}{7.2}\selectfont]{../../splat_editor/render/grid_renderer.py}
\subsubsection*{Khai báo mức module (dòng 15--19)}
\VerbatimInput[firstline=15,lastline=19,firstnumber=15,numbers=left,numbersep=3pt,frame=lines,framesep=2pt,rulecolor=\color{gray!60},breaklines=true,fontsize=\fontsize{6.6}{7.6}\selectfont]{../../splat_editor/render/grid_renderer.py}
\subsubsection{\texorpdfstring{\code{class GridRenderer}}{GridRenderer}}
{\footnotesize\textit{Dòng 22--51}}\par
Các phương thức: \code{__init__}, \code{draw}.

\VerbatimInput[firstline=22,lastline=22,firstnumber=22,numbers=left,numbersep=3pt,frame=lines,framesep=2pt,rulecolor=\color{gray!60},breaklines=true,fontsize=\fontsize{6.6}{7.6}\selectfont]{../../splat_editor/render/grid_renderer.py}
\subsubsection{\texorpdfstring{\code{GridRenderer.__init__(self)}}{GridRenderer.__init__}}
{\footnotesize\textit{Dòng 23--26}}\par
\VerbatimInput[firstline=23,lastline=26,firstnumber=23,numbers=left,numbersep=3pt,frame=lines,framesep=2pt,rulecolor=\color{gray!60},breaklines=true,fontsize=\fontsize{6.6}{7.6}\selectfont]{../../splat_editor/render/grid_renderer.py}
\subsubsection{\texorpdfstring{\code{GridRenderer.draw(self, view, projection, camera_distance)}}{GridRenderer.draw}}
{\footnotesize\textit{Dòng 28--51}}\par
\VerbatimInput[firstline=28,lastline=51,firstnumber=28,numbers=left,numbersep=3pt,frame=lines,framesep=2pt,rulecolor=\color{gray!60},breaklines=true,fontsize=\fontsize{6.6}{7.6}\selectfont]{../../splat_editor/render/grid_renderer.py}
\subsection{\texorpdfstring{\code{splat_editor/render/light_renderer.py}}{splat_editor/render/light_renderer.py}}
{\footnotesize 114 dòng}\par
\textbf{Vai trò.} Vẽ mặt trời, các khối đèn (quả cầu phát sáng) và đom đóm (hạt) quanh đèn.

\paragraph{Mô tả đầu file và cách dùng.} Phần chú thích đầu file do tác giả viết, nêu ý tưởng và ví dụ sử dụng:
\VerbatimInput[firstline=1,lastline=11,firstnumber=1,numbers=left,numbersep=3pt,frame=lines,framesep=2pt,rulecolor=\color{gray!60},breaklines=true,fontsize=\fontsize{6.2}{7.2}\selectfont]{../../splat_editor/render/light_renderer.py}
\subsubsection*{Khai báo mức module (dòng 12--17)}
\VerbatimInput[firstline=12,lastline=17,firstnumber=12,numbers=left,numbersep=3pt,frame=lines,framesep=2pt,rulecolor=\color{gray!60},breaklines=true,fontsize=\fontsize{6.6}{7.6}\selectfont]{../../splat_editor/render/light_renderer.py}
\subsubsection{\texorpdfstring{\code{class LightRenderer}}{LightRenderer}}
{\footnotesize\textit{Dòng 20--114}}\par
Các phương thức: \code{__init__}, \code{draw_sun}, \code{draw}, \code{_sphere}.

\VerbatimInput[firstline=20,lastline=20,firstnumber=20,numbers=left,numbersep=3pt,frame=lines,framesep=2pt,rulecolor=\color{gray!60},breaklines=true,fontsize=\fontsize{6.6}{7.6}\selectfont]{../../splat_editor/render/light_renderer.py}
\subsubsection{\texorpdfstring{\code{LightRenderer.__init__(self)}}{LightRenderer.__init__}}
{\footnotesize\textit{Dòng 21--32}}\par
\VerbatimInput[firstline=21,lastline=32,firstnumber=21,numbers=left,numbersep=3pt,frame=lines,framesep=2pt,rulecolor=\color{gray!60},breaklines=true,fontsize=\fontsize{6.6}{7.6}\selectfont]{../../splat_editor/render/light_renderer.py}
\subsubsection{\texorpdfstring{\code{LightRenderer.draw_sun(self, view, projection, lighting)}}{LightRenderer.draw_sun}}
{\footnotesize\textit{Dòng 34--52}}\par
\textbf{Mô tả trong mã nguồn.} Vẽ mặt trời TRƯỚC các gaussian: nó là phần tử xa nhất nên các gaussian (đã sắp xa -\textgreater{} gần) tự blend đè lên nó, không cần kiểm tra che khuất.

\VerbatimInput[firstline=34,lastline=52,firstnumber=34,numbers=left,numbersep=3pt,frame=lines,framesep=2pt,rulecolor=\color{gray!60},breaklines=true,fontsize=\fontsize{6.6}{7.6}\selectfont]{../../splat_editor/render/light_renderer.py}
\subsubsection{\texorpdfstring{\code{LightRenderer.draw(self, view, projection, lighting, time, unit, physics=None)}}{LightRenderer.draw}}
{\footnotesize\textit{Dòng 54--102}}\par
\textbf{Mô tả trong mã nguồn.} unit: chiều cao người chơi (đơn vị cảnh) để chọn kích thước quả cầu.

\VerbatimInput[firstline=54,lastline=102,firstnumber=54,numbers=left,numbersep=3pt,frame=lines,framesep=2pt,rulecolor=\color{gray!60},breaklines=true,fontsize=\fontsize{6.6}{7.6}\selectfont]{../../splat_editor/render/light_renderer.py}
\subsubsection{\texorpdfstring{\code{LightRenderer._sphere(self, view, projection, center, radius, halo, color, brightness)}}{LightRenderer._sphere}}
{\footnotesize\textit{Dòng 104--114}}\par
\VerbatimInput[firstline=104,lastline=114,firstnumber=104,numbers=left,numbersep=3pt,frame=lines,framesep=2pt,rulecolor=\color{gray!60},breaklines=true,fontsize=\fontsize{6.6}{7.6}\selectfont]{../../splat_editor/render/light_renderer.py}
\subsection{\texorpdfstring{\code{splat_editor/render/avatar_renderer.py}}{splat_editor/render/avatar_renderer.py}}
{\footnotesize 190 dòng}\par
\textbf{Vai trò.} Vẽ nhân vật người chơi từ 6 hộp (đầu, thân, hai tay, hai chân) với skin $64\times64$. Có sẵn bộ giải mã PNG tự viết (\code{load_png}) nên không cần thư viện ảnh.

\paragraph{Mô tả đầu file và cách dùng.} Phần chú thích đầu file do tác giả viết, nêu ý tưởng và ví dụ sử dụng:
\VerbatimInput[firstline=1,lastline=13,firstnumber=1,numbers=left,numbersep=3pt,frame=lines,framesep=2pt,rulecolor=\color{gray!60},breaklines=true,fontsize=\fontsize{6.2}{7.2}\selectfont]{../../splat_editor/render/avatar_renderer.py}
%% ---- merged math: s:texture
\subsubsection{Cơ sở toán: Texture}\label{s:texture}
Các công thức dưới đây giải thích cách dựng UV từ skin, hộp nhân vật, shader avatar (\code{avatar.vert}, \code{avatar.frag}) và bộ giải mã PNG \code{load_png} trong file này.
\begin{itemize}
\item \textbf{Skin $64\times64$} kiểu chữ thập Minecraft, \code{avatar_renderer._uv_region(face, ox, oy, w, h, d)} ($w,h,d$ là rộng, cao, sâu của hộp). Mặt trước $(ox+d,\,oy+d)$; sau $(ox+2d+w,\,oy+d)$; phải $(ox,\,oy+d,\,d\times h)$; trái $(ox+d+w,\,oy+d)$; trên $(ox+d,\,oy)$; dưới $(ox+d+w,\,oy)$. $\mathrm{UV}=\mathrm{pixel}/64$. Bốn đỉnh mỗi mặt: $(r_x,r_y)$, $(r_x+r_w,r_y)$, $(r_x+r_w,r_y+r_h)$, $(r_x,r_y+r_h)$.
\item Hộp: nửa kích thước $h_x=w/2+\mathrm{inflate}$, \ldots (lớp áo ``inflate'' phồng ra để tránh z-fighting). Đỉnh $=$ tâm $+(\pm1)\cdot$nửa kích thước; pháp tuyến là 6 hướng trục.
\item \textbf{Avatar shader:} $\mathrm{shade}=\mathrm{bright}\cdot(0.62+0.38\max(n\cdot l,0))$ (ambient 0.62 + Lambert 0.38); $n_{\mathrm{world}}=\mathrm{normalize}(\mathrm{mat3}(M)\,n)$; bỏ texel $\alpha<0.5$ (lớp áo trong suốt).
\item \textbf{Giải mã PNG tự viết} (\code{load_png}): zlib + bộ lọc dòng, với $a$ là trái, $b$ là trên, $c$ là trên-trái và mọi phép cộng đều mod 256: 1 Sub: $x+a$; 2 Up: $x+b$; 3 Average: $x+\lfloor(a+b)/2\rfloor$; 4 \textbf{Paeth}: $p=a+b-c$, chọn trong $\{a,b,c\}$ gần $p$ nhất ($p_a=|p-a|$, $p_b$, $p_c$; ưu tiên $a\le b\le c$).
\item Texel buffer: dữ liệu gaussian đọc bằng \code{texelFetch} (số nguyên, không nội suy).
\end{itemize}

\subsubsection*{Khai báo mức module (dòng 14--49)}
\VerbatimInput[firstline=14,lastline=49,firstnumber=14,numbers=left,numbersep=3pt,frame=lines,framesep=2pt,rulecolor=\color{gray!60},breaklines=true,fontsize=\fontsize{6.6}{7.6}\selectfont]{../../splat_editor/render/avatar_renderer.py}
\subsubsection{\texorpdfstring{\code{load_png(path)}}{load_png}}
{\footnotesize\textit{Dòng 52--102}}\par
\textbf{Mô tả trong mã nguồn.} Đọc PNG 8 bit RGB / RGBA không xen kẽ -\textgreater{} mảng (h, w, 4) uint8 (không cần thư viện ảnh).

\textbf{Giải thích.} Giải mã PNG 8 bit RGB/RGBA không xen kẽ bằng \code{zlib} và 5 bộ lọc dòng (None, Sub, Up, Average, Paeth), không cần thư viện ảnh.

\textbf{Công thức.}
Bộ lọc dòng PNG ($a$ trái, $b$ trên, $c$ trên-trái, mọi phép cộng $\bmod256$):
\[
\text{Sub}:x+a,\quad\text{Up}:x+b,\quad\text{Avg}:x+\Big\lfloor\tfrac{a+b}{2}\Big\rfloor,\quad
\text{Paeth}:x+\arg\min_{u\in\{a,b,c\}}|a+b-c-u|
\]


\VerbatimInput[firstline=52,lastline=102,firstnumber=52,numbers=left,numbersep=3pt,frame=lines,framesep=2pt,rulecolor=\color{gray!60},breaklines=true,fontsize=\fontsize{6.6}{7.6}\selectfont]{../../splat_editor/render/avatar_renderer.py}
\subsubsection{\texorpdfstring{\code{_uv_region(face, ox, oy, w, h, d)}}{_uv_region}}
{\footnotesize\textit{Dòng 105--108}}\par
\textbf{Mô tả trong mã nguồn.} Ô skin của một mặt hộp theo kiểu trải chữ thập của Minecraft: (x, y, rộng, cao).

\textbf{Giải thích.} Ô skin của một mặt hộp theo kiểu trải chữ thập của Minecraft: $(x,y,\text{rộng},\text{cao})$.

\textbf{Công thức.}
\[
\mathrm{UV}=\frac{\mathrm{pixel}}{64},\qquad
\text{trước }(o_x+d,\,o_y+d),\ \text{sau }(o_x+2d+w,\,o_y+d),\ \text{trên }(o_x+d,\,o_y)
\]


\VerbatimInput[firstline=105,lastline=108,firstnumber=105,numbers=left,numbersep=3pt,frame=lines,framesep=2pt,rulecolor=\color{gray!60},breaklines=true,fontsize=\fontsize{6.6}{7.6}\selectfont]{../../splat_editor/render/avatar_renderer.py}
\subsubsection{\texorpdfstring{\code{_box(ox, oy, centre, size, inflate)}}{_box}}
{\footnotesize\textit{Dòng 111--121}}\par
\textbf{Mô tả trong mã nguồn.} Đỉnh (x, y, z, nx, ny, nz, u, v) của một hộp: 6 mặt x 4 đỉnh.

\VerbatimInput[firstline=111,lastline=121,firstnumber=111,numbers=left,numbersep=3pt,frame=lines,framesep=2pt,rulecolor=\color{gray!60},breaklines=true,fontsize=\fontsize{6.6}{7.6}\selectfont]{../../splat_editor/render/avatar_renderer.py}
\subsubsection{\texorpdfstring{\code{build_mesh()}}{build_mesh}}
{\footnotesize\textit{Dòng 124--137}}\par
\textbf{Mô tả trong mã nguồn.} (đỉnh float32 (n, 8), chỉ số uint32, [(offset byte, số chỉ số)] cho từng khối theo LIMBS).

\VerbatimInput[firstline=124,lastline=137,firstnumber=124,numbers=left,numbersep=3pt,frame=lines,framesep=2pt,rulecolor=\color{gray!60},breaklines=true,fontsize=\fontsize{6.6}{7.6}\selectfont]{../../splat_editor/render/avatar_renderer.py}
\subsubsection{\texorpdfstring{\code{class AvatarRenderer}}{AvatarRenderer}}
{\footnotesize\textit{Dòng 140--190}}\par
Các phương thức: \code{__init__}, \code{draw}.

\VerbatimInput[firstline=140,lastline=140,firstnumber=140,numbers=left,numbersep=3pt,frame=lines,framesep=2pt,rulecolor=\color{gray!60},breaklines=true,fontsize=\fontsize{6.6}{7.6}\selectfont]{../../splat_editor/render/avatar_renderer.py}
\subsubsection{\texorpdfstring{\code{AvatarRenderer.__init__(self)}}{AvatarRenderer.__init__}}
{\footnotesize\textit{Dòng 141--166}}\par
\VerbatimInput[firstline=141,lastline=166,firstnumber=141,numbers=left,numbersep=3pt,frame=lines,framesep=2pt,rulecolor=\color{gray!60},breaklines=true,fontsize=\fontsize{6.6}{7.6}\selectfont]{../../splat_editor/render/avatar_renderer.py}
\subsubsection{\texorpdfstring{\code{AvatarRenderer.draw(self, view, projection, matrices, brightness=(1.0, 1.0, 1.0), light...)}}{AvatarRenderer.draw}}
{\footnotesize\textit{Dòng 168--190}}\par
\textbf{Mô tả trong mã nguồn.} matrices: 6 ma trận model theo LIMBS (core.avatar\_pose.limb\_matrices).

\VerbatimInput[firstline=168,lastline=190,firstnumber=168,numbers=left,numbersep=3pt,frame=lines,framesep=2pt,rulecolor=\color{gray!60},breaklines=true,fontsize=\fontsize{6.6}{7.6}\selectfont]{../../splat_editor/render/avatar_renderer.py}
\subsection{\texorpdfstring{\code{shaders/}}{shaders/}}
Mã GLSL được giải thích riêng ở Chương~\ref{ch:shader}.


\clearpage
\section{Thư mục ui: giao diện}\label{ch:ui}
Thư mục \code{ui} là \textbf{giao diện}: cửa sổ GLFW, Dear ImGui, menu, bảng bên trái, thanh trạng thái, xử lý chuột/bàn phím.
Thư mục này hầu như không có toán; nó chỉ gọi lại \code{Editor}, \code{Viewport}, \code{Lighting}, \code{Physics}. Ý tưởng chính là mẫu
\textbf{Command}: mọi hành động (menu, phím tắt, nút) khai báo một lần trong \code{ActionRegistry} nên không bao giờ lệch nhau.

\subsection{\texorpdfstring{\code{splat_editor/ui/window.py}}{splat_editor/ui/window.py}}
{\footnotesize 66 dòng}\par
\textbf{Vai trò.} Cửa sổ GLFW với ngữ cảnh OpenGL 3.3 core. Bọc các hàm GLFW (kích thước, vị trí con trỏ, phím, chuột, \code{swap_buffers}).

\paragraph{Mô tả đầu file và cách dùng.} Phần chú thích đầu file do tác giả viết, nêu ý tưởng và ví dụ sử dụng:
\VerbatimInput[firstline=1,lastline=13,firstnumber=1,numbers=left,numbersep=3pt,frame=lines,framesep=2pt,rulecolor=\color{gray!60},breaklines=true,fontsize=\fontsize{6.2}{7.2}\selectfont]{../../splat_editor/ui/window.py}
\subsubsection*{Khai báo mức module (dòng 14--14)}
\VerbatimInput[firstline=14,lastline=14,firstnumber=14,numbers=left,numbersep=3pt,frame=lines,framesep=2pt,rulecolor=\color{gray!60},breaklines=true,fontsize=\fontsize{6.6}{7.6}\selectfont]{../../splat_editor/ui/window.py}
\subsubsection{\texorpdfstring{\code{class Window}}{Window}}
{\footnotesize\textit{Dòng 17--66}}\par
Các phương thức: \code{__init__}, \code{should_close}, \code{close}, \code{size}, \code{framebuffer_size}, \code{cursor_pos}, \code{is_key_down}, \code{is_mouse_down}, \code{poll_events}, \code{swap_buffers}, \code{destroy}.

\VerbatimInput[firstline=17,lastline=17,firstnumber=17,numbers=left,numbersep=3pt,frame=lines,framesep=2pt,rulecolor=\color{gray!60},breaklines=true,fontsize=\fontsize{6.6}{7.6}\selectfont]{../../splat_editor/ui/window.py}
\subsubsection{\texorpdfstring{\code{Window.__init__(self, title, width, height)}}{Window.__init__}}
{\footnotesize\textit{Dòng 18--30}}\par
\VerbatimInput[firstline=18,lastline=30,firstnumber=18,numbers=left,numbersep=3pt,frame=lines,framesep=2pt,rulecolor=\color{gray!60},breaklines=true,fontsize=\fontsize{6.6}{7.6}\selectfont]{../../splat_editor/ui/window.py}
\subsubsection{\texorpdfstring{\code{Window.should_close(self)}}{Window.should_close}}
{\footnotesize\textit{Dòng 32--34}}\par
\VerbatimInput[firstline=32,lastline=34,firstnumber=32,numbers=left,numbersep=3pt,frame=lines,framesep=2pt,rulecolor=\color{gray!60},breaklines=true,fontsize=\fontsize{6.6}{7.6}\selectfont]{../../splat_editor/ui/window.py}
\subsubsection{\texorpdfstring{\code{Window.close(self)}}{Window.close}}
{\footnotesize\textit{Dòng 36--37}}\par
\VerbatimInput[firstline=36,lastline=37,firstnumber=36,numbers=left,numbersep=3pt,frame=lines,framesep=2pt,rulecolor=\color{gray!60},breaklines=true,fontsize=\fontsize{6.6}{7.6}\selectfont]{../../splat_editor/ui/window.py}
\subsubsection{\texorpdfstring{\code{Window.size(self)}}{Window.size}}
{\footnotesize\textit{Dòng 39--41}}\par
\VerbatimInput[firstline=39,lastline=41,firstnumber=39,numbers=left,numbersep=3pt,frame=lines,framesep=2pt,rulecolor=\color{gray!60},breaklines=true,fontsize=\fontsize{6.6}{7.6}\selectfont]{../../splat_editor/ui/window.py}
\subsubsection{\texorpdfstring{\code{Window.framebuffer_size(self)}}{Window.framebuffer_size}}
{\footnotesize\textit{Dòng 43--45}}\par
\VerbatimInput[firstline=43,lastline=45,firstnumber=43,numbers=left,numbersep=3pt,frame=lines,framesep=2pt,rulecolor=\color{gray!60},breaklines=true,fontsize=\fontsize{6.6}{7.6}\selectfont]{../../splat_editor/ui/window.py}
\subsubsection{\texorpdfstring{\code{Window.cursor_pos(self)}}{Window.cursor_pos}}
{\footnotesize\textit{Dòng 47--49}}\par
\VerbatimInput[firstline=47,lastline=49,firstnumber=47,numbers=left,numbersep=3pt,frame=lines,framesep=2pt,rulecolor=\color{gray!60},breaklines=true,fontsize=\fontsize{6.6}{7.6}\selectfont]{../../splat_editor/ui/window.py}
\subsubsection{\texorpdfstring{\code{Window.is_key_down(self, key)}}{Window.is_key_down}}
{\footnotesize\textit{Dòng 51--52}}\par
\VerbatimInput[firstline=51,lastline=52,firstnumber=51,numbers=left,numbersep=3pt,frame=lines,framesep=2pt,rulecolor=\color{gray!60},breaklines=true,fontsize=\fontsize{6.6}{7.6}\selectfont]{../../splat_editor/ui/window.py}
\subsubsection{\texorpdfstring{\code{Window.is_mouse_down(self, button)}}{Window.is_mouse_down}}
{\footnotesize\textit{Dòng 54--55}}\par
\VerbatimInput[firstline=54,lastline=55,firstnumber=54,numbers=left,numbersep=3pt,frame=lines,framesep=2pt,rulecolor=\color{gray!60},breaklines=true,fontsize=\fontsize{6.6}{7.6}\selectfont]{../../splat_editor/ui/window.py}
\subsubsection{\texorpdfstring{\code{Window.poll_events()}}{Window.poll_events}}
{\footnotesize\textit{Dòng 57--59}}\par
\VerbatimInput[firstline=57,lastline=59,firstnumber=57,numbers=left,numbersep=3pt,frame=lines,framesep=2pt,rulecolor=\color{gray!60},breaklines=true,fontsize=\fontsize{6.6}{7.6}\selectfont]{../../splat_editor/ui/window.py}
\subsubsection{\texorpdfstring{\code{Window.swap_buffers(self)}}{Window.swap_buffers}}
{\footnotesize\textit{Dòng 61--62}}\par
\VerbatimInput[firstline=61,lastline=62,firstnumber=61,numbers=left,numbersep=3pt,frame=lines,framesep=2pt,rulecolor=\color{gray!60},breaklines=true,fontsize=\fontsize{6.6}{7.6}\selectfont]{../../splat_editor/ui/window.py}
\subsubsection{\texorpdfstring{\code{Window.destroy(self)}}{Window.destroy}}
{\footnotesize\textit{Dòng 64--66}}\par
\VerbatimInput[firstline=64,lastline=66,firstnumber=64,numbers=left,numbersep=3pt,frame=lines,framesep=2pt,rulecolor=\color{gray!60},breaklines=true,fontsize=\fontsize{6.6}{7.6}\selectfont]{../../splat_editor/ui/window.py}
\subsection{\texorpdfstring{\code{splat_editor/ui/gui.py}}{splat_editor/ui/gui.py}}
{\footnotesize 58 dòng}\par
\textbf{Vai trò.} Ngữ cảnh Dear ImGui gắn với cửa sổ GLFW: \code{begin_frame}, \code{end_frame}, \code{draw}, \code{shutdown}, và cờ \code{wants_mouse} / \code{wants_text} để biết ImGui đang chiếm chuột/bàn phím.

\paragraph{Mô tả đầu file và cách dùng.} Phần chú thích đầu file do tác giả viết, nêu ý tưởng và ví dụ sử dụng:
\VerbatimInput[firstline=1,lastline=19,firstnumber=1,numbers=left,numbersep=3pt,frame=lines,framesep=2pt,rulecolor=\color{gray!60},breaklines=true,fontsize=\fontsize{6.2}{7.2}\selectfont]{../../splat_editor/ui/gui.py}
\subsubsection*{Khai báo mức module (dòng 20--23)}
\VerbatimInput[firstline=20,lastline=23,firstnumber=20,numbers=left,numbersep=3pt,frame=lines,framesep=2pt,rulecolor=\color{gray!60},breaklines=true,fontsize=\fontsize{6.6}{7.6}\selectfont]{../../splat_editor/ui/gui.py}
\subsubsection{\texorpdfstring{\code{class Gui}}{Gui}}
{\footnotesize\textit{Dòng 26--58}}\par
Các phương thức: \code{__init__}, \code{io}, \code{wants_mouse}, \code{wants_text}, \code{begin_frame}, \code{end_frame}, \code{draw}, \code{shutdown}.

\VerbatimInput[firstline=26,lastline=26,firstnumber=26,numbers=left,numbersep=3pt,frame=lines,framesep=2pt,rulecolor=\color{gray!60},breaklines=true,fontsize=\fontsize{6.6}{7.6}\selectfont]{../../splat_editor/ui/gui.py}
\subsubsection{\texorpdfstring{\code{Gui.__init__(self, window)}}{Gui.__init__}}
{\footnotesize\textit{Dòng 27--32}}\par
\VerbatimInput[firstline=27,lastline=32,firstnumber=27,numbers=left,numbersep=3pt,frame=lines,framesep=2pt,rulecolor=\color{gray!60},breaklines=true,fontsize=\fontsize{6.6}{7.6}\selectfont]{../../splat_editor/ui/gui.py}
\subsubsection{\texorpdfstring{\code{Gui.io(self)}}{Gui.io}}
{\footnotesize\textit{Dòng 34--36}}\par
\VerbatimInput[firstline=34,lastline=36,firstnumber=34,numbers=left,numbersep=3pt,frame=lines,framesep=2pt,rulecolor=\color{gray!60},breaklines=true,fontsize=\fontsize{6.6}{7.6}\selectfont]{../../splat_editor/ui/gui.py}
\subsubsection{\texorpdfstring{\code{Gui.wants_mouse(self)}}{Gui.wants_mouse}}
{\footnotesize\textit{Dòng 38--40}}\par
\VerbatimInput[firstline=38,lastline=40,firstnumber=38,numbers=left,numbersep=3pt,frame=lines,framesep=2pt,rulecolor=\color{gray!60},breaklines=true,fontsize=\fontsize{6.6}{7.6}\selectfont]{../../splat_editor/ui/gui.py}
\subsubsection{\texorpdfstring{\code{Gui.wants_text(self)}}{Gui.wants_text}}
{\footnotesize\textit{Dòng 42--44}}\par
\VerbatimInput[firstline=42,lastline=44,firstnumber=42,numbers=left,numbersep=3pt,frame=lines,framesep=2pt,rulecolor=\color{gray!60},breaklines=true,fontsize=\fontsize{6.6}{7.6}\selectfont]{../../splat_editor/ui/gui.py}
\subsubsection{\texorpdfstring{\code{Gui.begin_frame(self)}}{Gui.begin_frame}}
{\footnotesize\textit{Dòng 46--48}}\par
\VerbatimInput[firstline=46,lastline=48,firstnumber=46,numbers=left,numbersep=3pt,frame=lines,framesep=2pt,rulecolor=\color{gray!60},breaklines=true,fontsize=\fontsize{6.6}{7.6}\selectfont]{../../splat_editor/ui/gui.py}
\subsubsection{\texorpdfstring{\code{Gui.end_frame(self)}}{Gui.end_frame}}
{\footnotesize\textit{Dòng 50--52}}\par
\textbf{Mô tả trong mã nguồn.} Kết thúc khung hình ImGui (chưa vẽ ra màn hình).

\VerbatimInput[firstline=50,lastline=52,firstnumber=50,numbers=left,numbersep=3pt,frame=lines,framesep=2pt,rulecolor=\color{gray!60},breaklines=true,fontsize=\fontsize{6.6}{7.6}\selectfont]{../../splat_editor/ui/gui.py}
\subsubsection{\texorpdfstring{\code{Gui.draw(self)}}{Gui.draw}}
{\footnotesize\textit{Dòng 54--55}}\par
\VerbatimInput[firstline=54,lastline=55,firstnumber=54,numbers=left,numbersep=3pt,frame=lines,framesep=2pt,rulecolor=\color{gray!60},breaklines=true,fontsize=\fontsize{6.6}{7.6}\selectfont]{../../splat_editor/ui/gui.py}
\subsubsection{\texorpdfstring{\code{Gui.shutdown(self)}}{Gui.shutdown}}
{\footnotesize\textit{Dòng 57--58}}\par
\VerbatimInput[firstline=57,lastline=58,firstnumber=57,numbers=left,numbersep=3pt,frame=lines,framesep=2pt,rulecolor=\color{gray!60},breaklines=true,fontsize=\fontsize{6.6}{7.6}\selectfont]{../../splat_editor/ui/gui.py}
\subsection{\texorpdfstring{\code{splat_editor/ui/theme.py}}{splat_editor/ui/theme.py}}
{\footnotesize 74 dòng}\par
\textbf{Vai trò.} Font và màu sắc của giao diện ImGui.

\paragraph{Mô tả đầu file và cách dùng.} Phần chú thích đầu file do tác giả viết, nêu ý tưởng và ví dụ sử dụng:
\VerbatimInput[firstline=1,lastline=7,firstnumber=1,numbers=left,numbersep=3pt,frame=lines,framesep=2pt,rulecolor=\color{gray!60},breaklines=true,fontsize=\fontsize{6.2}{7.2}\selectfont]{../../splat_editor/ui/theme.py}
\subsubsection*{Khai báo mức module (dòng 8--10)}
\VerbatimInput[firstline=8,lastline=10,firstnumber=8,numbers=left,numbersep=3pt,frame=lines,framesep=2pt,rulecolor=\color{gray!60},breaklines=true,fontsize=\fontsize{6.6}{7.6}\selectfont]{../../splat_editor/ui/theme.py}
\subsubsection{\texorpdfstring{\code{class Theme}}{Theme}}
{\footnotesize\textit{Dòng 13--74}}\par
Các phương thức: \code{apply}, \code{_load_font}.

\VerbatimInput[firstline=13,lastline=47,firstnumber=13,numbers=left,numbersep=3pt,frame=lines,framesep=2pt,rulecolor=\color{gray!60},breaklines=true,fontsize=\fontsize{6.6}{7.6}\selectfont]{../../splat_editor/ui/theme.py}
\subsubsection{\texorpdfstring{\code{Theme.apply(cls)}}{Theme.apply}}
{\footnotesize\textit{Dòng 49--66}}\par
\VerbatimInput[firstline=49,lastline=66,firstnumber=49,numbers=left,numbersep=3pt,frame=lines,framesep=2pt,rulecolor=\color{gray!60},breaklines=true,fontsize=\fontsize{6.6}{7.6}\selectfont]{../../splat_editor/ui/theme.py}
\subsubsection{\texorpdfstring{\code{Theme._load_font(cls)}}{Theme._load_font}}
{\footnotesize\textit{Dòng 68--74}}\par
\VerbatimInput[firstline=68,lastline=74,firstnumber=68,numbers=left,numbersep=3pt,frame=lines,framesep=2pt,rulecolor=\color{gray!60},breaklines=true,fontsize=\fontsize{6.6}{7.6}\selectfont]{../../splat_editor/ui/theme.py}
\subsection{\texorpdfstring{\code{splat_editor/ui/actions.py}}{splat_editor/ui/actions.py}}
{\footnotesize 70 dòng}\par
\textbf{Vai trò.} Mẫu Command cho giao diện: \code{ActionRegistry} là nơi duy nhất khai báo nhãn, phím tắt và hàm xử lý. Menu, nút bấm, phím tắt và bảng Keyboard Shortcuts
đều đọc cùng danh sách này.

\paragraph{Mô tả đầu file và cách dùng.} Phần chú thích đầu file do tác giả viết, nêu ý tưởng và ví dụ sử dụng:
\VerbatimInput[firstline=1,lastline=17,firstnumber=1,numbers=left,numbersep=3pt,frame=lines,framesep=2pt,rulecolor=\color{gray!60},breaklines=true,fontsize=\fontsize{6.2}{7.2}\selectfont]{../../splat_editor/ui/actions.py}
\subsubsection*{Khai báo mức module (dòng 18--19)}
\VerbatimInput[firstline=18,lastline=19,firstnumber=18,numbers=left,numbersep=3pt,frame=lines,framesep=2pt,rulecolor=\color{gray!60},breaklines=true,fontsize=\fontsize{6.6}{7.6}\selectfont]{../../splat_editor/ui/actions.py}
\subsubsection{\texorpdfstring{\code{class Action}}{Action}}
{\footnotesize\textit{Dòng 22--33}}\par
\VerbatimInput[firstline=22,lastline=33,firstnumber=22,numbers=left,numbersep=3pt,frame=lines,framesep=2pt,rulecolor=\color{gray!60},breaklines=true,fontsize=\fontsize{6.6}{7.6}\selectfont]{../../splat_editor/ui/actions.py}
\subsubsection{\texorpdfstring{\code{class ActionRegistry}}{ActionRegistry}}
{\footnotesize\textit{Dòng 36--70}}\par
Các phương thức: \code{__init__}, \code{add}, \code{__getitem__}, \code{__iter__}, \code{trigger}, \code{handle_key}, \code{run_pending}.

\VerbatimInput[firstline=36,lastline=36,firstnumber=36,numbers=left,numbersep=3pt,frame=lines,framesep=2pt,rulecolor=\color{gray!60},breaklines=true,fontsize=\fontsize{6.6}{7.6}\selectfont]{../../splat_editor/ui/actions.py}
\subsubsection{\texorpdfstring{\code{ActionRegistry.__init__(self)}}{ActionRegistry.__init__}}
{\footnotesize\textit{Dòng 37--39}}\par
\VerbatimInput[firstline=37,lastline=39,firstnumber=37,numbers=left,numbersep=3pt,frame=lines,framesep=2pt,rulecolor=\color{gray!60},breaklines=true,fontsize=\fontsize{6.6}{7.6}\selectfont]{../../splat_editor/ui/actions.py}
\subsubsection{\texorpdfstring{\code{ActionRegistry.add(self, name, action)}}{ActionRegistry.add}}
{\footnotesize\textit{Dòng 41--42}}\par
\VerbatimInput[firstline=41,lastline=42,firstnumber=41,numbers=left,numbersep=3pt,frame=lines,framesep=2pt,rulecolor=\color{gray!60},breaklines=true,fontsize=\fontsize{6.6}{7.6}\selectfont]{../../splat_editor/ui/actions.py}
\subsubsection{\texorpdfstring{\code{ActionRegistry.__getitem__(self, name)}}{ActionRegistry.__getitem__}}
{\footnotesize\textit{Dòng 44--45}}\par
\VerbatimInput[firstline=44,lastline=45,firstnumber=44,numbers=left,numbersep=3pt,frame=lines,framesep=2pt,rulecolor=\color{gray!60},breaklines=true,fontsize=\fontsize{6.6}{7.6}\selectfont]{../../splat_editor/ui/actions.py}
\subsubsection{\texorpdfstring{\code{ActionRegistry.__iter__(self)}}{ActionRegistry.__iter__}}
{\footnotesize\textit{Dòng 47--48}}\par
\VerbatimInput[firstline=47,lastline=48,firstnumber=47,numbers=left,numbersep=3pt,frame=lines,framesep=2pt,rulecolor=\color{gray!60},breaklines=true,fontsize=\fontsize{6.6}{7.6}\selectfont]{../../splat_editor/ui/actions.py}
\subsubsection{\texorpdfstring{\code{ActionRegistry.trigger(self, name)}}{ActionRegistry.trigger}}
{\footnotesize\textit{Dòng 50--58}}\par
\VerbatimInput[firstline=50,lastline=58,firstnumber=50,numbers=left,numbersep=3pt,frame=lines,framesep=2pt,rulecolor=\color{gray!60},breaklines=true,fontsize=\fontsize{6.6}{7.6}\selectfont]{../../splat_editor/ui/actions.py}
\subsubsection{\texorpdfstring{\code{ActionRegistry.handle_key(self, key, ctrl)}}{ActionRegistry.handle_key}}
{\footnotesize\textit{Dòng 60--66}}\par
\textbf{Mô tả trong mã nguồn.} Chạy hành động gắn với phím. Trả về True nếu có hành động khớp.

\textbf{Giải thích.} Tìm hành động có phím tắt khớp và chạy; trả \code{True} nếu có.

\VerbatimInput[firstline=60,lastline=66,firstnumber=60,numbers=left,numbersep=3pt,frame=lines,framesep=2pt,rulecolor=\color{gray!60},breaklines=true,fontsize=\fontsize{6.6}{7.6}\selectfont]{../../splat_editor/ui/actions.py}
\subsubsection{\texorpdfstring{\code{ActionRegistry.run_pending(self)}}{ActionRegistry.run_pending}}
{\footnotesize\textit{Dòng 68--70}}\par
\VerbatimInput[firstline=68,lastline=70,firstnumber=68,numbers=left,numbersep=3pt,frame=lines,framesep=2pt,rulecolor=\color{gray!60},breaklines=true,fontsize=\fontsize{6.6}{7.6}\selectfont]{../../splat_editor/ui/actions.py}
\subsection{\texorpdfstring{\code{splat_editor/ui/input.py}}{splat_editor/ui/input.py}}
{\footnotesize 136 dòng}\par
\textbf{Vai trò.} \code{InputController} nhận sự kiện chuột/bàn phím: chuyển cho ImGui trước, phần còn lại điều khiển camera (xoay, kéo, zoom, bay), cọ chọn và phím tắt.

\paragraph{Mô tả đầu file và cách dùng.} Phần chú thích đầu file do tác giả viết, nêu ý tưởng và ví dụ sử dụng:
\VerbatimInput[firstline=1,lastline=16,firstnumber=1,numbers=left,numbersep=3pt,frame=lines,framesep=2pt,rulecolor=\color{gray!60},breaklines=true,fontsize=\fontsize{6.2}{7.2}\selectfont]{../../splat_editor/ui/input.py}
\subsubsection*{Khai báo mức module (dòng 17--17)}
\VerbatimInput[firstline=17,lastline=17,firstnumber=17,numbers=left,numbersep=3pt,frame=lines,framesep=2pt,rulecolor=\color{gray!60},breaklines=true,fontsize=\fontsize{6.6}{7.6}\selectfont]{../../splat_editor/ui/input.py}
\subsubsection{\texorpdfstring{\code{class InputController}}{InputController}}
{\footnotesize\textit{Dòng 20--136}}\par
Các phương thức: \code{__init__}, \code{brush}, \code{looking}, \code{_set_drag}, \code{reset}, \code{on_mouse_move}, \code{on_mouse_button}, \code{on_scroll}, \code{on_drop}, \code{on_key}, \code{update}.

\VerbatimInput[firstline=20,lastline=20,firstnumber=20,numbers=left,numbersep=3pt,frame=lines,framesep=2pt,rulecolor=\color{gray!60},breaklines=true,fontsize=\fontsize{6.6}{7.6}\selectfont]{../../splat_editor/ui/input.py}
\subsubsection{\texorpdfstring{\code{InputController.__init__(self, app)}}{InputController.__init__}}
{\footnotesize\textit{Dòng 21--34}}\par
\VerbatimInput[firstline=21,lastline=34,firstnumber=21,numbers=left,numbersep=3pt,frame=lines,framesep=2pt,rulecolor=\color{gray!60},breaklines=true,fontsize=\fontsize{6.6}{7.6}\selectfont]{../../splat_editor/ui/input.py}
\subsubsection{\texorpdfstring{\code{InputController.brush(self)}}{InputController.brush}}
{\footnotesize\textit{Dòng 36--38}}\par
\VerbatimInput[firstline=36,lastline=38,firstnumber=36,numbers=left,numbersep=3pt,frame=lines,framesep=2pt,rulecolor=\color{gray!60},breaklines=true,fontsize=\fontsize{6.6}{7.6}\selectfont]{../../splat_editor/ui/input.py}
\subsubsection{\texorpdfstring{\code{InputController.looking(self)}}{InputController.looking}}
{\footnotesize\textit{Dòng 40--43}}\par
\textbf{Mô tả trong mã nguồn.} Đang giữ chuột phải: nhìn quanh + bay kiểu Minecraft.

\VerbatimInput[firstline=40,lastline=43,firstnumber=40,numbers=left,numbersep=3pt,frame=lines,framesep=2pt,rulecolor=\color{gray!60},breaklines=true,fontsize=\fontsize{6.6}{7.6}\selectfont]{../../splat_editor/ui/input.py}
\subsubsection{\texorpdfstring{\code{InputController._set_drag(self, button)}}{InputController._set_drag}}
{\footnotesize\textit{Dòng 45--50}}\par
\VerbatimInput[firstline=45,lastline=50,firstnumber=45,numbers=left,numbersep=3pt,frame=lines,framesep=2pt,rulecolor=\color{gray!60},breaklines=true,fontsize=\fontsize{6.6}{7.6}\selectfont]{../../splat_editor/ui/input.py}
\subsubsection{\texorpdfstring{\code{InputController.reset(self)}}{InputController.reset}}
{\footnotesize\textit{Dòng 52--55}}\par
\textbf{Mô tả trong mã nguồn.} Gọi sau khi mở hộp thoại hệ thống (cửa sổ mất focus, chuột có thể đã nhả).

\VerbatimInput[firstline=52,lastline=55,firstnumber=52,numbers=left,numbersep=3pt,frame=lines,framesep=2pt,rulecolor=\color{gray!60},breaklines=true,fontsize=\fontsize{6.6}{7.6}\selectfont]{../../splat_editor/ui/input.py}
\subsubsection{\texorpdfstring{\code{InputController.on_mouse_move(self, window, x, y)}}{InputController.on_mouse_move}}
{\footnotesize\textit{Dòng 58--74}}\par
\VerbatimInput[firstline=58,lastline=74,firstnumber=58,numbers=left,numbersep=3pt,frame=lines,framesep=2pt,rulecolor=\color{gray!60},breaklines=true,fontsize=\fontsize{6.6}{7.6}\selectfont]{../../splat_editor/ui/input.py}
\subsubsection{\texorpdfstring{\code{InputController.on_mouse_button(self, window, button, action, mods)}}{InputController.on_mouse_button}}
{\footnotesize\textit{Dòng 76--99}}\par
\VerbatimInput[firstline=76,lastline=99,firstnumber=76,numbers=left,numbersep=3pt,frame=lines,framesep=2pt,rulecolor=\color{gray!60},breaklines=true,fontsize=\fontsize{6.6}{7.6}\selectfont]{../../splat_editor/ui/input.py}
\subsubsection{\texorpdfstring{\code{InputController.on_scroll(self, window, dx, dy)}}{InputController.on_scroll}}
{\footnotesize\textit{Dòng 101--102}}\par
\VerbatimInput[firstline=101,lastline=102,firstnumber=101,numbers=left,numbersep=3pt,frame=lines,framesep=2pt,rulecolor=\color{gray!60},breaklines=true,fontsize=\fontsize{6.6}{7.6}\selectfont]{../../splat_editor/ui/input.py}
\subsubsection{\texorpdfstring{\code{InputController.on_drop(self, window, paths)}}{InputController.on_drop}}
{\footnotesize\textit{Dòng 105--106}}\par
\VerbatimInput[firstline=105,lastline=106,firstnumber=105,numbers=left,numbersep=3pt,frame=lines,framesep=2pt,rulecolor=\color{gray!60},breaklines=true,fontsize=\fontsize{6.6}{7.6}\selectfont]{../../splat_editor/ui/input.py}
\subsubsection{\texorpdfstring{\code{InputController.on_key(self, window, key, scancode, action, mods)}}{InputController.on_key}}
{\footnotesize\textit{Dòng 109--114}}\par
\VerbatimInput[firstline=109,lastline=114,firstnumber=109,numbers=left,numbersep=3pt,frame=lines,framesep=2pt,rulecolor=\color{gray!60},breaklines=true,fontsize=\fontsize{6.6}{7.6}\selectfont]{../../splat_editor/ui/input.py}
\subsubsection{\texorpdfstring{\code{InputController.update(self, dt)}}{InputController.update}}
{\footnotesize\textit{Dòng 116--136}}\par
\textbf{Mô tả trong mã nguồn.} Bay kiểu Minecraft sáng tạo (giữ phím), gọi mỗi khung hình.

W A S D: đi ngang, Space: lên, Shift: xuống, Ctrl: bay nhanh (khi giữ chuột phải).

\textbf{Giải thích.} Mỗi khung hình khi giữ chuột phải: đọc W/A/S/D/Q/E và gọi \code{OrbitCamera.fly} (hoặc \code{Application.walk} ở chế độ sinh tồn).

\VerbatimInput[firstline=116,lastline=136,firstnumber=116,numbers=left,numbersep=3pt,frame=lines,framesep=2pt,rulecolor=\color{gray!60},breaklines=true,fontsize=\fontsize{6.6}{7.6}\selectfont]{../../splat_editor/ui/input.py}
\subsection{\texorpdfstring{\code{splat_editor/ui/layout.py}}{splat_editor/ui/layout.py}}
{\footnotesize 53 dòng}\par
\textbf{Vai trò.} Bố cục màn hình: menu trên cùng, bảng bên trái, thanh trạng thái dưới cùng, phần còn lại là vùng 3D.

\paragraph{Mô tả đầu file và cách dùng.} Phần chú thích đầu file do tác giả viết, nêu ý tưởng và ví dụ sử dụng:
\VerbatimInput[firstline=1,lastline=21,firstnumber=1,numbers=left,numbersep=3pt,frame=lines,framesep=2pt,rulecolor=\color{gray!60},breaklines=true,fontsize=\fontsize{6.2}{7.2}\selectfont]{../../splat_editor/ui/layout.py}
\subsubsection*{Khai báo mức module (dòng 22--22)}
\VerbatimInput[firstline=22,lastline=22,firstnumber=22,numbers=left,numbersep=3pt,frame=lines,framesep=2pt,rulecolor=\color{gray!60},breaklines=true,fontsize=\fontsize{6.6}{7.6}\selectfont]{../../splat_editor/ui/layout.py}
\subsubsection{\texorpdfstring{\code{class Layout}}{Layout}}
{\footnotesize\textit{Dòng 25--53}}\par
Các phương thức: \code{__init__}, \code{update}, \code{_body_height}, \code{side_panel}, \code{viewport}, \code{status_bar}.

\VerbatimInput[firstline=25,lastline=26,firstnumber=25,numbers=left,numbersep=3pt,frame=lines,framesep=2pt,rulecolor=\color{gray!60},breaklines=true,fontsize=\fontsize{6.6}{7.6}\selectfont]{../../splat_editor/ui/layout.py}
\subsubsection{\texorpdfstring{\code{Layout.__init__(self)}}{Layout.__init__}}
{\footnotesize\textit{Dòng 28--32}}\par
\VerbatimInput[firstline=28,lastline=32,firstnumber=28,numbers=left,numbersep=3pt,frame=lines,framesep=2pt,rulecolor=\color{gray!60},breaklines=true,fontsize=\fontsize{6.6}{7.6}\selectfont]{../../splat_editor/ui/layout.py}
\subsubsection{\texorpdfstring{\code{Layout.update(self, width, height, menu_height, status_height)}}{Layout.update}}
{\footnotesize\textit{Dòng 34--36}}\par
\textbf{Giải thích.} Tính ba hình chữ nhật (bảng trái, thanh trạng thái, vùng 3D) từ kích thước cửa sổ, chiều cao menu và thanh trạng thái.

\VerbatimInput[firstline=34,lastline=36,firstnumber=34,numbers=left,numbersep=3pt,frame=lines,framesep=2pt,rulecolor=\color{gray!60},breaklines=true,fontsize=\fontsize{6.6}{7.6}\selectfont]{../../splat_editor/ui/layout.py}
\subsubsection{\texorpdfstring{\code{Layout._body_height(self)}}{Layout._body_height}}
{\footnotesize\textit{Dòng 38--40}}\par
\VerbatimInput[firstline=38,lastline=40,firstnumber=38,numbers=left,numbersep=3pt,frame=lines,framesep=2pt,rulecolor=\color{gray!60},breaklines=true,fontsize=\fontsize{6.6}{7.6}\selectfont]{../../splat_editor/ui/layout.py}
\subsubsection{\texorpdfstring{\code{Layout.side_panel(self)}}{Layout.side_panel}}
{\footnotesize\textit{Dòng 42--44}}\par
\VerbatimInput[firstline=42,lastline=44,firstnumber=42,numbers=left,numbersep=3pt,frame=lines,framesep=2pt,rulecolor=\color{gray!60},breaklines=true,fontsize=\fontsize{6.6}{7.6}\selectfont]{../../splat_editor/ui/layout.py}
\subsubsection{\texorpdfstring{\code{Layout.viewport(self)}}{Layout.viewport}}
{\footnotesize\textit{Dòng 46--49}}\par
\VerbatimInput[firstline=46,lastline=49,firstnumber=46,numbers=left,numbersep=3pt,frame=lines,framesep=2pt,rulecolor=\color{gray!60},breaklines=true,fontsize=\fontsize{6.6}{7.6}\selectfont]{../../splat_editor/ui/layout.py}
\subsubsection{\texorpdfstring{\code{Layout.status_bar(self)}}{Layout.status_bar}}
{\footnotesize\textit{Dòng 51--53}}\par
\VerbatimInput[firstline=51,lastline=53,firstnumber=51,numbers=left,numbersep=3pt,frame=lines,framesep=2pt,rulecolor=\color{gray!60},breaklines=true,fontsize=\fontsize{6.6}{7.6}\selectfont]{../../splat_editor/ui/layout.py}
\subsection{\texorpdfstring{\code{splat_editor/ui/menu_bar.py}}{splat_editor/ui/menu_bar.py}}
{\footnotesize 47 dòng}\par
\textbf{Vai trò.} Thanh menu File/Edit/View/Help sinh từ \code{ActionRegistry}.

\paragraph{Mô tả đầu file và cách dùng.} Phần chú thích đầu file do tác giả viết, nêu ý tưởng và ví dụ sử dụng:
\VerbatimInput[firstline=1,lastline=7,firstnumber=1,numbers=left,numbersep=3pt,frame=lines,framesep=2pt,rulecolor=\color{gray!60},breaklines=true,fontsize=\fontsize{6.2}{7.2}\selectfont]{../../splat_editor/ui/menu_bar.py}
\subsubsection*{Khai báo mức module (dòng 8--8)}
\VerbatimInput[firstline=8,lastline=8,firstnumber=8,numbers=left,numbersep=3pt,frame=lines,framesep=2pt,rulecolor=\color{gray!60},breaklines=true,fontsize=\fontsize{6.6}{7.6}\selectfont]{../../splat_editor/ui/menu_bar.py}
\subsubsection{\texorpdfstring{\code{class MainMenuBar}}{MainMenuBar}}
{\footnotesize\textit{Dòng 11--47}}\par
Các phương thức: \code{__init__}, \code{draw}, \code{_draw_item}.

\VerbatimInput[firstline=11,lastline=21,firstnumber=11,numbers=left,numbersep=3pt,frame=lines,framesep=2pt,rulecolor=\color{gray!60},breaklines=true,fontsize=\fontsize{6.6}{7.6}\selectfont]{../../splat_editor/ui/menu_bar.py}
\subsubsection{\texorpdfstring{\code{MainMenuBar.__init__(self, actions)}}{MainMenuBar.__init__}}
{\footnotesize\textit{Dòng 23--24}}\par
\VerbatimInput[firstline=23,lastline=24,firstnumber=23,numbers=left,numbersep=3pt,frame=lines,framesep=2pt,rulecolor=\color{gray!60},breaklines=true,fontsize=\fontsize{6.6}{7.6}\selectfont]{../../splat_editor/ui/menu_bar.py}
\subsubsection{\texorpdfstring{\code{MainMenuBar.draw(self)}}{MainMenuBar.draw}}
{\footnotesize\textit{Dòng 26--40}}\par
\textbf{Mô tả trong mã nguồn.} Vẽ menu, trả về chiều cao (pixel) của thanh menu.

\VerbatimInput[firstline=26,lastline=40,firstnumber=26,numbers=left,numbersep=3pt,frame=lines,framesep=2pt,rulecolor=\color{gray!60},breaklines=true,fontsize=\fontsize{6.6}{7.6}\selectfont]{../../splat_editor/ui/menu_bar.py}
\subsubsection{\texorpdfstring{\code{MainMenuBar._draw_item(self, name)}}{MainMenuBar._draw_item}}
{\footnotesize\textit{Dòng 42--47}}\par
\VerbatimInput[firstline=42,lastline=47,firstnumber=42,numbers=left,numbersep=3pt,frame=lines,framesep=2pt,rulecolor=\color{gray!60},breaklines=true,fontsize=\fontsize{6.6}{7.6}\selectfont]{../../splat_editor/ui/menu_bar.py}
\subsection{\texorpdfstring{\code{splat_editor/ui/side_panel.py}}{splat_editor/ui/side_panel.py}}
{\footnotesize 39 dòng}\par
\textbf{Vai trò.} Bảng cố định bên trái, gồm các mục thu gọn được.

\paragraph{Mô tả đầu file và cách dùng.} Phần chú thích đầu file do tác giả viết, nêu ý tưởng và ví dụ sử dụng:
\VerbatimInput[firstline=1,lastline=7,firstnumber=1,numbers=left,numbersep=3pt,frame=lines,framesep=2pt,rulecolor=\color{gray!60},breaklines=true,fontsize=\fontsize{6.2}{7.2}\selectfont]{../../splat_editor/ui/side_panel.py}
\subsubsection*{Khai báo mức module (dòng 8--11)}
\VerbatimInput[firstline=8,lastline=11,firstnumber=8,numbers=left,numbersep=3pt,frame=lines,framesep=2pt,rulecolor=\color{gray!60},breaklines=true,fontsize=\fontsize{6.6}{7.6}\selectfont]{../../splat_editor/ui/side_panel.py}
\subsubsection{\texorpdfstring{\code{class SidePanel}}{SidePanel}}
{\footnotesize\textit{Dòng 14--39}}\par
Các phương thức: \code{__init__}, \code{draw}.

\VerbatimInput[firstline=14,lastline=17,firstnumber=14,numbers=left,numbersep=3pt,frame=lines,framesep=2pt,rulecolor=\color{gray!60},breaklines=true,fontsize=\fontsize{6.6}{7.6}\selectfont]{../../splat_editor/ui/side_panel.py}
\subsubsection{\texorpdfstring{\code{SidePanel.__init__(self, app)}}{SidePanel.__init__}}
{\footnotesize\textit{Dòng 19--21}}\par
\VerbatimInput[firstline=19,lastline=21,firstnumber=19,numbers=left,numbersep=3pt,frame=lines,framesep=2pt,rulecolor=\color{gray!60},breaklines=true,fontsize=\fontsize{6.6}{7.6}\selectfont]{../../splat_editor/ui/side_panel.py}
\subsubsection{\texorpdfstring{\code{SidePanel.draw(self, rect)}}{SidePanel.draw}}
{\footnotesize\textit{Dòng 23--39}}\par
\VerbatimInput[firstline=23,lastline=39,firstnumber=23,numbers=left,numbersep=3pt,frame=lines,framesep=2pt,rulecolor=\color{gray!60},breaklines=true,fontsize=\fontsize{6.6}{7.6}\selectfont]{../../splat_editor/ui/side_panel.py}
\subsection{\texorpdfstring{\code{splat_editor/ui/sections.py}}{splat_editor/ui/sections.py}}
{\footnotesize 248 dòng}\par
\textbf{Vai trò.} Các mục của bảng trái: Layers, Transform, Selection, View. Mỗi mục là một lớp con của \code{PanelSection} (mẫu Template Method).

\paragraph{Mô tả đầu file và cách dùng.} Phần chú thích đầu file do tác giả viết, nêu ý tưởng và ví dụ sử dụng:
\VerbatimInput[firstline=1,lastline=20,firstnumber=1,numbers=left,numbersep=3pt,frame=lines,framesep=2pt,rulecolor=\color{gray!60},breaklines=true,fontsize=\fontsize{6.2}{7.2}\selectfont]{../../splat_editor/ui/sections.py}
\subsubsection*{Khai báo mức module (dòng 21--27)}
\VerbatimInput[firstline=21,lastline=27,firstnumber=21,numbers=left,numbersep=3pt,frame=lines,framesep=2pt,rulecolor=\color{gray!60},breaklines=true,fontsize=\fontsize{6.6}{7.6}\selectfont]{../../splat_editor/ui/sections.py}
\subsubsection{\texorpdfstring{\code{class PanelSection(ABC)}}{PanelSection}}
{\footnotesize\textit{Dòng 30--54}}\par
Các phương thức: \code{__init__}, \code{editor}, \code{scene}, \code{header}, \code{is_visible}, \code{draw}.

\VerbatimInput[firstline=30,lastline=31,firstnumber=30,numbers=left,numbersep=3pt,frame=lines,framesep=2pt,rulecolor=\color{gray!60},breaklines=true,fontsize=\fontsize{6.6}{7.6}\selectfont]{../../splat_editor/ui/sections.py}
\subsubsection{\texorpdfstring{\code{PanelSection.__init__(self, app)}}{PanelSection.__init__}}
{\footnotesize\textit{Dòng 33--34}}\par
\VerbatimInput[firstline=33,lastline=34,firstnumber=33,numbers=left,numbersep=3pt,frame=lines,framesep=2pt,rulecolor=\color{gray!60},breaklines=true,fontsize=\fontsize{6.6}{7.6}\selectfont]{../../splat_editor/ui/sections.py}
\subsubsection{\texorpdfstring{\code{PanelSection.editor(self)}}{PanelSection.editor}}
{\footnotesize\textit{Dòng 36--38}}\par
\VerbatimInput[firstline=36,lastline=38,firstnumber=36,numbers=left,numbersep=3pt,frame=lines,framesep=2pt,rulecolor=\color{gray!60},breaklines=true,fontsize=\fontsize{6.6}{7.6}\selectfont]{../../splat_editor/ui/sections.py}
\subsubsection{\texorpdfstring{\code{PanelSection.scene(self)}}{PanelSection.scene}}
{\footnotesize\textit{Dòng 40--42}}\par
\VerbatimInput[firstline=40,lastline=42,firstnumber=40,numbers=left,numbersep=3pt,frame=lines,framesep=2pt,rulecolor=\color{gray!60},breaklines=true,fontsize=\fontsize{6.6}{7.6}\selectfont]{../../splat_editor/ui/sections.py}
\subsubsection{\texorpdfstring{\code{PanelSection.header(self)}}{PanelSection.header}}
{\footnotesize\textit{Dòng 44--47}}\par
\textbf{Mô tả trong mã nguồn.} Nhãn tiêu đề; phần sau '\#\#\#' giữ ID cố định khi chữ thay đổi.

\VerbatimInput[firstline=44,lastline=47,firstnumber=44,numbers=left,numbersep=3pt,frame=lines,framesep=2pt,rulecolor=\color{gray!60},breaklines=true,fontsize=\fontsize{6.6}{7.6}\selectfont]{../../splat_editor/ui/sections.py}
\subsubsection{\texorpdfstring{\code{PanelSection.is_visible(self)}}{PanelSection.is_visible}}
{\footnotesize\textit{Dòng 49--50}}\par
\VerbatimInput[firstline=49,lastline=50,firstnumber=49,numbers=left,numbersep=3pt,frame=lines,framesep=2pt,rulecolor=\color{gray!60},breaklines=true,fontsize=\fontsize{6.6}{7.6}\selectfont]{../../splat_editor/ui/sections.py}
\subsubsection{\texorpdfstring{\code{PanelSection.draw(self)}}{PanelSection.draw}}
{\footnotesize\textit{Dòng 52--54}}\par
\VerbatimInput[firstline=52,lastline=54,firstnumber=52,numbers=left,numbersep=3pt,frame=lines,framesep=2pt,rulecolor=\color{gray!60},breaklines=true,fontsize=\fontsize{6.6}{7.6}\selectfont]{../../splat_editor/ui/sections.py}
\subsubsection{\texorpdfstring{\code{class LayersSection(PanelSection)}}{LayersSection}}
{\footnotesize\textit{Dòng 59--120}}\par
Các phương thức: \code{header}, \code{draw}, \code{_draw_list}.

\VerbatimInput[firstline=59,lastline=61,firstnumber=59,numbers=left,numbersep=3pt,frame=lines,framesep=2pt,rulecolor=\color{gray!60},breaklines=true,fontsize=\fontsize{6.6}{7.6}\selectfont]{../../splat_editor/ui/sections.py}
\subsubsection{\texorpdfstring{\code{LayersSection.header(self)}}{LayersSection.header}}
{\footnotesize\textit{Dòng 63--65}}\par
\VerbatimInput[firstline=63,lastline=65,firstnumber=63,numbers=left,numbersep=3pt,frame=lines,framesep=2pt,rulecolor=\color{gray!60},breaklines=true,fontsize=\fontsize{6.6}{7.6}\selectfont]{../../splat_editor/ui/sections.py}
\subsubsection{\texorpdfstring{\code{LayersSection.draw(self)}}{LayersSection.draw}}
{\footnotesize\textit{Dòng 67--83}}\par
\VerbatimInput[firstline=67,lastline=83,firstnumber=67,numbers=left,numbersep=3pt,frame=lines,framesep=2pt,rulecolor=\color{gray!60},breaklines=true,fontsize=\fontsize{6.6}{7.6}\selectfont]{../../splat_editor/ui/sections.py}
\subsubsection{\texorpdfstring{\code{LayersSection._draw_list(self)}}{LayersSection._draw_list}}
{\footnotesize\textit{Dòng 85--120}}\par
\VerbatimInput[firstline=85,lastline=120,firstnumber=85,numbers=left,numbersep=3pt,frame=lines,framesep=2pt,rulecolor=\color{gray!60},breaklines=true,fontsize=\fontsize{6.6}{7.6}\selectfont]{../../splat_editor/ui/sections.py}
\subsubsection{\texorpdfstring{\code{class TransformSection(PanelSection)}}{TransformSection}}
{\footnotesize\textit{Dòng 125--192}}\par
Các phương thức: \code{is_visible}, \code{draw}, \code{_draw_level}.

\VerbatimInput[firstline=125,lastline=127,firstnumber=125,numbers=left,numbersep=3pt,frame=lines,framesep=2pt,rulecolor=\color{gray!60},breaklines=true,fontsize=\fontsize{6.6}{7.6}\selectfont]{../../splat_editor/ui/sections.py}
\subsubsection{\texorpdfstring{\code{TransformSection.is_visible(self)}}{TransformSection.is_visible}}
{\footnotesize\textit{Dòng 129--130}}\par
\VerbatimInput[firstline=129,lastline=130,firstnumber=129,numbers=left,numbersep=3pt,frame=lines,framesep=2pt,rulecolor=\color{gray!60},breaklines=true,fontsize=\fontsize{6.6}{7.6}\selectfont]{../../splat_editor/ui/sections.py}
\subsubsection{\texorpdfstring{\code{TransformSection.draw(self)}}{TransformSection.draw}}
{\footnotesize\textit{Dòng 132--154}}\par
\VerbatimInput[firstline=132,lastline=154,firstnumber=132,numbers=left,numbersep=3pt,frame=lines,framesep=2pt,rulecolor=\color{gray!60},breaklines=true,fontsize=\fontsize{6.6}{7.6}\selectfont]{../../splat_editor/ui/sections.py}
\subsubsection{\texorpdfstring{\code{TransformSection._draw_level(self, t)}}{TransformSection._draw_level}}
{\footnotesize\textit{Dòng 156--192}}\par
\textbf{Mô tả trong mã nguồn.} Chỉnh cho model nằm ngang: xoay theo hướng đang nhìn, quanh điểm camera đang nhìn vào.

Roll: quanh trục nhìn tới (nằm ngang) -\textgreater{} đường chân trời nghiêng trái / phải trên màn hình. Pitch: quanh trục ngang của màn hình -\textgreater{} mặt sàn chúi về phía mình / ngửa ra xa.

\VerbatimInput[firstline=156,lastline=192,firstnumber=156,numbers=left,numbersep=3pt,frame=lines,framesep=2pt,rulecolor=\color{gray!60},breaklines=true,fontsize=\fontsize{6.6}{7.6}\selectfont]{../../splat_editor/ui/sections.py}
\subsubsection{\texorpdfstring{\code{class SelectionSection(PanelSection)}}{SelectionSection}}
{\footnotesize\textit{Dòng 197--227}}\par
Các phương thức: \code{draw}.

\VerbatimInput[firstline=197,lastline=199,firstnumber=197,numbers=left,numbersep=3pt,frame=lines,framesep=2pt,rulecolor=\color{gray!60},breaklines=true,fontsize=\fontsize{6.6}{7.6}\selectfont]{../../splat_editor/ui/sections.py}
\subsubsection{\texorpdfstring{\code{SelectionSection.draw(self)}}{SelectionSection.draw}}
{\footnotesize\textit{Dòng 201--227}}\par
\VerbatimInput[firstline=201,lastline=227,firstnumber=201,numbers=left,numbersep=3pt,frame=lines,framesep=2pt,rulecolor=\color{gray!60},breaklines=true,fontsize=\fontsize{6.6}{7.6}\selectfont]{../../splat_editor/ui/sections.py}
\subsubsection{\texorpdfstring{\code{class ViewSection(PanelSection)}}{ViewSection}}
{\footnotesize\textit{Dòng 232--248}}\par
Các phương thức: \code{draw}.

\VerbatimInput[firstline=232,lastline=234,firstnumber=232,numbers=left,numbersep=3pt,frame=lines,framesep=2pt,rulecolor=\color{gray!60},breaklines=true,fontsize=\fontsize{6.6}{7.6}\selectfont]{../../splat_editor/ui/sections.py}
\subsubsection{\texorpdfstring{\code{ViewSection.draw(self)}}{ViewSection.draw}}
{\footnotesize\textit{Dòng 236--248}}\par
\VerbatimInput[firstline=236,lastline=248,firstnumber=236,numbers=left,numbersep=3pt,frame=lines,framesep=2pt,rulecolor=\color{gray!60},breaklines=true,fontsize=\fontsize{6.6}{7.6}\selectfont]{../../splat_editor/ui/sections.py}
\subsection{\texorpdfstring{\code{splat_editor/ui/world_sections.py}}{splat_editor/ui/world_sections.py}}
{\footnotesize 133 dòng}\par
\textbf{Vai trò.} Hai mục của bảng trái cho ánh sáng và sinh tồn: \code{ModeSection} (chế độ sáng tạo/sinh tồn) và \code{LightingSection} (mặt trời, đèn, bóng).

\paragraph{Mô tả đầu file và cách dùng.} Phần chú thích đầu file do tác giả viết, nêu ý tưởng và ví dụ sử dụng:
\VerbatimInput[firstline=1,lastline=6,firstnumber=1,numbers=left,numbersep=3pt,frame=lines,framesep=2pt,rulecolor=\color{gray!60},breaklines=true,fontsize=\fontsize{6.2}{7.2}\selectfont]{../../splat_editor/ui/world_sections.py}
\subsubsection*{Khai báo mức module (dòng 7--11)}
\VerbatimInput[firstline=7,lastline=11,firstnumber=7,numbers=left,numbersep=3pt,frame=lines,framesep=2pt,rulecolor=\color{gray!60},breaklines=true,fontsize=\fontsize{6.6}{7.6}\selectfont]{../../splat_editor/ui/world_sections.py}
\subsubsection{\texorpdfstring{\code{class ModeSection(PanelSection)}}{ModeSection}}
{\footnotesize\textit{Dòng 14--52}}\par
Các phương thức: \code{draw}.

\VerbatimInput[firstline=14,lastline=16,firstnumber=14,numbers=left,numbersep=3pt,frame=lines,framesep=2pt,rulecolor=\color{gray!60},breaklines=true,fontsize=\fontsize{6.6}{7.6}\selectfont]{../../splat_editor/ui/world_sections.py}
\subsubsection{\texorpdfstring{\code{ModeSection.draw(self)}}{ModeSection.draw}}
{\footnotesize\textit{Dòng 18--52}}\par
\VerbatimInput[firstline=18,lastline=52,firstnumber=18,numbers=left,numbersep=3pt,frame=lines,framesep=2pt,rulecolor=\color{gray!60},breaklines=true,fontsize=\fontsize{6.6}{7.6}\selectfont]{../../splat_editor/ui/world_sections.py}
\subsubsection{\texorpdfstring{\code{class LightingSection(PanelSection)}}{LightingSection}}
{\footnotesize\textit{Dòng 55--134}}\par
Các phương thức: \code{draw}, \code{_draw_sun}, \code{_draw_lights}.

\VerbatimInput[firstline=55,lastline=57,firstnumber=55,numbers=left,numbersep=3pt,frame=lines,framesep=2pt,rulecolor=\color{gray!60},breaklines=true,fontsize=\fontsize{6.6}{7.6}\selectfont]{../../splat_editor/ui/world_sections.py}
\subsubsection{\texorpdfstring{\code{LightingSection.draw(self)}}{LightingSection.draw}}
{\footnotesize\textit{Dòng 59--86}}\par
\VerbatimInput[firstline=59,lastline=86,firstnumber=59,numbers=left,numbersep=3pt,frame=lines,framesep=2pt,rulecolor=\color{gray!60},breaklines=true,fontsize=\fontsize{6.6}{7.6}\selectfont]{../../splat_editor/ui/world_sections.py}
\subsubsection{\texorpdfstring{\code{LightingSection._draw_sun(self, lighting)}}{LightingSection._draw_sun}}
{\footnotesize\textit{Dòng 88--99}}\par
\VerbatimInput[firstline=88,lastline=99,firstnumber=88,numbers=left,numbersep=3pt,frame=lines,framesep=2pt,rulecolor=\color{gray!60},breaklines=true,fontsize=\fontsize{6.6}{7.6}\selectfont]{../../splat_editor/ui/world_sections.py}
\subsubsection{\texorpdfstring{\code{LightingSection._draw_lights(self, lighting)}}{LightingSection._draw_lights}}
{\footnotesize\textit{Dòng 101--133}}\par
\VerbatimInput[firstline=101,lastline=133,firstnumber=101,numbers=left,numbersep=3pt,frame=lines,framesep=2pt,rulecolor=\color{gray!60},breaklines=true,fontsize=\fontsize{6.6}{7.6}\selectfont]{../../splat_editor/ui/world_sections.py}
\subsection{\texorpdfstring{\code{splat_editor/ui/status_bar.py}}{splat_editor/ui/status_bar.py}}
{\footnotesize 53 dòng}\par
\textbf{Vai trò.} Thanh trạng thái dưới cùng: thông báo gần nhất bên trái, thống kê bên phải.

\paragraph{Mô tả đầu file và cách dùng.} Phần chú thích đầu file do tác giả viết, nêu ý tưởng và ví dụ sử dụng:
\VerbatimInput[firstline=1,lastline=8,firstnumber=1,numbers=left,numbersep=3pt,frame=lines,framesep=2pt,rulecolor=\color{gray!60},breaklines=true,fontsize=\fontsize{6.2}{7.2}\selectfont]{../../splat_editor/ui/status_bar.py}
\subsubsection*{Khai báo mức module (dòng 9--9)}
\VerbatimInput[firstline=9,lastline=9,firstnumber=9,numbers=left,numbersep=3pt,frame=lines,framesep=2pt,rulecolor=\color{gray!60},breaklines=true,fontsize=\fontsize{6.6}{7.6}\selectfont]{../../splat_editor/ui/status_bar.py}
\subsubsection{\texorpdfstring{\code{class StatusBar}}{StatusBar}}
{\footnotesize\textit{Dòng 12--53}}\par
Các phương thức: \code{__init__}, \code{height}, \code{draw}, \code{_fit}.

\VerbatimInput[firstline=12,lastline=16,firstnumber=12,numbers=left,numbersep=3pt,frame=lines,framesep=2pt,rulecolor=\color{gray!60},breaklines=true,fontsize=\fontsize{6.6}{7.6}\selectfont]{../../splat_editor/ui/status_bar.py}
\subsubsection{\texorpdfstring{\code{StatusBar.__init__(self, app)}}{StatusBar.__init__}}
{\footnotesize\textit{Dòng 18--19}}\par
\VerbatimInput[firstline=18,lastline=19,firstnumber=18,numbers=left,numbersep=3pt,frame=lines,framesep=2pt,rulecolor=\color{gray!60},breaklines=true,fontsize=\fontsize{6.6}{7.6}\selectfont]{../../splat_editor/ui/status_bar.py}
\subsubsection{\texorpdfstring{\code{StatusBar.height(self)}}{StatusBar.height}}
{\footnotesize\textit{Dòng 21--22}}\par
\VerbatimInput[firstline=21,lastline=22,firstnumber=21,numbers=left,numbersep=3pt,frame=lines,framesep=2pt,rulecolor=\color{gray!60},breaklines=true,fontsize=\fontsize{6.6}{7.6}\selectfont]{../../splat_editor/ui/status_bar.py}
\subsubsection{\texorpdfstring{\code{StatusBar.draw(self, rect)}}{StatusBar.draw}}
{\footnotesize\textit{Dòng 24--43}}\par
\VerbatimInput[firstline=24,lastline=43,firstnumber=24,numbers=left,numbersep=3pt,frame=lines,framesep=2pt,rulecolor=\color{gray!60},breaklines=true,fontsize=\fontsize{6.6}{7.6}\selectfont]{../../splat_editor/ui/status_bar.py}
\subsubsection{\texorpdfstring{\code{StatusBar._fit(text, width)}}{StatusBar._fit}}
{\footnotesize\textit{Dòng 45--53}}\par
\textbf{Mô tả trong mã nguồn.} Cắt bớt chữ (thêm '...') cho vừa `width` pixel.

\VerbatimInput[firstline=45,lastline=53,firstnumber=45,numbers=left,numbersep=3pt,frame=lines,framesep=2pt,rulecolor=\color{gray!60},breaklines=true,fontsize=\fontsize{6.6}{7.6}\selectfont]{../../splat_editor/ui/status_bar.py}
\subsection{\texorpdfstring{\code{splat_editor/ui/viewport_overlay.py}}{splat_editor/ui/viewport_overlay.py}}
{\footnotesize 108 dòng}\par
\textbf{Vai trò.} Vẽ chồng lên vùng 3D: gợi ý khi cảnh trống, nhãn và vòng tròn cọ, tọa độ người chơi, đồng hồ.

\paragraph{Mô tả đầu file và cách dùng.} Phần chú thích đầu file do tác giả viết, nêu ý tưởng và ví dụ sử dụng:
\VerbatimInput[firstline=1,lastline=7,firstnumber=1,numbers=left,numbersep=3pt,frame=lines,framesep=2pt,rulecolor=\color{gray!60},breaklines=true,fontsize=\fontsize{6.2}{7.2}\selectfont]{../../splat_editor/ui/viewport_overlay.py}
\subsubsection*{Khai báo mức module (dòng 8--10)}
\VerbatimInput[firstline=8,lastline=10,firstnumber=8,numbers=left,numbersep=3pt,frame=lines,framesep=2pt,rulecolor=\color{gray!60},breaklines=true,fontsize=\fontsize{6.6}{7.6}\selectfont]{../../splat_editor/ui/viewport_overlay.py}
\subsubsection{\texorpdfstring{\code{class ViewportOverlay}}{ViewportOverlay}}
{\footnotesize\textit{Dòng 13--108}}\par
Các phương thức: \code{__init__}, \code{draw}, \code{_draw_empty_hint}, \code{_draw_brush_hint}, \code{_draw_brush_cursor}, \code{_draw_survival}, \code{_draw_coords}, \code{_draw_clock}.

\VerbatimInput[firstline=13,lastline=19,firstnumber=13,numbers=left,numbersep=3pt,frame=lines,framesep=2pt,rulecolor=\color{gray!60},breaklines=true,fontsize=\fontsize{6.6}{7.6}\selectfont]{../../splat_editor/ui/viewport_overlay.py}
\subsubsection{\texorpdfstring{\code{ViewportOverlay.__init__(self, app)}}{ViewportOverlay.__init__}}
{\footnotesize\textit{Dòng 21--22}}\par
\VerbatimInput[firstline=21,lastline=22,firstnumber=21,numbers=left,numbersep=3pt,frame=lines,framesep=2pt,rulecolor=\color{gray!60},breaklines=true,fontsize=\fontsize{6.6}{7.6}\selectfont]{../../splat_editor/ui/viewport_overlay.py}
\subsubsection{\texorpdfstring{\code{ViewportOverlay.draw(self, rect)}}{ViewportOverlay.draw}}
{\footnotesize\textit{Dòng 24--38}}\par
\VerbatimInput[firstline=24,lastline=38,firstnumber=24,numbers=left,numbersep=3pt,frame=lines,framesep=2pt,rulecolor=\color{gray!60},breaklines=true,fontsize=\fontsize{6.6}{7.6}\selectfont]{../../splat_editor/ui/viewport_overlay.py}
\subsubsection{\texorpdfstring{\code{ViewportOverlay._draw_empty_hint(self, rect)}}{ViewportOverlay._draw_empty_hint}}
{\footnotesize\textit{Dòng 40--49}}\par
\VerbatimInput[firstline=40,lastline=49,firstnumber=40,numbers=left,numbersep=3pt,frame=lines,framesep=2pt,rulecolor=\color{gray!60},breaklines=true,fontsize=\fontsize{6.6}{7.6}\selectfont]{../../splat_editor/ui/viewport_overlay.py}
\subsubsection{\texorpdfstring{\code{ViewportOverlay._draw_brush_hint(self, rect)}}{ViewportOverlay._draw_brush_hint}}
{\footnotesize\textit{Dòng 51--59}}\par
\VerbatimInput[firstline=51,lastline=59,firstnumber=51,numbers=left,numbersep=3pt,frame=lines,framesep=2pt,rulecolor=\color{gray!60},breaklines=true,fontsize=\fontsize{6.6}{7.6}\selectfont]{../../splat_editor/ui/viewport_overlay.py}
\subsubsection{\texorpdfstring{\code{ViewportOverlay._draw_brush_cursor(self, rect)}}{ViewportOverlay._draw_brush_cursor}}
{\footnotesize\textit{Dòng 61--68}}\par
\VerbatimInput[firstline=61,lastline=68,firstnumber=61,numbers=left,numbersep=3pt,frame=lines,framesep=2pt,rulecolor=\color{gray!60},breaklines=true,fontsize=\fontsize{6.6}{7.6}\selectfont]{../../splat_editor/ui/viewport_overlay.py}
\subsubsection{\texorpdfstring{\code{ViewportOverlay._draw_survival(self, rect)}}{ViewportOverlay._draw_survival}}
{\footnotesize\textit{Dòng 70--85}}\par
\VerbatimInput[firstline=70,lastline=85,firstnumber=70,numbers=left,numbersep=3pt,frame=lines,framesep=2pt,rulecolor=\color{gray!60},breaklines=true,fontsize=\fontsize{6.6}{7.6}\selectfont]{../../splat_editor/ui/viewport_overlay.py}
\subsubsection{\texorpdfstring{\code{ViewportOverlay._draw_coords(self, rect)}}{ViewportOverlay._draw_coords}}
{\footnotesize\textit{Dòng 87--97}}\par
\textbf{Mô tả trong mã nguồn.} Tọa độ XYZ của người đang đứng (vị trí camera) ở góc dưới trái vùng 3D.

\VerbatimInput[firstline=87,lastline=97,firstnumber=87,numbers=left,numbersep=3pt,frame=lines,framesep=2pt,rulecolor=\color{gray!60},breaklines=true,fontsize=\fontsize{6.6}{7.6}\selectfont]{../../splat_editor/ui/viewport_overlay.py}
\subsubsection{\texorpdfstring{\code{ViewportOverlay._draw_clock(self, rect, lighting)}}{ViewportOverlay._draw_clock}}
{\footnotesize\textit{Dòng 99--108}}\par
\VerbatimInput[firstline=99,lastline=108,firstnumber=99,numbers=left,numbersep=3pt,frame=lines,framesep=2pt,rulecolor=\color{gray!60},breaklines=true,fontsize=\fontsize{6.6}{7.6}\selectfont]{../../splat_editor/ui/viewport_overlay.py}
\subsection{\texorpdfstring{\code{splat_editor/ui/shortcuts_window.py}}{splat_editor/ui/shortcuts_window.py}}
{\footnotesize 68 dòng}\par
\textbf{Vai trò.} Cửa sổ Keyboard Shortcuts (F1), lấy danh sách phím trực tiếp từ \code{ActionRegistry}.

\paragraph{Mô tả đầu file và cách dùng.} Phần chú thích đầu file do tác giả viết, nêu ý tưởng và ví dụ sử dụng:
\VerbatimInput[firstline=1,lastline=8,firstnumber=1,numbers=left,numbersep=3pt,frame=lines,framesep=2pt,rulecolor=\color{gray!60},breaklines=true,fontsize=\fontsize{6.2}{7.2}\selectfont]{../../splat_editor/ui/shortcuts_window.py}
\subsubsection*{Khai báo mức module (dòng 9--9)}
\VerbatimInput[firstline=9,lastline=9,firstnumber=9,numbers=left,numbersep=3pt,frame=lines,framesep=2pt,rulecolor=\color{gray!60},breaklines=true,fontsize=\fontsize{6.6}{7.6}\selectfont]{../../splat_editor/ui/shortcuts_window.py}
\subsubsection{\texorpdfstring{\code{class ShortcutsWindow}}{ShortcutsWindow}}
{\footnotesize\textit{Dòng 12--68}}\par
Các phương thức: \code{__init__}, \code{toggle}, \code{draw}, \code{_table}.

\VerbatimInput[firstline=12,lastline=28,firstnumber=12,numbers=left,numbersep=3pt,frame=lines,framesep=2pt,rulecolor=\color{gray!60},breaklines=true,fontsize=\fontsize{6.6}{7.6}\selectfont]{../../splat_editor/ui/shortcuts_window.py}
\subsubsection{\texorpdfstring{\code{ShortcutsWindow.__init__(self, actions)}}{ShortcutsWindow.__init__}}
{\footnotesize\textit{Dòng 30--32}}\par
\VerbatimInput[firstline=30,lastline=32,firstnumber=30,numbers=left,numbersep=3pt,frame=lines,framesep=2pt,rulecolor=\color{gray!60},breaklines=true,fontsize=\fontsize{6.6}{7.6}\selectfont]{../../splat_editor/ui/shortcuts_window.py}
\subsubsection{\texorpdfstring{\code{ShortcutsWindow.toggle(self)}}{ShortcutsWindow.toggle}}
{\footnotesize\textit{Dòng 34--35}}\par
\VerbatimInput[firstline=34,lastline=35,firstnumber=34,numbers=left,numbersep=3pt,frame=lines,framesep=2pt,rulecolor=\color{gray!60},breaklines=true,fontsize=\fontsize{6.6}{7.6}\selectfont]{../../splat_editor/ui/shortcuts_window.py}
\subsubsection{\texorpdfstring{\code{ShortcutsWindow.draw(self)}}{ShortcutsWindow.draw}}
{\footnotesize\textit{Dòng 37--55}}\par
\VerbatimInput[firstline=37,lastline=55,firstnumber=37,numbers=left,numbersep=3pt,frame=lines,framesep=2pt,rulecolor=\color{gray!60},breaklines=true,fontsize=\fontsize{6.6}{7.6}\selectfont]{../../splat_editor/ui/shortcuts_window.py}
\subsubsection{\texorpdfstring{\code{ShortcutsWindow._table(self, title, rows)}}{ShortcutsWindow._table}}
{\footnotesize\textit{Dòng 57--68}}\par
\VerbatimInput[firstline=57,lastline=68,firstnumber=57,numbers=left,numbersep=3pt,frame=lines,framesep=2pt,rulecolor=\color{gray!60},breaklines=true,fontsize=\fontsize{6.6}{7.6}\selectfont]{../../splat_editor/ui/shortcuts_window.py}
\subsection{\texorpdfstring{\code{splat_editor/ui/widgets.py}}{splat_editor/ui/widgets.py}}
{\footnotesize 50 dòng}\par
\textbf{Vai trò.} Các widget ImGui nhỏ dùng lại: nút toàn chiều ngang, hàng nút, nút nguy hiểm, tooltip.

\paragraph{Mô tả đầu file và cách dùng.} Phần chú thích đầu file do tác giả viết, nêu ý tưởng và ví dụ sử dụng:
\VerbatimInput[firstline=1,lastline=12,firstnumber=1,numbers=left,numbersep=3pt,frame=lines,framesep=2pt,rulecolor=\color{gray!60},breaklines=true,fontsize=\fontsize{6.2}{7.2}\selectfont]{../../splat_editor/ui/widgets.py}
\subsubsection*{Khai báo mức module (dòng 13--15)}
\VerbatimInput[firstline=13,lastline=15,firstnumber=13,numbers=left,numbersep=3pt,frame=lines,framesep=2pt,rulecolor=\color{gray!60},breaklines=true,fontsize=\fontsize{6.6}{7.6}\selectfont]{../../splat_editor/ui/widgets.py}
\subsubsection{\texorpdfstring{\code{full_width_button(label)}}{full_width_button}}
{\footnotesize\textit{Dòng 18--19}}\par
\VerbatimInput[firstline=18,lastline=19,firstnumber=18,numbers=left,numbersep=3pt,frame=lines,framesep=2pt,rulecolor=\color{gray!60},breaklines=true,fontsize=\fontsize{6.6}{7.6}\selectfont]{../../splat_editor/ui/widgets.py}
\subsubsection{\texorpdfstring{\code{button_row(labels)}}{button_row}}
{\footnotesize\textit{Dòng 22--32}}\par
\textbf{Mô tả trong mã nguồn.} Một hàng nút chia đều chiều ngang. Trả về chỉ số nút được bấm hoặc None.

\VerbatimInput[firstline=22,lastline=32,firstnumber=22,numbers=left,numbersep=3pt,frame=lines,framesep=2pt,rulecolor=\color{gray!60},breaklines=true,fontsize=\fontsize{6.6}{7.6}\selectfont]{../../splat_editor/ui/widgets.py}
\subsubsection{\texorpdfstring{\code{danger_button(label)}}{danger_button}}
{\footnotesize\textit{Dòng 35--41}}\par
\textbf{Mô tả trong mã nguồn.} Nút toàn chiều ngang màu đỏ (thao tác xoá).

\VerbatimInput[firstline=35,lastline=41,firstnumber=35,numbers=left,numbersep=3pt,frame=lines,framesep=2pt,rulecolor=\color{gray!60},breaklines=true,fontsize=\fontsize{6.6}{7.6}\selectfont]{../../splat_editor/ui/widgets.py}
\subsubsection{\texorpdfstring{\code{tooltip(text)}}{tooltip}}
{\footnotesize\textit{Dòng 44--46}}\par
\textbf{Mô tả trong mã nguồn.} Hiện tooltip cho widget vừa vẽ khi rê chuột lên.

\VerbatimInput[firstline=44,lastline=46,firstnumber=44,numbers=left,numbersep=3pt,frame=lines,framesep=2pt,rulecolor=\color{gray!60},breaklines=true,fontsize=\fontsize{6.6}{7.6}\selectfont]{../../splat_editor/ui/widgets.py}
\subsubsection{\texorpdfstring{\code{color_u32(rgba)}}{color_u32}}
{\footnotesize\textit{Dòng 49--50}}\par
\VerbatimInput[firstline=49,lastline=50,firstnumber=49,numbers=left,numbersep=3pt,frame=lines,framesep=2pt,rulecolor=\color{gray!60},breaklines=true,fontsize=\fontsize{6.6}{7.6}\selectfont]{../../splat_editor/ui/widgets.py}
\subsection{\texorpdfstring{\code{splat_editor/ui/file_dialog.py}}{splat_editor/ui/file_dialog.py}}
{\footnotesize 37 dòng}\par
\textbf{Vai trò.} Hộp thoại chọn file của hệ điều hành bằng \code{tkinter}.

\paragraph{Mô tả đầu file và cách dùng.} Phần chú thích đầu file do tác giả viết, nêu ý tưởng và ví dụ sử dụng:
\VerbatimInput[firstline=1,lastline=8,firstnumber=1,numbers=left,numbersep=3pt,frame=lines,framesep=2pt,rulecolor=\color{gray!60},breaklines=true,fontsize=\fontsize{6.2}{7.2}\selectfont]{../../splat_editor/ui/file_dialog.py}
\subsubsection*{Khai báo mức module (dòng 9--10)}
\VerbatimInput[firstline=9,lastline=10,firstnumber=9,numbers=left,numbersep=3pt,frame=lines,framesep=2pt,rulecolor=\color{gray!60},breaklines=true,fontsize=\fontsize{6.6}{7.6}\selectfont]{../../splat_editor/ui/file_dialog.py}
\subsubsection{\texorpdfstring{\code{class FileDialog}}{FileDialog}}
{\footnotesize\textit{Dòng 13--37}}\par
Các phương thức: \code{_run}, \code{open_file}, \code{open_files}, \code{save_file}.

\VerbatimInput[firstline=13,lastline=14,firstnumber=13,numbers=left,numbersep=3pt,frame=lines,framesep=2pt,rulecolor=\color{gray!60},breaklines=true,fontsize=\fontsize{6.6}{7.6}\selectfont]{../../splat_editor/ui/file_dialog.py}
\subsubsection{\texorpdfstring{\code{FileDialog._run(cls, ask, **options)}}{FileDialog._run}}
{\footnotesize\textit{Dòng 16--24}}\par
\VerbatimInput[firstline=16,lastline=24,firstnumber=16,numbers=left,numbersep=3pt,frame=lines,framesep=2pt,rulecolor=\color{gray!60},breaklines=true,fontsize=\fontsize{6.6}{7.6}\selectfont]{../../splat_editor/ui/file_dialog.py}
\subsubsection{\texorpdfstring{\code{FileDialog.open_file(cls)}}{FileDialog.open_file}}
{\footnotesize\textit{Dòng 26--28}}\par
\VerbatimInput[firstline=26,lastline=28,firstnumber=26,numbers=left,numbersep=3pt,frame=lines,framesep=2pt,rulecolor=\color{gray!60},breaklines=true,fontsize=\fontsize{6.6}{7.6}\selectfont]{../../splat_editor/ui/file_dialog.py}
\subsubsection{\texorpdfstring{\code{FileDialog.open_files(cls)}}{FileDialog.open_files}}
{\footnotesize\textit{Dòng 30--32}}\par
\VerbatimInput[firstline=30,lastline=32,firstnumber=30,numbers=left,numbersep=3pt,frame=lines,framesep=2pt,rulecolor=\color{gray!60},breaklines=true,fontsize=\fontsize{6.6}{7.6}\selectfont]{../../splat_editor/ui/file_dialog.py}
\subsubsection{\texorpdfstring{\code{FileDialog.save_file(cls, initial='scene.ply')}}{FileDialog.save_file}}
{\footnotesize\textit{Dòng 34--37}}\par
\VerbatimInput[firstline=34,lastline=37,firstnumber=34,numbers=left,numbersep=3pt,frame=lines,framesep=2pt,rulecolor=\color{gray!60},breaklines=true,fontsize=\fontsize{6.6}{7.6}\selectfont]{../../splat_editor/ui/file_dialog.py}

\clearpage
\section{Thư mục render/shaders: vertex và fragment shader}\label{ch:shader}
Chương này giải thích \textbf{toàn bộ mã GLSL} trong \code{splat_editor/render/shaders/}. Mỗi chương trình shader gồm một cặp
\textbf{vertex shader} (\code{.vert}) và \textbf{fragment shader} (\code{.frag}); riêng pass 1 của bộ vẽ gaussian và bản đồ bóng chỉ dùng vertex shader.
Các file đều dùng \code{\#version 330 core} (OpenGL 3.3, tương thích máy yếu).

\subsection*{Vertex shader và fragment shader làm gì}
\begin{itemize}
\item \textbf{Vertex shader} chạy \emph{một lần cho mỗi đỉnh}. Nhiệm vụ chính là ghi \code{gl_Position} (vị trí đỉnh trong không gian clip) và các biến \code{out} sẽ được nội suy.
\item Giữa hai giai đoạn, GPU làm bước \textbf{rasterization}: ghép các đỉnh thành hình (điểm, tam giác), chia chúng thành các pixel (fragment) và \textbf{nội suy} biến \code{out} của vertex shader.
\item \textbf{Fragment shader} chạy \emph{một lần cho mỗi pixel} của hình. Nhận biến \code{in} đã nội suy, ghi màu ra \code{fragColor}. Lệnh \code{discard} bỏ pixel đó.
\end{itemize}

\subsection*{Từ khóa GLSL xuất hiện nhiều lần}
\begin{center}\small
\begin{tabular}{@{}p{0.2\textwidth}p{0.72\textwidth}@{}}\toprule
Từ khóa & Ý nghĩa\\\midrule
\code{in} / \code{out} & Đầu vào / đầu ra của giai đoạn. \code{out} của vertex nối với \code{in} cùng tên của fragment.\\
\code{layout(location=n)} & Số thứ tự thuộc tính đỉnh, khớp với \code{glVertexAttribPointer} phía Python.\\
\code{uniform} & Hằng số cho cả lần vẽ (ma trận, màu, thời gian\ldots), được đặt bằng \code{ShaderProgram.set_*}.\\
\code{flat} & Không nội suy; lấy giá trị của đỉnh đầu (dùng cho số nguyên như màu đóng gói).\\
\code{samplerBuffer}, \code{usamplerBuffer} & Texture buffer: mảng lớn đọc ngẫu nhiên bằng \code{texelFetch(buf, i)}; bản \code{u} đọc ra số nguyên không dấu.\\
\code{texelFetch} & Đọc đúng một texel theo chỉ số nguyên, \textbf{không nội suy}.\\
\code{gl_VertexID}, \code{gl_PointSize}, \code{gl_PointCoord} & Số thứ tự đỉnh, kích thước điểm (pixel) và tọa độ trong điểm.\\
\code{mat3(m)}, \code{mat4(c0,c1,c2,c3)} & Ma trận GLSL điền theo \textbf{cột}; \code{mat3(model)} lấy khối $3\times3$ trên trái.\\
\code{uintBitsToFloat} & Đọc lại bit của uint như float (dữ liệu đóng gói trong texel \code{RGBA32UI}).\\
\code{fwidth(x)} & Độ biến thiên của $x$ giữa các pixel kề nhau (đạo hàm màn hình), dùng khử răng cưa.\\
\bottomrule
\end{tabular}
\end{center}

\subsection*{Hai pass của bộ vẽ gaussian}
Mỗi khung hình \code{SplatRenderer} chạy hai chương trình nối tiếp (xem Chương 4):
\begin{enumerate}
\item \textbf{Pass 1} \code{splat_prep.vert} (transform feedback, không có fragment shader): mỗi gaussian đã chọn chạy đúng 1 lần, ghi ellipse 2D và màu vào một buffer 32 byte/gaussian.
\item \textbf{Pass 2} \code{splat.vert} + \code{splat.frag}: đọc buffer đó theo thứ tự xa $\to$ gần, vẽ mỗi gaussian một hình vuông và trộn alpha.
\end{enumerate}
Cách chia này giúp phần tính nặng (ma trận, SH, Phong, bóng) chỉ làm \textbf{một lần cho mỗi gaussian}, thay vì 4 lần cho 4 đỉnh của hình vuông.

\subsection*{Bố cục dữ liệu gaussian trên GPU (\code{uData})}
Mỗi gaussian là 2 texel RGBA32UI (8 uint32 $=32$ byte):
\begin{center}\small
\begin{tabular}{@{}lll@{}}\toprule
Texel & Kênh & Nội dung\\\midrule
0 & \code{xyz} & tâm (float32, đọc bằng \code{uintBitsToFloat})\\
0 & \code{w} & màu đỏ, xanh lá (hai half-float)\\
1 & \code{x} & màu xanh dương và độ đục $\alpha$ (hai half-float)\\
1 & \code{y}, \code{z[0:16]} & scale $x,y$ và scale $z$ (half-float)\\
1 & \code{z[16:32]} & chỉ số layer\\
1 & \code{w} & quaternion $4\times8$ bit\\\bottomrule
\end{tabular}
\end{center}


\subsection{\texorpdfstring{\code{splat_prep.vert}}{splat_prep.vert}}
{\footnotesize Giai đoạn: Vertex (transform feedback) \;$\cdot$\; 215 dòng \;$\cdot$\; \code{splat_editor/render/shaders/splat_prep.vert}}\par
\textbf{Vai trò.} Shader nặng nhất chương trình. Nó \textbf{không vẽ gì}; kết quả được ghi vào buffer rồi pass 2 đọc lại. Đầu vào duy nhất là \code{aSplatIndex}, vị trí gaussian trên GPU (thứ tự Morton);
mọi dữ liệu khác được tra từ texture buffer. Toán học: \S\ref{s:ewa} (EWA), \S\ref{s:phong} (Phong, bóng đổ), \S\ref{s:sun} (mặt trời, bản đồ bóng).

%% ---- merged math: s:ewa
\subsubsection{Cơ sở toán: EWA Splatting, chiếu ellipsoid 3D thành ellipse 2D}\label{s:ewa}
Mục này giải thích toán của pass 1 (\code{splat_prep.vert}); phần chiếu và trộn nằm ở \code{splat.vert}, \code{splat.frag} (mục \ref{s:blend}).

\paragraph{Hiệp phương sai 3D.}
$\Sig = R S S^{\top} R^{\top} = M M^{\top}$, với $M=R\,\mathrm{diag}(s)$ (shader: \code{M = R * mat3(scale...)}, \code{sigma = M * transpose(M)}). $\Sig$ đối xứng bán xác định dương; trục ellipsoid là cột của $R$, bán trục là $s$.

\paragraph{Jacobian của phép chiếu.}
Điểm ảnh: $x_{px}=f_{px}v_x/(-v_z)$, $y_{px}=f_{px}v_y/(-v_z)$. Đạo hàm theo $(v_x,v_y,v_z)$ tại tâm:
\[
J=\begin{bmatrix}\frac{f}{-v_z}&0&\frac{f v_x}{v_z^2}\\0&\frac{f}{-v_z}&\frac{f v_y}{v_z^2}\\0&0&0\end{bmatrix},\qquad
f_{px}=\tfrac H2 P_{11}=\frac{H/2}{\tan(\mathrm{fov}/2)}
\]
(\code{splat_renderer}: \code{uFocal = height/2 * projection[1,1]}). Phép chia cho $c_w$ là phi tuyến, nên dùng \textbf{xấp xỉ tuyến tính bậc nhất} tại tâm (khai triển Taylor) để ellipsoid vẫn cho ra ellipse.

\paragraph{Hiệp phương sai 2D.}
\[
\Sig' = J\,W\,\Sig\,W^{\top}J^{\top}.
\]
Trong shader gộp thêm ma trận layer $sA$: \code{T = J * mat3(ModelView) * linear}, \code{cov = T Σ Tᵀ} ($\Sig$ ở tọa độ file; \code{linear = s·A} đã gồm co giãn nên không cần $s^2$ riêng). Lấy khối $2\times2$: $a=\mathrm{cov}_{00}$, $b=\mathrm{cov}_{01}$, $c=\mathrm{cov}_{11}$.

\textbf{Bộ lọc thông thấp:} $a\mathrel{+}=0.3$, $c\mathrel{+}=0.3$ (tương đương cộng $0.3I$), bảo đảm gaussian rộng ít nhất khoảng 1 pixel, chống răng cưa như EWA gốc và 3DGS.

\paragraph{Trị riêng và vector riêng của ma trận 2x2.}
\begin{align*}
\mathrm{mid}&=\tfrac{a+c}{2}, & \mathrm{radius}&=\sqrt{\left(\tfrac{a-c}{2}\right)^2+b^2},\\
\lambda_1&=\mathrm{mid}+\mathrm{radius}, & \lambda_2&=\max(\mathrm{mid}-\mathrm{radius},\,0.1)\ \text{(kẹp chống suy biến)},\\
\mathrm{axis}_1&=\mathrm{normalize}(b,\lambda_1-a), & \mathrm{axis}_2&=(-\mathrm{axis}_{1,y},\mathrm{axis}_{1,x}).
\end{align*}
Nếu $|b|<10^{-6}$ thì dùng trục $x$ hoặc $y$. Kiểm tra: $(A-\lambda_1 I)v=0\Rightarrow(a-\lambda_1)v_1+bv_2=0\Rightarrow v\propto(b,\lambda_1-a)$. Bán trục ellipse là $\sqrt{\lambda_1},\sqrt{\lambda_2}$ (độ lệch chuẩn theo hai hướng riêng).

\paragraph{Kích thước hình vuông (billboard).}
\begin{itemize}
\item $\mathrm{extent}=\min\!\big(\sqrt{2\ln(255\alpha)},\,3\big)$, từ phương trình $\alpha e^{-r^2/2}=1/255\Rightarrow r^2=2\ln(255\alpha)$. Gaussian mờ cho hình vuông nhỏ; tối đa $3\sigma$.
\item $\mathrm{major}=\min(\mathrm{extent}\cdot\sqrt{\lambda_1},1024)\cdot\mathrm{axis}_1\cdot(2/\mathrm{viewport})$, \code{minor} tương tự với $\lambda_2$. Nhân $2/W$, $2/H$ để đổi pixel sang NDC (NDC rộng 2).
\item \textbf{Loại bỏ (cull):} $\alpha<1/255$; $c_w\le0$ (sau camera); $|\mathrm{ndc}_x|$ hoặc $|\mathrm{ndc}_y|>1.3$ (đệm 30\% vì ellipse có thể lấn vào màn hình).
\end{itemize}

%% ---- merged math: s:phong
\subsubsection{Cơ sở toán: Ánh sáng Phong và bóng đổ}\label{s:phong}
Các hàm \code{shade} và \code{shadowFactor} trong \code{splat_prep.vert} tính ánh sáng cho từng gaussian. Quỹ đạo mặt trời, máy trạng thái ngày--đêm và thiết lập bản đồ bóng (phía CPU, \code{core/lighting.py}) nằm ở mục \ref{s:sun}.

\paragraph{Pháp tuyến gaussian.}
$n$ là trục ngắn nhất của ellipsoid: $\mathrm{axis}=\arg\min|s|$, $n=\mathrm{normalize}(\mathrm{linear}\cdot R[\mathrm{axis}])$ (cột $R$ tương ứng, đưa sang thế giới), lật nếu $n\cdot\mathrm{toCam}<0$. Với gaussian tròn (không có mặt phẳng): $\mathrm{roundness}=\mathrm{smoothstep}(0.35,0.9,s_{\min}/s_{\max})$, $n=\mathrm{normalize}(\mathrm{mix}(n,(0,1,0),\mathrm{roundness}))$. Hàm $\mathrm{smoothstep}(a,b,x)$: $t=\mathrm{clip}((x-a)/(b-a),0,1)$, kết quả $t^2(3-2t)$ (đa thức Hermite bậc 3).

\paragraph{Mô hình Phong.}
\begin{align*}
w(n,l)&=\mathrm{clip}\!\left(\frac{n\cdot l+0.35}{1.35},0,1\right)\quad\text{(wrap diffuse: $n\cdot l=-0.35\mapsto0$, $1\mapsto1$)}\\
I_a&=c\,L_a\\
I_d&=c\sum_j L_j\,w(n,l_j)\,S_j\\
I_s&=k_s\sum_j L_j\max(r_j\cdot v,0)^m S_j,\quad r=\mathrm{reflect}(-l,n)=2(n\cdot l)n-l,\ \text{chỉ khi }n\cdot l>0\\
c_{\mathrm{out}}&=I_a+I_d+I_s\quad\text{(cuối cùng clamp 0..1)}
\end{align*}
$v=\mathrm{normalize}(\mathrm{camPos}-p)$; $m=$ \code{shininess} (mặc định 24); $k_s=$ \code{specular} (0.25).

\paragraph{Đèn điểm: suy giảm.}
$\mathrm{att}=\mathrm{clamp}(1-d^2/\mathrm{range}^2,0,1)^2$ (bình phương nên mượt về 0 tại $d=\mathrm{range}$). Tối đa 8 đèn, không đổ bóng. Màu $=$ màu $\times$ cường độ.

\subsubsection{Khai báo đầu vào, uniform và biến ra}
{\footnotesize\textit{Dòng 1--42}}\par
\textbf{Giải thích.} Ba nhóm: (1) dữ liệu gaussian đọc từ texture buffer (\code{uData}, \code{uState}, \code{uSh}, \code{uLayers}); (2) uniform camera (\code{uModelView}, \code{uProjection}, \code{uViewport}, \code{uFocal}, \code{uCamPos});
(3) uniform ánh sáng và bóng. Ba biến \code{out} là kết quả pass 1: \code{vCenterMajor} (tâm và trục lớn, NDC), \code{vMinorExtent} (trục nhỏ và bán kính hình vuông tính theo $\sigma$), \code{vColor} (màu RGBA8 đóng gói).

\VerbatimInput[firstline=1,lastline=42,firstnumber=1,numbers=left,numbersep=3pt,frame=lines,framesep=2pt,rulecolor=\color{gray!60},breaklines=true,fontsize=\fontsize{6.6}{7.6}\selectfont]{../../splat_editor/render/shaders/splat_prep.vert}
\subsubsection{Hàm tiện ích: half-float, đóng gói màu, hệ số SH}
{\footnotesize\textit{Dòng 44--60}}\par
\textbf{Giải thích.} \code{halfToFloat} giải half-float IEEE-754 thủ công (GLSL 3.30 chưa có \code{unpackHalf2x16}): $e=0$ là số subnormal $m\cdot2^{-24}$, còn lại $(m+1024)\cdot2^{e-25}$. \code{packColor} nén RGBA vào một \code{uint}
(mỗi kênh 8 bit, cộng $0.5$ để làm tròn). \code{shCoeff} đọc một hệ số SH 8 bit rồi giải lượng tử $(\mathrm{texel}\cdot255-128)\cdot uShScale$.

\textbf{Công thức.}
\[
v=(-1)^{s}\cdot\begin{cases}m\cdot2^{-24}&e=0\\(m+1024)\cdot2^{e-25}&e\neq0\end{cases},\qquad
h\approx(\mathrm{texel}\cdot255-128)\cdot\mathrm{uShScale}
\]


\VerbatimInput[firstline=44,lastline=60,firstnumber=44,numbers=left,numbersep=3pt,frame=lines,framesep=2pt,rulecolor=\color{gray!60},breaklines=true,fontsize=\fontsize{6.6}{7.6}\selectfont]{../../splat_editor/render/shaders/splat_prep.vert}
\subsubsection{Màu theo hướng nhìn: shColor}
{\footnotesize\textit{Dòng 62--84}}\par
\textbf{Giải thích.} Cộng phần SH bậc 1, 2, 3 vào màu theo hướng nhìn $d=(x,y,z)$. Mỗi dòng là một hàm cơ sở (hằng số $C_1,C_2,C_3$) nhân hệ số tương ứng; \code{if (uShDegree < 2) return c;} dừng sớm khi người dùng chọn bậc thấp, tiết kiệm tính toán.

\textbf{Công thức.}
\[
c_{SH}(d)=\sum_{k=1}^{15}Y_k(d)\,h_k,\qquad d=\frac{A^{\top}(p_w-e)}{\|A^{\top}(p_w-e)\|}
\]


\VerbatimInput[firstline=62,lastline=84,firstnumber=62,numbers=left,numbersep=3pt,frame=lines,framesep=2pt,rulecolor=\color{gray!60},breaklines=true,fontsize=\fontsize{6.6}{7.6}\selectfont]{../../splat_editor/render/shaders/splat_prep.vert}
\subsubsection{Bóng đổ: shadowFactor}
{\footnotesize\textit{Dòng 86--98}}\par
\textbf{Giải thích.} Dịch điểm cần tô dọc pháp tuyến hai texel (để gaussian không tự đổ bóng lên mình), chiếu vào không gian ánh sáng bằng \code{uLightVP}, đổi NDC sang $[0,1]$, rồi lấy \textbf{9 mẫu} quanh điểm (PCF $3\times3$)
và so độ sâu với \code{uShadowBias}. Kết quả $1$ là sáng, $0$ là bóng, giá trị giữa là mép bóng mềm. Ngoài bản đồ thì coi là sáng.

\textbf{Công thức.}
\[
s=\frac{\mathrm{VP}_{light}\,(p+2\,\mathrm{texel}\,n)}{w}\cdot0.5+0.5,\qquad
S=\frac19\sum_{i,j=-1}^{1}\Big[\,s_z-\mathrm{bias}\le D\Big(s_{xy}+\frac{(i,j)}{2048}\Big)\Big]
\]


\VerbatimInput[firstline=86,lastline=98,firstnumber=86,numbers=left,numbersep=3pt,frame=lines,framesep=2pt,rulecolor=\color{gray!60},breaklines=true,fontsize=\fontsize{6.6}{7.6}\selectfont]{../../splat_editor/render/shaders/splat_prep.vert}
\subsubsection{Chiếu sáng Phong: shade}
{\footnotesize\textit{Dòng 100--122}}\par
\textbf{Giải thích.} Ba phần: \textbf{ambient} (\code{uAmbient}); \textbf{diffuse} dạng \emph{wrap} $\mathrm{clamp}((n\cdot l+0.35)/1.35,0,1)$ để làm mềm pháp tuyến ước lượng nhiễu; \textbf{specular} $\max(r\cdot v,0)^m$ với $r=\mathrm{reflect}(-l,n)$.
Mặt trời nhân thêm hệ số bóng; đèn điểm nhân hệ số suy giảm $\mathrm{att}=\mathrm{clamp}(1-d^2/\mathrm{range}^2,0,1)^2$. Màu cuối $=\mathrm{albedo}\cdot\mathrm{diffuse}+\mathrm{spec}\cdot k_s$.

\textbf{Công thức.}
\[
w(n,l)=\mathrm{clip}\Big(\frac{n\cdot l+0.35}{1.35},0,1\Big),\qquad r=2(n\cdot l)\,n-l
\]
\[
c_{out}=c\Big(L_a+\sum_jL_j\,w(n,l_j)\,S_j\Big)+k_s\sum_jL_j\max(r_j\cdot v,0)^m\,S_j,\qquad
\mathrm{att}=\mathrm{clamp}\Big(1-\frac{d^2}{\mathrm{range}^2},0,1\Big)^2
\]


\VerbatimInput[firstline=100,lastline=122,firstnumber=100,numbers=left,numbersep=3pt,frame=lines,framesep=2pt,rulecolor=\color{gray!60},breaklines=true,fontsize=\fontsize{6.6}{7.6}\selectfont]{../../splat_editor/render/shaders/splat_prep.vert}
\subsubsection{Hàm cull}
{\footnotesize\textit{Dòng 124--128}}\par
\textbf{Giải thích.} Gán hình vuông suy biến nằm ngoài màn hình ($\mathrm{ndc}=(2,2)$, trục bằng 0, màu 0) để pass 2 không tô pixel nào. Cách này loại gaussian mà không phải đổi số lượng đầu ra của transform feedback.

\VerbatimInput[firstline=124,lastline=128,firstnumber=124,numbers=left,numbersep=3pt,frame=lines,framesep=2pt,rulecolor=\color{gray!60},breaklines=true,fontsize=\fontsize{6.6}{7.6}\selectfont]{../../splat_editor/render/shaders/splat_prep.vert}
\subsubsection{main (1): đọc dữ liệu, loại bỏ, ma trận layer}
{\footnotesize\textit{Dòng 130--156}}\par
\textbf{Giải thích.} Đọc cờ trạng thái (đã xóa thì cull). Đọc hai texel gaussian. Lấy chỉ số layer từ nửa cao của \code{t1.z}, dựng \code{mat4} từ 4 texel \code{uLayers}, đưa tâm sang tọa độ thế giới ($M\,p$), giải $\alpha$.
Tính \code{extent} $=\min(\sqrt{2\ln(255\alpha)},3)$ và chiếu tâm bằng \code{uProjection*uModelView}; loại gaussian sau camera ($w\le0$) hoặc ngoài màn hình quá 30\%.

\textbf{Công thức.}
\[
c_{world}=M\begin{bmatrix}p_f\\1\end{bmatrix},\quad \mathrm{clip}=P\,V\,c_{world},\quad \mathrm{ndc}=\frac{\mathrm{clip}_{xyz}}{\mathrm{clip}_w},\quad
\mathrm{extent}=\min\big(\sqrt{2\ln(255\alpha)},\,3\big)
\]
Giải từ $\alpha\,e^{-r^2/2}=\tfrac1{255}$. Cull nếu $w\le0$, $|\mathrm{ndc}_{x,y}|>1.3$ hoặc $\alpha<\tfrac1{255}$.


\VerbatimInput[firstline=130,lastline=156,firstnumber=130,numbers=left,numbersep=3pt,frame=lines,framesep=2pt,rulecolor=\color{gray!60},breaklines=true,fontsize=\fontsize{6.6}{7.6}\selectfont]{../../splat_editor/render/shaders/splat_prep.vert}
\subsubsection{main (2): hiệp phương sai 3D rồi 2D (EWA)}
{\footnotesize\textit{Dòng 158--179}}\par
\textbf{Giải thích.} Giải mã quaternion 4$\times$8 bit và chuẩn hóa; dựng $R$; $M=R\,\mathrm{diag}(s)$; $\Sigma=MM^\top$. Dựng Jacobian $J$ của phép chiếu tại tâm; $T=J\cdot W\cdot(sA)$ gộp luôn ma trận layer; $\Sigma'=T\Sigma T^\top$.
Lấy ba phần tử $a,b,c$ của khối $2\times2$ và cộng $0.3$ vào đường chéo (bộ lọc thông thấp, ellipse không nhỏ hơn khoảng 1 pixel).

\textbf{Công thức.}
\[
\Sigma=MM^\top,\quad M=R\,\mathrm{diag}(s),\qquad
J=\begin{bmatrix}\dfrac{f}{-v_z}&0&\dfrac{f\,v_x}{v_z^2}\\[2mm]0&\dfrac{f}{-v_z}&\dfrac{f\,v_y}{v_z^2}\\0&0&0\end{bmatrix},\quad f=\frac H2P_{11}
\]
\[
T=J\,W\,(sA),\qquad \Sigma'=T\,\Sigma\,T^{\top}=\begin{bmatrix}a&b\\b&c\end{bmatrix}+0.3\,I
\]


\VerbatimInput[firstline=158,lastline=179,firstnumber=158,numbers=left,numbersep=3pt,frame=lines,framesep=2pt,rulecolor=\color{gray!60},breaklines=true,fontsize=\fontsize{6.6}{7.6}\selectfont]{../../splat_editor/render/shaders/splat_prep.vert}
\subsubsection{main (3): trị riêng, trục ellipse, màu và đầu ra}
{\footnotesize\textit{Dòng 181--215}}\par
\textbf{Giải thích.} Giải trị riêng của $\begin{bmatrix}a&b\\b&c\end{bmatrix}$ bằng công thức đóng ($\lambda_{1,2}=\mathrm{mid}\pm\mathrm{radius}$), vector riêng $(b,\lambda_1-a)$, trục nhỏ vuông góc. Đổi pixel sang NDC (nhân $2/\mathrm{viewport}$).
Tính màu: cộng SH (hướng nhìn đưa ngược về tọa độ file bằng $\mathrm{linear}^\top$), nếu bật ánh sáng thì lấy pháp tuyến là trục ngắn nhất của ellipsoid rồi \code{shade}. Cuối cùng pha màu chọn (\code{mix(color, uSelectColor, 0.7*selected)}) và ghi ba biến đầu ra.

\textbf{Công thức.}
\[
\mathrm{mid}=\frac{a+c}{2},\quad \mathrm{radius}=\sqrt{\Big(\frac{a-c}{2}\Big)^2+b^2},\quad
\lambda_1=\mathrm{mid}+\mathrm{radius},\quad \lambda_2=\max(\mathrm{mid}-\mathrm{radius},0.1)
\]
\[
\hat a_1=\frac{(b,\ \lambda_1-a)}{\|(b,\lambda_1-a)\|},\quad \hat a_2=(-\hat a_{1,y},\hat a_{1,x}),\quad
\mathrm{major}=\min(\mathrm{extent}\sqrt{\lambda_1},1024)\,\hat a_1\cdot\frac2{\mathrm{viewport}}
\]
Pháp tuyến: $n=\mathrm{normalize}(\mathrm{linear}\cdot R_{[\arg\min|s|]})$; pha màu chọn $c\leftarrow\mathrm{mix}(c,c_{sel},0.7\,\mathrm{selected})$.


\VerbatimInput[firstline=181,lastline=215,firstnumber=181,numbers=left,numbersep=3pt,frame=lines,framesep=2pt,rulecolor=\color{gray!60},breaklines=true,fontsize=\fontsize{6.6}{7.6}\selectfont]{../../splat_editor/render/shaders/splat_prep.vert}
\subsection{\texorpdfstring{\code{splat.vert}}{splat.vert}}
{\footnotesize Giai đoạn: Vertex \;$\cdot$\; 20 dòng \;$\cdot$\; \code{splat_editor/render/shaders/splat.vert}}\par
\textbf{Vai trò.} Vẽ mỗi gaussian bằng \textbf{instancing}: 4 đỉnh $(\pm1,\pm1)$ lặp lại cho từng instance; \code{aRank} là chỉ số vào kết quả pass 1, đã xếp xa $\to$ gần.

\subsubsection{Toàn bộ shader}
{\footnotesize\textit{Dòng 1--20}}\par
\textbf{Giải thích.} \code{aCorner} là góc hình vuông. \code{uPrep} chứa kết quả pass 1 (2 texel/gaussian). \code{offset = cornerx*major + cornery*minor} kéo hình vuông theo hai trục ellipse; \code{gl_Position} là tâm cộng offset.
Màu RGBA8 được tách từ uint bằng dịch bit. \code{vCorner = aCorner*extent} là tọa độ trong hệ trục ellipse theo đơn vị $\sigma$; nội suy giữa các đỉnh cho fragment shader biết pixel cách tâm bao nhiêu $\sigma$.

\textbf{Công thức.}
\[
\mathrm{gl\_Position}_{xy}=\mathrm{center}+x_c\,\mathrm{major}+y_c\,\mathrm{minor},\qquad (x_c,y_c)\in\{\pm1\}^2,\qquad \mathrm{vCorner}=(x_c,y_c)\cdot\mathrm{extent}
\]


\VerbatimInput[firstline=1,lastline=20,firstnumber=1,numbers=left,numbersep=3pt,frame=lines,framesep=2pt,rulecolor=\color{gray!60},breaklines=true,fontsize=\fontsize{6.6}{7.6}\selectfont]{../../splat_editor/render/shaders/splat.vert}
%% ---- merged math: s:blend
\subsubsection{Cơ sở toán: hiển thị gaussian (pass 2, đỉnh)}\label{s:blend}
Mỗi instance vẽ 4 góc $(\pm1,\pm1)$: \code{pos_ndc = center + cx*major + cy*minor}; \code{vCorner = (cx,cy)*extent} (đơn vị $\sigma$ trong hệ trục ellipse). Fragment shader và phép trộn alpha xem tiếp ở \code{splat.frag} ngay sau đây; kích thước \code{major}, \code{minor}, \code{extent} lấy từ mục \ref{s:ewa}.

\subsection{\texorpdfstring{\code{splat.frag}}{splat.frag}}
{\footnotesize Giai đoạn: Fragment \;$\cdot$\; 12 dòng \;$\cdot$\; \code{splat_editor/render/shaders/splat.frag}}\par
\textbf{Vai trò.} Biến hình vuông thành gaussian mờ dần.

\subsubsection{Toàn bộ shader}
{\footnotesize\textit{Dòng 1--12}}\par
\textbf{Giải thích.} $r^2=|\mathrm{vCorner}|^2$ là bình phương khoảng cách tới tâm (đơn vị $\sigma$). Bỏ pixel nếu $r^2>9$ (ngoài $3\sigma$). $\alpha_{px}=\alpha\,e^{-r^2/2}$ là mật độ Gauss 2D có đỉnh bằng $\alpha$; bỏ nếu nhỏ hơn $1/255$.
Xuất \code{(rgb*alpha, alpha)}, tức \textbf{premultiplied alpha}, khớp với \code{glBlendFunc(GL_ONE, GL_ONE_MINUS_SRC_ALPHA)}; ta được công thức trộn ``over'' $C=c\,a+(1-a)C_{dst}$.

\textbf{Công thức.}
\[
r^2=\|\mathrm{vCorner}\|^2,\qquad a=\alpha\,e^{-r^2/2},\qquad \mathrm{fragColor}=(a\cdot\mathrm{rgb},\ a)\qquad(\text{bỏ nếu }r^2>9\text{ hoặc }a<\tfrac1{255})
\]


\VerbatimInput[firstline=1,lastline=12,firstnumber=1,numbers=left,numbersep=3pt,frame=lines,framesep=2pt,rulecolor=\color{gray!60},breaklines=true,fontsize=\fontsize{6.6}{7.6}\selectfont]{../../splat_editor/render/shaders/splat.frag}
%% ---- merged math: s:blend-frag
\subsubsection{Cơ sở toán: fragment và trộn alpha ``over''}\label{s:blend-frag}
\paragraph{Fragment (\code{splat.frag}).}
\begin{itemize}
\item $r^2=|\mathrm{vCorner}|^2$; bỏ nếu $r^2>9$ (ngoài $3\sigma$).
\item $a=\alpha\exp(-r^2/2)$ (mật độ Gauss 2D chuẩn hóa về đỉnh bằng $\alpha$). Bỏ nếu $a<1/255$.
\item Xuất $(\mathrm{rgb}\cdot a,\,a)$, tức \textbf{premultiplied alpha}.
\end{itemize}

\paragraph{Trộn ``over'' (\code{glBlendFunc(GL_ONE, GL_ONE_MINUS_SRC_ALPHA)}).}
\[
C_i = c_i a_i + (1-a_i)\,C_{i-1}
\]
(với rgb đã nhân $a$: $C=\mathrm{src.rgb}+(1-\mathrm{src.a})\,\mathrm{dst}$). Phép này không giao hoán, nên phải vẽ \textbf{từ xa đến gần}. Depth test tắt (vật trong suốt không được che nhau).

\subsection{\texorpdfstring{\code{shadow.vert}}{shadow.vert}}
{\footnotesize Giai đoạn: Vertex \;$\cdot$\; 34 dòng \;$\cdot$\; \code{splat_editor/render/shaders/shadow.vert}}\par
\textbf{Vai trò.} Vẽ \textbf{tâm} mỗi gaussian như một điểm vài pixel vào bản đồ độ sâu nhìn từ mặt trời. Làm thế rất rẻ so với dựng ellipse.

\subsubsection{Toàn bộ shader}
{\footnotesize\textit{Dòng 1--34}}\par
\textbf{Giải thích.} Giải half-float (\code{halfToFloat}) để lấy $\alpha$; gaussian đã xóa hoặc quá mờ ($\alpha<0.3$) bị đẩy ra ngoài (\code{gl_Position = (2,2,2,1)}) nên không đổ bóng. Dựng ma trận layer, đưa tâm sang thế giới rồi nhân \code{uLightVP}.
\code{gl_PointSize = uPointSize} quyết định mỗi điểm phủ bao nhiêu pixel của bản đồ bóng (tránh lỗ hổng).

\textbf{Công thức.}
\[
\mathrm{gl\_Position}=\mathrm{VP}_{light}\,M\begin{bmatrix}p_f\\1\end{bmatrix}\qquad(\text{bỏ qua nếu }\alpha<0.3\text{ hoặc đã xóa})
\]


\VerbatimInput[firstline=1,lastline=34,firstnumber=1,numbers=left,numbersep=3pt,frame=lines,framesep=2pt,rulecolor=\color{gray!60},breaklines=true,fontsize=\fontsize{6.6}{7.6}\selectfont]{../../splat_editor/render/shaders/shadow.vert}
\subsection{\texorpdfstring{\code{shadow.frag}}{shadow.frag}}
{\footnotesize Giai đoạn: Fragment \;$\cdot$\; 3 dòng \;$\cdot$\; \code{splat_editor/render/shaders/shadow.frag}}\par
\textbf{Vai trò.} Fragment shader \textbf{rỗng}: bản đồ bóng chỉ cần giá trị độ sâu do phần cứng tự ghi, không cần màu.

\subsubsection{Toàn bộ shader}
{\footnotesize\textit{Dòng 1--3}}\par
\textbf{Giải thích.} \code{main} không làm gì. Độ sâu của điểm được GPU ghi vào depth attachment, sau đó \code{shadowFactor} trong pass 1 so sánh với nó.

\VerbatimInput[firstline=1,lastline=3,firstnumber=1,numbers=left,numbersep=3pt,frame=lines,framesep=2pt,rulecolor=\color{gray!60},breaklines=true,fontsize=\fontsize{6.6}{7.6}\selectfont]{../../splat_editor/render/shaders/shadow.frag}
\subsection{\texorpdfstring{\code{grid.vert}}{grid.vert}}
{\footnotesize Giai đoạn: Vertex \;$\cdot$\; 8 dòng \;$\cdot$\; \code{splat_editor/render/shaders/grid.vert}}\par
\textbf{Vai trò.} Vẽ lưới mà \textbf{không cần vertex buffer}: tạo tam giác lớn phủ toàn màn hình từ \code{gl_VertexID}.

\subsubsection{Toàn bộ shader}
{\footnotesize\textit{Dòng 1--8}}\par
\textbf{Giải thích.} $p=((\mathrm{id}\ll1)\,\&\,2,\ \mathrm{id}\,\&\,2)\cdot2-1$: \code{id=0,1,2} cho $(-1,-1)$, $(3,-1)$, $(-1,3)$. Tam giác này phủ kín hình vuông $[-1,1]^2$ nên mỗi pixel màn hình chạy fragment shader đúng một lần. \code{vNdc} chuyển tọa độ NDC xuống fragment.

\textbf{Công thức.}
\[
p=\big((\mathrm{id}\ll1)\,\&\,2,\ \mathrm{id}\,\&\,2\big)\cdot2-1:\quad \mathrm{id}=0,1,2\ \mapsto\ (-1,-1),\ (3,-1),\ (-1,3)
\]


\VerbatimInput[firstline=1,lastline=8,firstnumber=1,numbers=left,numbersep=3pt,frame=lines,framesep=2pt,rulecolor=\color{gray!60},breaklines=true,fontsize=\fontsize{6.6}{7.6}\selectfont]{../../splat_editor/render/shaders/grid.vert}
%% ---- merged math: grid_vert
\subsubsection{Cơ sở toán: tam giác phủ toàn màn hình}\label{s:math-grid-vert}
Chỉ số đỉnh \code{gl_VertexID} (ký hiệu $\mathrm{id}$) được đổi thành tọa độ $p=((\mathrm{id}\ll1)\,\&\,2,\;\mathrm{id}\,\&\,2)\cdot2-1$, nên $\mathrm{id}=0,1,2$ cho $(-1,-1)$, $(3,-1)$, $(-1,3)$.

\subsection{\texorpdfstring{\code{grid.frag}}{grid.frag}}
{\footnotesize Giai đoạn: Fragment \;$\cdot$\; 44 dòng \;$\cdot$\; \code{splat_editor/render/shaders/grid.frag}}\par
\textbf{Vai trò.} Với mỗi pixel, \textbf{bắn một tia} từ camera, giao với mặt phẳng $y=h$ và tô đường lưới tại điểm giao.

\subsubsection{Uniform và hàm gridLine}
{\footnotesize\textit{Dòng 1--18}}\par
\textbf{Giải thích.} \code{uCamRot} (cột right/up/back của camera), \code{uCamPos}, \code{uProjScale} (hai phần tử chéo của ma trận chiếu), \code{uHeight}, \code{uCell}, \code{uBlend}, \code{uFade}.
\code{gridLine}: $c=p/\mathrm{cell}$, \code{w = fwidth(c)} là độ rộng 1 pixel theo đơn vị ô, $g=|\mathrm{fract}(c-\tfrac12)-\tfrac12|/w$, đường $=1-\min(\min(g_x,g_y),1)$: đường luôn dày khoảng 1 pixel dù zoom thế nào (khử răng cưa theo đạo hàm).

\textbf{Công thức.}
\[
c=\frac{p}{\mathrm{cell}},\quad w=\mathrm{fwidth}(c),\quad g=\frac{\big|\mathrm{fract}(c-\tfrac12)-\tfrac12\big|}{w},\quad \mathrm{line}=1-\min\big(\min(g_x,g_y),1\big)
\]


\VerbatimInput[firstline=1,lastline=18,firstnumber=1,numbers=left,numbersep=3pt,frame=lines,framesep=2pt,rulecolor=\color{gray!60},breaklines=true,fontsize=\fontsize{6.6}{7.6}\selectfont]{../../splat_editor/render/shaders/grid.frag}
\subsubsection{Bắn tia và giao với mặt sàn}
{\footnotesize\textit{Dòng 20--25}}\par
\textbf{Giải thích.} Hướng tia $\mathrm{dir}=\mathrm{CamRot}\cdot(x/P_{00},\,y/P_{11},\,-1)$ (ngược phép chiếu). Giao điểm $t=(h-e_y)/\mathrm{dir}_y$, $p=e+t\,\mathrm{dir}$. Bỏ pixel nếu tia song song sàn ($|\mathrm{dir}_y|<10^{-6}$) hoặc sàn nằm sau camera ($t\le0$).

\textbf{Công thức.}
\[
\hat d=C_{rot}\Big(\frac{x_{ndc}}{P_{00}},\ \frac{y_{ndc}}{P_{11}},\ -1\Big),\qquad t=\frac{h-e_y}{\hat d_y},\qquad p=e+t\,\hat d
\]


\VerbatimInput[firstline=20,lastline=25,firstnumber=20,numbers=left,numbersep=3pt,frame=lines,framesep=2pt,rulecolor=\color{gray!60},breaklines=true,fontsize=\fontsize{6.6}{7.6}\selectfont]{../../splat_editor/render/shaders/grid.frag}
\subsubsection{Đường lưới, trục tọa độ, mờ dần}
{\footnotesize\textit{Dòng 27--44}}\par
\textbf{Giải thích.} Ô nhỏ nhân $(1-\mathrm{blend})$ và ô lớn $10\times$ tạo hiệu ứng chuyển mức chi tiết liên tục khi zoom. Trục $X$ (đỏ, $z=0$) và trục $Z$ (xanh, $x=0$) tô đè lên. Mờ theo khoảng cách $1-\mathrm{smoothstep}(fade,3\,fade,dist)$ và theo góc nhìn
(mờ khi nhìn gần song song mặt sàn). Xuất premultiplied alpha.

\textbf{Công thức.}
\[
\alpha=\max(0.35\,\mathrm{minor}\,(1-\mathrm{blend}),\ 0.7\,\mathrm{major})\cdot\big(1-\mathrm{smoothstep}(f,3f,\mathrm{dist})\big)\cdot\mathrm{smoothstep}(0,0.15,|\hat d_y|)
\]


\VerbatimInput[firstline=27,lastline=44,firstnumber=27,numbers=left,numbersep=3pt,frame=lines,framesep=2pt,rulecolor=\color{gray!60},breaklines=true,fontsize=\fontsize{6.6}{7.6}\selectfont]{../../splat_editor/render/shaders/grid.frag}
%% ---- merged math: grid_frag
\subsubsection{Cơ sở toán: lưới sàn}\label{s:math-grid-frag}
Tia từ camera: $\mathrm{dir}=\mathrm{CamRot}\cdot(\mathrm{ndc}_x/P_{00},\,\mathrm{ndc}_y/P_{11},\,-1)$ (ngược phép chiếu). \textbf{Giao tia--mặt phẳng $y=h$:} $t=(h-e_y)/\mathrm{dir}_y$, $p=e+t\cdot\mathrm{dir}$ (bỏ nếu $|\mathrm{dir}_y|<10^{-6}$ hoặc $t\le0$).

Đường lưới: $c=p/\mathrm{cell}$, $w=\mathrm{fwidth}(c)$ (đạo hàm màn hình, tức độ rộng 1 pixel), $g=|\mathrm{fract}(c-\tfrac12)-\tfrac12|/w$, $\mathrm{line}=1-\min(\min(g_x,g_y),1)$ (khử răng cưa theo đạo hàm). Mức chi tiết liên tục: $\mathrm{level}=\log_{10}(\mathrm{size})$, $\mathrm{cell}=10^{\lfloor\mathrm{level}\rfloor-1}$, $\mathrm{blend}=\mathrm{level}-\lfloor\mathrm{level}\rfloor$; lưới lớn $=10\,\mathrm{cell}$. Mờ theo khoảng cách $1-\mathrm{smoothstep}(\mathrm{fade},3\,\mathrm{fade},\mathrm{dist})$ và theo góc nhìn $\mathrm{smoothstep}(0,.15,|\mathrm{dir}_y|)$.

Vị trí camera được lấy từ ma trận view: $e=-R^{\top}t$ (dùng trong \code{splat_renderer.py}, \code{light_renderer.py} và \code{grid_renderer.py}).

\subsection{\texorpdfstring{\code{light_gizmo.vert}}{light_gizmo.vert}}
{\footnotesize Giai đoạn: Vertex \;$\cdot$\; 18 dòng \;$\cdot$\; \code{splat_editor/render/shaders/light_gizmo.vert}}\par
\textbf{Vai trò.} Billboard luôn hướng về camera.

\subsubsection{Toàn bộ shader}
{\footnotesize\textit{Dòng 1--18}}\par
\textbf{Giải thích.} Đưa tâm đèn sang không gian view rồi \textbf{cộng offset trong mặt phẳng màn hình} (\code{viewPos.xy += corner*radius*halo}) nên hình vuông luôn vuông góc với tia nhìn. \code{vPos} là tọa độ tính theo đơn vị bán kính cầu, \code{uHalo} là tỉ lệ quầng sáng.

\textbf{Công thức.}
\[
v=V\,c,\qquad v_{xy}\leftarrow v_{xy}+x_c\cdot R\cdot\mathrm{halo},\qquad \mathrm{gl\_Position}=P\,v
\]


\VerbatimInput[firstline=1,lastline=18,firstnumber=1,numbers=left,numbersep=3pt,frame=lines,framesep=2pt,rulecolor=\color{gray!60},breaklines=true,fontsize=\fontsize{6.6}{7.6}\selectfont]{../../splat_editor/render/shaders/light_gizmo.vert}
%% ---- merged math: light_gizmo
\subsubsection{Cơ sở toán: quả cầu đèn}\label{s:math-light-gizmo}
Billboard: \code{view.xy += corner*R*halo}. Trên hình tròn $r\le1$: $z=\sqrt{1-r^2}$ (cao độ mặt cầu), cho độ sáng $0.7+0.5z$. Bên ngoài là quầng $h=\big(1-\frac{r-1}{\mathrm{halo}-1}\big)^2\cdot0.65$. Bán kính hiển thị $\max(0.12\,\mathrm{unit},\,0.012\,\mathrm{khoảng\ cách})$ (đủ lớn trên màn hình). Phần fragment ở \code{light_gizmo.frag} dùng $r$, $z$ và quầng này để tô màu.

\subsection{\texorpdfstring{\code{light_gizmo.frag}}{light_gizmo.frag}}
{\footnotesize Giai đoạn: Fragment \;$\cdot$\; 21 dòng \;$\cdot$\; \code{splat_editor/render/shaders/light_gizmo.frag}}\par
\textbf{Vai trò.} Biến hình vuông thành quả cầu phát sáng cộng quầng.

\subsubsection{Toàn bộ shader}
{\footnotesize\textit{Dòng 1--21}}\par
\textbf{Giải thích.} Trong hình tròn $r\le1$: $z=\sqrt{1-r^2}$ là cao độ mặt cầu, màu $=\mathrm{uColor}\cdot(0.7+0.5z)$ (sáng nhất ở giữa), pha với màu tối $0.12$ theo \code{uBrightness} (đèn tắt thì mờ). Ngoài hình tròn là \textbf{quầng}:
$h=\big(1-\tfrac{r-1}{\mathrm{halo}-1}\big)^2\cdot0.65\cdot\mathrm{brightness}$, xuất với alpha $0$ nên được \emph{cộng thêm} vào cảnh (additive).

\textbf{Công thức.}
\[
r\le1:\ z=\sqrt{1-r^2},\ c=c_0\,(0.7+0.5z);\qquad
r>1:\ h=\Big(1-\frac{r-1}{\mathrm{halo}-1}\Big)^2\cdot0.65\cdot\mathrm{brightness}
\]


\VerbatimInput[firstline=1,lastline=21,firstnumber=1,numbers=left,numbersep=3pt,frame=lines,framesep=2pt,rulecolor=\color{gray!60},breaklines=true,fontsize=\fontsize{6.6}{7.6}\selectfont]{../../splat_editor/render/shaders/light_gizmo.frag}
\subsection{\texorpdfstring{\code{firefly.vert}}{firefly.vert}}
{\footnotesize Giai đoạn: Vertex \;$\cdot$\; 22 dòng \;$\cdot$\; \code{splat_editor/render/shaders/firefly.vert}}\par
\textbf{Vai trò.} Hạt bay quanh đèn, \textbf{toàn bộ chuyển động tính trên GPU} từ \code{gl_VertexID} và thời gian, không cần buffer vị trí.

\subsubsection{Toàn bộ shader}
{\footnotesize\textit{Dòng 1--22}}\par
\textbf{Giải thích.} \code{phase = i*2.399963} (góc vàng) cho các hạt pha lệch đều. Offset là tổ hợp \code{sin}/\code{cos} với tần số hơi khác nhau theo từng hạt (\code{0.4+0.13 sin(1.7 i)}) nên quỹ đạo không lặp cứng; nhân bán kính $(0.35+0.65\,\mathrm{fract}(0.618i))$ để phân bố đều theo bán kính (dãy Weyl, $0.618$ là nghịch đảo tỉ lệ vàng).
\code{gl_PointSize = uPixel} đặt kích thước hạt; \code{vTwinkle} là hệ số nhấp nháy $0.5+0.5\sin(2t+3\varphi)$.

\textbf{Công thức.}
\[
\varphi_i=2.399963\,i\ (=\pi(3-\sqrt5)\,i),\qquad
\mathrm{offset}_i=\begin{bmatrix}\sin(\omega_it+\varphi_i)\\0.6\sin(0.7t+1.3\varphi_i)+0.2\\\cos(\omega_i't+0.7\varphi_i)\end{bmatrix}R\,\big(0.35+0.65\,\mathrm{fract}(0.618\,i)\big)
\]
\[
\mathrm{twinkle}=0.5+0.5\sin(2t+3\varphi_i)
\]


\VerbatimInput[firstline=1,lastline=22,firstnumber=1,numbers=left,numbersep=3pt,frame=lines,framesep=2pt,rulecolor=\color{gray!60},breaklines=true,fontsize=\fontsize{6.6}{7.6}\selectfont]{../../splat_editor/render/shaders/firefly.vert}
%% ---- merged math: firefly
\subsubsection{Cơ sở toán: đom đóm}\label{s:math-firefly}
Góc vàng $i\cdot2.399963$ ($=\pi(3-\sqrt5)$) cho pha. Offset $=(\sin(\omega t+\varphi),\;0.6\sin(0.7t+1.3\varphi)+0.2,\;\cos(\omega't+0.7\varphi))\cdot R\cdot(0.35+0.65\,\mathrm{fract}(0.618\,i))$, trong đó $0.618$ là nghịch đảo tỉ lệ vàng (dãy Weyl phân bố đều bán kính). Nhấp nháy $0.5+0.5\sin(2t+3\varphi)$. Điểm mờ: $a=\mathrm{clamp}(1-d,0,1)^2\cdot\mathrm{twinkle}$ (phần mờ này tính ở \code{firefly.frag}).

\subsection{\texorpdfstring{\code{firefly.frag}}{firefly.frag}}
{\footnotesize Giai đoạn: Fragment \;$\cdot$\; 12 dòng \;$\cdot$\; \code{splat_editor/render/shaders/firefly.frag}}\par
\textbf{Vai trò.} Biến điểm vuông thành chấm sáng mờ dần.

\subsubsection{Toàn bộ shader}
{\footnotesize\textit{Dòng 1--12}}\par
\textbf{Giải thích.} \code{gl_PointCoord} là tọa độ trong điểm $[0,1]^2$; $d=2\,|\mathrm{coord}-0.5|$ là khoảng cách tới tâm. $a=\mathrm{clamp}(1-d,0,1)^2\cdot\mathrm{twinkle}$. Xuất alpha $0$ để cộng thêm vào cảnh như quầng đèn.

\textbf{Công thức.}
\[
d=2\,\|\mathrm{coord}-0.5\|,\qquad a=\mathrm{clamp}(1-d,0,1)^2\cdot\mathrm{twinkle}
\]


\VerbatimInput[firstline=1,lastline=12,firstnumber=1,numbers=left,numbersep=3pt,frame=lines,framesep=2pt,rulecolor=\color{gray!60},breaklines=true,fontsize=\fontsize{6.6}{7.6}\selectfont]{../../splat_editor/render/shaders/firefly.frag}
\subsection{\texorpdfstring{\code{avatar.vert}}{avatar.vert}}
{\footnotesize Giai đoạn: Vertex \;$\cdot$\; 14 dòng \;$\cdot$\; \code{splat_editor/render/shaders/avatar.vert}}\par
\textbf{Vai trò.} Vẽ 6 hộp cơ thể; mỗi hộp có ma trận \code{uModel} riêng nên có thể xoay quanh khớp.

\subsubsection{Toàn bộ shader}
{\footnotesize\textit{Dòng 1--14}}\par
\textbf{Giải thích.} Ba thuộc tính đỉnh: vị trí, pháp tuyến, UV (khớp bố cục \code{build_mesh}). Pháp tuyến đổi sang thế giới bằng \code{mat3(uModel)}; \code{gl_Position = uProj*uView*uModel*pos}. UV và pháp tuyến truyền xuống fragment.

\textbf{Công thức.}
\[
n_w=\mathrm{normalize}(M_{3\times3}\,n),\qquad \mathrm{gl\_Position}=P\,V\,M\,p
\]


\VerbatimInput[firstline=1,lastline=14,firstnumber=1,numbers=left,numbersep=3pt,frame=lines,framesep=2pt,rulecolor=\color{gray!60},breaklines=true,fontsize=\fontsize{6.6}{7.6}\selectfont]{../../splat_editor/render/shaders/avatar.vert}
\subsection{\texorpdfstring{\code{avatar.frag}}{avatar.frag}}
{\footnotesize Giai đoạn: Fragment \;$\cdot$\; 14 dòng \;$\cdot$\; \code{splat_editor/render/shaders/avatar.frag}}\par
\textbf{Vai trò.} Tô skin và chiếu sáng Lambert đơn giản.

\subsubsection{Toàn bộ shader}
{\footnotesize\textit{Dòng 1--14}}\par
\textbf{Giải thích.} Lấy mẫu skin $64\times64$ bằng \code{texture}; \code{discard} nếu $\alpha<0.5$ (lớp áo trong suốt). $\mathrm{shade}=\mathrm{bright}\,(0.62+0.38\max(n\cdot l,0))$: 0.62 là ambient, 0.38 là Lambert.

\textbf{Công thức.}
\[
\mathrm{shade}=\mathrm{bright}\,\big(0.62+0.38\max(n\cdot l,0)\big),\qquad \mathrm{fragColor}=(\mathrm{texel}_{rgb}\cdot\mathrm{shade},\ 1)
\]


\VerbatimInput[firstline=1,lastline=14,firstnumber=1,numbers=left,numbersep=3pt,frame=lines,framesep=2pt,rulecolor=\color{gray!60},breaklines=true,fontsize=\fontsize{6.6}{7.6}\selectfont]{../../splat_editor/render/shaders/avatar.frag}

\clearpage
%====================================================================
\section{Bảng tra nhanh: công thức và file}\label{ch:quickref}
%====================================================================
Bảng dưới liệt kê từng nhóm công thức và nơi nó được cài đặt; cột cuối trỏ tới mục
\textit{Cơ sở toán} tương ứng nằm ngay trong chương của file đó.

\begin{center}\small
\begin{tabular}{@{}p{0.34\textwidth}p{0.46\textwidth}p{0.1\textwidth}@{}}\toprule
Chủ đề & File:hàm & Mục\\\midrule
sigmoid, exp, $C_0$, chuẩn hóa quaternion & \code{io/codec.py: _decode_block, encode} & \S\ref{s:codec}\\
Euler, Rodrigues, quaternion & \code{core/math3d.py} & \S\ref{s:quat}, \S\ref{s:layer}\\
$M=[sA\,|\,t]$, xoay quanh pivot & \code{core/layer.py: world_matrix}; \code{core/transform.py: rotate_about} & \S\ref{s:layer}\\
view, projection, ortho & \code{core/math3d.py: look_at, perspective, orthographic} & \S\ref{s:cam}\\
orbit, pan, fly, frame & \code{render/camera.py} & \S\ref{s:cam-orbit}\\
EWA $\Sig'=JW\Sig W^{\top}J^{\top}$, eigen, extent & \code{render/shaders/splat_prep.vert: main} & \S\ref{s:ewa}\\
Gauss và trộn over & \code{render/shaders/splat.frag}; \code{splat_renderer.py: _draw_splats} & \S\ref{s:blend}, \S\ref{s:blend-frag}\\
SH cơ sở và xoay & \code{core/sh.py}; \code{splat_prep.vert: shColor} & \S\ref{s:shrot}\\
Phong, suy giảm, wrap & \code{splat_prep.vert: shade} & \S\ref{s:phong}\\
Mặt trời, ngày--đêm, bóng & \code{core/lighting.py} & \S\ref{s:sun}\\
PCF & \code{splat_prep.vert: shadowFactor} & \S\ref{s:phong}, \S\ref{s:sun}\\
Morton, radix, frustum, budget & \code{core/lod.py; splat_renderer.py: adapt_budget} & \S\ref{s:lod}, \S\ref{s:budget}\\
Sắp xếp độ sâu 16 bit & \code{render/sorter.py} & \S\ref{s:sort}\\
Cọ capsule, viewport & \code{editor/brush.py} & \S\ref{s:brush}, \S\ref{s:viewport}\\
Voxel, khóa 63 bit, trọng lực & \code{core/physics.py} & \S\ref{s:phys}\\
Hoạt họa và ma trận khớp & \code{core/avatar_pose.py} & \S\ref{s:anim}\\
Grid ray--plane, fwidth & \code{render/shaders/grid.frag}; \code{grid_renderer.py} & \S\ref{s:math-grid-frag}, \S\ref{s:math-grid-vert}\\
Skin UV, Paeth PNG & \code{render/avatar_renderer.py} & \S\ref{s:texture}\\
Half-float, lượng tử hóa & \code{splat_renderer.py: pack}; shaders; quaternion 8 bit & \S\ref{s:quant}, \S\ref{s:quat}\\
Hoàn tác (Command, 20 bước), bit trạng thái & \code{core/history.py} & \S\ref{s:math-undo}\\
Quả cầu đèn (billboard, quầng) & \code{render/shaders/light_gizmo.*} & \S\ref{s:math-light-gizmo}\\
Đom đóm (góc vàng, Weyl) & \code{render/shaders/firefly.*} & \S\ref{s:math-firefly}\\\bottomrule
\end{tabular}
\end{center}

\end{document}
